\documentclass[11pt]{article}
\usepackage[a4paper,margin=28mm]{geometry}
\usepackage[T1]{fontenc}
\usepackage{lmodern,amsmath,amssymb,amsthm,mathtools,booktabs,longtable}
\usepackage[round,authoryear]{natbib}
\usepackage{xurl}
\usepackage[colorlinks=true,allcolors=blue]{hyperref}
\hypersetup{hypertexnames=false}
\usepackage{microtype}
\usepackage{needspace}
\setlength{\emergencystretch}{2em}
\newcommand{\Q}{\mathbb Q}
\newcommand{\Z}{\mathbb Z}
\newcommand{\F}{\mathbb F}
\newcommand{\C}{\mathbb C}
\newcommand{\Cl}{\operatorname{Cl}}
\newcommand{\Spec}{\operatorname{Spec}}
\newcommand{\Hom}{\operatorname{Hom}}
\newcommand{\Lie}{\operatorname{Lie}}
\newcommand{\loc}{\operatorname{loc}}
\newcommand{\inv}{\operatorname{inv}}
\newcommand{\divisor}{\operatorname{div}}
\newcommand{\red}{\operatorname{red}}
\newcommand{\tors}{\operatorname{tors}}
\newcommand{\et}{\mathrm{\acute et}}
\numberwithin{equation}{section}
\newcommand{\A}{\mathbb A}
\newcommand{\G}{\mathbb G}
\newcommand{\wh}{\widehat}
\newcommand{\Sha}{\operatorname{Sha}}
\newcommand{\Sec}{\operatorname{Sec}}
\newcommand{\Sel}{\operatorname{Sel}}
\newcommand{\Gal}{\operatorname{Gal}}
\newcommand{\Jac}{\operatorname{Jac}}
\newcommand{\ab}{\mathrm{ab}}
\newcommand{\sol}{\mathrm{sol}}
\newcommand{\pc}{\mathrm{pc}}
\newcommand{\workbenchurl}{https://github.com/KokunoYumeto/k-theory-and-galois-groups-workbench}
\newcommand{\recordurl}{https://github.com/KokunoYumeto/mathematics-workbenches}
\newtheorem{theorem}{Theorem}[section]
\newtheorem{proposition}[theorem]{Proposition}
\newtheorem{lemma}[theorem]{Lemma}
\newtheorem{corollary}[theorem]{Corollary}
\theoremstyle{definition}
\newtheorem{definition}[theorem]{Definition}
\newtheorem{remark}[theorem]{Remark}
\title{Fermat's theorem for Selmer sections in solvable quotients\\
\large Grothendieck's programme, after Kummer, McCallum, P\'al, Stoll, and Sa\"idi}
\author{}
\date{}
\begin{document}
\maketitle
\begin{abstract}
Grothendieck asked whether sections of the arithmetic fundamental group
of a Fermat curve arise from its evident boundary points. We study the
point-or-cusp Selmer form of this question in geometric quotients.
For every odd exponent divisible by three, five, or a regular prime,
every Selmer section in every finite derived quotient and in the maximal
prosolvable quotient is cuspidal. The section set is exactly three
disjoint packets, each parametrized by the full completed Kummer group.
Kummer's cyclotomic unit argument excludes first-case local centres;
a primary refinement of McCallum's supported-class argument isolates
the cusps. A global unit then identifies all local centres, and
Sa\"idi's localization theorem gives the global section.
For every regular prime greater than five, the full geometric
pro-prime quotient and the stated central and derived quotients have
the same three-packet classification throughout the prime-power tower.
At irregular primes, a cyclotomic class-rank bound gives one fixed set
of at most the exponent prime containing every compact local centre;
this applies at thirty-seven, fifty-nine, and sixty-seven.
We also prove the full geometrically pro-three classification at
exponents \(3^r\), \(5\cdot3^r\), and \(7\cdot3^r\), for \(r\ge1\),
using explicit arithmetic sources. The remaining pro-three exponents
and the complete irregular-prime classification remain open.
Separate applications of Stoll's finite-descent criterion provide
proper-curve results at exponents five, eleven, thirteen, and nineteen,
with explicit Jacobian factors, analytic certificates, and the precise
published arithmetic inputs. Two families of noncuspidal pro-three
sections fail the local Selmer condition at three. P\'al's finite-action
theorem also separates different coordinate-boundary components in
the proper pro-three point map of every Fermat curve over finite
fields of characteristic prime to three and the exponent. Boundary
fibres have a uniform cyclic-block description; they are singletons
whenever three divides the exponent.
The mathematics and text were produced by GPT-6 Astra Pro in ChatGPT
Pro mode, selected principal statements were separately re-proved by
the same model, and no human has checked this paper.
\end{abstract}

\section{Conventions and scope}

In \emph{La Longue Marche \`a travers la th\'eorie de Galois},
\citet[\S51.16, fols.~733--736v]{Grothendieck1981} formulates the
Fermat problem through sections of the arithmetic fundamental sequence.
On folio~734 he asks whether all sections come from the points at
infinity, explaining that this would establish the desired Fermat
conclusion without first establishing the full anabelian programme.
His calculation of an abelian approximation on folios~736--736v ends
with the observation that this approximation is too coarse and that
finer ones are needed. The quotient distinctions below address that
specific issue; they do not identify arbitrary sections with Selmer
sections.

The nearest full-group statement is
\citet[Theorem~A and \S5]{Steklov2026}, which asserts cuspidality for
full profinite Selmer sections at prime exponents at least seven.
Its argument uses residual Legendre representations and modular
finite descent. The cuspidal divisibility assertion used in
Proposition~5.10 and the exponent-seven endpoint in Theorem~5.14
require separate corrections. Independently of those issues, a
classification in the full fundamental group does not classify
sections in a geometric quotient: a quotient section need not have
a full global lift. The results below establish cuspidality in the
stated quotients with full local point-or-cusp witnesses and without
assuming such a lift.

The lower-central Galois modules of Fermat curves have also been studied
explicitly. \citet[Theorem~1.1 and Corollary~1.2; author version,
pp.~2--3]{DavisPriesWickelgren2024} determine the defining graded
relation for the proper Fermat curve and, under the stated Vandiver
hypothesis, the second graded module modulo the prime exponent.
\citet[Theorem~A, equivalently Theorem~3.5]{DuqueRoseroPries2024}
give the corresponding classifying element for the totally ramified
cyclic quotients $v^n=u(1-u)^k$.
Those results describe graded Galois modules. The full pro-$3$ section
argument below additionally uses local geometric inertia, positive
Frobenius weights, and localization on an arithmetic Galois source.

At exponent thirteen we use the classical Hurwitz--Klein decomposition.
\citet[Theorem~6.2 and Corollary~6.3, pp.~89--90]{Lim1992}
identifies the cyclotomic decomposition, while
\citet[Proposition~2.1 and Corollary~2.2, p.~942]{Tzermias2004}
gives the rational complementary factor and its rank parity.
Together with the rank-one result of Gross--Rohrlich, these already
give a rank-zero fourfold. Section~\ref{subsec:thirteen} gives explicit
equations and a central-value certificate. The arithmetic step in
Section~\ref{subsec:lm-lz-rational} uses the ordinary family constructions
and abstract Euler-system arguments of \citet{LoefflerZerbes2025},
\citet{LoefflerRockwoodZerbes2024}, and \citet{LoefflerZerbes2026},
with the precise \citet{MazurRubin2004} hypotheses verified directly.
In this finite-descent application, the identification of the rational
Fermat points uses
\citet[Theorem~0.5, p.~448]{Wiles1995}. The regular-prime section theorem
below has a separate local proof and does not require that
exponent-thirteen arithmetic application.

For every positive integer $n$, the notation $X_n$ means the all-plus
curve $X^n+Y^n+Z^n=0$. For odd $n$, changing the sign of $Z$ identifies
it with $X^n+Y^n=Z^n$. For even $n$, the all-plus curve has no real point.

Let $k$ be a number field, $X/k$ a smooth proper geometrically connected
curve, $D\subset X$ a reduced nonempty divisor, and $U=X\setminus D$.
Write
\[
1\longrightarrow\Delta\longrightarrow\Pi\longrightarrow G_k
\longrightarrow1
\]
for its arithmetic fundamental exact sequence. For a finite set $S$ of finite places, $G_{k,S}$ denotes the Galois
group of the maximal extension unramified outside $S$; ramification at
archimedean places is allowed. All subgroups generated
by commutators below are closed. For a profinite group $R$, its ordinary
lower central series is $\gamma_1R=R$ and
$\gamma_{m+1}R=\overline{[R,\gamma_mR]}$. Put $\Delta^{[0]}=\Delta$ and
$\Delta^{[m+1]}=[\Delta^{[m]},\Delta^{[m]}]$, and define
\[
\Delta_m=\Delta/\Delta^{[m]},\qquad
\Delta^{\sol}=\Delta/\bigcap_{m\geq1}\Delta^{[m]},\qquad
\Delta^{(3)}=\text{the maximal pro-$3$ quotient of }\Delta.
\]
Thus $\Delta_1$ is abelian and $\Delta_2$ is metabelian.
The corresponding arithmetic group is obtained by dividing $\Pi$ by
the indicated geometric kernel; the quotient $G_k$ is retained in full.
In particular a geometrically pro-$3$ arithmetic group is not the maximal
pro-$3$ quotient of $\Pi$.

\begin{definition}\label{def:point-cusp-selmer}
For any of these arithmetic quotients, a \emph{point-or-cusp Selmer
section} is a continuous section, considered modulo geometric
conjugacy, whose restriction at every place $v$ is the image either of
a point section from $U(k_v)$ or of a section of the decomposition group
of a $k_v$-rational point of $D$. We write $\Sec^{\pc}_{\Sel}$ for this
set. A global section is \emph{cuspidal} if it factors, up to geometric
conjugacy, through the image of the decomposition group of one point
of $D(k)$.
\end{definition}

The local condition in this definition concerns the quotient itself;
it does not require a lift of the global section to the full fundamental
group. A nonzero rational tangent at $Q\in D(k)$ chooses an origin among
the splittings of its punctured formal neighbourhood. These splittings
are parametrized by
\[
H^1(k,\wh\Z(1))=\varprojlim_N k^\times/(k^\times)^N.
\]
They induce the cusp packet maps in the geometric quotients.
Their injection and disjointness are proved under the hypotheses of
each theorem below. For the pro-$3$ group replace $\wh\Z(1)$ by
$\Z_3(1)$. These are completed Kummer groups, not merely the images
of ordinary elements of $k^\times$.

This convention differs from the affine $S$-Selmer convention of
\citet[Definitions 1.5 and 2.1]{BKL2026}, which requires ordinary local
points, with integrality outside $S$. On the proper curve $X$, both
conventions give ordinary local points. At real places we record
connected components; complex places are omitted from the modified
adelic product, as required by \citet[\S1]{StollErrata2017}.
We write $X(\mathbb A_k)_\bullet^{\mathrm{f\text{-}ab}}$ for the
modified adelic points surviving finite abelian descent: for every
finite étale torsor under a finite commutative $k$-group scheme,
some twist defined over $k$ has local lifts of the given adelic
point at every place \citep{Stoll2007}.

\subsection{Main results}

A prime \(p\) is \emph{regular} if
\(p\nmid h(\Q(\zeta_p))\). The first theorem answers Grothendieck's
Fermat question for point-or-cusp Selmer sections in the stated
all-prime quotients. Its local condition is exactly
Definition~\ref{def:point-cusp-selmer}, with the full arithmetic
quotient \(G_\Q\).

\begin{theorem}[Fermat sections at exponents with a regular prime factor]
\label{thm:regular-solvable-intro}
Let \(n\ge3\) be odd and divisible by \(3\), by \(5\), or by a regular
prime \(p>5\). For every finite derived quotient
\(\Delta/\Delta^{[d]}\), \(d\ge1\), of the geometric fundamental
group of \(U_n\), and for its maximal prosolvable quotient, every
point-or-cusp Selmer section is cuspidal. The section set is the
disjoint union of the three rational cusp packets at
\[
 [0:1:-1],\qquad [1:0:-1],\qquad [1:-1:0],
\]
each a torsor under \(H^1(\Q,\widehat\Z(1))\).
This includes the geometric abelian and metabelian quotients.
The input section need not lift globally to the full fundamental
group, and its ramification need not have finite support.
\end{theorem}

Thus the cusps determine both the centres and every possible section
parameter, including completed Kummer classes with infinite support.
The first prime not covered by this theorem is \(37\).
The theorem does not identify the finite-abelian-descent adelic set
of every proper Fermat curve with its rational points; that stronger
proper-curve assertion is obtained in the separate applications below.

The general transfer theorem and its split-prime variant explain those
proper-curve applications. The later statements give the single-prime
classification, the finite local bounds at irregular primes, and the
distinct conclusions for geometrically pro-\(3\) groups.

\begin{theorem}[Finite descent and every derived quotient]
\label{thm:transfer}
Suppose $g(X)\geq1$, $\#D(\bar k)\geq2$, and
$X(\A_k)_\bullet^{\mathrm{f\text{-}ab}}=X(k)$. For every $m\geq1$
and for the maximal geometrically prosolvable quotient,
\begin{equation}
\Sec^{\pc}_{\Sel}(U/k)=
 U(k)\ \sqcup\ \coprod_{Q\in D(k)}H^1(k,\wh\Z(1)),
\label{eq:classification}
\end{equation}
where the maps on the right are the ordinary point maps and the cusp
packet maps, respectively. In particular, if $X(k)\subset D(k)$ every
such Selmer section is cuspidal.
\end{theorem}

\begin{theorem}[A completely split-prime descent criterion]
\label{thm:split-prime-descent}
Let $X/\Q$ be a smooth proper geometrically connected curve with a
nonconstant morphism $f:X\to A$ to an abelian variety. Suppose
$A(\Q)$ is finite and $f(X)_{\bar\Q}$ is not a torsion translate
of an elliptic subvariety. If a finite Galois extension $L/\Q$
satisfies
\[
T_\ell\Sha(\Q,A)=0
\quad\text{for all but finitely many primes $\ell$ splitting
completely in $L$},
\]
then $X(\A_\Q)_\bullet^{\mathrm{f\text{-}ab}}=X(\Q)$.
Thus Theorem~\ref{thm:transfer} applies to every affine open with
at least two geometric boundary points. No rank or
Tate--Shafarevich hypothesis over $L$ is required.
\end{theorem}


\begin{theorem}[The exponent-thirteen split-prime arithmetic input]
\label{thm:thirteen-selmer-intro}
Let $\mathcal K:A^4B+B^4C+C^4A=0$, let $c[A:B:C]=[C:A:B]$,
and put $B=\mathcal K/\langle c\rangle$.
Define the rational abelian fourfold
$P_{13}=(1-c)\Jac(\mathcal K)$, and put
\[
 K=\Q(\zeta_{13})^{\langle3\rangle},\qquad
 L=K\Q(\zeta_3).
\]
There is a nonconstant rational morphism $X_{13}\to P_{13}$ whose
image generates $P_{13}$ by differences. Moreover,
\[
 P_{13}(\Q)\simeq\Z/13\Z,\qquad
 2.632<L(P_{13}/\Q,1)<2.636.
\]
For all but finitely many rational primes $p$ splitting completely
in the finite Galois extension $L/\Q$,
\[
 H^1_f(\Q,V_pP_{13})=0,\qquad T_p\Sha(P_{13}/\Q)=0,
 \qquad V_pP_{13}=T_pP_{13}\otimes_{\Z_p}\Q_p.
\]
Here $H^1_f$ is the rational Bloch--Kato Selmer group. Consequently
$X_n(\A_\Q)_\bullet^{\mathrm{f\text{-}ab}}=X_n(\Q)$ for every
positive multiple $n$ of $13$. No finiteness of the entire group
$\Sha(P_{13}/\Q)$ is asserted.
\end{theorem}

\begin{theorem}[The exponent-nineteen arithmetic input]
\label{thm:nineteen-selmer-intro}
Let $C$ be the smooth projective curve $y^{19}=x(1-x)^7$, put
$J=\Jac(C)$, and let
\[
 c(x,y)=\left(\frac1{1-x},-\frac{y^7}{(1-x)^3}\right),
 \qquad e=\frac{1+c+c^2}{3}.
\]
Let $A=\Jac(C/\langle c\rangle)$ and let $P$ be the image of
$1-e$ in the rational isogeny category. Then
\[
 J\sim_\Q A\times P,\qquad \dim A=3,\quad\dim P=6,
 \qquad \operatorname{rank}A(\Q)=1,\quad
 \operatorname{rank}P(\Q)=0.
\]
The central derivative satisfies
\[
 9.96<L'(J/\Q,1)<9.99.
\]
The groups $\Sha(J/\Q)$, $\Sha(A/\Q)$ and $\Sha(P/\Q)$ are finite.
There is a nonconstant rational morphism $X_{19}\to P$ whose image
generates $P$ by differences. Consequently
$X_n(\A_\Q)_\bullet^{\mathrm{f\text{-}ab}}=X_n(\Q)$ for every
positive multiple $n$ of $19$.
\end{theorem}

\begin{theorem}[Compact localization and full compact section images]
\label{thm:compact-localization-sections-intro}
Let $J_n=\Jac(X_n)$, and let $\Sel(T_3J_n)$ denote the integral
Selmer group with completed Mordell--Weil Kummer images at every
finite place. The localization map
\[
 \Sel(T_3J_n)\longrightarrow H^1(\Q_3,T_3J_n)
\]
is injective for
\[
 \begin{gathered}
 n=3^r\quad(r\ge1),\qquad
 n=5\cdot3^r\ \hbox{or}\ 7\cdot3^r\quad(r\ge0),\qquad n=11.
 \end{gathered}
\]
At $n=11$, the group $\Sha(J_{11}/\Q)$ is finite,
$\operatorname{rank}J_{11}(\Q)=6$, and
$\Sel(T_3J_{11})\simeq\Z_3^6$. More generally, for every rational
prime $\ell$ the map
\[
 \Sel(T_\ell J_{11})\longrightarrow H^1(\Q_\ell,T_\ell J_{11})
\]
is injective integrally, including at $\ell=11$.
For $a=5,7$, let $N_a=\pi_1(U_{a,\overline\Q})^{\mathrm{pro}\text{-}3}$,
let $I_a$ be the closed normal subgroup generated by its boundary inertia,
and let $R$ be a stable quotient of $N_a$ retaining its abelianization.
The compact quotient is
$R^{\mathrm{cp}}=N_a/(\ker(N_a\to R)I_a)$.
Every point-or-cusp Selmer section in $R$ has compact image globally
conjugate to the image of the proper point section at a rational
coordinate cusp supplied by its localization at $3$.
This includes the full compact pro-$3$ image when $R=N_a$.
These conclusions concern compact localization and compact section images;
they do not enlarge the affine pro-$3$ classification in the next theorem.
\end{theorem}

\begin{theorem}[Unrestricted pro-$3$ Selmer sections]
\label{thm:unrestrictedthreefamiliesintro}
Let $n=3^r$, $n=5\cdot3^r$, or $n=7\cdot3^r$, with $r\ge1$, and put
$N_n=\pi_1(U_{n,\overline\Q})^{\mathrm{pro}\text{-}3}$.
Let $R$ be a $G_\Q$-stable quotient of $N_n$ such that
$N_n^{\mathrm{ab}}\to R^{\mathrm{ab}}$ is an isomorphism and every
closed ordinary lower-central factor
$\gamma_mR/\gamma_{m+1}R$ is torsion-free over $\Z_3$.
The point-or-cusp Selmer sections in the corresponding arithmetic
quotient are exactly the three disjoint rational cusp packets,
each parametrized by $H^1(\Q,\Z_3(1))$.

In particular this holds for the full group $N_n$, every quotient
$N_n/\gamma_cN_n$ with $c\ge2$, and every finite-derived quotient
$N_n/N_n^{[d]}$ with $d\ge1$, including the maximal geometrically
metabelian pro-$3$ quotient. Here $N_n^{[d]}$ is the closed derived series.
No finite ramification hypothesis or global lift to the full group is assumed.

For all three exponent families, writing $J_n=\Jac(X_n)$, the compact
integral Selmer localization is injective:
\[
\ker\bigl(\Sel(T_3J_n)\longrightarrow H^1(\Q_3,T_3J_n)\bigr)=0.
\]
Here the local conditions defining $\Sel(T_3J_n)$ are the completed
Mordell--Weil Kummer images at all finite places. No rank or
Tate--Shafarevich hypothesis for the higher $J_n$ is required.
\end{theorem}

\begin{theorem}[Explicit restricted pro-$3$ cases]
\label{thm:restrictedintro}
Let $n=3^r$, $r\geq1$, and let
$N=\pi_1((U_n)_{\bar\Q})^{(3)}$. Fix a rational coordinate tangent
with leading coefficient one, extending over $\Z[1/3]$. Let $M=N/J$
where $J$ is closed, normal, Galois stable, and contained in $[N,N]$;
in particular $M=N$ is allowed. Take any of
\[
S=\{3\},\qquad S=\{3,7\},\qquad
S=\{3,q\}\quad(q\text{ prime},\ q\equiv2\text{ or }5\pmod9).
\]
Every section of $M\rtimes G_\Q$ whose entire relative cocycle is
trivial on inertia outside $S$, and whose restriction at $3$ is
point-or-cusp geometric, is globally cuspidal. Such sections form
three disjoint cusp packets, each parametrized by
$\bigoplus_{q\in S}\Z_3\kappa_3(q)$. They satisfy the point-or-cusp
condition at every place.
\end{theorem}

\begin{theorem}[Arbitrary pro-$3$ sections and the local obstruction]
\label{thm:arbitraryintro}
For every $n=3^r$, $r\geq1$, there are continuum many pairwise
geometrically nonconjugate sections of the full geometrically pro-$3$
arithmetic fundamental group of $U_n/\Q$ whose entire cocycles relative to a standard integral
tangent are unramified outside $3$ and whose compact classes in
$H^1(\Q,T_3E)$, for $E:v^2+v=u^3-7$, have infinite order. They are noncuspidal, and each
restriction at $3$ is neither point geometric nor cusp geometric.
In particular none is a point-or-cusp Selmer section.
Their noncuspidality, pairwise nonconjugacy, and local obstruction
persist in every Galois-stable geometric quotient retaining the
geometric abelianization. A second family of continuum many sections
has zero class in $H^1(\Q,T_3\Jac(X_n))$, with the same ramification,
nonconjugacy, and local obstruction; its two boundary valuation
coordinates can be chosen as $(1,b)$ for $b\in3+9\Z_3$.
\end{theorem}

\begin{theorem}[Local cuspidality at three]
\label{thm:localthreeintro}
Let $n\ge3$ be odd and divisible by at least one of $3$, $5$, $7$,
or $11$. Put $N=\pi_1((U_n)_{\bar\Q})^{(3)}$, and let $M=N/J$,
where $J$ is closed, normal, Galois stable, and contained in $[N,N]$.
For every point-or-cusp Selmer section of the arithmetic quotient
with geometric group $M$, its localization at $3$ is cuspidal at
a rational coordinate cusp. This includes the geometric abelian
quotient and the full geometric pro-$3$ group. Every full local
geometric representative at $3$ has a cusp as its center.
This local theorem alone does not prove global conjugacy to a cusp
section; the global results for the three families above require
their additional arithmetic arguments.
\end{theorem}

\begin{theorem}[A conditional mixed-exponent criterion]
\label{thm:realconditionalmixedintro}
Let $n=3^r q^s$, where $r,s\ge1$ and $q$ is an odd prime with
$q\ne3$. Assume $r=1$ or $q\equiv2,4,5,7\pmod9$, and put
$S=\{3,q\}$ and $J_n=\Jac(X_n)$.
Write $\Sel(T_3J_n)$ for the integral $3$-adic Selmer group, whose
local conditions are the Kummer images of $J_n(\Q_v)^{\wedge_3}$
at every finite place. Assume
\[
\ker\bigl(\Sel(T_3J_n)\longrightarrow
                 H^1(\Q_3,T_3J_n)\bigr)=0.\tag{$\mathrm C_n$}
\]
Then the point-or-cusp Selmer sections of the geometrically pro-$3$
abelian quotient of $U_n/\Q$ are exactly the three disjoint rational
cusp packets, each parametrized by $H^1(\Q,\Z_3(1))$.

More generally, let $R$ be a $G_\Q$-stable quotient of the geometric
pro-$3$ group retaining its abelianization, with finite free ordinary
lower-central factors $L_m=\gamma_mR/\gamma_{m+1}R$.
The same conclusion holds in $R$ if, for every $m\ge2$,
\[
L_m^{G_{\Q_3}}=0,\qquad
\ker\bigl(H^1(G_{\Q,S},L_m)\longrightarrow H^1(\Q_3,L_m)\bigr)=0.
\tag{$\mathrm G_m$}
\]
The conditions $(\mathrm C_n)$ and $(\mathrm G_m)$ are proved here
when $q=5$ or $7$ and $s=1$, giving the mixed families in
Theorem~\ref{thm:unrestrictedthreefamiliesintro}. For the other exponents
they remain explicit hypotheses. Rank $J_n(\Q)=0$ together with
$T_3\Sha(J_n/\Q)=0$ is sufficient for $(\mathrm C_n)$.

The exponent condition is exact for the boundary localization used here.
For any odd prime $q\ne3$ and $r,s\ge1$, write the boundary Tate module
as
$B_n=\Z_3(1)[D_n(\bar\Q)]/\Z_3(1)\sum_{Q\in D_n}e_Q$.
Choose distinct rational cusps $P,Q$, eliminate $e_Q$ using the
diagonal relation, and split off $I_P=\Z_3(1)e_P$ to obtain
$B_n=I_P\oplus M_n$. Then
\[
 H^1(G_{\Q,S},M_n)\longrightarrow H^1(\Q_3,M_n)
 \text{ is injective}\quad\Longleftrightarrow\quad
 r=1\ \text{or}\ q\equiv2,4,5,7\pmod9.
\]
Failure of this cohomological injectivity does not by itself produce a
noncuspidal Selmer section.
\end{theorem}


\begin{theorem}[Regular-prime nonabelian localization]
\label{thm:regular-nonabelian-localization-intro}
Let $p$ be an odd regular prime, meaning $p\nmid h(\Q(\zeta_p))$,
let $n=p^r$ with $r\ge1$, and use a standard integral rational cusp
tangent on $U_n$. Write $N_n^{(p)}$ and $C_n^{(p)}$ for the full
geometric pro-$p$ fundamental groups of $U_n$ and $X_n$.
For $M=N_n^{(p)}$, $C_n^{(p)}$, or any finitely generated
$G_\Q$-stable quotient, the entire map of pointed sets
\[
 H^1(G_{\Q,\{p\}},M)\longrightarrow H^1(\Q_p,M)
\]
is injective. Proper point-local Selmer sections of $C_n^{(p)}$
and its stable quotients are determined by their localization at $p$.
The integral maps
\[
 H^1(G_{\Q,\{p\}},T_pJ_n)\longrightarrow H^1(\Q_p,T_pJ_n),
 \qquad \Sel(\Q,T_pJ_n)\longrightarrow H^1(\Q_p,T_pJ_n)
\]
are injective, with no rank or Tate--Shafarevich finiteness hypothesis.

\end{theorem}

\begin{theorem}[The compact local centre at a regular prime]
\label{thm:regular-primary-local-intro}
Let \(p>5\) be regular, put \(J_p=\Jac(X_p)\), and fix a rational
coordinate cusp \(O\). Suppose
\[
 t\in H^1(\Q,T_pJ_p),\qquad R_v\in X_p(\Q_v)\quad(v<\infty),
 \qquad
 \loc_v(t)=\kappa_{p,v}([R_v-O])
\]
for every finite place, where \(\kappa_{p,v}\) is the inverse-limit
\(p\)-primary Kummer map. Then every permitted \(R_p\) is one of the
three rational coordinate cusps. Consequently the same local-centre
conclusion holds for any quotient section mapping to the proper
geometrically abelian pro-\(p\) quotient and having full local
point-or-cusp representatives.
\end{theorem}

\begin{theorem}[Geometric pro-prime cusp classification at regular primes]
\label{thm:regular-primary-packets-intro}
Let \(p>5\) be regular, let \(n=p^r\) with \(r\ge1\), and put
\(N_n^{(p)}=\pi_1(U_{n,\overline\Q})^{\mathrm{pro}\text{-}p}\).
Let \(R\) be a \(G_\Q\)-stable quotient of \(N_n^{(p)}\) inducing
an isomorphism on abelianization, and suppose that every closed
ordinary lower-central factor of \(R\) is torsion-free over \(\Z_p\).
The point-or-cusp Selmer sections of its arithmetic extension are
exactly the three disjoint rational cusp packets, each a torsor
under \(H^1(\Q,\Z_p(1))\).
This includes the full geometric pro-\(p\) group, every finite-derived
quotient, and every quotient by \(\gamma_c\), \(c\ge2\), of the
ordinary lower central series. No global full section lift or
finite ramification support is assumed.
\end{theorem}

\begin{corollary}[The classified prosolvable fibre at seventeen]
\label{cor:classified-seventeen-fibre-intro}
Let \(n=17^r\), \(r\ge1\), fix a rational cusp \(P\), and fix a
parameter \(\beta_{17}\in H^1(\Q,\Z_{17}(1))\).
Extend it by zero at the other coefficient primes to
\(\widehat\beta\in H^1(\Q,\widehat\Z(1))\), and use the full cusp
section \(t_{P,\widehat\beta}\) to define all actions. Put
\[
 K_{17}=\ker\bigl(N_n^{\mathrm{sol}}\longrightarrow N_n^{(17)}\bigr).
\]
The natural map identifies \(H^1(\Q,K_{17})\) bijectively with the
fibre over the distinguished class of
\[
 H^1(\Q,N_n^{\mathrm{sol}})
       \longrightarrow H^1(\Q,N_n^{(17)}).
\]
Inside this fibre, the point-or-cusp Selmer subset is exactly the
image of the cusp parameters
\[
 H^1\left(\Q,\prod_{\ell\ne17}\Z_\ell(1)\right).
\]
Equivalently, the Selmer lifts of the fixed pro-\(17\) cusp section
form a torsor under that prime-to-\(17\) completed Kummer group.
This identifies a subset of \(H^1(\Q,K_{17})\), not that whole
nonabelian cohomology set.
\end{corollary}

\begin{theorem}[Finite cyclotomic first-case obstructions]
\label{thm:first-case-obstructions-intro}
Let \(p\ge5\), \(K=\Q(\zeta_p)\), and
\(r^-=\dim_{\F_p}\Cl(K)[p]^-\).
In the model \(F_p:X^p+Y^p=Z^p\), put
\(\Delta_p=\pi_1(F_{p,\overline\Q})\).
Suppose a section of
\[
 \pi_1(F_p/\Q)\big/
 \overline{[\Delta_p,\Delta_p]\Delta_p^p}\longrightarrow G_\Q
\]
has at every finite place a representative that is the image of the
full proper point section of \(R_v\in F_p(\Q_v)\).
If \((r^-+2)(r^-+3)-2<p\), every permitted \(R_p\) reduces to a
rational coordinate cusp.

More generally, if \(R_p\) has primitive unit coordinates, use the
signed coordinates \(a,b,c\) with \(a+b+c=0\) modulo \(p\), and for
each ordered pair put \(e=b/(a+b)\in\F_p\).
Define
\[
 C_m(E)=[W^m]\log(1-E+E\exp W)\in\F_p[E],
 \qquad 3\le m\le p-2,
\]
with the series truncated through degree \(p-2\).
For every odd such \(m\), the direct character convention
\(\sigma_i(\zeta_p)=\zeta_p^i\), \(\omega(\sigma_i)=i\), gives
\[
 \Cl(K)[p]^{(\omega^m)}=0\ \Longrightarrow\ C_m(e)=0,
 \qquad B_{p-m}C_m(e)=0\quad\hbox{in }\F_p.
\]
Here \(B_j\) is the \(j\)-th Bernoulli number.
In particular \(p\nmid\operatorname{num}(B_{p-3})\) excludes a
first-case centre. All these assertions apply to every compatible
quotient section with the same full local representatives.
\end{theorem}

\begin{theorem}[One fixed finite set of compact local centres]
\label{thm:finite-local-centres-intro}
Let \(p\ge5\), put \(J_n=\Jac(X_n)\), and use the coordinate cusp
\(O=[0:1:-1]\) on each \(X_n\). Suppose
\[
 r=\dim_{\F_p}\Cl(\Q(\zeta_p))/p\Cl(\Q(\zeta_p))
                 <\frac{p+3}{8}.
\]
There is a set \(S_p\subset X_p(\Q_p)\), containing the three rational
coordinate cusps and of size at most \(p\), with the following
uniform property. For every
\[
 t\in H^1(\Q,T_pJ_p),\quad R_v\in X_p(\Q_v),\quad
 \loc_v(t)=\kappa_{p,v}([R_v-O])\quad(v<\infty),
\]
one has \(R_p\in S_p\); the set is independent of the class and of
all its local representatives. At \(p=37,59,67\), these hypotheses
hold and \(S_p\) can be chosen inside the three cusp residue discs.

For odd \(n=pm\), put \(d=\gcd(m,p-1)\).
There is similarly one fixed set \(S_{n,p}\subset X_n(\Q_p)\)
containing every centre of a compact \(T_pJ_n\) datum of this form,
and every \(\Q_p\)-rational coordinate boundary point, such that
\[
 |S_{n,p}|\le(p-3)d^2+3d.
\]
If \(p\) is the least prime factor of \(n\), then \(d=1\), so the
bound is \(p\). These are finite-set bounds; equality with the set
of rational cusps is not asserted.
\end{theorem}

P\'al proves the section conjecture for prime-order finite actions on
complex curves in the pro-prime group
\citep[Theorem~5.2, author pp.~12--13]{FCPal2014}.
The following application uses his order-three theorem and the
coordinate symmetries of Fermat curves. It separates components of
the boundary over finite fields, where unrestricted point injectivity
for a single geometric prime is false
\citep[Remark~10(i), pp.~742--743]{FCHoshi2009}.

\begin{theorem}[Finite-field boundary fibres]
\label{thm:finitefield-boundary-intro}
Let \(n\ge3\), let \(k=\F_Q\) have characteristic prime to \(3n\),
and let \(C_n:X^n+Y^n+Z^n=0\). Write \(D_X,D_Y,D_Z\) for the
coordinate divisors and \(D\) for their disjoint union. Put
\[
 N=\pi_1(C_{n,\bar k})^{\mathrm{pro}\text{-}3},\qquad
 \Pi=\pi_1(C_n)/\ker\bigl(\pi_1(C_{n,\bar k})\longrightarrow N\bigr),
\]
keeping full \(G_k\) as arithmetic quotient, and let \(\kappa\)
be the proper point-section map modulo \(N\)-conjugacy.

Points in different coordinate components have distinct classes.
If \(D(k)\ne\varnothing\), put \(d=\gcd(n,Q-1)\).
Each component has \(d\) points, and there is one divisor \(h\mid d\)
such that its fibres are cosets of \(\mu_h(k)\) in its regular
\(\mu_d(k)\)-action. There are exactly \(3d/h\) boundary classes,
and \(3\nmid h\). If \(3\mid n\), then \(h=1\): the boundary map
is injective at every allowed finite field.

For every prime exponent \(a\ge5\), the nonsplit case
\(Q\not\equiv1\pmod a\) gives three distinct boundary classes.
In the split case, either all \(3a\) classes are distinct or the
three whole coordinate components are precisely the fibres.
For odd \(n\) prime to three, the sufficient injectivity primes
\(q\nmid3n\), \(\gcd(n,q-1)=1\), have natural density
\[
 \prod_{p\mid n}\frac{p-2}{p-1}.
\]
This is \(3/4\) at \(n=5\) and \(5/6\) at \(n=7\).
For \(3\mid n\), all good primes give injectivity.
\end{theorem}

The fibres are intersected with the coordinate boundary; ordinary
proper points are not excluded. For odd \(n\), that boundary is
nonempty. For even \(n\), it is nonempty exactly when
\(-1\in(k^\times)^n\). The split same-component alternative for
prime exponents different from three remains undecided.
In particular, the theorem does not establish density-one separation
or the affine pro-three classification at the base exponents five
and seven. Its coefficient-prime conclusion comes from the exact
order \(n\) of every coordinate-point divisor difference; the
remaining component separation uses Pál's theorem.

\begin{theorem}[A finite-descent criterion allowing positive rank]
\label{thm:positive-rank-split-intro}
Let $X/\Q$ be a smooth proper curve of genus at least one with
$O\in X(\Q)$, let $f:X\to A$ be a morphism to an abelian variety,
and let $Z\subset X$ be a finite $\Q$-subscheme. Suppose that a
finite Galois extension $L/\Q$ has the following property: for all
but finitely many rational primes $\ell$ splitting completely in
$L$, one has $T_\ell\Sha(A/\Q)=0$ and
\[
 \left\{x\in X(\Q_\ell):
 \log_\ell(f(x)-f(O))\in
 \operatorname{span}_{\Q_\ell}\log_\ell A(\Q)\right\}
       \subset Z(\Q_\ell).
\]
Then $X(\A_\Q)_\bullet^{\mathrm{f\text{-}ab}}=X(\Q)\subset Z(\Q)$.
The local logarithmic condition has not been proved for a
sufficient set of split primes at exponent seventeen.
\end{theorem}


\begin{theorem}[Open Frey images in the prosolvable quotient]
\label{thm:prosolvable-frey-open-intro}
Let $n\ge5$ be odd, and let $s$ be a point-or-cusp Selmer section
of the geometrically prosolvable quotient of $\pi_1(U_n)$.
In the model $X^n+Y^n=Z^n$, the Legendre family
$E_\lambda:y^2=x(x-1)(x-\lambda)$, $\lambda=(X/Z)^n$,
gives representations
$\rho_{s,\ell}:G_\Q\to\operatorname{GL}_2(\Z_\ell)$,
$\ell=2,3$, with determinant $\chi_\ell$.
If $\rho_{s,3}\bmod9$ has a stable cyclic subgroup of order nine,
then $s$ is globally cuspidal.
If $s$ is not globally cuspidal, both representations are absolutely
irreducible and have open images, and their joint image is open in
$\operatorname{GL}_2(\Z_2)\times\operatorname{GL}_2(\Z_3)$.
At a set of rational primes $q$ of Dirichlet density one, its local
representatives are interior and integral and both representations
have the same Frobenius polynomial $T^2-a_qT+q$, with $a_q\in\Z$
and complex roots of absolute value $q^{1/2}$.
No common finite ramification set, modularity, or classification of
this simultaneous open-image case is asserted. The exponent-three
case is already classified by Theorem~\ref{thm:regular-solvable-intro}.
\end{theorem}

At exponent one, an ordinary point gives a noncuspidal Selmer section
in every quotient considered here; see Proposition~\ref{prop:one}.
For even plus exponents the real curve has no points, so the all-place
point-or-cusp Selmer set is empty for every quotient. The unrestricted
pro-$3$ theorem here covers $3^r$, $5\cdot3^r$, and $7\cdot3^r$ for $r\ge1$.
For odd exponents
outside Theorem~\ref{thm:regular-solvable-intro}, the general transfer theorem still
applies whenever its finite-abelian-descent hypothesis is established.
At exponent thirty-one, Section~\ref{subsec:thirtyone-gate} identifies a
complementary tenfold with rational rank zero or two and explains the
remaining analytic and Selmer hypotheses for that abelian variety.
The section classification at \(31\) already follows from
Theorem~\ref{thm:regular-solvable-intro}; this arithmetic question
is separate.

At exponent nineteen, a rank-one theorem of
\citet[Theorem~1.6(2), author p.~3]{McCallumTzermias2004} supplies
an unconditional rank-zero complementary sixfold. The explicit
central-derivative certificate below and
\citet[Theorem~A, p.~28]{Zhang2001} supply its full
Tate--Shafarevich finiteness. The genus-three CM quotient used in
the calculation is the published quartic $X_{15}$ of
\citet[author pp.~2 and 22; Table~1, case~15, p.~5]{KilicerEtAl2018};
its rational identification with this Fermat quotient is already given by
\citet[Remark~7.2, p.~12]{Tafazolian2026}.

\subsection{Proof overview}

At a regular coefficient prime, proper finite Kummer torsors turn
each local representative into a global cyclotomic virtual unit.
The class group and the minus eigenspaces then exclude a first-case
representative at that prime. McCallum's supported cohomology classes,
paired with the finite levels of a single compact Tate-module class,
isolate the rational cusp inside each remaining residue disc.
Regularity also makes the local Galois group dominate the relevant
arithmetic pro-prime source, so the compact class equals that of this
rational cusp globally.

For an all-prime quotient, one coordinate unit has valuation zero at
every finite prime. Its global completed Kummer class is therefore
a sign, and its full local value at the coefficient prime fixes that
sign. This forces every finite-place representative to be the same
cusp. Filling that puncture identifies a unique global cusp parameter,
and Sa\"idi's density-one localization theorem identifies the section
in every finite-derived quotient and in the maximal prosolvable group.
Odd multiplication is injective on the global completed Kummer group,
which extends the result to arbitrary odd multiples of the exponent.

For a geometric pro-prime quotient, compact cancellation instead pins
every boundary reduction to that cusp. Boundary cohomology determines
one global primary parameter. Positive weights, local inertia, and
the regular-prime arithmetic source then kill the successive
lower-central comparison classes. Only those comparison classes
are shown unramified away from the coefficient prime.

At an irregular prime, the same proper torsors give class-rank and
odd-character restrictions. A supported tower selected independently
of the Selmer datum yields one fixed nonzero differential on a wild
cyclic quotient. Finite reciprocity makes its integral vanish at all
permitted local centres, and McCallum's zero bound gives a fixed
finite set. Identifying that set with the cusps needs a further
prescribed local pairing, stated precisely at the end.



Finite abelian descent first forces the local compact centers of a
Selmer section to coincide with one rational point. If that point is a
cusp, filling just that puncture turns the remaining abelian obstruction
into a class that vanishes everywhere locally. Sa\"idi's localization
theorem makes it vanish globally. The resulting global inertia
parameter determines every local cusp section, and Sa\"idi's higher
derived or prosolvable theorem then identifies the global section.

The Fermat applications supply Stoll's hypothesis through a nonconstant
map to a rank-zero abelian variety. The cubic and septic targets are
classical modular elliptic curves. For the quintic, three explicit
maps to a genus-two curve yield the Jacobian isogeny, and exact Euler
coefficients with rigorous Bessel-function tail bounds certify central
nonvanishing. At exponent eleven the classical genus-five factor is
handled by a Mellin integral whose tail is bounded using the
Dedekind zeta function of the cyclotomic field.
Nekov\'a\v r's theorem supplies the arithmetic finiteness in both cases.
The split-prime variant instead combines Manin--Mumford with local
completion at the residue characteristic. Base change makes the
completely split primes have density one among prime ideals, and
Stoll's finite-support theorem then recovers the rational point.
At exponent thirteen the classical rank-zero Prym has two cubic twists
of a cyclic-quartic CM surface as its Tate constituents. Split-prime
CM inertia supplies both the rank-one cokernel element and the
cohomological vanishing required by Mazur--Rubin. Ordinary families,
the completed explicit reciprocity law, and a uniform leading-term
specialization then give rational Bloch--Kato vanishing. The argument
uses characteristic-zero specialization, so it does not require the
large-image integral period normalization that fails for these CM
representations. The split-prime descent criterion completes this
application, using Wiles's theorem only to identify the rational centers.
At exponent nineteen, a rational order-three projector separates the
known rank-one part from a rank-zero sixfold. Three degree-six
Hecke functions replace the larger genus-nine function; exact norm
witnesses determine their Euler factors and a local Gauss sum fixes
the complex root number. Positive Mellin kernels and explicit tails
certify the central derivative. Zhang's rank-one theorem over the
odd-degree real field then gives full Tate--Shafarevich finiteness,
which descends to the rational sixfold.

At exponent eleven, the nine cyclic quotient maps give an integral
three-adic decomposition into three rank-zero factors and six copies
of a rank-one factor. A derivative certificate and Zhang's theorem
give finite Tate--Shafarevich group for the latter. Each completed
rank-one Mordell--Weil group injects locally because its generator
has nonzero three-adic logarithm. This proves compact localization
for the whole Jacobian, while the affine boundary argument remains
a separate requirement.

For the pure-power case of the unrestricted pro-$3$ result, the compact
 cubic first identifies the local center at $3$ with a rational cusp
 at every exponent $3^r$.
 The difference in the higher compact Jacobian is unramified away $3$
 and zero locally at $3$, so arithmetic-source localization kills it.
 Tangent-line divisors make every boundary difference $3$-primary torsion;
 this pins boundary reductions to the same cusp at all other primes.
 Boundary cohomology then gives one global cusp parameter, possibly
 with infinite support. Positive Frobenius weights force fixed points
 to vanish in every central quotient, and arithmetic localization does
 the same at $3$. Thus a local conjugator lies in the current central
 term. Its next additive difference is unramified outside $3$ and locally
 zero at $3$, so it vanishes globally. Induction and compactness prove
 the assertion directly in the retained quotient. The free-polynilpotent
 theorem and an explicit completion argument supply torsion-free grades
 for every finite-derived quotient.

For $5\cdot3^r$, the quintic Jacobian isogeny controls the geometric
action over $K=\Q(A[3])$. CM reciprocity determines this degree-eighty
field and its discriminant. A cyclic subgroup of order five divides
a hypothetical localization kernel into two cases; unconditional
Odlyzko bounds exclude both. The genuine $3$-power cover to $U_5$
passes the action to every $U_{5\cdot3^r}$. This proves compact and
graded localization, while source surjectivity gives the missing
fixed-point vanishing at three. The same central induction then applies.
An independent explicit-formula calculation uses the forced splitting
of the prime above three to prove that this residual field has no
three-primary ideal classes, giving a second source argument.
For the septic family, comparison of CM characters controls all
Jacobian factors over the residual field. Ray classes then reduce
source dominance to its three-primary class group. Complex conjugation
and the ambiguous class number formula reduce its vanishing to a
degree-$56$ field. Exact principal-ideal witnesses force enough small
primes to split in any unramified cubic extension. Their positive
contributions in Odlyzko's unconditional explicit formula make such
an extension impossible. This removes the class-group condition and
proves the septic family with the same central induction.

For the broader conditional mixed-exponent result, compact Selmer localization
first removes the compact difference. The exact order of a cusp
difference then restricts ramification of the remaining boundary
class to the two bad primes. The maximal real subfield of each
cyclotomic boundary orbit has one prime above $q$. Norms, abelian
Leopoldt, and a rank-one anti-invariant unit calculation prove
$S$-unit localization injective. A dimension count shows that this
boundary range is exact. Higher quotients require
separate fixed-point and graded localization hypotheses.

For the further local conclusion, two maps to the same proper
hyperelliptic quotient have zero $3$-adic Selmer image. Its zero
local Kummer locus consists of four explicit points. An ordinary
Fermat point would have to map to the unique noncuspidal one under
both maps, and their coordinate equations are inconsistent. The
septic case uses a separate genus-three central-value certificate.

For the finite-field boundary theorem, equality of two point sections
is a Galois-invariant geometric path. The finite-field centralizer
lemma makes this path unique, so invariance under an automorphism
group produces a coherent splitting of its geometric action
extension. P\'al's theorem gives no splitting class for
the coordinate three-cycle when three divides the exponent, and
otherwise exactly two classes, exchanged by a transposition.
A splitting for the whole coordinate permutation group is therefore
impossible. Two twisted transpositions then separate different
coordinate components, including for even exponents. The cyclic
diagonal action makes every remaining fibre a subgroup coset.
A tangent-line pole calculation proves that every difference of
distinct coordinate points has exact order equal to the exponent
in the Jacobian. Its three-primary part is nonzero when three
divides the exponent, giving full boundary injectivity in that case.

For the separate one-place criterion, explicit cyclotomic Kummer
calculations show that the local Galois group at $3$ surjects onto the
global arithmetic source through which the whole section cocycle
factors. One local conjugation into cusp inertia then works globally.
This last argument explains exactly why the ramification condition is
on the complete cocycle.
The abstract noncuspidal sections come from the rank-one compact
cohomology of a free pro-$3$ arithmetic source. A free pair exchanged
by complex conjugation allows its generator values to lift directly
to the full geometric group equivariantly. Their
infinite compact classes then force failure of local geometry at $3$.
A second proof lifts prescribed boundary cocycles using
$\operatorname{cd}_3\Gamma=1$ and detects the local obstruction by
coordinate valuations, even when compact abelianization is zero.

The stronger base-five and base-seven compact assertion uses the same
arithmetic sources without abelianizing the cocycle. Proper good reduction
kills its whole inertia outside the two bad primes. A single local
conjugator then identifies the global compact section. The remaining
affine question concerns fibres of the proper pro-$3$ point map over
finite fields; general singleton-prime point injectivity cannot be used.

For the Frey images, the order-nine modular argument first gives
density-one good fibres. A rank-one Hodge--Tate character is algebraic
and has a fixed Hecke weight. Potentially Tate reduction at two
contradicts weight one if the three-adic connected monodromy is a torus.
Algebraic companions of rank-one characters then transfer the exclusion
to the two-adic representation. Lie algebras give individual openness,
and the distinct residue characteristics give joint openness.

\begingroup
\small
\setcounter{tocdepth}{2}
\tableofcontents
\endgroup
\clearpage

\section{A descent-to-cusps transfer theorem}

The following published inputs are used, with their original scopes.
First, \citet[Theorem 8.6]{Stoll2007} proves
\begin{equation}
X(\A_k)_\bullet^{\mathrm{f\text{-}ab}}=X(k)
\label{eq:stoll}
\end{equation}
if there is a nonconstant $k$-morphism $X\to A$ to an abelian variety
such that $A(k)$ is finite and $\Sha(k,A)_{\mathrm{div}}=0$.
His erratum repairs the real-place step in the proof through
Proposition 8.5 by using the density-one form of Theorem 8.2.
Second, \citet[Proposition 3.1 and Theorem B]{Saidi2022} proves injectivity
of localization on $H^1(k,\Delta_m)$ for every $m\geq1$, and on
$H^1(k,\Delta^{\sol})$, at a set of finite places of density one.
The action is obtained from any fixed full section; here rational
tangential sections supply such a section.

The full-group local-to-global cuspidality theorem of
\citet[Theorem 2.2, pp.~3--4]{Porowski2025} also fills one rational
cusp and passes through a cuspidally central extension. That theorem
starts from local cuspidality on a density-one set. Here finite
abelian descent first identifies the common center, and the affine
abelian exact sequence is used inside each actual derived quotient.
Porowski's theorem does not itself assert the single-prime quotient
conclusion considered below.

\begin{lemma}[Filling one rational cusp]\label{lem:boundary}
Suppose $g(X)\geq1$, $\#D(\bar k)\geq2$, and $Q\in D(k)$.
The peripheral map
\[
H^1(k,\wh\Z(1))\longrightarrow H^1(k,\Delta_1)
\]
at $Q$ is injective, as are its analogues at every finite completion.
Moreover, if every finite localization of a class in $H^1(k,\Delta_1)$
belongs to the peripheral image at $Q$, the global class belongs to
that image.
\end{lemma}

\begin{proof}
Put $W=U\cup\{Q\}$, which remains affine. The homology sequence obtained
by filling the puncture $Q$ is
\begin{equation}
0\longrightarrow I_Q\longrightarrow\Delta_1\longrightarrow
\pi_1(W_{\bar k})^{\ab}\longrightarrow0,
\qquad I_Q\simeq\wh\Z(1).
\label{eq:boundary}
\end{equation}
To verify injectivity on the left, the boundary part of $\Delta_1$ is
$(\bigoplus_{d\in D(\bar k)}\wh\Z(1)e_d)/\wh\Z(1)\sum_d e_d$.
The vector $e_Q$ remains primitive because at least one other geometric
puncture exists. Quotienting by it gives precisely the boundary part
for $W$, and compact homology is unchanged. This proves the exact
sequence over $k$ without requiring a second rational puncture.

The right-hand term is the full Tate module of the generalized
Jacobian of $W$. Its invariants over $k$, and over every finite
completion, vanish: rational torsion is finite for the toric part and
for the abelian quotient, hence also for their extension. Equivalently,
one may apply \citet[Lemma 2.3]{Saidi2022}. The cohomology sequence now
proves the asserted global and local injections. Finally a class
locally peripheral at $Q$ maps to zero in $H^1(k_v,\pi_1(W_{\bar k})^{\ab})$
at every finite place. This image is globally zero by
\citet[Proposition 3.2]{Saidi2022}, applied to $W$ with its full section
at the ordinary point $Q$. Exactness gives a unique global class in
$H^1(k,I_Q)$, as required.
\end{proof}

\begin{proof}[Proof of Theorem~\ref{thm:transfer}]
Let $s$ be a section of one of the stated quotients satisfying the
local condition. Choose its local point or cusp realizations, and let
$x_v\in X(k_v)$ be their compact centers. The induced section of the
geometrically abelian fundamental sequence of $X$ shows that the
modified adelic point $x=(x_v)_v$ survives every finite abelian torsor
on $X$. One may see this directly by restricting a finite torsor to
the section and twisting by the resulting global cocycle: the chosen
local points then lift everywhere. Equivalently this is the abelian
case of the section--descent comparison used in
\citet[\S2.1, Theorem 2.7]{BKL2026}.
Equation~\eqref{eq:stoll} now gives one $P\in X(k)$ with $x_v=P$ at
every finite place.

If $P\in U(k)$, every finite localization is the point section at
$P$. Apply Sa\"idi's localization theorem with the action defined by
the actual full point section $s_P$. The difference class is locally
trivial, so the theorem gives $s=s_P$.
Suppose instead that $P=Q\in D(k)$. The chosen local realization at
every finite place must be a cusp section at this same $Q$; an ordinary
point has its own distinct compact center. Choose a rational tangent
at $Q$ and let $s_Q$ be its full section. The difference between
$s_{\ab}$ and $(s_Q)_{\ab}$ is locally peripheral at $Q$ everywhere.

Lemma~\ref{lem:boundary} produces a unique
$b\in H^1(k,\wh\Z(1))$ such that $s_{\ab}$ is the abelianization of
the global cusp section $s_{Q,b}$. At every finite place the original
local section $s_v$ is a cusp section at $Q$. Injectivity of its
peripheral abelian map, again Lemma~\ref{lem:boundary}, says its local
parameter is exactly $b_v$. Hence $s_v$ and $(s_{Q,b})_v$ are conjugate
in the stated quotient at every finite place. Apply localization with
the action defined by the constructed full section $s_{Q,b}$.
The difference class is locally trivial, so Sa\"idi's Proposition~3.1,
or Theorem~B in the prosolvable case, gives global conjugacy.

Conversely every point section and every global cusp section satisfies
the local condition. Distinct centers are separated after compact
abelianization by local full Tate-module Kummer injectivity, and
parameters at one center are separated by Lemma~\ref{lem:boundary}.
This establishes both the equality and the disjointness in
\eqref{eq:classification}.
\end{proof}

\begin{remark}\label{rem:transfer-scope}
The theorem does not use an unrestricted local section conjecture.
Nor does it assert that the corresponding statement follows for an
arbitrary quotient or from an arbitrary abelian section criterion.
The proof uses the full geometric abelianization and a specific
localization theorem for the actual derived quotients. These points
are essential when comparing with a single-prime quotient.
\end{remark}

\subsection{Descent from completely split primes}
\label{subsec:split-prime-descent}

\begin{proof}[Proof of Theorem~\ref{thm:split-prime-descent}]
Take $x=(x_v)\in X(\A_\Q)_\bullet^{\mathrm{f\text{-}ab}}$.
Functoriality gives $f(x)\in A(\A_\Q)_\bullet^{\mathrm{f\text{-}ab}}$.
For every rational prime $\ell$, its componentwise $\ell$-adic
Kummer image comes from $\Sel(T_\ell A)$: pull back the
multiplication-by-$\ell^r$ torsors and take a compatible inverse
limit of the finite Selmer sets. The exact sequence is
\[
0\longrightarrow A(\Q)^{\wedge_\ell}
\longrightarrow\Sel(T_\ell A)
\longrightarrow T_\ell\Sha(\Q,A)\longrightarrow0.
\]
At each completely split prime satisfying the hypothesis, this
Selmer class is induced by some $a_\ell\in A(\Q)[\ell^\infty]$.
The difference $f(x_\ell)-a_\ell$ has zero $\ell$-adic Kummer
class in $A(\Q_\ell)$. A compact abelian $\ell$-adic Lie group
has a finite-index open subgroup isomorphic to $\Z_\ell^d$;
therefore the kernel of its pro-$\ell$ completion is precisely
its finite prime-to-$\ell$ torsion. Consequently $f(x_\ell)$
is torsion. This assertion does not require good reduction.

By Manin--Mumford, in the form of
\citet[Theorem~1.1$^*$, p.~151]{PilaZannier2008}, the hypothesis
on the image curve implies that
\[
f(X)(\bar\Q)\cap A(\bar\Q)_{\rm tors}
\]
is finite and Galois stable. Its inverse image is a finite reduced
closed subscheme $Z\subset X$ defined over $\Q$, since $f$ is
nonconstant. Thus $x_\ell\in Z(\Q_\ell)$ at all the indicated
split primes.

Finite abelian descent survives base change by
\citet[Proposition~5.16, p.~370]{Stoll2007}; the abelian assertion
is explicitly retained in \citet[\S2(4)]{StollErrata2017}. Hence
\[
x_L\in X(\A_L)_\bullet^{\mathrm{f\text{-}ab}}.
\]
At every place $v\mid\ell$ above one of the indicated completely
split primes, $(x_L)_v\in Z(L_v)$. These places have Dirichlet
density one among the finite places of $L$. Indeed, outside
ramification, a prime ideal has residue degree one exactly when
its rational prime splits completely, since $L/\Q$ is Galois.
The sum of $\operatorname{Norm}(v)^{-s}$ over residue degrees at
least two stays bounded as $s\to1^+$, whereas the sum over all
prime ideals diverges logarithmically. Finitely many exceptions
do not change the density.

A nonconstant map from $X$ to an abelian variety implies
$g(X)\ge1$. Therefore
\citet[Theorem~8.2, pp.~383--384]{Stoll2007}, in its
density-one finite-support form, gives $x_L=P$ for some
$P\in Z(L)$. The base-changed adelic point is Galois invariant.
The diagonal map $X(L)\to X(\A_L)_\bullet$ is injective, as
one sees in any finite completion, so $P\in Z(\Q)$.

For each finite rational prime $p$, choose $v\mid p$.
Equality over $L_v$ and the injectivity
$X(\Q_p)\to X(L_v)$ give $x_p=P$. Base change to a field
with complex places can discard the real component. Apply
Stoll's Theorem~8.2 a second time over $\Q$, to the singleton
$\{P\}$ and all finite places. This restores equality in the
entire modified adelic space. The validity of using finite
places is also explicit in \citet[\S3]{StollErrata2017}.
The converse follows because rational points survive descent.
\end{proof}

\begin{remark}[Geometric and arithmetic hypotheses]
\label{rem:split-prime-scope}
The geometric hypothesis holds if $\dim A\ge2$ and the
homomorphism $\Jac(X)\to A$ induced by $f$ is surjective:
differences of an elliptic translate generate only that
elliptic subvariety. Genus at least two of $X$ alone does not
exclude an elliptic image. Over any number field, the same
initial argument followed directly by Stoll's Theorem~8.2
works if the places whose residue characteristic $\ell$
satisfies $T_\ell\Sha(k,A)=0$ have density one. The completely
split criterion avoids imposing that density-one condition
over $\Q$ itself. Neither formulation proves the required
arithmetic vanishing for a new abelian variety.
\end{remark}


\section{Cyclotomic constraints and supported local conditions}
\label{sec:irreg-setup}

Finite torsors on the proper Fermat curve allow the classical first-case
arguments to use the local representatives of quotient sections.
Kummer's unit argument and its cyclotomic setting are presented in
\citet[Theorem~1.1 and Proposition~1.5, pp.~1--6]{Washington1982}.
Herbrand's class-group implication is
\citet[\S4, pp.~423--431]{Herbrand1932}; we use the direct character
convention recalled by
\citet[Theorem~1.1 and the following paragraph, p.~151]{Ribet1976}.
The constructions below establish the passage from those arithmetic
inputs to sections without assuming rational Fermat coordinates.

Throughout, \(p\geq5\) is prime. Let
\[
 F_p:\ X^p+Y^p=Z^p,\qquad
 D_p=F_p\cap\{XYZ=0\},\qquad U_p=F_p\setminus D_p.
\]
Here \(F_p\) is the minus-sign model of \(X_p\), identified with
the all-plus model by changing the sign of \(Z\); the notation
\(U_p\) is used through that identification. Finite fields are
written \(\F_p\). Its three rational coordinate cusps are
\[
 Q_0=(0:1:1),\qquad Q_1=(1:0:1),\qquad Q_\infty=(1:-1:0).
\]
Write \(J_p=\Jac(F_p)\) and \(j_p(R)=[R-Q_0]\).

Let \(\Delta_p=\pi_1(F_{p,\overline{\Q}})\). Define
\[
 \Pi_p^{\ab,p}
 =
 \pi_1(F_p/\Q)/
 \overline{[\Delta_p,\Delta_p]\Delta_p^p}.
\]
This quotient retains the full arithmetic quotient \(G_\Q\).
We consider a geometric quotient
\[
 \pi_1(U_p/\Q)\twoheadrightarrow\Pi\twoheadrightarrow G_\Q
\]
equipped with a compatible map \(\Pi\to\Pi_p^{\ab,p}\), and a continuous
section \(s:G_\Q\to\Pi\). At each finite place \(v\), suppose that its
restriction is, up to geometric conjugacy, induced by a full local
point section or by a full section of the decomposition group of a
\(\Q_v\)-rational point of \(D_p\). Such a representative has a proper
centre \(R_v\in F_p(\Q_v)\). The arguments below require only that the
image in \(\Pi_p^{\ab,p}\) be represented by the full proper point
section of \(R_v\).

All assertions concern every permitted choice of the \(R_v\). In
particular, uniqueness of a local point from a finite abelian
cohomology class is not assumed. A centre at \(p\) is called
\emph{first case} if its primitive integral homogeneous coordinates
are all units in \(\Z_p\).

Set
\[
 K=\Q(\zeta_p),\qquad L=\Q_p(\zeta_p),\qquad
 \Delta=\Gal(K/\Q),\qquad
 \sigma_i(\zeta_p)=\zeta_p^i,\quad
 \omega(\sigma_i)=i\in\F_p^\times.
\]
For an \(\F_p[\Delta]\)-module, a superscript \((\omega^m)\)
always denotes the direct eigenspace on which \(\sigma_i\) acts
as multiplication by \(i^m\).

\subsection{Proper Kummer torsors and virtual units}
\label{sec:irreg-torsors}

Define the virtual-unit space
\[
 V=\left\{[\alpha]\in K^\times/K^{\times p}:
             v(\alpha)\equiv0\pmod p
             \text{ for every finite place }v\right\}.
\]
Let \(E_K=\mathcal O_K^\times\), and let a bar denote complex conjugation.

\begin{lemma}\label{lem:irreg-virtual}
There is an exact sequence of \(\F_p[\Delta]\)-modules
\begin{equation}
 0\longrightarrow E_K/E_K^p
   \longrightarrow V
   \longrightarrow\Cl(K)[p]\longrightarrow0.
 \label{eq:irreg-virtual}
\end{equation}
Consequently, if \(r^-=\dim_{\F_p}\Cl(K)[p]^-\), then
\begin{equation}
 \dim_{\F_p}V^-=1+r^-.
 \label{eq:irreg-minus-dim}
\end{equation}
For odd \(m\) with \(3\leq m\leq p-2\),
\begin{equation}
 V^{(\omega^m)}\simeq\Cl(K)[p]^{(\omega^m)}.
 \label{eq:irreg-odd-components}
\end{equation}
\end{lemma}

\begin{proof}
If \((\alpha)=\mathfrak a^p\), send \([\alpha]\) to
\([\mathfrak a]\). Changing \(\alpha\) by a \(p\)-th power changes
\(\mathfrak a\) by a principal ideal. Every class killed by \(p\)
is obtained in this way, and the kernel consists of units times
\(p\)-th powers. This proves \eqref{eq:irreg-virtual}, including its
equivariance for the natural actions on elements and ideals.

For \(\epsilon\in E_K\), the unit
\(\epsilon/\overline\epsilon\) has absolute value one under every
complex embedding and is therefore a root of unity. The roots of
unity in \(K\) are \(\mu_{2p}\). Since \(2\) is invertible modulo \(p\),
the odd part of \(E_K/E_K^p\) is the line generated by \(\zeta_p\);
this line has character \(\omega\). It is nonzero, since a \(p\)-th
root of \(\zeta_p\) in \(K\) would be a root of unity of order divisible
by \(p^2\). Finally, \(\#\Delta=p-1\) is prime to \(p\), so taking
eigenspaces is exact. This gives both conclusions.
\end{proof}

Use signed coordinate functions \(A,B,C\), obtained by permuting
\(X,Y,-Z\), so that
\[
 A^p+B^p+C^p=0.
\]
For \(i\in\Z/p\Z\), let \(P_i\) be the point where \(C=0\) and
\(A+\zeta_p^iB=0\), and put
\[
 g_i=\frac{A+\zeta_p^iB}{A+B}.
\]

\begin{lemma}\label{lem:irreg-proper-torsors}
The section supplies classes
\[
 \xi_i\in H^1(K,\mu_p)=K^\times/K^{\times p}
\]
satisfying \(\xi_i=\sigma_i\xi_1\) for \(i\neq0\).
If its chosen \(p\)-local centre is first case, then every \(\xi_i\)
belongs to \(V\). If that centre has signed coordinates \(a,b,c\) and
\(e=b/(a+b)\), then
\begin{equation}
 \loc_{\mathfrak p}(\xi_i)
  =[1-e+e\zeta_p^i]\in L^\times/L^{\times p}.
 \label{eq:irreg-local-values}
\end{equation}
Here \(\mathfrak p\) is the unique prime of \(K\) above \(p\).
\end{lemma}

\begin{proof}
The tangent line \(A+\zeta_p^iB=0\) meets the Fermat curve in the
divisor \(pP_i\). Thus
\begin{equation}
 \divisor(g_i)=p(P_i-P_0).
 \label{eq:irreg-divisor}
\end{equation}
The equation \(W^p=g_i\) consequently extends to an unramified
\(\mu_p\)-torsor on the proper curve: locally, a zero or pole of
order \(pm\) is removed by rescaling \(W\) by the \(m\)-th power of a
local parameter. Its geometric monodromy is abelian of exponent \(p\),
so the torsor factors through the stipulated proper quotient. Pulling
it back along the restricted section over \(K\) gives \(\xi_i\).

For a finite place of \(K\) away from \(p\), the standard Fermat model
is smooth and proper, and the sections \(P_i\) are disjoint. The
divisor identity \eqref{eq:irreg-divisor} has no vertical component:
neither tangent line contains a fibre. The same local construction,
now with \(p\) invertible, extends the torsor to a finite etale torsor
over the integral model. Proper local points extend integrally, so
their evaluations are unramified. At a zero or pole this means
evaluation of the extended torsor, not of an undefined rational
function. A cusp section projects to the point section of its
proper centre, so it has the same property. Hence the valuation of
\(\xi_i\) is divisible by \(p\) away from \(\mathfrak p\).

At a first-case centre, \(a+b\equiv-c\not\equiv0\pmod p\).
Evaluation gives \eqref{eq:irreg-local-values}, and every displayed
value is a unit because \(\zeta_p\equiv1\pmod{\mathfrak p}\). Thus
\(\xi_i\in V\).

Finally, conjugating \(g_1\) by \(\sigma_i\) gives \(g_i\). The Kummer
isomorphism is natural when \(\mu_p\) carries its natural Galois
action. Pullback along the rational arithmetic section therefore
gives \(\xi_i=\sigma_i\xi_1\). Geometric conjugacy changes the
corresponding abelian cocycle by a coboundary only. There is no
identification with a constant coefficient module that could shift
the character.
\end{proof}

The same construction is available for each signed coordinate pair
on the original curve. No action of coordinate permutations on the
kernel of an arbitrary geometric quotient is required.

\subsection{A relative-class-rank criterion}
\label{sec:irreg-rank}

We give the local derivative argument in the precise form needed
here. Its classical antecedent is Eichler's method, as presented in
\cite[Section 3.2, pp.~107--108]{Gras2010}. The explicit sufficient
inequality below is used without an optimality assertion.

\begin{lemma}\label{lem:irreg-derivative}
Let \(F(T)\) be a rational function over \(\Z_p\), with a denominator
that is a unit at \(T=1\). Suppose \(F(1)=1\) and
\(F(\zeta_p)=b^p\) for \(b\in\mathcal O_L^\times\). Then
\begin{equation}
 \frac{F'(T)}{F(T)}=0
 \quad\text{in}\quad
 \F_p[T,\text{denominators}^{-1}]/(T-1)^p.
 \label{eq:irreg-derivative}
\end{equation}
\end{lemma}

\begin{proof}
Write \(b=B(\zeta_p)\), where \(B\in\Z_p[T]\) and
\(\deg B\leq p-2\). Clearing denominators and dividing by the monic
polynomial \(\Phi_p\) gives
\[
 F(T)-B(T)^p=\Phi_p(T)H(T)
\]
in a localization with denominator a unit at \(1\).
The residue of \(b\) is one, so \(B(1)\equiv1\pmod p\). Evaluation at
\(1\) gives
\[
 pH(1)=1-B(1)^p\equiv0\pmod{p^2}.
\]
Thus \(\overline H\) has a factor \(T-1\). Since
\(\Phi_p(T)\equiv(T-1)^{p-1}\pmod p\), we have
\(F-B^p=0\) modulo \((p,(T-1)^p)\). Differentiation preserves
\((T-1)^p\) in characteristic \(p\) and kills \(B^p\). The image of
\(F\) is a unit, giving \eqref{eq:irreg-derivative}.
\end{proof}

\begin{lemma}\label{lem:irreg-k-independent}
Suppose \(e,1-e,1-2e\in\Z_p^\times\), and let \(k\geq1\) satisfy
\begin{equation}
 k(k+1)-2<p.
 \label{eq:irreg-k-bound}
\end{equation}
Then the classes
\[
 d_i=
 \left[\frac{1-e+e\zeta_p^i}{1-e+e\zeta_p^{-i}}\right],
 \qquad 1\leq i\leq k,
\]
are linearly independent in \(L^\times/L^{\times p}\).
\end{lemma}

\begin{proof}
The bound implies \(k<p\). Suppose
\(\prod_i d_i^{\lambda_i}=1\), with
\(\lambda_i\in\{0,\ldots,p-1\}\) not all zero. The product of the
displayed unit representatives is therefore \(b^p\) for a unit
\(b\in\mathcal O_L^\times\). Apply
Lemma~\ref{lem:irreg-derivative} to
\[
 G(T)=\prod_{i=1}^k
 \left(\frac{1-e+eT^i}{1-e+eT^{-i}}\right)^{\lambda_i}.
\]
Its logarithmic derivative at \(1\) is
\(2e\sum_i i\lambda_i\), so \(\sum_i i\lambda_i\equiv0\pmod p\).
In particular \(\zeta_p^{\sum_i i\lambda_i}=1\) exactly. The identity
\[
 \frac{1-e+eT^i}{1-e+eT^{-i}}
   =T^i\frac{1-e+eT^i}{e+(1-e)T^i}
\]
then permits a second application to
\[
 F(T)=\prod_i
 \left(\frac{1-e+eT^i}{e+(1-e)T^i}\right)^{\lambda_i}.
\]
Writing
\(D_i(T)=(1-e+eT^i)(e+(1-e)T^i)\), we obtain
\[
 (2e-1)\sum_i\frac{\lambda_i iT^{i-1}}{D_i(T)}
       =0\pmod{p,(T-1)^p}.
\]
Let \(I\) be the nonempty support of the \(\lambda_i\).
All \(D_i(1)=1\), and \(2e-1\) is a unit, so
\[
 P(T)=
 \sum_{i\in I}\lambda_i iT^{i-1}
       \prod_{\substack{j\in I\\j\neq i}}D_j(T)
\]
is divisible by \((T-1)^p\) over \(\F_p\). On the other hand,
\[
 \deg P\leq2\sum_{i\in I}i-\min(I)-1
          \leq k(k+1)-2<p.
\]
Hence \(P=0\). For \(i_0=\min(I)\), its coefficient of \(T^{i_0-1}\)
is
\[
 \lambda_{i_0}i_0[e(1-e)]^{|I|-1}\neq0,
\]
and no other summand contributes at that degree. This contradiction
proves the lemma.
\end{proof}

\begin{theorem}\label{thm:irreg-rank}
Under the hypotheses of Section~\ref{sec:irreg-setup}, if
\begin{equation}
 (r^-+2)(r^-+3)-2<p,
 \qquad r^-=\dim_{\F_p}\Cl(K)[p]^-,
 \label{eq:irreg-rank-bound}
\end{equation}
then no permitted \(p\)-local centre is first case.
\end{theorem}

\begin{proof}
Assume otherwise and choose signed primitive unit coordinates
\(a,b,c\). Their nonzero residues satisfy \(a+b+c=0\), and, since
\(p\geq5\), are not all equal. Choose an unequal pair \(a,b\).
Then \(e=b/(a+b)\) satisfies the unit hypotheses of
Lemma~\ref{lem:irreg-k-independent}.

Lemma~\ref{lem:irreg-proper-torsors} gives
\(\xi_i\in V\), with \(\overline{\xi_i}=\xi_{-i}\). Thus
\(\xi_i/\xi_{-i}\in V^-\), and their localizations are the classes
\(d_i\). Take \(k=r^-+2\). By \eqref{eq:irreg-rank-bound} these \(k\)
local classes are independent; the global classes are therefore
independent as well. This contradicts
\(\dim_{\F_p}V^-=r^-+1\).
\end{proof}

\begin{corollary}\label{cor:irreg-residue}
Under \eqref{eq:irreg-rank-bound}, every permitted \(p\)-local
centre reduces to one of \(Q_0,Q_1,Q_\infty\).
\end{corollary}

\begin{proof}
For primitive integral coordinates the reduction of the equation is
\(X+Y=Z\). If they are not all units, exactly one is zero in the
residue field, giving one of the three displayed points. Moreover
\(\mu_p(\Q_p)=\{1\}\), so the only \(\Q_p\)-rational coordinate
boundary points themselves are the three rational cusps.
\end{proof}

\subsection{Odd-character Kummer congruences}
\label{sec:irreg-character}

In \(\F_p[E,W]/(W^{p-1})\), define
\begin{equation}
 C_j(E)=[W^j]\log(1-E+E\exp W),
 \qquad 1\leq j\leq p-2,
 \label{eq:irreg-C}
\end{equation}
where both formal series are truncated through degree \(p-2\).
All denominators are invertible in \(\F_p\).

\begin{theorem}\label{thm:irreg-characters}
Suppose that a section as in Section~\ref{sec:irreg-setup} has a
first-case \(p\)-local centre. For each ordered pair of its signed
coordinates, put \(e=b/(a+b)\bmod p\). For every odd integer
\(3\leq m\leq p-2\),
\begin{equation}
 \Cl(K)[p]^{(\omega^m)}=0
       \quad\Longrightarrow\quad C_m(e)=0.
 \label{eq:irreg-character-zero}
\end{equation}
In particular, without a class-component hypothesis,
\begin{equation}
 B_{p-m}\,C_m(e)=0\quad\text{in }\F_p.
 \label{eq:irreg-Bernoulli-congruence}
\end{equation}
\end{theorem}

\begin{proof}
By Lemmas~\ref{lem:irreg-virtual} and
\ref{lem:irreg-proper-torsors}, the vanishing class component implies
\(V^{(\omega^m)}=0\), and the section supplies
\(\xi_i=\sigma_i\xi_1\in V\). The direct \(\omega^m\)-projector is
\[
 e_m=\frac1{p-1}\sum_{i=1}^{p-1}i^{-m}\sigma_i.
\]
Choose integers \(c_i\) lifting its coefficients. The equality
\(e_m\xi_1=0\) means
\[
 \prod_{i=1}^{p-1}\xi_i^{c_i}=1
       \quad\text{in }K^\times/K^{\times p}.
\]
After localization at \(\mathfrak p\), it gives
\begin{equation}
 \prod_{i=1}^{p-1}(1-e+e\zeta_p^i)^{c_i}=\gamma^p
 \label{eq:irreg-projector-product}
\end{equation}
for a unit \(\gamma\in\mathcal O_L^\times\) with residue one.
Here and until reduction, \(e\) denotes its original value in
\(\Z_p\).

Set \(\lambda=\zeta_p-1\). Since \(v_\lambda(p)=p-1\) and
\(m+1\leq p-1\), there is an isomorphism
\begin{equation}
 \mathcal O_L/(\lambda^{m+1})
       \simeq R_m:=\F_p[t]/(t^{m+1}),\qquad
 \lambda\longmapsto t.
 \label{eq:irreg-trunc-ring}
\end{equation}
Indeed, \(\mathcal O_L=\Z_p[\lambda]\), and the Eisenstein equation
\(\Phi_p(1+\lambda)=0\) imposes no further relation among
\(1,\lambda,\ldots,\lambda^m\) in this quotient.
Writing \(\gamma=1+\lambda u\), the binomial formula shows that
\(\gamma^p\equiv1\pmod{\lambda^p}\). Thus the right side of
\eqref{eq:irreg-projector-product} equals \(1\) in \(R_m\).

Put
\[
 w=\log(1+t)=\sum_{r=1}^m\frac{(-1)^{r+1}t^r}{r}\in R_m.
\]
Truncated logarithm and exponential are inverse isomorphisms
between the multiplicative group \(1+(t)\) and the additive group
\((t)\). This follows by reducing their usual formal identities:
every denominator occurring through degree \(m\) has prime factors
less than \(p\). In particular \((1+t)^i=\exp(iw)\), and
\[
 \log\bigl(1-e+e(1+t)^i\bigr)
       =\sum_{j=1}^m C_j(e)i^j w^j.
\]
For \(1\leq j\leq m\),
\[
 \sum_{i=1}^{p-1}c_i i^j
       =\frac1{p-1}\sum_{i\in\F_p^\times}i^{j-m}
       =
       \begin{cases}
          0,&j<m,\\
          1,&j=m.
       \end{cases}
\]
The nonzero exponents \(j-m\) have absolute value at most \(p-3\),
so none is a multiple of \(p-1\). The logarithm of
\eqref{eq:irreg-projector-product} therefore reduces to
\(C_m(e)w^m=0\). Since \(w^m=t^m\neq0\), this proves
\eqref{eq:irreg-character-zero}.

For the Bernoulli assertion, use Herbrand's implication
\begin{equation}
 \Cl(K)[p]^{(\omega^m)}\neq0
       \quad\Longrightarrow\quad
       p\mid\operatorname{num}(B_{p-m}).
 \label{eq:irreg-Herbrand}
\end{equation}
The direct character convention in \eqref{eq:irreg-Herbrand} is that
of \cite[Theorem~1.1 and following paragraph, p.~151]{Ribet1976}; see also \cite[Section~6.3, pp.~100--102, Theorem~6.17]{Washington1982};
the projector there is written using \(\sigma_i^{-1}\) and agrees
with ours after reindexing. See also
\cite[Theorem 3.7 and the following paragraph, pp.~104--105]{Gras2010}.
If \(B_{p-m}\) is nonzero modulo \(p\), the contrapositive of
\eqref{eq:irreg-Herbrand} and
\eqref{eq:irreg-character-zero} give \(C_m(e)=0\).
Otherwise \eqref{eq:irreg-Bernoulli-congruence} is automatic.

The Bernoulli number has an unambiguous reduction: for
\(1\leq n\leq p-3\), its recurrence
\[
 \sum_{j=0}^n\binom{n+1}{j}B_j=0,\qquad B_0=1,
\]
shows inductively that \(B_n\in\Z_{(p)}\), since \(n+1<p\).
\end{proof}

\begin{remark}\label{rem:irreg-endpoint}
The proof includes \(m=p-2\). In that case
\(\mathcal O_L/(\lambda^{m+1})\) has characteristic \(p\) and length
\(p-1\), and no denominator divisible by \(p\) occurs. The logarithm
used above is a polynomial on a nilpotent ideal. It is not the
\(p\)-adic logarithm of \(\zeta_p\), which is zero. The reflected
character \(\omega^{1-m}\) in the usual notation for Mirimanoff
polynomials is likewise not the direct character of \(V\) used
here. The classical rational-point congruences are treated in
\cite[Theorems 3.5 and 3.7]{Gras2010}; the proper-torsor construction
above establishes their applicability under the stated quotient
hypotheses.
\end{remark}

\begin{corollary}\label{cor:irreg-cubic}
If \(\Cl(K)[p]^{(\omega^3)}=0\), no section under consideration has a
first-case \(p\)-local centre. In particular it suffices that
\begin{equation}
 p\nmid\operatorname{num}(B_{p-3}).
 \label{eq:irreg-cubic-condition}
\end{equation}
\end{corollary}

\begin{proof}
Expansion gives
\[
 C_1(E)=E,\qquad C_2(E)=\frac{E(1-E)}2,\qquad
 C_3(E)=\frac{E(1-E)(1-2E)}6.
\]
At a first-case centre, \(e\) and \(1-e\) are nonzero. Thus
Theorem~\ref{thm:irreg-characters} forces \(e=1/2\), or \(a=b\), for
each signed coordinate pair. All three nonzero signed residues
would be equal and have sum zero, impossible for \(p\geq5\).
Herbrand gives the stated sufficient Bernoulli condition.
\end{proof}

\begin{corollary}\label{cor:irreg-multiples}
Let \(n\) be an odd multiple of \(p\). The rank criterion, the
odd-character congruences, and the cubic corollary apply to sections
on \(U_n\) whose geometric quotient retains the pullbacks of the
proper mod-\(p\) abelian torsors on \(F_p\), with full \(G_\Q\) and
full local point-or-cusp representatives. In particular they apply
to its geometrically maximal prosolvable arithmetic fundamental
group.
\end{corollary}

\begin{proof}
The coordinate-power map
\[
 F_n\longrightarrow F_p,\qquad
 (X:Y:Z)\longmapsto
 (X^{n/p}:Y^{n/p}:Z^{n/p})
\]
takes boundary to boundary and first-case \(p\)-local points to
first-case points. Pullback retains the finite proper torsors used
in the proofs, and the full local representatives map to full
local representatives. Their geometric monodromy is abelian, so
these pullbacks factor through the maximal prosolvable quotient.
\end{proof}


The preceding rank criterion and character congruences, with the
cubic consequence, prove
Theorem~\ref{thm:first-case-obstructions-intro}.

\subsection{Local Brauer support at seventeen, after McCallum}
\label{subsec:seventeen-local-brauer}

This subsection proves a single-primary local condition using
McCallum's cusp-disc argument, then proves full nonabelian localization
over a regular-prime arithmetic source.  Write $F_p:X^p+Y^p=Z^p$; for odd $p$ this is isomorphic over
$\mathbb Q$ to $X_p$ by changing the sign of $Z$.  Let $\mathcal F_p$
be its proper plane model over $\mathbb Z$. The rational boundary
$D_p(\mathbb Q)$ consists of its three rational coordinate cusps.

For a proper curve $C/\mathbb Q$ with proper model $\mathcal C$, put
\[
 \operatorname{Br}_{\{p\}}(\mathcal C)
 =\{a\in\operatorname{Br}(C):
 a_v\text{ extends over }\mathcal C_{\mathbb Z_v}
 \text{ for every finite }v\ne p,\quad a_\infty=0\}.
\]
Its local zero set consists of the $x\in C(\mathbb Q_p)$ satisfying
$\operatorname{inv}_p a(x)=0$ for every such $a$.  Evaluation outside
$p$ vanishes by properness and $\operatorname{Br}(\mathbb Z_v)=0$.
Thus this zero set contains the projection of the Brauer--Manin set,
and every rational point belongs to it
\citep[\S1, pp.~393--394]{McCallum2001Brauer}.

\begin{proposition}[The regular-prime cusp residue classes]
\label{prop:regular-primary-brauer-cusp-discs}
Let $p>5$ be regular, meaning that the class number of
$\mathbb Q(\zeta_p)$ is prime to $p$.  In each residue class of a
rational coordinate cusp, the local zero set
$F_p(\mathbb Q_p)^{\operatorname{Br}_{\{p\}}(\mathcal F_p)[p^\infty]}$
contains just that cusp.  In particular,
\[
 F_{17}(\mathbb Q_{17})^{\operatorname{Br}_{\{17\}}(\mathcal F_{17})[17^\infty]}
       =D_{17}(\mathbb Q).
\]
\end{proposition}
\begin{proof}
We first explain the extension of McCallum's regular-prime argument
from rational points to this local zero set.  Put
$K=\mathbb Q(\zeta_p)$, let $\lambda$ be its unique prime over $p$,
and set $\pi=1-\zeta_p$ and $\Delta=\operatorname{Gal}(K/\mathbb Q)$.
Choose a wild cyclic quotient $C$ of $F_p$, with rational branch
cusps $O_0,O_1,O_\infty$, and write $A=\operatorname{Jac}(C)$.
Here wild means that, for branch exponents $a+b+c=0$,
$\gamma=a^ab^bc^c$ satisfies $\gamma^{p-1}\not\equiv1\pmod {p^2}$.
Such a quotient exists by
\citet[Lemma~9, pp.~580--581]{McCallum1994}.
Use $O_0$ as Abel--Jacobi base point and write $j:C\to A$.

The regularity hypothesis supplies classes
\[
 d_k\in H^1(K,A)[\pi^k]\qquad(k\ge1)
\]
which vanish locally away from $\lambda$, satisfy
$\pi^{k-1}d_k=d_1$, and whose first member pairs nontrivially with
$(A(K_\lambda)/\pi A(K_\lambda))^\Delta$.
These are exactly the classes of
\citet[Lemma~10 and Theorem~5, pp.~581--582]{McCallum1994}.
To specify where regularity is used, the construction begins with
the Kummer class of $\zeta_p$ and lifts it through
$H^1(\mathcal O_K[1/p],A[\pi^k])$.
The obstruction $H^2(\mathcal O_K[1/p],\mu_p)$ vanishes: its
Picard contribution vanishes by regularity, and its Brauer
contribution vanishes because the only allowed finite invariant
is at $\lambda$, all infinite places are complex, and the global
sum of invariants is zero.

Let $B$ be the common annihilator of these local classes in
$A(\mathbb Q_p)$:
\[
 B=\{P:\langle\operatorname{res}P,d_{k,\lambda}\rangle=0
                     \text{ for every }k\ge1\}.
\]
The classes $d_k$ are $p$-primary, because $p$ is a unit times
$\pi^{p-1}$. The canonical principal polarization of $A$, defined
 over $\Q$, identifies these with classes for the dual abelian variety.
The global classes $e_k=\operatorname{cor}_{K/\mathbb Q}d_k$
vanish at every place other than $p$.  Compatibility of the Tate
pairing with restriction and corestriction gives
\[
 \langle P,e_{k,p}\rangle_{\mathbb Q_p}
       =\langle\operatorname{res}P,d_{k,\lambda}\rangle_{K_\lambda}.
\]
There is no degree factor in this identity.  The relation between
the Tate pairing and Brauer evaluation, including the choice of
representatives with zero evaluations away from $p$, is
\citet[Lemma~1 and equation~(5), author pp.~4--5]{McCallum2001Brauer}.
Normalize the Brauer representative $a_k$ of $e_k$ by
$a_k(O_0)=0$. The rational base point splits the constants in
$\operatorname{Br}(C)$, and Hochschild--Serre identifies its quotient
by constants with $H^1(\Q,A)$. If $p^me_k=0$, then $p^ma_k$ is
constant and evaluates to zero at $O_0$, so $p^ma_k=0$. At every
$v\ne p$, the vanishing of $e_{k,v}$ makes $a_{k,v}$ constant;
normalization makes it the zero class. Thus these are genuinely
$p$-primary supported Brauer classes, not merely primary modulo constants.
It follows that
\begin{equation}\label{eq:regular-supported-local-annihilator}
 j\bigl(C(\mathbb Q_p)^{\operatorname{Br}_{\{p\}}(\mathcal C)[p^\infty]}\bigr)
       \subset B.
\end{equation}
The quotient map extends to proper models away from $p$, so Brauer
functoriality applies to it as in that paper's equation~(3).

Consider the natural map
\[
 \chi:A(\mathbb Q_p)/pA(\mathbb Q_p)
                \longrightarrow A(K_\lambda)/\pi A(K_\lambda).
\]
Its image is the one-dimensional $\Delta$-invariant subspace,
and it is nonzero
\citep[Corollary~2 and Proposition~5, pp.~574--575]{McCallum1994}.
We claim that the saturation $B^{\mathrm{sat}}$ lies in
$\ker\chi$.  If $rP\in B$ for a nonzero integer $r$, set
$t=v_\pi(r)$ and pair $rP$ with $d_{t+1,\lambda}$.
Writing $r=\pi^tu$ and moving the endomorphism across the pairing
reduces the result to a nonzero scalar multiple of the pairing
of $\chi(P)$ with $d_{1,\lambda}$.  Indeed the Rosati dual of
$\pi$ is a unit times $\pi$, and every unit multiplier reduces
modulo $\pi$ to a scalar in $\mathbb F_p$.
Since $P$ is fixed by $\Delta$, its class remains in the invariant
line.  The nontrivial pairing on that line forces $\chi(P)=0$.
This is the algebra in
\citet[Theorem~3, pp.~575--576]{McCallum1994}; equation~\eqref{eq:regular-supported-local-annihilator}
replaces the rational-point reciprocity condition in its proof.

Here $B$ is closed.  After quotienting torsion, $A(\mathbb Q_p)$
is a free $\mathbb Z_p$-module of rank $g=\dim A$, and the
saturation of a closed submodule is again closed.  Since
$B^{\mathrm{sat}}\bmod p\subset\ker\chi$, elementary lattice
linear algebra gives a saturated subgroup $M$ containing
$B^{\mathrm{sat}}$, with $M/A(\mathbb Q_p)_{\mathrm{tors}}$
of rank $g-1$ and image modulo $p$ equal to $\ker\chi$.
For example, extend a basis of the saturated submodule to a
basis of the free module and lift the nonzero functional
$\chi$ on the resulting quotient.

The logarithmic pairing with invariant differentials has torsion
kernel on the left.  Its annihilator of $M$ is a line; choose
a primitive element $\eta$ in the integral N\'eron-differential
lattice on this line and put $\omega=j^*\eta$.
The integral pairing and the identification of N\'eron and
dualizing differential lattices are
\citet[Propositions~2, 3 and~7, pp.~569--570, 578]{McCallum1994}.
Thus the nonzero reduction of $\omega$ belongs to the subspace
$V$ annihilating $\ker\chi$ in that paper.  By its
Theorem~4, p.~579, no nonzero member of $V$ vanishes at all
three branch cusps.

The Coleman integral $\lambda_\omega$, normalized at $O_0$,
vanishes on $j^{-1}(B)$ and hence on the local Brauer zero set.
It vanishes at all three rational branch cusps because their
Abel--Jacobi differences are torsion.  At a branch cusp where
the reduced differential is nonzero, the integral has at most
one zero in that residue disc by
\citet[Lemma~7, pp.~579--580]{McCallum1994}.
Consequently that cusp is the only point of the local Brauer
zero set in its disc.

The quotient $F_p\to C$ is injective on $\mathbb Q_p$-points:
the quotient coordinates determine the $p$th powers of the
Fermat coordinates, and $\mu_p(\mathbb Q_p)=\{1\}$.
The three cusp residue classes map to the three branch cusp
classes, as follows from the same explicit formulas used in
\citet[Lemma~8, p.~580]{McCallum1994}.
Therefore one coordinate-cusp disc on $F_p$ contains just its
cusp in the Brauer zero set.  Rational coordinate permutations,
with the signs appropriate to $X^p+Y^p=Z^p$, preserve the
model away from $p$ and permute the three cusps transitively.
They preserve the local Brauer zero set, proving the first
assertion for all three discs.

For $p=17$, regularity follows directly from Kummer's criterion,
in the full-class-number form of
\citet[Theorem~5.34, p.~78]{Washington1982}.  The relevant
Bernoulli numbers, in increasing even degree, are
\[
 B_2,B_4,B_6,B_8,B_{10},B_{12},B_{14}
 =\frac16,-\frac1{30},\frac1{42},-\frac1{30},
   \frac5{66},-\frac{691}{2730},\frac76.
\]
None has numerator divisible by seventeen.  The wild quotient can be taken to be
$y^{17}=x(1-x)$: here $\gamma=1/4$ and
$4^{16}\equiv154\not\equiv1\pmod {289}$.

It remains to exclude other residue classes.  A primitive
integral point with three unit coordinates would give
$x^{17}+y^{17}=1$ with $x+y\equiv1\pmod {17}$.
Seventeenth powers modulo $289$ depend only on residues
modulo $17$.  For residues $a=0,\ldots,16$, the errors
$a^{17}+((1-a)\bmod17)^{17}-1$ modulo $289$ are, in order,
\[
 \begin{gathered}
 0,0,153,68,102,119,170,34,34,\\
 68,34,34,170,119,102,68,153.
 \end{gathered}
\]
Only $a=0,1$ give zero.  Thus every $\mathbb Q_{17}$-point
reduces to a coordinate cusp.  The first assertion now gives
the stated equality.
\end{proof}

\begin{lemma}[Finite-primary reciprocity]\label{lem:finite-primary-reciprocity}
Let $p$ be an odd prime, let $A/k$ be an abelian variety over a number
field, and let $z\in H^1(k,T_pA)$ have local Kummer representatives
$P_v\in A(k_v)^{\wedge_p}$. For every
$d\in H^1(k,A^\vee)[p^\infty]$,
\[
 \sum_v\langle P_v,d_v\rangle_v=0.
\]
In particular a class supported at one place pairs to zero there.
\end{lemma}
\begin{proof}
Choose $m$ with $p^md=0$. The Kummer sequence supplies a global
lift $b\in H^1(k,A^\vee[p^m])$ of $d$. Reduction of $z$ gives
$z_m\in H^1(k,A[p^m])$. Their Weil cup product lies in
$H^2(k,\mu_{p^m})=\operatorname{Br}(k)[p^m]$.
Its local invariant is, up to one fixed sign convention,
$\langle P_v,d_v\rangle_v$ by the finite Kummer diagram of
\citet[Theorem~I.3.2, Corollary~I.3.4 and Remark~I.3.5,
pp.~41--44]{Milne2006ADT}. Global Brauer reciprocity makes
its finite invariant sum zero; equivalently use that source's
Theorem~I.4.10(b), pp.~56--57, with all places and
$M=A[p^m]$, whose Cartier dual is $A^\vee[p^m]$.
These statements include the residue-characteristic-$p$ place.
The real-place terms vanish because $p$ is odd.
Where $d_v=0$ the pairing is zero even if $b_v$ is a nonzero
Kummer class. This finite-level argument needs no inverse-limit
duality theorem and no finiteness assumption on $\Sha$.
\end{proof}

\begin{corollary}[The primary local cusp condition]
\label{cor:seventeen-primary-local}
Let $z\in H^1(\Q,T_{17}J_{17})$ be locally the Kummer class
of $[x_v-O]$, with $x_v\in X_{17}(\Q_v)$ and a fixed rational
coordinate cusp $O$. Every chosen $x_{17}$ is a rational
coordinate cusp. Consequently this holds for every quotient
Selmer section retaining the compact $T_{17}J_{17}$ class.
The same cusp-disc assertion holds for every regular $p>5$;
if every $\Q_p$-point reduces to the coordinate boundary, the
entire local representative is a rational cusp.
\end{corollary}
\begin{proof}
Push $z$ to the Jacobian of a wild cyclic quotient used in
Proposition~\ref{prop:regular-primary-brauer-cusp-discs}, and
restrict to $K=\Q(\zeta_p)$. Apply
Lemma~\ref{lem:finite-primary-reciprocity} to every $d_k$, using
the rational principal polarization. There is only one place
of $K$ above $p$ and the other local classes vanish. The chosen
local point therefore belongs to exactly the common annihilator
$B$ used in that proposition. Pulling back the classes and
permuting the coordinate cusps gives its cusp-disc conclusion.
At seventeen the displayed residue calculation exhausts all
local points. A quotient section supplies $z$ by filling the
boundary and abelianizing; a full local cusp witness fills to
its center, so every chosen witness has the required Kummer class.
No all-prime descent or global full-section lift is involved.
The older finite-descent consequence also follows, since
\citet[Corollary~7.3, p.~379]{Stoll2007} identifies proper finite
abelian descent with the Brauer--Manin condition.
\end{proof}

\subsection{Nonabelian localization at every odd regular prime}
\label{subsec:regular-nonabelian}

\begin{proposition}[The regular-prime arithmetic source]
\label{prop:regular-nonabelian-localization}
Let $p$ be odd with $p\nmid h(K)$, where $K=\Q(\zeta_p)$,
and let $n=p^r$, $r\ge1$. Put
$N_n=\pi_1(U_{n,\overline\Q})^{\mathrm{pro}\text{-}p}$ and
$C_n=\pi_1(X_{n,\overline\Q})^{\mathrm{pro}\text{-}p}$.
For a standard rational tangent the actual actions on these
and all their stable quotients factor through
$\Gamma=\Gal(K_{\{p\}}(p)/\Q)$, where $K_{\{p\}}(p)$ is the
maximal pro-$p$ extension unramified outside the sole prime above $p$.
Their $G_K$-images are pro-$p$, and $G_{\Q_p}\to\Gamma$ is surjective.
For every finitely generated pro-$p$ group $M$ with such an action,
\[
 H^1(G_{\Q,\{p\}},M)\longrightarrow H^1(\Q_p,M)
\]
is injective as a map of nonabelian cohomology sets.
Localization is consequently injective on proper point-local
Selmer sections for every stable quotient of $C_n$, and on
$H^1(G_{\Q,\{p\}},T_pJ_n)$ and its compact Selmer subgroup.
\end{proposition}
\begin{proof}
In the minus-sign model, put $t=x^n$, $y^n=1-t$. The geometric
pro-$p$ group $\Pi$ of $\mathbb P^1-\{0,1,\infty\}$ at its
unit tangent is free of rank two, with
$\Pi^{\mathrm{ab}}=\Z_p(1)^2$. The coordinate cover has group
$(\Z/n)^2$ and full geometric group
\[
 N_n=\ker(\Pi\longrightarrow\Pi^{\mathrm{ab}}/n\Pi^{\mathrm{ab}}).
\]
Indeed, the core of the kernel of any finite $p$-quotient of the
covering subgroup has $p$-group kernel and $p$-group deck quotient;
thus the induced topology is the full pro-$p$ topology. Riemann
existence supplies this covering comparison
\citep[Expos\'e~XII, Theorem~5.1, p.~251]{SGA1}.
The branch $x=t^{1/n}$, $y=(1-t)^{1/n}$ with leading coefficient
and constant term one gives compatible rational tangential splittings.
The upstairs action is the restriction of the actual action on $\Pi$.

On $\Pi/\Phi(\Pi)$ the action is cyclotomic, so $G_K$ acts
trivially. The automorphism kernel on this Frattini quotient is
pro-$p$: for a finite $p$-group, automorphisms trivial on the
Frattini quotient act freely on the $p$-power set of tuples lifting
a fixed generating basis. Pass to characteristic finite quotients.
Thus the whole $G_K$-image is pro-$p$, also on subgroups and stable
quotients. The based action is unramified outside $p$ by
\citet[Definitions~17.4--17.6, Theorem~17.7 and Corollary~17.8,
pp.~49--51]{Wojtkowiak2009}: the unit tangent extends nontrivially
on the smooth pair there. This is an actual action statement.
Filling the punctures gives the proper action at the rational center.

Let $P=\Gal(K_{\{p\}}(p)/K)$ and let $D$ be the image of the
sole local group $G_{K_\lambda}$. If $D$ were proper, a finite
$p$-quotient and a maximal subgroup would give a cyclic degree-$p$
extension of $K$, unramified outside $\lambda$ and split there.
It would be everywhere unramified, contradicting $p\nmid h(K)$
by class field theory. Hence $D=P$. Total ramification of $p$ in
$K/\Q$ makes $G_{\Q_p}$ surject onto $\Gal(K/\Q)$ as well, proving
surjectivity onto $\Gamma$. The pro-$p$ unramified action over
$K$ factors through $P$, and hence the full action factors through
$\Gamma$.

For any cocycle $c$ on $G_{\Q,\{p\}}$, its restriction to $G_K$
has graph $g\mapsto(c(g),\rho(g))$ in the pro-$p$ semidirect
product $M\rtimes\rho(G_K)$. It factors through $P$. Thus $c$
and its action factor through $\Gamma$. If two classes agree at
$p$, conjugate one cocycle globally by the same local conjugator.
The cocycles now agree on the surjective image of $G_{\Q_p}$,
so agree globally. This proves injection of the entire pointed set.
The identical graph argument applies to finite free $\Z_p$-modules
obtained by subquotients and tensor operations from $N_n^{\mathrm{ab}}$.

For the proper group, smooth proper specialization at $q\ne p$
identifies the pointed arithmetic sequence with the pullback of
the special-fibre sequence
\citep[Expos\'e~X, Theorem~3.8 and Corollary~3.9, p.~217]{SGA1}.
Each proper local point extends and gives a cocycle trivial on
inertia. Filling a local cusp gives its proper point section.
Since the actual coefficient inertia is trivial, conjugacy preserves
this property. Proper Selmer cocycles therefore descend to
$G_{\Q,\{p\}}$. Compare two of them using the same fixed rational
cusp as origin, and apply the proved injection. This also covers
all stable proper quotients. Abelianization is $T_pJ_n$
\citep[Expos\'e~IX, \S5.8, p.~194]{SGA1}, proving the integral
compact assertions, without a rank or $\Sha$ hypothesis.
\end{proof}

\begin{remark}[Why one place does not give the all-prime theorem]
\label{rem:regular-all-prime-limitation}
The primary local condition and nonabelian localization do not imply
$X_{17}(\mathbb A_{\mathbb Q})_\bullet^{f\text{-}\mathrm{ab}}
=X_{17}(\mathbb Q)$.  At positive Mordell--Weil rank, a
single finite place cannot detect the full profinite
Mordell--Weil completion.  Indeed $A(\mathbb Q_q)$ has an
open pro-$q$ subgroup with finite quotient $G$.  For every
prime $\ell\ne q$ not dividing $|G|$, its $\ell$-adic
completion is zero, whereas $A(\mathbb Q)^{\wedge}$ contains
$\mathbb Z_\ell^{\operatorname{rank}A(\mathbb Q)}$.
The same obstruction applies to any fixed finite collection
of local places.

There is also a curve-level distinction.  Given a rational
point $O$ on a smooth proper curve, fix a good prime $q\ne p$
and choose a different $Q_q$ sufficiently close to $O$.
The Abel--Jacobi difference lies in the pro-$q$ formal
subgroup, where multiplication by $p$ is an automorphism.
The adelic point equal to $O$ except for the component $Q_q$
therefore has trivial compact $p$-adic Kummer class at every
place and survives every $p$-primary finite abelian torsor.
It need not be diagonal.  This is a limitation of the
single-primary condition, not a counterexample to full
finite abelian descent.  Stoll's Theorem~8.2 requires finite
support at a density-one set of places; the one place
seventeen does not meet that hypothesis after a finite
number-field extension.
\end{remark}

\begin{proposition}[A sufficient criterion allowing positive rank]
\label{prop:positive-rank-split-prime-criterion}
Let $X/\mathbb Q$ be a smooth proper curve of genus at least
one, with rational base point $O$.  Let $f:X\to A$ be a
morphism to an abelian variety and $Z\subset X$ a finite
$\mathbb Q$-subscheme.  Suppose that there is a finite
Galois extension $L/\mathbb Q$ such that, for all but finitely
many rational primes $\ell$ split completely in $L$,
\[
 T_\ell\Sha(A/\mathbb Q)=0
\]
and
\[
 \{x\in X(\mathbb Q_\ell):
   \log_\ell(f(x)-f(O))\in
   \operatorname{span}_{\mathbb Q_\ell}
                         \log_\ell A(\mathbb Q)\}
                 \subset Z(\mathbb Q_\ell).
\]
Then
$X(\mathbb A_{\mathbb Q})_\bullet^{f\text{-}\mathrm{ab}}
=X(\mathbb Q)\subset Z(\mathbb Q)$.
\end{proposition}
\begin{proof}
The image in $A$ of a finite-descent adelic point has a
compatible compact $\ell$-adic Selmer class.  The inverse-limit
Kummer sequence and the vanishing of $T_\ell\Sha$ identify
that class with an element of $A(\mathbb Q)^{\wedge_\ell}$.
Its logarithm at $\ell$ lies in the stated rational
Mordell--Weil span.  The local hypothesis therefore places
$x_\ell$ in $Z(\mathbb Q_\ell)$ at every such split prime.

After base change to $L$, the degree-one primes have density
one; outside the ramified primes they lie above precisely
the rational primes splitting completely in $L$.
\citet[Theorem~8.2, pp.~383--384]{Stoll2007} makes the
modified adelic point diagonal in $Z(L)$.  Compatibility
with the original $\mathbb Q$-adeles makes this point
$\mathbb Q$-rational: at a split prime, all its conjugate
components equal the same original local point.
Applying the same theorem over $\mathbb Q$ to this rational
singleton recovers the real components as well. Rational points
survive finite abelian descent and $Z$ is a subscheme of $X$, so
the inclusion is exactly the stated equality. Thus the conclusion
concerns the full modified adelic set, proving
Theorem~\ref{thm:positive-rank-split-intro}.
\end{proof}

The local hypothesis of this criterion has not been proved
at a density-one set of split primes for $X_{17}$.
Proposition~\ref{prop:regular-primary-brauer-cusp-discs} determines
the component at seventeen, while
Corollary~\ref{cor:seventeen-no-rankzero-quotient} explains
why a rank-zero abelian target cannot complete the argument.

\section{The local cusp condition at every regular prime}
\label{sec:all-regular-local}

The first-case obstruction and McCallum's cusp-disc calculation
together determine the entire local representative of a compact
Selmer datum at its coefficient prime.

\begin{corollary}[The regular-prime local cusp condition]
\label{cor:all-regular-primary-local}
Let \(p>5\) be regular, and let
\(t\in H^1(\Q,T_pJ_p)\) have localizations
\(\kappa_{p,v}([R_v-Q_0])\) for points \(R_v\in F_p(\Q_v)\)
at every finite place. Then every permitted \(R_p\) is one of
the three rational coordinate cusps. The same conclusion holds
for every quotient Selmer section that maps to the proper
geometrically abelian pro-\(p\) quotient and has full local
point-or-cusp representatives.
\end{corollary}
\begin{proof}
Reduce \(t\) modulo \(p\) and use the proper point section at
\(Q_0\) as origin. This gives a section of the proper geometric
mod-\(p\) abelian quotient with the prescribed full proper local
point representatives. It satisfies the hypotheses of
Theorem~\ref{thm:irreg-rank}. Regularity gives \(r^-=0\), and its
inequality is \(4<p\). Thus every \(R_p\) belongs to a coordinate
cusp residue disc. The regular-prime cusp-disc conclusion of
Corollary~\ref{cor:seventeen-primary-local}, proved by finite
primary reciprocity against McCallum's supported classes,
then identifies \(R_p\) with the cusp itself.
A quotient section gives precisely this compact class by
abelianizing after filling the boundary. A full local cusp
representative fills to its proper centre, so the same argument
applies to every choice of its local representatives. This proves
Theorem~\ref{thm:regular-primary-local-intro}.
\end{proof}

This is a restriction on the local centres of global Selmer data.
It is not a statement that every point of \(F_p(\Q_p)\) lies in a
coordinate cusp disc. For example, at \(p=7\), if \(\omega\) is
a nontrivial cube root of unity in \(\Q_7\), then
\([1:\omega:\omega^2]\) is a first-case point of the all-plus
Fermat curve. The finite global torsors exclude it as a centre
of the data under consideration.

\section{Regular-prime Selmer sections, after McCallum}
\label{sec:regular-prime-sections}

The regular-prime local cusp condition combines with arithmetic-source
localization to classify the complete primary section. Throughout
this section \(p>5\) is regular and \(n=p^r\), with \(r\ge1\).
We use the all-plus model; changing the sign of \(Z\) identifies it
with the coordinates in Section~\ref{sec:irreg-setup}. Put
\(N_n=\pi_1(U_{n,\overline\Q})^{\mathrm{pro}\text{-}p}\).

\begin{lemma}[Compact cancellation and boundary centers]
\label{lem:all-regular-compact-boundary}
Let $s$ be a quotient Selmer section retaining $N_n^{\mathrm{ab}}$.
There is a unique rational coordinate cusp $P$ such that every
full local witness at $p$ is a cusp witness at $P$ and the global
compact abelian class equals that of $P$. At every finite $q\ne p$,
each local cusp witness has center $P$, and each ordinary local
witness reducing to the boundary reduces to $P$.
\end{lemma}
\begin{proof}
The power map $X_n\to X_p$ sends the compact abelian class to a
locally Kummer class. Corollary~\ref{cor:all-regular-primary-local} forces its local center at
$p$ to be a rational coordinate cusp. The fibre upstairs is boundary,
and $\Q_p$ has no nontrivial $p$-power roots of unity. Its only
$\Q_p$-point is the corresponding rational cusp $P$ of $X_n$.
An ordinary point cannot map into this boundary, so the full local
witness is cuspidal. Relative to the standard tangent at $P$, its
compact class $a\in H^1(\Q,T_pJ_n)$ is zero at $p$ and unramified
elsewhere by proper good reduction.
Proposition~\ref{prop:regular-nonabelian-localization} gives $a=0$.
For a full section its entire compact pro-$p$ image is the proper
point section $P$ by the same proposition; a smaller quotient
has only its own compact point image asserted.

For any geometric coordinate cusps $Q,P$, tangent-line ratios
have divisor $n(Q-P)$. Thus $[Q-P]$ is $p$-primary torsion. If it
were zero for distinct cusps, a function with divisor $Q-P$ would
give a degree-one map from the positive-genus curve to $\mathbb P^1$.
Hence distinct cusps have distinct classes. This calculation agrees
with the exact cusp-order statement of
\citet[p.~191, before Theorem~A]{Coleman1986}.
The $p$-primary torsion of $J_n(\Q_q)$ injects into its full
$p$-primary completion at every finite $q$: its torsion is finite
and its local group has an open additive lattice
\citep[I.3.3, p.~41]{Milne2006ADT}. Injection at an individual
finite Kummer level is not required. This proves uniqueness of
$P$ and forces any local cusp center $Q$ with zero compact
difference to equal $P$. If an ordinary point $x$ reduces to the
boundary at $q\ne p$, that boundary point lifts uniquely to
$Q\in D_n(\Q_q)$, since the boundary is finite \'etale and disjoint
there. The difference $[x-Q]$ is in a pro-$q$ formal group, on
which multiplication by $p$ is invertible. Its Kummer class is zero,
so again $Q=P$. This treats all locally occurring geometric
boundary components, without requiring interior witnesses to equal $P$.
\end{proof}

\begin{lemma}[The global abelian cusp parameter]
\label{lem:all-regular-abelian-packet}
The abelian Selmer section is the image of one global cusp section
$t_{P,\beta}$, with unique $\beta\in H^1(\Q,\Z_p(1))$. Its parameter
equals that of every full local cusp witness at $p$.
At $q\ne p$, the whole geometric inertia of a full local pro-$p$
witness is, up to conjugacy, $i_P(v_q(\beta)\tau_q)$, where
$\tau_q$ is the tame pro-$p$ character.
\end{lemma}
\begin{proof}
The inertia sequence is
\[
0\longrightarrow B_n\longrightarrow N_n^{\mathrm{ab}}
 \longrightarrow T_pJ_n\longrightarrow0,\qquad
B_n=\Z_p(1)[D_n(\overline\Q)]/\Z_p(1)\sum_Qe_Q.
\]
The compact class is zero relative to $P$.
The Tate module has no fixed vectors over $\Q$ or a finite
completion, since local torsion is finite. Thus the abelian class
has a unique preimage $b\in H^1(\Q,B_n)$, and boundary cohomology
injects locally too. Eliminate a second rational cusp generator
$e_{R_0}$ from the boundary relation. This gives an integral
$G_\Q$-equivariant splitting $B_n=I_P\oplus M_n$, where
$I_P=\Z_p(1)e_P$. Write $b=(\beta,b_M)$. The local cusp at $p$
gives $(b_M)_p=0$.

At $q\ne p$, an ordinary local witness reducing inside $U_n$
has zero geometric inertia. If it meets a boundary component,
its inertia is the tame character times its positive intersection
multiplicity. A local cusp has inertia given by its completed
Kummer valuation. Indeed, for a parameter $u$ in the strictly
henselian punctured normal disc, the pro-$p$ covers are
$u^{1/p^m}$; pullback gives the Kummer class of $u(x)$.
The same calculation applies to a cusp parameter.
Tame specialization and boundary inertia are
\citet[Expos\'e~XIII, Corollary~2.9, \S2.10 and Lemma~2.11,
pp.~288--290]{SGA1}; changing paths conjugates inertia.
By Lemma~\ref{lem:all-regular-compact-boundary}, every occurring
boundary center is $P$. Hence $b_M$ is zero on inertia away
from $p$. Coefficient inertia is trivial, so coboundaries do
not change this statement about the actual cocycle.
The module $M_n$ inherits the pro-$p$ arithmetic-source action;
Proposition~\ref{prop:regular-nonabelian-localization} gives $b_M=0$.

The summand $I_P$ and fixed-point vanishing make
$H^1(\Q,I_P)\to H^1(\Q,N_n^{\mathrm{ab}})$ injective, and likewise
at each finite place. Thus $\beta$ and its local parameter are
unique. This injection also identifies the exponent of the full
inertia formula with $v_q(\beta)$. If it is zero, the whole
inertia is zero. No finite-support restriction has been imposed.
\end{proof}

\begin{lemma}[Free pro-$p$ groups and positive-weight fixed points]
\label{lem:regular-group-fixed-points}
The group $N_n$ is free pro-$p$ of rank $n^2+1$, and each cusp
inertia is generated by a free basis element. Its ordinary
lower-central factors are finite free $\Z_p$-modules and form
the free graded Lie algebra. For nonzero $b\in\Z_p$,
$N_{N_n}(I_P^b)=I_P$.
If a quotient $R$ retains $N_n^{\mathrm{ab}}$ and all its ordinary
factors $L_m=\gamma_mR/\gamma_{m+1}R$ are torsion-free, then for
$t=t_{P,\beta}$ and every finite place $v$,
\[
 L_m^{G_{\Q_v}}=0,\qquad
 (R/\gamma_mR)^{t(G_{\Q_v})}=1,\qquad R^{t(G_\Q)}=1.
\]
The full group, its finite ordinary-central quotients and its
finite-derived quotients satisfy the factor hypothesis.
\end{lemma}
\begin{proof}
The punctured-surface presentation has $2g(X_n)$ handle generators,
$3n$ inertia generators and one relation. Eliminate a boundary
generator different from the chosen one. The free rank is
$2g+3n-1=n^2+1$; the presentation and comparison are
\citet[Expos\'e~XIII, Corollary~2.12, pp.~290--291;
Expos\'e~XII, Theorem~5.1]{SGA1}.
For every prime, the ordinary factors of a finite-rank free
pro-$p$ group are finite free with Witt ranks
\citep[\S3.5, p.~71, after Theorem~3.5.1]{LubotzkySegal2002}.
The natural free-Lie surjection has equal ranks, hence is an
isomorphism. Thus $(N_n^{\mathrm{ab}})^{\otimes m}\twoheadrightarrow L_m$.
Every pro-$p$ group's ordinary lower-central intersection is
trivial: a nonidentity element survives a finite $p$-group,
which is nilpotent.

A normalizer of $I_P^b=\langle a^b\rangle$ acts by a unit on this
procyclic group. Abelianization, where $a$ is primitive and $b\ne0$,
forces the unit to be one. In the free Lie algebra over $\Q_p$,
the centralizer of $A$ is $\Q_pA$: in the free associative algebra,
a homogeneous element commuting with $A$ is a multiple of $A^m$;
for $m>1$ this is not primitive for the characteristic-zero
coproduct. The degree-one class of a group centralizer is therefore
an integral multiple of $A$. Remove the corresponding power of
$a$. Every higher class then commutes with $bA$ and vanishes by
torsion-freeness. The trivial intersection proves the normalizer
assertion, also when $b$ is divisible by $p$.

Inner actions are trivial on $L_m$, so the cusp action there is
the original standard-tangent action. At $q\ne p$, Frobenius
eigenvalues on $T_pJ_n$ have absolute value $q^{1/2}$ and those
on the boundary are $q$ times roots of unity
\citep[Theorem~(1.6), p.~276]{Deligne1974}.
The tensor surjection gives absolute values $q^{w/2}$ on
$L_m\otimes\Q_p$, with $m\le w\le2m$; none is one.
Torsion-freeness proves fixed-point vanishing. At $p$ the graded
action factors through $\Gamma$, dominated by $G_{\Q_p}$.
Its fixed vectors are global fixed vectors, hence zero at any
good prime. Filtering $R/\gamma_mR$ gives its trivial fixed
subgroup, even for an infinitely ramified cusp parameter.
The trivial intersection also gives $R^{t(G_\Q)}=1$.

For finite-derived quotients,
\citet[\S3, Theorem~3.1, p.~101; definition p.~91]{Shmelkin1964}
gives free abelian ordinary factors for free solvable discrete
groups. For a finite nilpotent truncation, choose collected Malcev
coordinates. Its integer-valued binomial multiplication and
inversion formulas extend to $\Z_p$; density proves their laws.
Every map to a finite $p$-group is continuous: if $p^a$ is an
exponent of the target, expressions with all coordinates divisible
by $p^a$ map to one. The resulting compact pro-$p$ group is
therefore the full completion, with central factors obtained by
tensoring with $\Z_p$. Apply this to $S_d/\gamma_{c+1}S_d$,
where $F_0$ is the discrete free group on $n^2+1$ generators and
$S_d=F_0/F_0^{[d]}$ is free solvable. Its completion is
$(N_n/N_n^{[d]})/\gamma_{c+1}(N_n/N_n^{[d]})$, because both
classify the same finite $p$-quotients of the prescribed derived
length and nilpotency class. Taking $c\ge m$ proves the assertion.
No derived quotient of an open subgroup is assumed to embed in
a derived quotient downstairs.
\end{proof}

\begin{theorem}[Regular-prime cusp packets]
\label{thm:all-regular-prime-packets}
Let \(p>5\) be regular and let \(n=p^r\), \(r\ge1\).
For every stable quotient $R$ of $N_n$ retaining abelianization
and having torsion-free ordinary lower-central factors, its
point-or-cusp Selmer set is exactly the three disjoint rational
coordinate-cusp packets, each injectively parametrized by
$H^1(\Q,\Z_p(1))$. This includes the full group, every finite
ordinary-central quotient retaining abelianization and every
finite-derived quotient.
\end{theorem}
\begin{proof}
Let $t=t_{P,\beta}$ be the abelian comparison section of
Lemma~\ref{lem:all-regular-abelian-packet}.
At $q\ne p$ with $v_q(\beta)\ne0$, project a full local witness
to the full group $N_n$. Its inertia is a conjugate of
$I_P^{v_q(\beta)}$. Arithmetic inertia is normal, so the section
image normalizes this subgroup. Conjugate it to $I_P^{v_q(\beta)}$.
The standard action preserves $I_P$, and the normalizer lemma
makes its normalizer in $N_n\rtimes G_{\Q_q}$ precisely
$I_P\rtimes G_{\Q_q}$. Thus the local pro-$p$ witness is
cuspidal, with parameter $\beta_q$ by the local abelian injection.
An ordinary point in a boundary disc need not be cuspidal in the
all-prime group; only its pro-$p$ image is used here.
At $p$ the witness is actually a cusp at $P$ with parameter
$\beta_p$. These give local conjugacy to $t$ in every quotient $R$.
At the other $q\ne p$, both whole geometric inertias are zero.

Conjugate globally to make $s=t$ modulo $\gamma_2R$.
Suppose agreement holds modulo $\gamma_mR$, $m\ge2$.
The next difference is a cocycle $z_m$ in $L_m$, where the action
of $t$ equals the standard-tangent action.
At a place of local conjugacy, the image of a local conjugator
in $R/\gamma_mR$ is fixed by the common section and hence is one.
The conjugator lies in $\gamma_mR$, so its image in $L_m$
makes the local class of $z_m$ zero; in particular $(z_m)_p=0$.
At $q\ne p$ with nonzero cusp valuation, its inertia commutator
lies in $[\gamma_mR,R]\subset\gamma_{m+1}R$. Thus $z_m$ is zero
on inertia. At zero cusp valuation both whole inertias already
vanish. Prior global conjugations preserve this because the
original coefficient inertia action is trivial.

The graded cocycle is unramified outside $p$ and zero locally at
$p$. Proposition~\ref{prop:regular-nonabelian-localization} makes
its global class zero. Lift a correcting coboundary to $\gamma_mR$
and advance one step. For the original pair, the sets of conjugators
giving agreement modulo $\gamma_mR$ are nonempty, closed and nested
in the compact group $R$. A common element and the trivial central
intersection give agreement in $R$. Only the graded difference,
not the whole relative cusp cocycle, has been declared unramified.

The inertia summand injects on global abelian cohomology, so each
packet is injective in every retained quotient. Distinct centers
have distinct compact $p$-primary torsion classes.
Conversely any $\beta\in H^1(\Q,\Z_p(1))$ lifts to
$H^1(\Q,\widehat\Z(1))$ by setting all other prime components to zero.
The full global cusp section gives every required local witness,
including a real one. Odd-primary real cohomology is zero;
uniqueness of a chosen real witness is not asserted.
Parameters with infinite support are allowed, for example
valuations $b_{q_j}=p^j$ at distinct primes.
Surjectivity, injection and disjointness of the packets now prove
Theorem~\ref{thm:regular-primary-packets-intro}.
\end{proof}

\begin{corollary}[The classified prosolvable cohomology fibre]
\label{cor:classified-seventeen-kernel-fibre}
Let $n=17^r$, $r\ge1$, write $N_n^{\mathrm{sol}}$ for the
maximal geometric prosolvable group, and put
$K_{17}=\ker(N_n^{\mathrm{sol}}\to N_n)$.
Every prosolvable Selmer section has cuspidal image
$t_{P,\beta_{17}}$ in $N_n$.
Extend $\beta_{17}$ by zero at other primes to $\widehat\beta$,
and use its global prosolvable cusp section $\widehat t$ as origin.
Then
\[
 H^1(\Q,K_{17})\ \xrightarrow{\ \sim\ }\
 \operatorname{fib}_{1}\bigl(H^1(\Q,N_n^{\mathrm{sol}})
                    \longrightarrow H^1(\Q,N_n)\bigr),
\]
where all actions are induced by $\widehat t$.
Its point-or-cusp Selmer subset is exactly the image of
\[
 H^1\left(\Q,\prod_{\ell\ne17}\Z_\ell(1)\right)
\]
under the full cusp packet based at \(\widehat t\).
\end{corollary}
\begin{proof}
Every finite $17$-group is solvable, so the geometric projection
is surjective. Lift the conjugator identifying the projected
section with $t_{P,\beta_{17}}$. Its relative cocycle then lies
in $K_{17}$, proving surjectivity onto the fibre.
If two such cocycles are conjugate by $h\in N_n^{\mathrm{sol}}$,
its image in $N_n$ is fixed by the cusp action.
Lemma~\ref{lem:regular-group-fixed-points} makes that fixed subgroup
trivial. Hence $h\in K_{17}$ and the map is injective.
By Theorem~\ref{thm:regular-solvable-intro}, proved in the next
section, every Selmer section in this fibre is a full cusp section.
Its centre is \(P\), since distinct rational cusps have distinct
compact \(17\)-primary classes. The difference of its parameter
from \(\widehat\beta\) has zero \(17\)-component, and every such
prime-to-\(17\) parameter gives a Selmer section in this fibre.
The injectivity of the full cusp packet gives the asserted torsor
description. This proves
Corollary~\ref{cor:classified-seventeen-fibre-intro}.
It does not identify all of \(H^1(\Q,K_{17})\) with that subset.
\end{proof}

\section{All-prime Fermat sections at regular exponents}
\label{sec:allprime-regular-bridge}

One global affine unit turns the regular-primary local cusp condition
into a common centre at every finite place. Write
\[
 F_n:\ X^n+Y^n=Z^n,\qquad
 U_n=F_n\setminus\{XYZ=0\},\qquad
 \mathcal D_n=F_n\setminus U_n.
\]
Throughout this section \(n\geq3\) is odd.  The three rational
coordinate points are
\[
 P_0=(0:1:1),\qquad P_1=(1:0:1),\qquad
 P_\infty=(1:-1:0).
\]
The notation \(\mathcal D_n\) includes all \(3n\) geometric boundary
points, not just these three rational points.  Changing the sign of
\(Z\) identifies this model with the all-plus Fermat model.

Put \(\Delta_n=\pi_1(U_{n,\overline{\mathbb Q}})\), and define
\[
 \Delta_n^{[0]}=\Delta_n,\qquad
 \Delta_n^{[d+1]}=
 \overline{[\Delta_n^{[d]},\Delta_n^{[d]}]},\qquad
 R_{n,d}=\Delta_n/\Delta_n^{[d]}\quad(d\geq1).
\]
Also put
\[
 R_{n,\mathrm{sol}}
 =\Delta_n\big/\bigcap_{d\geq0}\Delta_n^{[d]}.
\]
Every arithmetic quotient considered below retains the entire group
\(G_{\mathbb Q}\).  Section classes are taken modulo conjugation by
the indicated geometric group.  The local condition means that,
at each finite place, the quotient section is the image of either
a full local point section of \(U_n\), or a full section of the
decomposition group of a \(\mathbb Q_v\)-rational point of
\(\mathcal D_n\).  At the real place the analogous point-or-cusp
condition is imposed.  A full global lift is not part of this
definition.

\begin{theorem}[All-prime cusp packets at regular exponents]
\label{thm:allprime-regular-bridge}
Suppose that \(n\) is divisible by \(3\), by \(5\), or by an odd
regular prime \(p>5\).  For every \(d\geq1\), and for the maximal
geometrically prosolvable quotient, every point-or-cusp Selmer
section is cuspidal.  After choosing one rational tangent at each
rational coordinate point, its section set is
\begin{equation}
 \coprod_{P\in\{P_0,P_1,P_\infty\}}
 H^1(\mathbb Q,\widehat{\mathbb Z}(1)).
 \label{eq:allprime-packets}
\end{equation}
Each displayed packet maps injectively to the section set, and
the three images are disjoint.  Every permitted full local
representative at a finite place has the same rational coordinate
centre as its global packet.

No common finite ramification set is required for the section or
its cusp parameter.
\end{theorem}

\subsection{Full Kummer completion and one fixed cusp}

We use the power completion
\begin{equation}
 \widehat{\mathbb Q^\times}
 =\varprojlim_{N\geq1}\mathbb Q^\times/(\mathbb Q^\times)^N
 =H^1(\mathbb Q,\widehat{\mathbb Z}(1)).
 \label{eq:allprime-power-completion}
\end{equation}
The inverse limit is over divisibility.  This notation does not
assert that the individual power quotients, or this inverse limit,
are compact.  The Kummer identification follows from the finite
Kummer sequences; the finite groups of rational roots of unity
satisfy the Mittag--Leffler condition on a cofinal factorial
subsystem.

\begin{lemma}[Valuations and local Kummer maps]
\label{lem:allprime-kummer}
The common kernel of all completed valuations
\[
 v_q:\widehat{\mathbb Q^\times}\longrightarrow\widehat{\mathbb Z}
 \qquad(q\text{ a rational prime})
\]
is \(\{1,-1\}\).  Multiplication by any odd positive integer is
injective on \(\widehat{\mathbb Q^\times}\).

For every finite place \(v\), the map from
\(\mathbb Q_v^\times\) to its full power completion is injective.
For an abelian variety \(A/\mathbb Q_v\), the map
\[
 A(\mathbb Q_v)\longrightarrow
 \varprojlim_N A(\mathbb Q_v)/N A(\mathbb Q_v)
\]
is injective as well.
\end{lemma}

\begin{proof}
Unique factorization gives
\[
 \mathbb Q^\times=\{1,-1\}\times
                     \bigoplus_q q^{\mathbb Z}.
\]
At every finite power level, a class whose valuation at each prime
is zero is represented by a sign.  Compatibility over the even
power levels makes that sign independent of the level.  This
proves the first assertion.

If \(\alpha^m=1\) and \(m\) is odd, then
\(m v_q(\alpha)=0\) for every \(q\).
The group \(\widehat{\mathbb Z}=\prod_\ell\mathbb Z_\ell\)
is torsion-free.  Thus all valuations vanish, so
\(\alpha\in\{1,-1\}\); oddness gives \(\alpha=1\).
In particular, for \(\epsilon\in\{1,-1\}\), the equation
\(\alpha^m=\epsilon\) has the unique solution
\(\alpha=\epsilon\).

The local multiplicative assertion follows from the usual
decomposition into valuation, finite roots of unity, and a
principal-unit group with an open \(\mathbb Z_v\)-lattice.
For the abelian assertion, \(A(\mathbb Q_v)\) is compact and has
an open subgroup \(\Lambda\simeq\mathbb Z_v^{\dim A}\);
see \citet[Lemma~I.3.3, p.~41]{Milne2006ADT}.
Choose a positive integer \(e\) annihilating
\(A(\mathbb Q_v)/\Lambda\).  If \(x\) is divisible by every
positive integer in \(A(\mathbb Q_v)\), then
\[
 x\in e v^r A(\mathbb Q_v)\subset v^r\Lambda
 \quad\text{for every }r,
\]
where \(v\) in this expression denotes the underlying rational
prime.  The intersection of these lattices is zero.
\end{proof}

In particular, full local Kummer cohomology distinguishes the
two rational signs.  A single odd-primary Kummer group need
not do so.  This distinction will be used at the prime \(p\).

\begin{lemma}[Recovering a full global cusp parameter]
\label{lem:allprime-one-cusp}
Let \(X/\mathbb Q\) be a smooth proper geometrically connected curve
of genus at least one, let \(D\) be a reduced finite divisor with at
least two geometric points, and put \(U=X\setminus D\).
Let \(P\in D(\mathbb Q)\).
Suppose a section of a finite full-geometric derived quotient, or
of the maximal geometrically prosolvable quotient, has at every
finite place a full local cusp representative at \(P\).
Then it is the image of a unique full global cusp section at \(P\),
where uniqueness refers to its
\(H^1(\mathbb Q,\widehat{\mathbb Z}(1))\) parameter after a
rational tangent has been chosen.
\end{lemma}

\begin{proof}
Choose a rational tangent at \(P\) and let \(t_P\) be its full section.
Filling \(P\) leaves an affine curve \(W=U\cup\{P\}\).
Lemma~\ref{lem:boundary} gives a unique
\(\beta\in H^1(\Q,\widehat\Z(1))\) whose peripheral image equals
the relative global abelian class of the given section.
Indeed, its image for \(W\) vanishes at every finite place,
and \citet[Proposition~3.2, author pp.~13--14]{Saidi2022} makes
that image globally zero. The lemma proves the global and local
peripheral injections, including their fixed-point hypothesis.

The parameter \(\beta\) defines a full cusp section \(t_{P,\beta}\).
At each finite place, the local peripheral injection identifies
its parameter with that of every full local cusp representative.
The two sections are therefore locally conjugate in the stated
quotient. Use the action from the constructed full section
\(t_{P,\beta}\), as required by the introduction of
\citet{Saidi2022}. His Proposition~3.1, author pp.~13--15,
makes the relative class globally zero at every derived depth
\(d\ge1\). His Theorem~B does the same in the prosolvable quotient.
Equivalently, the nonempty closed sets of conjugators giving
agreement at each derived depth are nested in the compact
geometric prosolvable group. A common conjugator gives agreement
in their inverse limit. This proves existence and uniqueness.
\end{proof}

The lemma imposes neither finite support on \(\beta\) nor a common
finite ramification set on its finite levels.  For example,
enumerate distinct rational primes \(q_j\), and specify completed
valuations \(v_{q_j}(\beta)=j!\), with all other valuations zero.
For each \(N\), only finitely many of these valuations are
nonzero modulo \(N\).  Their finite product defines a class in
\(\mathbb Q^\times/(\mathbb Q^\times)^N\), and these classes are
compatible as \(N\) varies.  They give a cusp parameter with
infinitely many nonzero valuations.

\subsection{The base regular prime}

\begin{proposition}[From one primary cusp to all finite centres]
\label{prop:allprime-base-regular}
Let \(p>5\) be regular.  Every section satisfying the local
condition on \(U_p\), in any quotient in
Theorem~\ref{thm:allprime-regular-bridge}, has all its full finite
local witnesses cuspidal at one rational coordinate point.
\end{proposition}

\begin{proof}
Choose arbitrary full local witnesses. Their projection to the
geometric pro-\(p\) abelian quotient satisfies the same local
condition. Corollary~\ref{cor:all-regular-primary-local} makes the
proper \(p\)-local centre a rational cusp \(P\).
Lemma~\ref{lem:all-regular-compact-boundary}, applied with \(n=p\),
then identifies the global compact \(T_pJ_p\)-class with that of
\(P\). At every \(q\ne p\), it also proves that a cusp witness
has centre \(P\), while an ordinary witness either reduces to
the open part or to \(P\). This use of the lemma concerns the
projection of the given quotient section and the same full
local witnesses; it assumes no full global lift.
A rational signed coordinate permutation lets us take
\(P=P_0=(0:1:1)\).
Put \(u=Y/Z\), a unit on \(U_p\) that extends as a unit at
\(P_0\), with \(u(P_0)=1\).

Pulling back the covers \(z^N=u\) along the quotient section
gives a single global class
\[
 \alpha\in H^1(\mathbb Q,\widehat{\mathbb Z}(1))
        =\widehat{\mathbb Q^\times}.
\]
All these covers have abelian geometric monodromy, so this
construction is available at every \(d\geq1\).
At a full local cusp section centred at \(P_0\), the value is
the Kummer class of \(u(P_0)=1\); inertia contributes nothing
because \(u\) extends as a unit there.
At an ordinary witness reducing to the open part, \(Y/Z\) is
a \(q\)-adic unit.
At an ordinary witness reducing to \(P_0\), both \(Y\) and \(Z\)
are units, so the same conclusion holds.
The \(p\)-local value is exactly \(1\).
We have proved
\[
 v_q(\alpha)=0\quad\text{for every rational prime }q,
 \qquad \alpha_p=1
 \quad\text{in }\widehat{\mathbb Q_p^\times}.
\]
Lemma~\ref{lem:allprime-kummer} first gives \(\alpha\in\{1,-1\}\).
Full local Kummer injectivity at \(p\) then gives \(\alpha=1\).

If a full witness at any finite \(q\) were an ordinary point
\(x_q\), the same local injectivity would imply
\(u(x_q)=1\).
The Fermat equation would then give \(X(x_q)=0\), contrary to
\(x_q\in U_p\).
Thus every finite witness is cuspidal at \(P_0\), as required.

The witnesses were arbitrary.  This proves the assertion for
every permitted choice, not just for a preferred system of
local representatives.
\end{proof}

Lemma~\ref{lem:allprime-one-cusp} completes the classification at
the base regular prime.

\subsection{Odd multiples without a good-reduction restriction}

\begin{proposition}[Odd-multiple transfer]
\label{prop:allprime-odd-transfer}
Let \(a\geq3\) be odd, and suppose that the conclusion of
Theorem~\ref{thm:allprime-regular-bridge} holds at exponent \(a\).
Then it holds at every exponent \(n=am\) with \(m\) odd.
\end{proposition}

\begin{proof}
The coordinate-power morphism
\[
 \phi_m:F_n\longrightarrow F_a,\qquad
 (X:Y:Z)\longmapsto(X^m:Y^m:Z^m)
\]
takes \(U_n\) to \(U_a\) and its boundary to the boundary.
Functoriality induces maps of all the stated finite-derived
arithmetic quotients and of their prosolvable limits.
The full arithmetic quotient \(G_{\mathbb Q}\) is unchanged.
There is no need to assert that a finite-derived quotient of a
covering subgroup embeds in the corresponding quotient below.

Push forward the section and each of its full local witnesses.
The resulting section at exponent \(a\) is globally cuspidal at
one rational coordinate point \(P_a\).
At each finite place \(v\), the proper full Tate-module Kummer
class of the image witness equals that of \(P_a\).
The Abel--Jacobi map is injective on \(F_a(\mathbb Q_v)\):
two distinct points with the same divisor class would give a
degree-one function on a curve of positive genus.
Its composition with full local Kummer completion remains
injective by Lemma~\ref{lem:allprime-kummer}.
Thus the image witness has proper centre exactly \(P_a\).
An ordinary point of \(U_n\) maps to an ordinary point of
\(U_a\); therefore every original finite witness is a cusp
witness at some
\[
 Q_v\in\mathcal D_n(\mathbb Q_v),\qquad \phi_m(Q_v)=P_a.
\]
This argument also applies at primes dividing \(m\).
It does not use good reduction there.

Choose signed rational coordinate permutations to make
\(P_a=P_0\).
Because \(m\) is odd, these permutations commute with the
coordinate-power maps.
Let \(u_n=Y/Z\) and \(u_a=Y/Z\).
Then \(\phi_m^*u_a=u_n^m\), and the \(u_a\)-Kummer class of the
pushed section is \(1\), since \(u_a\) extends as a unit of
value \(1\) at \(P_0\).
Consequently the global \(u_n\)-class satisfies
\[
 \alpha_n^m=1\quad\text{in }\widehat{\mathbb Q^\times}.
\]
Odd multiplication is injective, so \(\alpha_n=1\).

At every \(Q_v\) over \(P_0\), the function \(u_n\) extends as
a unit.  Full local Kummer injectivity gives \(u_n(Q_v)=1\).
Since \(X(Q_v)=0\), this says
\(Q_v=(0:1:1)\), the rational cusp \(P_0\) of \(F_n\).
Lemma~\ref{lem:allprime-one-cusp} now identifies the original
section with its unique full global cusp lift.
\end{proof}

\begin{proof}[Proof of Theorem~\ref{thm:allprime-regular-bridge}]
At a regular prime \(p>5\), use
Proposition~\ref{prop:allprime-base-regular} and
Lemma~\ref{lem:allprime-one-cusp}.
For the bases \(3\) and \(5\), the established arithmetic
inputs give
\begin{equation}
 F_a(\mathbb A_{\mathbb Q})_\bullet^{f\text{-}\mathrm{ab}}
       =F_a(\mathbb Q)
       =\{P_0,P_1,P_\infty\}\qquad(a=3,5).
 \label{eq:allprime-small-base-input}
\end{equation}
For \(a=3\), use the elliptic modular model \(X_0(27)\) and
\citet[Corollary~8.8, together with Theorem~8.6]{Stoll2007}.
For \(a=5\), the nonconstant map to the rank-zero abelian
variety with finite Tate--Shafarevich group in
Proposition~\ref{prop:arithmetic} supplies
\citet[Theorem~8.6, p.~385]{Stoll2007}.
The rational point lists and the modified real-place
convention are those established with these arithmetic
inputs; Stoll's real-place erratum is included.

For completeness, the step from
\eqref{eq:allprime-small-base-input} to sections uses no global
full lift.
A quotient section retaining proper geometric abelianization
makes its proper local centres survive every finite abelian
torsor: pull the torsor back to the section and twist by that
global class; the full local witnesses provide local lifts.
The displayed equality therefore makes all finite centres one
rational cusp.
Lemma~\ref{lem:allprime-one-cusp} gives the required full cusp
parameter and the section equality.
Proposition~\ref{prop:allprime-odd-transfer} now treats every
odd multiple in the theorem.

Every element of
\(H^1(\mathbb Q,\widehat{\mathbb Z}(1))\) defines a full
global cusp section, hence satisfies all the required local
conditions, including the real one.
Within one centre, injectivity follows from
Lemma~\ref{lem:boundary}.
If two different rational centres gave conjugate quotient
sections, their proper full Tate-module Kummer classes would
coincide at any finite place.
Local Kummer injectivity and injectivity of the Abel--Jacobi
map would make the two points equal, a contradiction.
This proves injectivity and disjointness of all three packets,
and hence Theorem~\ref{thm:regular-solvable-intro}.

No density or integrality condition was added to the local
hypothesis.  The only density-one theorem used is Sa\"idi's
localization theorem, applied after equality at every finite
place has already been proved.
\end{proof}

\begin{remark}[Scope of the bridge]
\label{rem:allprime-bridge-scope}
The proof uses all primes in geometric abelianization twice:
to select the rational sign of a global affine unit, and to
recover the full inertia parameter at a fixed cusp.
It does not assert diagonality of the finite-abelian-descent adelic set of the
proper curve at every regular exponent.
It is a theorem for the specified finite-derived and
prosolvable arithmetic quotients, with full local witnesses;
it does not deduce such a theorem for an arbitrary geometric
quotient or for a single-primary affine quotient.
\end{remark}

\section{Finite sets of compact local centres, after McCallum}
\label{sec:irreg-Tate}

McCallum bounds the local zero set of a suitable Coleman integral on
a wild cyclic quotient. We show that one fixed integral contains
the centres of every coherent compact Tate-Selmer datum. The
finite-level reciprocity argument is the step that supplies this
uniformity for integral cohomology classes.

For an abelian variety \(B/\Q_v\), let
\[
 \kappa_{p,v}:B(\Q_v)\longrightarrow H^1(\Q_v,T_pB)
\]
be the inverse-limit Kummer map; the analogous global map is denoted
\(\kappa_p:B(\Q)\to H^1(\Q,T_pB)\).

\begin{definition}\label{def:irreg-data}
A \emph{compact \(p\)-Selmer datum on \(F_p\)} consists of
\[
 t\in H^1(\Q,T_pJ_p),\qquad
 R_v\in F_p(\Q_v)\quad(v<\infty),
\]
such that
\begin{equation}
 \loc_v(t)=\kappa_{p,v}\bigl(j_p(R_v)\bigr)
 \quad\text{for every finite }v.
 \label{eq:irreg-data}
\end{equation}
\end{definition}

This definition uses integral Tate-module cohomology. It does not
mean an arbitrary Selmer class on the Jacobian without curve-valued
local representatives. It imposes no common finite set of
ramification primes on \(t\).

A quotient section with full local representatives supplies such a
datum whenever it retains the proper geometrically abelian pro-\(p\)
quotient. Indeed, the base point \(Q_0\) splits that quotient, whose
geometric group is \(T_pJ_p\); sections then give classes in
\(H^1(\Q,T_pJ_p)\), and point sections give Kummer classes.
The geometric identification follows at each finite level from
the Kummer sequence on \(F_{p,\overline\Q}\), the identification of
its degree-zero Picard variety with \(J_p\), and the principal
polarization of \(J_p\).

Conversely, reducing a datum modulo \(p\) gives a section of
\(\Pi_p^{\ab,p}\) with the prescribed proper local point
representatives. Thus the finite first-case criteria above apply
to every datum in Definition~\ref{def:irreg-data}.

For a principally polarized abelian variety \(A/\Q\), write
\begin{equation}
 H^1_{\{p\}}(\Q,A)
 =
 \ker\left(H^1(\Q,A)\longrightarrow
                  \prod_{v\neq p}H^1(\Q_v,A)\right).
 \label{eq:irreg-supported}
\end{equation}
The product includes the infinite place. This is the group of
Weil--Chatelet classes supported at \(p\), not the usual
everywhere locally trivial Tate--Shafarevich group.

\begin{lemma}\label{lem:irreg-reciprocity}
Let \(A/\Q\) be principally polarized, and let
\(u\in H^1(\Q,T_pA)\) have localizations
\(\kappa_{p,v}(P_v)\), with \(P_v\in A(\Q_v)\) at every finite place.
For every
\(d\in H^1_{\{p\}}(\Q,A)[p^n]\),
\begin{equation}
 \langle P_p,d_p\rangle_p=0
 \label{eq:irreg-supported-pairing}
\end{equation}
under local Tate duality.
\end{lemma}

\begin{proof}
Use the principal polarization to identify \(A\) with its dual.
Let \(u_n\in H^1(\Q,A[p^n])\) be the finite image of \(u\), and choose
\(b_n\in H^1(\Q,A[p^n])\) mapping to \(d\). Such a lift exists by
the Kummer exact sequence. The Weil pairing gives a global class
\[
 u_n\smile b_n\in H^2(\Q,\mu_{p^n})
                 \subseteq\operatorname{Br}(\Q).
\]
Compatibility of finite Kummer duality with the local Tate pairing
gives
\[
 \inv_v\bigl((u_n\smile b_n)_v\bigr)
     =\langle P_v,d_v\rangle_v.
\]
A consistent choice of the order of the dual factors fixes the sign;
only its vanishing is used. This compatibility is the finite-level
diagram in the proof of
\cite[I, Theorem 3.2, pp.~41--43]{Milne2006ADT};
the perfect local pairing is
\cite[I, Corollary 3.4, pp.~43--44]{Milne2006ADT}.
It also shows that changing the lift \(b_n\) by a Kummer class does
not change the displayed value.

For \(v\neq p\), the value is zero because \(d_v=0\). At the real
place the \(p\)-primary term is zero, since \(p\) is odd.
Global reciprocity now makes the remaining value at \(p\) zero;
equivalently one may use the orthogonality in
\cite[I, Theorem 4.10(b), pp.~56--57]{Milne2006ADT}.
This argument is made at one finite level at a time and needs no
uniform finite-ramification hypothesis on \(u\).
\end{proof}

\begin{lemma}\label{lem:irreg-tower-functional}
Suppose a compatible tower
\begin{equation}
 d=(d_n)_{n\geq1},\qquad
 d_n\in H^1_{\{p\}}(\Q,A)[p^n],\qquad
 p\,d_{n+1}=d_n,
 \label{eq:irreg-tower}
\end{equation}
has nonzero localization at \(p\). It determines a nonzero
continuous homomorphism
\begin{equation}
 \ell_d:A(\Q_p)\longrightarrow\Z_p.
 \label{eq:irreg-integral-functional}
\end{equation}
Its kernel contains all torsion and is saturated for multiplication
by every nonzero integer. There is a nonzero invariant differential
\(\omega_d\) on \(A/\Q_p\) such that
\begin{equation}
 \ell_d(P)=\omega_d\bigl(\log_A(P)\bigr)
 \quad\text{for all }P\in A(\Q_p).
 \label{eq:irreg-log-functional}
\end{equation}
Every collection \(P_v\) occurring in
Lemma~\ref{lem:irreg-reciprocity} satisfies \(\ell_d(P_p)=0\).
\end{lemma}

\begin{proof}
For each \(n\), pair \(P\) with \(d_{n,p}\) and identify
\[
 p^{-n}\Z/\Z\xrightarrow{\sim}\Z/p^n\Z,
 \qquad a/p^n\longmapsto a\bmod p^n.
\]
Compatibility of the tower gives a compatible system of residues,
hence \(\ell_d(P)\in\Z_p\). Every finite coordinate is continuous.
Perfect local Tate duality shows that this homomorphism is nonzero:
otherwise every \(d_{n,p}\) would vanish.

The target \(\Z_p\) is torsion-free, proving the assertions about
torsion and saturation. Choose an open subgroup \(H\subset A(\Q_p)\)
on which the logarithm identifies \(H\) with a full
\(\Z_p\)-lattice in \(\Lie(A)\). Such a subgroup exists by
\cite[I, Lemma 3.3, p.~41]{Milne2006ADT}.
The restriction of \(\ell_d\) becomes a continuous additive map on
that lattice, hence a \(\Z_p\)-linear map. Extend it uniquely to a
\(\Q_p\)-linear functional on \(\Lie(A)\).
This extension is nonzero: a homomorphism from the finite quotient
\(A(\Q_p)/H\) to \(\Z_p\) must vanish.

For arbitrary \(P\), choose a positive integer \(N\) with \(NP\in H\).
Both the logarithm and \(\ell_d\) commute with multiplication by
\(N\); division by \(N\) in \(\Q_p\) proves
\eqref{eq:irreg-log-functional}. The functional is the invariant
differential \(\omega_d\). Finally,
Lemma~\ref{lem:irreg-reciprocity} makes every finite coordinate of
\(\ell_d(P_p)\) zero.
\end{proof}

\subsection{Wild cyclic quotients}
\label{sec:irreg-wild}

Let \(a,b,c\) be integers with \(a+b+c=0\) and
\(p\nmid abc\). Choose \(a,b>0\) in their residue classes and set
\(c=-a-b\). Let \(C_{a,b,c}\) be the smooth projective curve with
affine equation
\begin{equation}
 y^p=(-1)^c x^a(1-x)^b.
 \label{eq:irreg-cyclic}
\end{equation}
It has genus \((p-1)/2\), by Riemann--Hurwitz for its degree-\(p\)
map to the \(x\)-line, totally ramified over \(0,1,\infty\).
Denote the unique points over these three values by
\(O_0,O_1,O_\infty\). They are rational. Put
\[
 A=\Jac(C_{a,b,c}),\qquad j(R)=[R-O_0].
\]
There is a morphism
\begin{equation}
 f:F_p\longrightarrow C_{a,b,c},\qquad
 x=X^pZ^{-p},\qquad y=X^aY^b(-Z)^c.
 \label{eq:irreg-quotient-map}
\end{equation}
The rational functions define the morphism on the smooth projective
curves by extension from their common affine domain.

\begin{lemma}\label{lem:irreg-quotient-free}
The morphism \(f\) is the quotient by the free action
\begin{equation}
 \mu_p\times F_p\longrightarrow F_p,\qquad
 \zeta\cdot(X:Y:Z)=(\zeta^bX:\zeta^{-a}Y:Z).
 \label{eq:irreg-action}
\end{equation}
In particular, \(f:F_p(\Q_p)\to C_{a,b,c}(\Q_p)\) is injective.
Moreover \(j(O_0),j(O_1),j(O_\infty)\) are torsion points.
\end{lemma}

\begin{proof}
The functions in \eqref{eq:irreg-quotient-map} are invariant under
\eqref{eq:irreg-action}. Over the \(x\)-line, the function field of
\(F_p\) has degree \(p^2\), whereas
\(\Q(x,y)/\Q(x)\) has degree \(p\), since the valuation \(a\) at
\(x=0\) is not divisible by \(p\). These invariants therefore give
the quotient field, and hence the quotient curve.

At a point with all coordinates nonzero, a projective stabilizer
must have \(\zeta^b=\zeta^{-a}=1\), hence \(\zeta=1\).
At \(X=0\) the stabilizer condition is \(\zeta^{-a}=1\);
at \(Y=0\) it is \(\zeta^b=1\); and at \(Z=0\) it is
\(\zeta^{a+b}=1\). Each again forces \(\zeta=1\), using
\(p\nmid abc\). Thus the action is free in every stratum and the
proper quotient is a \(\mu_p\)-torsor.

If two \(\Q_p\)-points have the same image, the unique geometric
element carrying one to the other is Galois-fixed, hence lies in
\(\mu_p(\Q_p)=\{1\}\). This proves injectivity, including at the
boundary.

Finally,
\[
 \divisor(x)=p(O_0-O_\infty),\qquad
 \divisor(1-x)=p(O_1-O_\infty).
\]
Their classes show that each of the three images under \(j\) is
killed by \(p\).
\end{proof}

Following \cite[pp.~395--398]{McCallum2001Brauer}, set
\[
 \gamma=a^a b^b c^c\in\Z_{(p)}^\times.
\]
The quotient is called \emph{wild} if
\begin{equation}
 \gamma^p\not\equiv\gamma\pmod{p^2}.
 \label{eq:irreg-wild}
\end{equation}

\begin{lemma}\label{lem:irreg-wild-exists}
For every prime \(p\geq5\), at least one of the quotients
\(C_{1,k,-k-1}\), \(1\leq k\leq p-2\), is wild.
\end{lemma}

\begin{proof}
For \(u\in\Z_{(p)}^\times\), define its Fermat quotient by
\[
 q_p(u)=\frac{u^{p-1}-1}{p}\pmod p.
\]
Multiplication modulo \(p^2\) shows that this is a homomorphism
from \(\Z_{(p)}^\times\) to the additive group \(\F_p\).
For \(\gamma_k=k^k(-k-1)^{-k-1}\), one has
\[
 q_p(\gamma_k)=kq_p(k)-(k+1)q_p(k+1).
\]
Here \(q_p(-1)=0\). Also
\((p-1)^{p-1}\equiv1+p\pmod{p^2}\), so \(q_p(p-1)=1\).
The telescoping sum is therefore
\[
 \sum_{k=1}^{p-2}q_p(\gamma_k)
   =q_p(1)-(p-1)q_p(p-1)=1\quad\text{in }\F_p.
\]
At least one summand is nonzero, which is exactly
\eqref{eq:irreg-wild}.
\end{proof}

We use two precise results of McCallum. In the following statements,
the class rank is the full \(p\)-class rank, not only its minus part:
\[
 r=\dim_{\F_p}\Cl(K)/p\Cl(K).
\]

\begin{proposition}[McCallum]\label{prop:irreg-McCallum}
For \(C=C_{a,b,c}\) and \(A=\Jac(C)\) as above:
\begin{enumerate}
\item If \(r<(p+3)/8\), there is a compatible tower
\[
 d_n\in H^1_{\{p\}}(\Q,A)[p^n],\qquad p\,d_{n+1}=d_n,
\]
whose localization at \(p\) is nonzero.
\item If \(C\) is wild, there is an affinoid
\(\mathcal X\subset C(\C_p)\) containing \(C(\Q_p)\) such that,
for every nonzero holomorphic differential \(\eta\) on \(C\),
the Coleman integral
\[
 I_\eta(R)=\int_{O_0}^R\eta
\]
has at most \(p\) zeros in \(\mathcal X\).
\end{enumerate}
\end{proposition}

These are \cite[Propositions 2 and 3, pp.~397--398]{McCallum2001Brauer}.
The construction of the tower is proved there in Section 3,
pp.~399--403; the analytic bound is proved in Section 4,
pp.~403--406. The affine region is the tube of the nonsingular
locus of the indicated regular model. Every \(\Q_p\)-point belongs
to it in the wild case. The zero bound applies to any nonzero
holomorphic differential; it does not require a complex
multiplication eigendifferential. The relation of the Coleman
integral to the abelian logarithm is given in the same paper,
Section 2, pp.~397--398.

\subsection{One fixed finite set of local centres}
\label{sec:irreg-fixed}

\begin{theorem}\label{thm:irreg-fixed}
Suppose
\begin{equation}
 r=\dim_{\F_p}\Cl(\Q(\zeta_p))/p\Cl(\Q(\zeta_p))
       <\frac{p+3}{8}.
 \label{eq:irreg-Mc-bound}
\end{equation}
There is one finite set \(S_p\subset F_p(\Q_p)\), independent of
the compact \(p\)-Selmer datum and of all its chosen local
representatives, such that
\begin{equation}
 \{Q_0,Q_1,Q_\infty\}\subseteq S_p,\qquad
 |S_p|\leq p,
 \label{eq:irreg-fixed-size}
\end{equation}
and every datum in Definition~\ref{def:irreg-data} has
\(R_p\in S_p\).

More precisely, one may choose any wild quotient \(C_{a,b,c}\);
a single supported tower for its Jacobian, fixed independently of
the datum, determines such a set.
\end{theorem}

\begin{proof}
Choose a wild quotient, which exists by
Lemma~\ref{lem:irreg-wild-exists}, and let \(A=\Jac(C_{a,b,c})\).
By Proposition~\ref{prop:irreg-McCallum}(1), choose once and for all
a supported tower \(d=(d_n)\) with nonzero \(p\)-localization.
Lemma~\ref{lem:irreg-tower-functional} gives the corresponding
nonzero invariant differential \(\omega_d\) on \(A\).
Put
\[
 \eta=j^*\omega_d.
\]
This is nonzero because the Abel--Jacobi map identifies invariant
differentials on the Jacobian with holomorphic differentials on
the curve. Functoriality of the abelian logarithm and Coleman
integration gives
\begin{equation}
 I_\eta(R)=\omega_d\bigl(\log_A(j(R))\bigr)
           =\ell_d(j(R)).
 \label{eq:irreg-fixed-integral}
\end{equation}

Let \(t,(R_v)\) be any compact \(p\)-Selmer datum on \(F_p\).
The induced homomorphism \(f_*:J_p\to A\), with the specified base
points, satisfies \(f_*j_p=jf\). The image \(u=f_*t\) therefore has
localizations
\[
 \loc_v(u)=\kappa_{p,v}(j(f(R_v))).
\]
Lemma~\ref{lem:irreg-reciprocity}, applied to every \(d_n\), and
\eqref{eq:irreg-fixed-integral} imply
\[
 I_\eta(f(R_p))=0.
\]
All choices of \(d,\omega_d,\eta\) preceded the arbitrary datum.

Define
\begin{equation}
 S_p=
 \{R\in F_p(\Q_p): I_\eta(f(R))=0\}.
 \label{eq:irreg-fixed-set}
\end{equation}
The image of \(F_p(\Q_p)\) lies in \(\mathcal X\), where
Proposition~\ref{prop:irreg-McCallum}(2) bounds the zeros of
\(I_\eta\) by \(p\). Lemma~\ref{lem:irreg-quotient-free} makes \(f\)
injective on \(\Q_p\)-points, so \(|S_p|\leq p\).

The three rational coordinate cusps map to \(O_0,O_1,O_\infty\).
Their images in \(A\) are torsion, and \(\ell_d\) kills torsion.
Thus all three belong to \(S_p\), proving \eqref{eq:irreg-fixed-size}.
\end{proof}

\begin{remark}\label{rem:irreg-uniform}
The set \(S_p\) need not be canonical: it depends on a fixed quotient
and a fixed supported tower. Its defining choices are independent
of the Selmer datum, which is the uniformity asserted in the
theorem. Neither a first-case hypothesis nor finiteness of a
Tate--Shafarevich group is used for the bound.
\end{remark}

\begin{corollary}\label{cor:irreg-fixed-discs}
In addition to \eqref{eq:irreg-Mc-bound}, suppose either
\eqref{eq:irreg-rank-bound} or the class-component hypothesis of
Corollary~\ref{cor:irreg-cubic}. Then \(S_p\) may be replaced by its
intersection with the three coordinate cusp residue discs while
retaining all conclusions of Theorem~\ref{thm:irreg-fixed}.
\end{corollary}

\begin{proof}
Reduction of every compact datum modulo \(p\) has the full proper
local representatives required for the first-case arguments.
Those arguments put every allowed centre in the three discs.
The intersection still contains the three cusps.
\end{proof}

\subsection{The primes 37, 59, and 67}
\label{sec:irreg-examples}

Masley's theorem gives class number one for every real abelian
field of prime-power conductor less than \(70\)
\cite[Theorem 3.2, p.~310]{Masley1978}. Newman's factored table gives
the following relative class numbers
\cite[Table 2, p.~218]{Newman1970}:
\begin{equation}
 \begin{array}{c|c|c}
 p&h^+(\Q(\zeta_p))&h^-(\Q(\zeta_p))\\ \hline
 37&1&37\\
 59&1&3\cdot59\cdot233\\
 67&1&67\cdot12739.
 \end{array}
 \label{eq:irreg-class-table}
\end{equation}
The cofactors at \(59\) and \(67\) are nonzero modulo the corresponding
prime:
\[
 699\equiv50\pmod{59},\qquad
 12739\equiv9\pmod{67}.
\]
Thus each \(p\)-Sylow class group has order \(p\). The norm to the
trivial real class group, followed by restriction, is \(1+\overline{\phantom{x}}\);
it follows that the \(p\)-class group is entirely minus. Therefore
\begin{equation}
 r=r^-=1\qquad(p=37,59,67).
 \label{eq:irreg-ranks-three}
\end{equation}
The full class groups at \(59\) and \(67\) have the additional factors
shown in \eqref{eq:irreg-class-table}.

For each of these primes one may take the same quotient
\begin{equation}
 C_p:\ y^p=x(1-x),\qquad
 f_p(X:Y:Z)=(X^p/Z^p,XY/Z^2),
 \label{eq:irreg-111-quotient}
\end{equation}
corresponding to \((a,b,c)=(1,1,-2)\). Here \(\gamma=1/4\).
Exact modular exponentiation gives
\begin{equation}
 \begin{array}{c|c|c}
 p&p^2&4^{p-1}\bmod p^2\\ \hline
 37&1369&75=1+2p\\
 59&3481&945=1+16p\\
 67&4489&1341=1+20p.
 \end{array}
 \label{eq:irreg-wild-table}
\end{equation}
Hence all three quotients are wild.

\begin{corollary}\label{cor:irreg-37-59-67}
For \(p=37,59,67\), there is one fixed set \(S_p\subset F_p(\Q_p)\)
containing the three rational cusps, contained in their three
coordinate residue discs, and satisfying
\[
 3\leq |S_p|\leq p,
\]
which contains the \(p\)-local centre of every compact
\(p\)-Selmer datum on \(F_p\).

For \(p=37\), one may use \(C_{37}:y^{37}=x(1-x)\), of genus \(18\),
to construct \(S_{37}\). In particular, the same \(S_{37}\) works for
every quotient section retaining the proper compact \(T_{37}J_{37}\)
projection and having the stipulated full local representatives.
\end{corollary}

\begin{proof}
Equation \eqref{eq:irreg-ranks-three} gives
\(r<(p+3)/8\), and the rank first-case criterion becomes \(10<p\).
Apply Theorem~\ref{thm:irreg-fixed} and
Corollary~\ref{cor:irreg-fixed-discs}, using the wild quotients in
\eqref{eq:irreg-111-quotient}. Their genus is \((p-1)/2\).
\end{proof}

The cardinality conclusions are upper bounds. They do not assert
\(S_p=\{Q_0,Q_1,Q_\infty\}\).

\subsubsection{Higher exponents}

\begin{corollary}\label{cor:irreg-higher-fixed}
Let \(n=pm\) be odd and assume \eqref{eq:irreg-Mc-bound}. Put
\(d=\gcd(m,p-1)\). There is one fixed finite set
\(S_{n,p}\subset F_n(\Q_p)\) containing every \(\Q_p\)-rational
coordinate boundary point, such that
\begin{equation}
 |S_{n,p}|\leq(p-3)d^2+3d.
 \label{eq:irreg-higher-bound}
\end{equation}
It contains the \(p\)-local centre of every datum consisting of
\[
 t\in H^1(\Q,T_pJ_n),\qquad R_v\in F_n(\Q_v),
 \qquad
 \loc_v(t)=\kappa_{p,v}([R_v-Q_0])
 \quad(v<\infty).
\]
If \(d=1\), the bound is \(p\). In particular \(d=1\) whenever
\(p\) is the least prime divisor of \(n\).
\end{corollary}

\begin{proof}
Let
\[
 \phi_m:F_n\longrightarrow F_p,\qquad
 (X:Y:Z)\longmapsto(X^m:Y^m:Z^m),
\]
and take \(S_{n,p}=\phi_m^{-1}(S_p)\) on \(\Q_p\)-points.
Functoriality on Jacobians pushes the given global Tate class and
the same local representatives to a compact \(p\)-Selmer datum on
\(F_p\). Hence every permitted centre belongs to this inverse image.

The roots of unity in \(\Q_p\) have order dividing \(p-1\), so
\(|\mu_m(\Q_p)|=d\). Above an ordinary point, normalize the nonzero
third coordinate to \(1\). Each of the other two coordinates has
at most \(d\) possible \(m\)-th roots, giving at most \(d^2\) points
in the fibre. Any such choices satisfy the exponent-\(n\) equation
automatically. Above \(Q_0\) or \(Q_1\), only one root choice remains,
giving exactly \(d\) points. Above \(Q_\infty\), the equation for
the coordinate ratio is \(u^m=-1\); since \(m\) is odd, its solutions
are \(-\mu_m(\Q_p)\), again exactly \(d\).

Every rational coordinate boundary point of \(F_n\) maps to one of
the three rational coordinate cusps of \(F_p\). Since those three
belong to \(S_p\), the inverse image contains them all, and
\[
 |S_{n,p}|
 \leq(|S_p|-3)d^2+3d
 \leq(p-3)d^2+3d.
\]
If \(p\) is the least prime divisor of \(n\), every prime dividing
\(m\) is at least \(p\), whereas every prime dividing \(p-1\) is
smaller than \(p\). Thus \(\gcd(m,p-1)=1\).
\end{proof}


The fixed-set theorem and its higher-exponent consequence prove
Theorem~\ref{thm:finite-local-centres-intro}.


\section{The quintic quotient and its Jacobian}

We construct the Jacobian quotient that supplies the rank-zero
target for finite descent. Set
\[
X_5:\ X^5+Y^5+Z^5=0,
\qquad C:\ v^2+v=u^5,
\qquad A=\Jac(C).
\]
Here $C$ means its smooth proper model; it has genus two and one point
$\infty$ above infinity. The plane quintic has genus six.

\begin{proposition}\label{prop:isogeny}
There is a rational isogeny $\Jac(X_5)\sim A^3$.
\end{proposition}
\begin{proof}
The three rational maps
\begin{equation}
f_Z=\left(-\frac{XY}{Z^2},\frac{X^5}{Z^5}\right),\quad
f_X=\left(-\frac{YZ}{X^2},\frac{Y^5}{X^5}\right),\quad
f_Y=\left(-\frac{ZX}{Y^2},\frac{Z^5}{Y^5}\right)
\label{eq:maps}
\end{equation}
land on $C$, since the Fermat equation makes $v(v+1)=u^5$.
They extend to morphisms between smooth proper curves and are
nonconstant. On $Z=1$ write $x=X/Z,y=Y/Z$, and put
$\eta=dx/y^4$. Adjunction gives the regular-differential basis
\[
\eta, x\eta, y\eta, x^2\eta, xy\eta, y^2\eta
\]
on $X_5$. A basis on $C$ is
$\omega_0=du/(2v+1)$ and $\omega_1=u\,du/(2v+1)$.
Using $dy=-x^4y^{-4}dx$, direct differentiation gives
\begin{equation}
\begin{array}{c|cc}
 &f^*\omega_0&f^*\omega_1\\\hline
f_Z&\eta&-xy\eta\\
f_X&x^2\eta&-y\eta\\
f_Y&y^2\eta&-x\eta.
\end{array}
\label{eq:differentials}
\end{equation}
For example $d(-xy)=(2x^5+1)dx/y^4$, proving the first row.
The other two rows follow by the same substitution, or by cyclic
permutation. Thus the pullback on cotangent spaces of
$(f_{Z,*},f_{X,*},f_{Y,*}):\Jac(X_5)\to A^3$ is an isomorphism.
A homomorphism of abelian varieties of equal dimension with invertible
differential has finite kernel and is surjective. It is therefore an
isogeny, defined over $\Q$. The determinant of the displayed
differential change of basis is not being used as an isogeny degree.
\end{proof}

\begin{proposition}[The integral quintic isogeny]
\label{prop:quintic-five-isogeny}
For the maps in \eqref{eq:maps}, put
\[
 \Phi=(f_{Z,*},f_{X,*},f_{Y,*}):\Jac(X_5)\longrightarrow A^3,
 \qquad
 \Psi=f_Z^*+f_X^*+f_Y^*:A^3\longrightarrow\Jac(X_5).
\]
Then
\[
 \Phi\Psi=[5]_{A^3},\qquad \Psi\Phi=[5]_{\Jac(X_5)}.
\]
In particular the kernels are killed by five, and
$T_3\Jac(X_5)\simeq(T_3A)^{\oplus3}$ as integral $G_\Q$-modules.
\end{proposition}
\begin{proof}
Each $f_i$ is the quotient by an order-five diagonal subgroup $H_i$.
For example the first map is invariant under
$(x,y)\mapsto(\zeta_5x,\zeta_5^{-1}y)$; its function field has
degree five because $x^5=v$ and $y=-u/x$.
The other two cases follow by cyclic permutation.
Consequently $f_{i,*}f_i^*=[5]$ and
$f_i^*f_{i,*}=\sum_{h\in H_i}h$.

The characters of the holomorphic differentials
$\omega_{a,b}=x^{a-1}y^{b-5}dx$ have $a,b\ge1$ and $a+b\le4$.
The three invariant pairs, in the order $Z,X,Y$, are
\[
 \{(1,1),(2,2)\},\qquad
 \{(1,2),(3,1)\},\qquad
 \{(1,3),(2,1)\}.
\]
They are disjoint and exhaust the six characters, as also follows
from \eqref{eq:differentials}. The $H_i$-trace is therefore five
on its own invariant subspace and zero on the other two. This
proves that $f_{i,*}f_j^*=0$ for $i\ne j$ and that the sum of
the three trace operators is $[5]$. These are equalities of
homomorphisms: a homomorphism of abelian varieties in characteristic
zero with zero differential is zero. More explicitly, the product
of distinct trace operators is zero; $f_{j,*}$ is surjective and
$f_i^*$ has finite kernel, so the mixed composition itself is zero.
The two displayed compositions
follow. Since five is a unit in $\Z_3$, the integral inverse of
$T_3\Phi$ is $5^{-1}T_3\Psi$.
\end{proof}

\begin{lemma}[An integral splitting of the quintic boundary]
\label{lem:quintic-boundary-split}
Let $D_5=X_5\setminus U_5$ and
$B_5=\Z_3(1)[D_5(\bar\Q)]/\Z_3(1)\sum_{Q\in D_5}e_Q$.
The sequence
\[
 0\longrightarrow B_5\longrightarrow
 \pi_1(U_{5,\bar\Q})^{\ab,(3)}
 \longrightarrow T_3\Jac(X_5)\longrightarrow0
\]
splits as a sequence of integral $G_\Q$-modules.
\end{lemma}
\begin{proof}
The generalized Jacobian extension by the boundary torus is
classified by the map from degree-zero divisors supported on $D_5$
to $\Jac(X_5)$, using the principal polarization.
For coordinate cusps $P,Q$, the quotient of their tangent-line
equations has divisor $5(Q-P)$. Thus this classifying map is killed
by five. Pullback of the extension by $[5]$ is trivial. Its splitting
is unique because a homomorphism from an abelian variety to a torus
is zero, and hence the splitting descends to $\Q$.
The resulting lift of $[5]$, divided by five on the $3$-adic Tate
module, gives the asserted splitting.
\end{proof}

These integral statements strengthen Proposition~\ref{prop:isogeny}.
For the pro-$3$ image argument below, a rational isogeny already
suffices: having pro-$3$ image is a topological property of the
rational Tate representation and is independent of a stable lattice.


\section{A central-value certificate}

\citet[p.~439, following Table 3]{Darmon2000} identifies $A$ with a CM
abelian surface attached to characters of $K=\Q(\zeta_5)$ of conductor
$(1-\zeta_5)^2$. He also identifies its base change to
$F=\Q(\sqrt5)$ with the two conjugate Hilbert forms of parallel weight
two and level $(\sqrt5)^3$. The conductor--discriminant formula for
induction gives the rational conductor $5^3\cdot5^2=3125$.
The associated completed Hasse--Weil function is, up to a nonzero
constant,
\begin{equation}
\Lambda(s)=3125^{s/2}(2\pi)^{-2s}\Gamma(s)^2 L(A/\Q,s),
\quad \Lambda(s)=W\Lambda(2-s),\quad W\in\{1,-1\}.
\label{eq:fe}
\end{equation}
The analytic continuation, entire property, and functional equation
here are the CM Hecke-character inputs; they are not inferred from
numerical data.

\begin{proposition}\label{prop:lvalue}
The root number in \eqref{eq:fe} is $+1$, and the enclosed central value is
\begin{equation}
\begin{gathered}
L(A/\Q,1)=0.7183136284205647705166034780506424+\varepsilon,\\
|\varepsilon|<6.57\cdot10^{-35}.
\end{gathered}
\label{eq:lvalue}
\end{equation}
In particular $L(A/\Q,1)>0$.
\end{proposition}

\begin{proof}
Write $L(A/\Q,s)=\sum_{n\geq1}a_nn^{-s}$ and
$c=4\pi/3125^{1/4}$. The inverse Mellin transform of $\Gamma(s)^2$
gives
\[
\theta(t)=2\sum_{n\geq1}a_nK_0(c\sqrt{nt}),\qquad
\Lambda(s)=\int_0^\infty\theta(t)t^{s-1}\,dt.
\]
The functional equation implies $\theta(t)=Wt^{-2}\theta(1/t)$.
Therefore $\theta(1)>0$ proves $W=1$; after this is known, integration
and $\int_z^\infty uK_0(u)\,du=zK_1(z)$ give
\begin{equation}
L(A/\Q,1)=2c\sum_{n\geq1}\frac{a_n}{\sqrt n}K_1(c\sqrt n).
\label{eq:central}
\end{equation}

All coefficients through $N=4096$ are computed exactly. At $p\ne5$
put $N_r=\#C(\F_{p^r})$, $t=p+1-N_1$, and
$b=(N_2-p^2-1+t^2)/2$. The Euler polynomial is
\begin{equation}
1-tT+bT^2-ptT^3+p^2T^4.
\end{equation}
Only $N_1$ is required when $p^2>N$, and $N_2$ is counted directly in
the quadratic finite field otherwise. The Euler factor at $5$ is one
because the ramified CM character has no inertia invariants. Inverting
these integer polynomials and using multiplicativity supplies the
coefficient list. For independent small checks the polynomials at
$3$ and $11$ are $1+9T^4$ and $1+T-9T^2+11T^3+121T^4$.

The Weil bounds give $|a_n|\leq d_4(n)\sqrt n\leq n^{7/2}$.
For real $z>0$, the integral formula for $K_\nu$ and
$\cosh u\geq1+u^2/2$ give
\begin{equation}
K_0(z)\leq e^{-z}\sqrt{\frac\pi{2z}},\qquad
K_1(z)\leq e^{-z}\sqrt{\frac\pi{2z}}e^{1/(2z)}.
\label{eq:kbound}
\end{equation}
In the second inequality, integrate the resulting Gaussian multiplied
by $\cosh u$. Since $x^a e^{-c\sqrt x}$ decreases for
$c\sqrt x>2a$, the omitted sums are bounded by integrals. With
$z_0=c\sqrt N$ and $B=\sqrt{\pi/(2c)}$, the two bounds are
\begin{align*}
|\theta(1)-\theta_N(1)|
&\leq4B c^{-17/2}\Gamma(17/2,z_0),\\
\left|L(A/\Q,1)-2c\sum_{n\leq N}\frac{a_n}{\sqrt n}K_1(c\sqrt n)\right|
&\leq4cB e^{1/(2z_0)}c^{-15/2}\Gamma(15/2,z_0).
\end{align*}
The monotonicity condition holds already at $N=4096$.
The companion program evaluates the finite sums and these bounds with
192-bit Arb interval arithmetic. It certifies
\[
\theta(1)\in[0.3300364476\mathbin\pm3.37\cdot10^{-11}],
\]
with omitted-tail bounds less than $1.674\cdot10^{-33}$ for $\theta$
and $4.374\cdot10^{-35}$ for the central sum. The first enclosure
determines the root sign, and the second gives \eqref{eq:lvalue}.
The intentionally loose bound on $d_4(n)$ is sufficient; no unproved
denominator bound or floating-point zero test is used.
\end{proof}

\subsection{An exact rational positivity check}

A shorter certificate of positivity uses only the first ten
Dirichlet coefficients and rational arithmetic. The sign can be
obtained without an interval evaluation of $\theta(1)$:
$C$ is isomorphic to $y^2=x^5+256$ via
$x=4u$, $y=16(2v+1)$. In
\citet[Proposition 2.2(4)--(6)]{StollYang2003}, the parameter $256=2^8$
has conductor exponent zero at $2$ and
$(256^4-1)/5\equiv4\pmod5$. Their formula gives conductor
$(1-\zeta_5)^2$ and root number $+1$.
Their Lemma 2.1 also verifies the trivial-central-character
normalization of the induced Hilbert form.

\begin{proposition}\label{prop:rationalpositive}
There is an exact rational certificate that $L(A/\Q,1)>1/20$.
\end{proposition}
\begin{proof}
For $p=2,3,7$ and $q=p,p^2$, the map $x\mapsto x^5$ is bijective
on $\F_q$. Thus $\#C(\F_q)=q+1$ and $P_p(T)=1+p^2T^4$.
The bad Euler factor at $5$ is one. Multiplicativity now gives
$a_1=1$ and $a_2=\cdots=a_{10}=0$.
Writing $z_n=c\sqrt n$, equation~\eqref{eq:central} becomes
\[
L(A/\Q,1)=2\sum_{n\geq1}\frac{a_n}{n}z_nK_1(z_n).
\]
The integral representation of $K_1$ and
\[
1+x/2-x^2/8\leq\sqrt{1+x}\leq1+x/2\qquad(x\geq0)
\]
give the global bounds
\[
\begin{split}
K_1(z)&\geq\sqrt{\frac\pi{2z}}e^{-z}
             \left(1+\frac3{8z}-\frac{15}{128z^2}\right),\\
K_1(z)&\leq\sqrt{\frac\pi{2z}}e^{-z}
             \left(1+\frac3{8z}\right).
\end{split}
\]
Indeed the remaining integral is the expectation of
$\sqrt{1+t/(2z)}$ against the density
$e^{-t}t^{1/2}/\Gamma(3/2)$, whose first two moments are $3/2$ and
$15/4$. These inequalities are valid for all positive $z$.

Use $|a_n|\leq d_4(n)\sqrt n$. The rational-arithmetic program
bounds the term at one below and the terms from $11$ through $400$
above using the displayed inequalities. It encloses $\pi$ by
Machin's formula and alternating arctangent sums, each square root
by rational numbers checked by squaring, and each exponential by
positive Taylor sums and a geometric remainder. Thus every
rounding direction is an exact rational inequality.
For the infinite tail note that $d_4(n)\leq n^2$; this follows
multiplicatively from $\binom{e+3}{3}\leq p^{2e}$.
The inequalities $c>5/3$ and $\sqrt{\pi c/2}<2$ give an upper
bound $8n^2e^{-(5/3)\sqrt n}$ for each absolute summand with $n>400$.
This decreases there, so its sum is at most
\[
16e^{-100/3}\sum_{j=0}^{5}
 \frac{5!}{(5-j)!}\,20^{5-j}\left(\frac35\right)^{j+1}.
\]
This last exponential is enclosed by the same positive Taylor sums.
The resulting exact inequalities, retained as executable assertions,
are
\[
2z_1K_1(z_1)>\frac7{10},\qquad
2\sum_{n\geq11}\frac{d_4(n)}{\sqrt n}z_nK_1(z_n)<\frac{13}{20}.
\]
They imply $L(A/\Q,1)>7/10-13/20=1/20$.
\end{proof}


\section{Arithmetic consequences and Fermat Selmer theorems}

\begin{proposition}\label{prop:arithmetic}
The groups $A(\Q)$ and $\Sha(\Q,A)$ are finite. More precisely
$A(\Q)\cong\Z/5\Z$ and
\begin{equation}
C(\Q)=\{\infty,(0,0),(0,-1)\}.
\label{eq:cpoints}
\end{equation}
Consequently $\Jac(X_5)(\Q)$ and $\Sha(\Q,\Jac(X_5))$ are finite.
\end{proposition}

\begin{proof}
We use the published Selmer theorem
\citet[Theorem B, p.~593; proof \S2.10.5, p.~641]{Nekovar2012}.
Its applicable additional hypothesis is
that the totally real base field has odd degree or that a finite local
component is not principal series. The relevant Hilbert forms over
$F=\Q(\sqrt5)$ have trivial central character and conductor exponent
three at $(\sqrt5)$. A principal series with trivial central character
is induced from a character and its inverse, and its conductor
exponent is twice that character's conductor exponent. Thus exponent
three excludes principal series and the required alternative holds.

The CM representation induced from $G_K$ is unchanged by twisting by
the quadratic character of $F/\Q$, because that character restricts
trivially to $G_K$. Artin formalism therefore gives
$L(A/F,s)=L(A/\Q,s)^2$. Proposition~\ref{prop:lvalue} makes each Hilbert
factor nonzero at the central point. Nekov\'a\v r's theorem gives zero
Bloch--Kato Selmer group with vector-space coefficients at every
coefficient prime, and zero Selmer group with divisible-torsion
coefficients at all but finitely many such primes. The classical
Selmer exact sequence, as used in his \S2.10.5, therefore gives finite
$A(F)$, finite primary parts of $\Sha(F,A)$ at every prime, and zero
primary parts outside a finite set. Thus the whole group $\Sha(F,A)$
is finite. Rational Mordell--Weil finiteness follows
by inclusion. Restriction on $\Sha$ has kernel contained in
$H^1(\Gal(F/\Q),A(F))$, a finite group, so it gives finiteness over $\Q$.
For comparison, \citet[Theorem 4.3.2, p.~588]{Zhang2010} states the
corresponding rank-zero conclusion while attributing it to work of
Tian--Zhang then listed as in preparation. The present dependency is
on Nekov\'a\v r's published theorem and proof, not that bibliographic
promise.

The exact counts in \eqref{eq:central}'s computation give
$\#A(\F_3)=10$ and $\#A(\F_{11})=125$. Prime-to-residue-characteristic
torsion injects into these good reductions, so the finite group
$A(\Q)$ has order dividing five. If $P_0=(0,0)$, then
$\operatorname{div}(v)=5P_0-5\infty$, and
$[P_0-\infty]$ is nonzero by injectivity of the Abel--Jacobi map.
It has order five and generates $A(\Q)$.
The hyperelliptic involution exchanges $P_0$ and $P_1=(0,-1)$, giving
$[P_1-\infty]=-[P_0-\infty]$. The two remaining nonzero group elements
cannot come from a point of $C$: an equality
$[P-\infty]=2[P_i-\infty]$ would give $P+\infty\sim2P_i$.
The divisor $2P_i$ is nonspecial of degree two, because $P_i$ is not
a Weierstrass point, so its linear system has just one effective
divisor. This would force $P+\infty=2P_i$, a contradiction.
This proves the point list \eqref{eq:cpoints}.
Finally Proposition~\ref{prop:isogeny} and isogeny invariance of rank
and of finiteness of $\Sha$ give the Jacobian claims.
\end{proof}

\subsection{The exponent-eleven quotient}
\label{subsec:eleven}

Put $C_{11}:v^2+v=u^{11}$ and $A_{11}=\Jac(C_{11})$.
This is the classical factor $J_{1,1,9}$: the substitution
$(x,y)=(-v,-u)$ identifies its function field with
$y^{11}=x(1-x)$. The map
\[
f_{11}:X_{11}\longrightarrow C_{11},\qquad
[X:Y:Z]\longmapsto(-XY/Z^2,X^{11}/Z^{11})
\]
is nonconstant and extends over the smooth proper models, by the
same direct identity as in \eqref{eq:maps}.
\citet[\S4, table on p.~219]{GrossRohrlich1978} records the finiteness
of $A_{11}(\Q)$ proved by Faddeev. Finiteness of this Mordell--Weil
group by itself does not imply the required condition on $\Sha$.
We establish that condition using the central value and the published
Hilbert-modular theorem.

\begin{proposition}\label{prop:eleven}
The groups $A_{11}(\Q)$ and $\Sha(\Q,A_{11})$ are finite,
$A_{11}(\Q)\simeq\Z/11\Z$, and
\[
C_{11}(\Q)=\{\infty,(0,0),(0,-1)\}.
\]
The central value satisfies
\begin{equation}
1.78<L(A_{11}/\Q,1)<1.80.
\label{eq:elevenvalue}
\end{equation}
\end{proposition}

\begin{proof}
We first fix the analytic normalization and certify nonvanishing.
\subsubsection{The analytic inputs}

Write $K=\Q(\zeta_{11})$, $F=K^+$, and
$\lambda=1-\zeta_{11}$. The Jacobi-sum description identifies
$L(A_{11}/\Q,s)$ with one Hecke $L$-function on $K$.
For a prime ideal $\mathfrak p\nmid11$, its value is
\[
j(\mathfrak p)=-\sum_{x\ne0,1}\chi_{\mathfrak p}(x)
                                      \chi_{\mathfrak p}(1-x),
\qquad |\sigma j(\mathfrak p)|=\sqrt{N\mathfrak p}
\]
for every complex embedding $\sigma$.
See \citet[\S1 and \S3]{GrossRohrlich1978} and
\citet[\S2.1--2.2, pp.~5--6]{ShuExpanded2021}.
Here and below the precise proposition numbering for Shu is that of
the expanded arXiv version; the root-number results are published in
\citet{Shu2021}.

The character has conductor $\lambda^2$ and the functional-equation
sign is $+1$. These inputs can be checked explicitly in
\citet[Proposition 3.12(2), p.~16; Theorem 3.13, pp.~17--18]{ShuExpanded2021}.
For $(r,s,t,\delta)=(1,1,9,1)$ put $q=(-2)^9$. Then
\[
v_{11}(q^{10}-1)=1,\qquad
 d=(q^{10}-1)/11\equiv
 9(2^{10}-1)/11\equiv1\pmod{11}.
\]
The conductor alternative is therefore $\lambda^2$, and the root-number
formula gives
\[
W=\left(\frac2{11}\right)
  \left[-\left(\frac{-2\cdot1\cdot1\cdot9\cdot d}{11}\right)\right]
  =(-1)(-1)=1.
\]
Since $|\operatorname{disc}K|=11^9$ and $N\lambda^2=11^2$, the rational
conductor is $11^{11}$. Thus, on putting
\[
B=\frac{\sqrt{11^{11}}}{(2\pi)^5},
\]
the completed function is
\begin{equation}
\Lambda_{11}(s)=B^s\Gamma(s)^5L(A_{11}/\Q,s),\qquad
\Lambda_{11}(s)=\Lambda_{11}(2-s).
\label{eq:elevenfe}
\end{equation}
Entire continuation and the functional equation are the cited Hecke
inputs. In particular the numerical calculation below is not being
used to guess a conductor or a sign.

\subsubsection{A finite central-value certificate}

Let $G_5$ be the density on $(0,\infty)$ of the product of five
independent exponential random variables of mean one, and put
$Q_5(x)=\int_x^\infty G_5(u)\,du$. Its Mellin transform is
$\Gamma(s)^5$. Write $L(A_{11}/\Q,s)=\sum a_nn^{-s}$.
Mellin inversion and \eqref{eq:elevenfe}, split at $t=1$, give
\begin{equation}
L(A_{11}/\Q,1)=2\sum_{n\geq1}\frac{a_n}{n}Q_5(n/B),
\qquad
Q_5(x)=\frac1{2\pi i}\int_{(c)}
             \frac{\Gamma(1+z)^5}{z}x^{-z}\,dz\quad(c>0).
\label{eq:elevenafe}
\end{equation}
For completeness the inverse Mellin series is
$\Theta(t)=\sum a_nG_5(nt/B)$, whose Mellin transform is
$\Lambda_{11}(s)$. The functional equation gives
$\Theta(t)=t^{-2}\Theta(1/t)$, so
$\Lambda_{11}(1)=2\int_1^\infty\Theta(t)\,dt$.
Substitution $u=nt/B$ yields the first formula in
\eqref{eq:elevenafe}. The positive-kernel bounds below justify absolute
convergence and the interchange of sums and integrals.

Let $b_K(n)$ count the integral ideals of $K$ of norm $n$.
The Hecke description and the triangle inequality imply
\[
|a_n|\leq \sqrt n\,b_K(n),\qquad
\sum_{n\geq1}b_K(n)n^{-\sigma}=\zeta_K(\sigma)\quad(\sigma>1).
\]
Markov's inequality for the product of the five exponential variables
uses its $k$-th moment $\Gamma(k+1)^5$ and gives
$Q_5(x)\leq\Gamma(k+1)^5x^{-k}$ for $k>0$.
Consequently, whenever $k+1/2>\sigma>1$,
\begin{equation}
2\sum_{n>N}\frac{|a_n|}{n}Q_5(n/B)
 \leq 2\Gamma(k+1)^5B^k
       N^{\sigma-k-1/2}\zeta_K(\sigma).
\label{eq:eleventail}
\end{equation}
This follows term by term from
$n^{-k-1/2}\leq N^{\sigma-k-1/2}n^{-\sigma}$ for $n>N$.

The companion program computes every required $a_n$ exactly through
$N=100000$. Here are all the ingredients of that finite calculation.
For a rational prime $\ell\ne11$, put
$f=\operatorname{ord}_{11}(\ell)$. If $f$ is even, the Euler polynomial
is
\[
P_\ell(T)=(1+\ell^{f/2}T^f)^{10/f};
\]
this is \citet[Lemma 1.1]{GrossRohrlich1978}.
If $f=5$, it is $1-s_\ell T^5+\ell^5T^{10}$, where a direct count
of $C_{11}(\F_{\ell^5})$ gives
$s_\ell=(\ell^5+1-\#C_{11}(\F_{\ell^5}))/5$.
Only $\ell=3,5$ occur in this case below the cutoff. The resulting
polynomials are
\[
P_3(T)=1-24T^5+243T^{10},\qquad
P_5(T)=1-90T^5+3125T^{10}.
\]
If $f=1$, choose a primitive root of $\F_\ell^\times$ and count the
residue classes of
$\log(x)+\log(1-x)\pmod{11}$. If these counts are $c_0,\ldots,c_{10}$,
then $j=-\sum c_i\zeta_{11}^i$.
The Euler polynomial is $\prod_{a=1}^{10}(1-\sigma_a(j)T)$;
its required coefficients follow from Newton's identities using
integer arithmetic in $\Z[\zeta_{11}]$.
The factor at $11$ is one. When $\ell^f>N$, its entire contribution
through the cutoff is the constant term. Inversion of the computed
local polynomials and multiplicativity produce 1077 nonzero
coefficients, beginning with $a_1=1$, $a_{23}=21$, $a_{67}=-23$,
$a_{89}=21$, and $a_{243}=24$.

The finite sum in \eqref{eq:elevenafe} is evaluated by one contour
integral on $\Re z=2$:
\begin{equation}
L_N=\frac2\pi\Re\int_0^\infty
 \frac{\Gamma(3+it)^5}{2+it}
 B^{2+it}\sum_{n\leq N}a_nn^{-3-it}\,dt.
\label{eq:elevencontour}
\end{equation}
The integral from zero to $20$ is enclosed with 96-bit Arb arithmetic
on twenty unit intervals. To bound the remainder, set
$S_N=\sum_{n\leq N}|a_n|n^{-3}$. The gamma integral gives
$|\Gamma(3+m+it)|\leq\Gamma(3+m)$; repeated use of its recurrence
therefore yields
$|\Gamma(3+it)|\leq\Gamma(3+m)|t|^{-m}$.
For every positive integer $m$, the contribution from $t\geq T$ is
at most
\begin{equation}
\frac2\pi B^2S_N\frac{\Gamma(3+m)^5}{5m\,T^{5m}}.
\label{eq:elevenvertical}
\end{equation}
With $T=20,m=17$ this is less than $1.538\cdot10^{-24}$.
The needed Dedekind-zeta value is enclosed using
\[
\zeta_K(\sigma)=\zeta(\sigma)
 \prod_{\substack{\chi\bmod11\\\chi\ne1}}L(\sigma,\chi),\qquad
L(\sigma,\chi)=11^{-\sigma}\sum_{a=1}^{10}
                           \chi(a)\zeta(\sigma,a/11).
\]
The separate Riemann-zeta factor includes the prime at $11$; replacing
it by the imprimitive principal Dirichlet factor would omit that prime.
Arb certifies $\zeta_K(11/10)<3.007268$.
The numerical enclosures are
\[
\begin{split}
L_N&=1.789861079603+\varepsilon_N,
                    \qquad |\varepsilon_N|<1.54\cdot10^{-13},\\
|L(A_{11}/\Q,1)-L_N|&<0.004239358.
\end{split}
\]
The latter uses \eqref{eq:eleventail} with $k=4,\sigma=11/10$ and
\eqref{eq:elevenvertical}. These bounds prove \eqref{eq:elevenvalue}.

\subsubsection{The arithmetic conclusion}

Automorphic induction of the Jacobi character from $K$ to $F$ gives
the parallel-weight-two Hilbert form attached to $A_{11}/F$;
this is the realization in \citet[\S2.2]{ShuExpanded2021}.
Its central character is trivial: the inducing character restricts to the
quadratic character of $K/F$ on the real idele group, so the central
character of its automorphic induction is the square of that character
\citep[\S3.4, p.~15]{ShuExpanded2021}. One can also check this normalization
is that the real multiplication $F\subset\operatorname{End}_F^0(A_{11})$
is fixed by Rosati, so each two-dimensional coefficient component of
the polarized Tate representation has determinant the cyclotomic
character. The corresponding weight-two form therefore has trivial
finite central character. Each coefficient conjugate has the same
Hecke $L$-function after reindexing the prime ideals of $K$ by the
corresponding automorphism; Jacobi sums satisfy
$j(\tau\mathfrak p)=\tau j(\mathfrak p)$.
Thus \eqref{eq:elevenvalue} establishes nonvanishing for every
coefficient component of $A_{11}/F$.

Now $[F:\Q]=5$ is odd. Apply
\citet[Theorems B and B$'$, alternative (a), p.~593]{Nekovar2012}.
The vector-space Selmer groups vanish at every coefficient prime,
and the divisible-torsion Selmer groups vanish outside a finite set.
The Selmer exact sequence gives finite $A_{11}(F)$ and finite
$\Sha(F,A_{11})$, exactly as in the proof of
Proposition~\ref{prop:arithmetic}. Restriction to $F$ has a finite
kernel on $\Sha$, contained in
$H^1(\Gal(F/\Q),A_{11}(F))$, so the two groups over $\Q$ are finite.
This proves the additional arithmetic input without a converse BSD
assumption. The torsion description of
\citet[Theorem 1.1]{GrossRohrlich1978} then gives
$A_{11}(\Q)\simeq\Z/11\Z$.

To recover the point set, put $P_0=(0,0)$ and $P_1=(0,-1)$.
The divisor of $v$ is $11P_0-11\infty$, and
$D=[P_0-\infty]$ generates $A_{11}(\Q)$.
The hyperelliptic involution gives $[P_1-\infty]=-D$.
At either $P_i$ the differentials
\[
\frac{du}{2v+1},\quad\frac{u\,du}{2v+1},\quad\ldots,
\quad\frac{u^4\,du}{2v+1}
\]
have vanishing orders $0,1,2,3,4$. Riemann--Roch therefore gives
$h^0(C_{11},dP_i)=1$ for $1\leq d\leq5$.
If a point $P$ represented $dD$, or $-dD$, with $2\leq d\leq5$,
then $P+(d-1)\infty\sim dP_i$ for one of the two $i$.
Its unique effective divisor is $dP_i$, a contradiction.
The only possible point classes are consequently $0,D,-D$, giving
the asserted three points. This completes the proof of
Proposition~\ref{prop:eleven}.
\end{proof}

\begin{remark}\label{rem:eleven-factor-only}
Only the factor $A_{11}$ is used. The whole Fermat Jacobian
$\Jac(X_{11})$ has positive Mordell--Weil rank: its rank is six by
\citet[Introduction and \S4]{GrossRohrlich1978}.
The target criterion in Stoll's theorem allows this distinction.
\end{remark}

\subsection{Compact Selmer localization for the Fermat curve of exponent eleven}
\label{subsec:eleven-compact-localization}

Put
\[
 A=\operatorname{Jac}(y^{11}=x(1-x)),\qquad
 B=\operatorname{Jac}(y^{11}=x(1-x)^2),\qquad
 J=\operatorname{Jac}(X^{11}+Y^{11}+Z^{11}=0),
\]
where smooth projective models are understood.  Both $A$ and $B$
have dimension five.  The rank-one assertion for $B(\mathbb Q)$ is
classical: it is the $(1,2,8)$ entry of the table in
\citet[\S4, p.~219]{GrossRohrlich1978}.  Their introduction also
records $\operatorname{rank}J(\mathbb Q)=6$.  The additional input
needed for compact Selmer localization is control of the full
Tate--Shafarevich group, rather than the rank alone.

\begin{proposition}[The rank-one cyclic factor at eleven]
\label{prop:eleven-rankone-arithmetic}
The abelian variety $B$ has rational Mordell--Weil rank one,
$B(\mathbb Q)_{\mathrm{tors}}\simeq\mathbb Z/11\mathbb Z$, and finite
$\Sha(B/\mathbb Q)$.  In the standard weight-one motivic normalization,
\[
 3.114<L'(B/\mathbb Q,1)<3.137.
\]
\end{proposition}
\begin{proof}
Write $K=\mathbb Q(\zeta_{11})$, $F=K^+$, and
$\lambda=1-\zeta_{11}$.  The Jacobi-sum character for the triple
$(1,2,8)$ is
\[
 j(\mathfrak p)=-\sum_{x\ne0,1}
       \chi_{\mathfrak p}(x)\chi_{\mathfrak p}(1-x)^2,
 \qquad\mathfrak p\nmid11.
\]
Its Hecke $L$-function is $L(B/\mathbb Q,s)$, and every complex
embedding of $j(\mathfrak p)$ has absolute value
$\sqrt{N\mathfrak p}$; see \citet[\S1, pp.~202--203]{GrossRohrlich1978}
and \citet[\S\S2.1--2.2, author pp.~5--6]{ShuExpanded2021}.

The conductor and sign can be checked without numerical inference.
For $r=1,s=2,t=8$ put $u=r^rs^s(t-11)^t$.  Direct arithmetic gives
\[
 u\equiv108\pmod{121},\qquad
 d=\frac{u^{10}-1}{11}\equiv10\pmod{11}.
\]
Thus \citet[Proposition~3.12(2), author p.~16]{ShuExpanded2021}
gives conductor $\lambda^2$.  The root-number formula in
\citet[Theorem~3.13, author pp.~17--18]{ShuExpanded2021} gives
\[
 W=\left(\frac2{11}\right)
       \left[-\left(\frac{-2rstd}{11}\right)\right]=-1,
\]
because $-2rstd\equiv10\pmod{11}$ and both displayed Legendre
symbols are $-1$.  The rational conductor is therefore
$11^9 N\lambda^2=11^{11}$.  On setting
\[
 a=\frac{\sqrt{11^{11}}}{(2\pi)^5},
\]
the entire completed function satisfies
\[
 \Lambda(s)=a^s\Gamma(s)^5L(B/\mathbb Q,s),\qquad
 \Lambda(s)=-\Lambda(2-s).
\]

We give the complete finite nonvanishing calculation.  Let $Y$ be the
product of five independent exponential random variables of mean one,
and define
\[
 H_5(x)=\mathbb E\bigl(\log(Y/x)_+\bigr)
       =\frac1{2\pi i}\int_{(c)}
          \frac{\Gamma(1+z)^5}{z^2}x^{-z}\,dz\quad(c>0).
\]
Here the subscript $+$ means the positive part.  Mellin inversion
follows either by integrating the tail distribution of $Y$ or from
the Mellin transform of the positive logarithm.  If
$L(B/\mathbb Q,s)=\sum a_n n^{-s}$, shifting the contour across
zero and using the odd functional equation gives
\[
 L'(B/\mathbb Q,1)
       =2\sum_{n\ge1}\frac{a_n}{n}H_5(n/a).
\]
Indeed $a^z\Gamma(1+z)^5L(1+z)$ is odd in $z$; the residue after
division by $z^2$ is $L'(1)$, and the two vertical integrals are
equal.  The bounds below justify the absolute convergence and
contour operations.  Since $L(1)=0$, differentiating the completed
function introduces no extra logarithmic term at the center.

Let $b_K(n)$ denote the number of integral ideals of $K$ of norm $n$.
The Hecke expansion gives $|a_n|\le\sqrt n\,b_K(n)$.  For every
$k>0$, the inequality $\log(v)_+\le v^k/k$ gives
\[
 H_5(x)\le\frac{\Gamma(k+1)^5}{k x^k}.
\]
Consequently, if $k+1/2>\sigma>1$, the omitted coefficients beyond
$N$ contribute at most
\begin{equation}
 \frac{2\Gamma(k+1)^5}{k}\,a^k
       N^{\sigma-k-1/2}\zeta_K(\sigma).
 \label{eq:eleven-rankone-tail}
\end{equation}

All coefficients through $N=50000$ are obtained by exact finite-field
counts.  For a prime $p\ne11$ put $f=\operatorname{ord}_{11}(p)$.
For even $f$, \citet[Lemma~1.1, p.~203]{GrossRohrlich1978} gives
\[
 P_p(T)=(1+p^{f/2}T^f)^{10/f}.
\]
For $f=5$, enumerate the logarithms in $\mathbb F_{p^5}^{\times}$
and count the classes of $\log(x)+2\log(1-x)$ modulo eleven.
Only $p=3,5$ are needed at the cutoff; the resulting polynomials are
\[
 P_3(T)=1+9T^5+243T^{10},\qquad
 P_5(T)=1+75T^5+3125T^{10}.
\]
The trace in $K$ is five times the trace over the quadratic field
fixed by $p$-Frobenius, so division by five in this computation is
exact.  At $f=1$, the same enumeration is in $\mathbb F_p$.
If $c_0,\ldots,c_{10}$ are the counts, then
$j=-\sum c_i\zeta_{11}^i$; taking the product over its ten
conjugates, with Newton's identities in integer arithmetic, gives
the required local coefficients.  The bad Euler factor at eleven
is one.  Primes with $p^f>N$ have no nonconstant contribution below
the cutoff.  Euler inversion and multiplicativity give 573 nonzero
coefficients, starting with
\[
 a_1=1,\quad a_{23}=-12,\quad a_{67}=-23,\quad
 a_{89}=-1,\quad a_{243}=-9.
\]

The finite sum is enclosed by the single contour integral
\[
 L'_N=\frac2\pi\Re\int_0^\infty
       \frac{\Gamma(3+it)^5}{(2+it)^2}
       a^{2+it}\sum_{n\le N}a_n n^{-3-it}\,dt.
\]
Put $S_N=\sum_{n\le N}|a_n|n^{-3}$.  From the gamma integral and
its recurrence, $|\Gamma(3+it)|\le\Gamma(3+m)|t|^{-m}$.
The part $t\ge T$ is therefore at most
\begin{equation}
 \frac2\pi a^2S_N
    \frac{\Gamma(3+m)^5}{(5m+1)T^{5m+1}}.
 \label{eq:eleven-rankone-vertical}
\end{equation}
Arb arithmetic at 128 bits encloses the integral on $[0,20]$ by
twenty unit intervals.  The Dedekind-zeta bound is obtained from
\[
 \zeta_K(\sigma)=\zeta(\sigma)
       \prod_{\substack{\chi\bmod11\\\chi\ne1}}L(\sigma,\chi),
 \qquad
 L(\sigma,\chi)=11^{-\sigma}\sum_{a=1}^{10}
                   \chi(a)\zeta(\sigma,a/11),
\]
so that the factor at eleven is included.  The resulting certified
bounds are
\[
 \begin{aligned}
 \zeta_K(11/10)&<3.007268,\\
 |L'_N-3.12547064977|&<3.14\cdot10^{-12},\\
 \text{vertical remainder}&<7.6\cdot10^{-26},\\
 \text{coefficient remainder}&<0.011188.
 \end{aligned}
\]
The last two lines use $T=20,m=17$ and $k=4$, respectively.
These explicit inequalities prove the asserted interval.

We now check the arithmetic theorem's precise hypotheses.
Automorphic induction of $j$ from $K$ to $F$ is a cuspidal
parallel-weight-two Hilbert form $f$: its two infinity exponents
differ by one, so the character is not a norm character.  Its finite
central character is trivial, since the inducing unitary character
restricts to $\eta_{K/F}$ on $\mathbb A_F^\times$, and the determinant
of induction multiplies this by $\eta_{K/F}$ once more
\citep[\S3.4, author p.~15]{ShuExpanded2021}.  One has
\[
 L(f/F,s)=L(B/\mathbb Q,s).
\]
Every coefficient conjugate of $f$ has the same $L$-function, since
conjugating a Jacobi character reindexes the prime ideals of the
abelian field $K$ while preserving their norms.

The CM type is primitive by
\citet[Lemma~1.6, pp.~205--207]{GrossRohrlich1978}: at eleven the
exception involving a nontrivial cube root of unity in
$(\mathbb Z/11\mathbb Z)^\times$ is impossible.  Thus $B$ is
absolutely simple.  The real CM subfield $F$ acts on $B_F$ and is
fixed by Rosati.  Its five coefficient components are the five
two-dimensional inductions of the coefficient conjugates of $j$.
The Hecke field of $f$ is contained in $F$ and is exactly $F$.
For an explicit check, choose $5$ as primitive element of
$\mathbb F_{23}^{\times}$ and send it to $\zeta_{11}$ under the
multiplicative character.  The exact Jacobi-sum coefficient vector is
\[
 (c_0,\ldots,c_{10})=(3,3,2,1,4,2,1,0,0,4,1).
\]
The coefficients of $\zeta_{11}^i+\zeta_{11}^{-i}$, $1\le i\le5$,
in the real trace of the Jacobi sum are therefore
$(4,6,1,4,3)$.  They are not all equal.  Since the sole rational
relation among one and these five cyclotomic pairs is
$1+\sum_{i=1}^5(\zeta_{11}^i+\zeta_{11}^{-i})=0$, this trace is
not rational.  The only proper subfield of $F$ is $\mathbb Q$,
so it generates $F$.  The attached-variety existence assertion
preceding Zhang's Theorem~A and the matching Tate representations
now identify $B_F$ with the abelian variety attached to $f$, up to
isogeny, by \citet[\S5, Corollary~2, p.~22]{Faltings1986}.

The real degree five is odd.  The verified simple zero therefore
meets every hypothesis of \citet[Theorem~A, p.~28; proof in
\S7.2, p.~143]{Zhang2001}, which gives finite full $\Sha(B/F)$.
There is no CM exclusion or additional finite-place hypothesis in
this odd-degree case.  The kernel of
$\Sha(B/\mathbb Q)\to\Sha(B/F)$ lies in
$H^1(\operatorname{Gal}(F/\mathbb Q),B(F))$, a finite group because
$B(F)$ is finitely generated.  Thus $\Sha(B/\mathbb Q)$ is finite.
The rational rank one is the cited classical table entry, and
\citet[Theorem~1.1, p.~202]{GrossRohrlich1978} gives the stated torsion.
\end{proof}

\begin{lemma}[An integral decomposition at eleven]
\label{lem:eleven-integral-isogeny}
There are rational homomorphisms
\[
 \Phi:J\longrightarrow A^3\times B^6,\qquad
 \Psi:A^3\times B^6\longrightarrow J
\]
with $\Phi\Psi=[11]$ and $\Psi\Phi=[11]$.  In particular,
\[
 T_3J\simeq(T_3A)^{\oplus3}\oplus(T_3B)^{\oplus6}
\]
integrally and $G_{\mathbb Q}$-equivariantly.
\end{lemma}
\begin{proof}
For $a=1,\ldots,9$ use the cyclic quotient
$C_a:y^{11}=x(1-x)^a$ and the rational map
\[
 x=-X^{11}/Z^{11},\qquad
 y=(-1)^{a+1}XY^a/Z^{a+1}.
\]
Its geometric deck subgroup is
$H_a=\{(\zeta_{11}^u,\zeta_{11}^v):u+av=0\}$.
The ninety characters in $H^1(X_{11})$ have indices
$(i,j)$ with $i,j,i+j\ne0$ modulo eleven.  They partition
into the nine lines $j=ai$.  The branch-exponent triples
for $a=1,5,9$ are equivalent to $(1,1,9)$ under permutation
and common nonzero scaling modulo eleven; those for
$a=2,3,4,6,7,8$ are equivalent to $(1,2,8)$.
These branch-coordinate and monomial transformations are defined
over $\mathbb Q$, so the product of the nine quotient Jacobians
is $A^3\times B^6$.

For the quotient maps $f_a$, pullback followed by trace satisfies
$f_a^*f_{a,*}=\sum_{h\in H_a}h$, while
$f_{a,*}f_a^*=[11]$.  The first operator acts by eleven
on its own character lines and by zero on every other line.
It follows that $f_{a,*}f_b^*=0$ for $a\ne b$ and
$\sum_a f_a^*f_{a,*}=[11]$.  Faithfulness on a rational Tate
module turns these into identities of abelian-variety
homomorphisms.  Taking the product of the traces and the
sum of the pullbacks gives $\Phi$ and $\Psi$.
Since eleven is a unit in $\mathbb Z_3$, these identities
give the integral Tate-module isomorphism.
\end{proof}

\begin{lemma}[Rank-one factors and integral compact localization]
\label{lem:rank-one-compact-localization}
Let $D/\mathbb Q$ be an abelian variety with finite
$\Sha(D/\mathbb Q)$. Suppose that $D$ is rationally isogenous to
a product $\prod_iD_i$ with
$\operatorname{rank}D_i(\mathbb Q)\le1$ for every $i$.
For every rational prime $\ell$, the integral compact localization
\[
 \operatorname{Sel}(\mathbb Q,T_\ell D)
       \longrightarrow H^1(\mathbb Q_\ell,T_\ell D)
\]
is injective. There is no restriction on the isogeny's degree.
\end{lemma}
\begin{proof}
The inverse-limit Kummer sequence and finite $\Sha$ identify the
source with $D(\mathbb Q)^{\wedge_\ell}$, including its finite
$\ell$-primary torsion. At the real place, $H^1(\mathbb R,D)$ is
killed by two and has zero Tate module, so the integral Tate-module
Kummer condition there is automatic, including when $\ell=2$.
Local Kummer theory embeds
$D(\mathbb Q_\ell)^{\wedge_\ell}$ into the target cohomology group.
It therefore suffices to prove injectivity on these completions.

The logarithm on $D_i(\mathbb Q_\ell)$ is injective on a sufficiently
small open formal subgroup, and a nonzero integer multiple of every
point belongs to that subgroup. Its kernel is consequently exactly
the torsion subgroup. A nontorsion rational point remains nontorsion
locally: an equality $[m]P=0$ over $\mathbb Q_\ell$ between rational
points already holds over $\mathbb Q$. Thus a generator of a rank-one
rational Mordell--Weil group has nonzero local logarithm, proving
injectivity of
\[
 D_i(\mathbb Q)\otimes\mathbb Q_\ell
       \longrightarrow\operatorname{Lie}(D_i)(\mathbb Q_\ell).
\]
For rank zero the source is zero. The product has the same
injectivity, and a rational isogeny induces isomorphisms on the
rationalized Mordell--Weil spaces and on the local Lie algebras.
Hence the asserted localization for $D$ is injective after
inverting $\ell$.

Any remaining integral kernel is therefore contained in
$D(\mathbb Q)[\ell^\infty]$. No nonzero such point disappears in
$D(\mathbb Q_\ell)^{\wedge_\ell}$. Indeed, if
$P\in\bigcap_n\ell^nD(\mathbb Q_\ell)$, choose $Q_n$ with
$\ell^nQ_n=P$. When $P$ is torsion, every $Q_n$ is torsion.
The local torsion subgroup is finite, since it injects into the
finite quotient by a sufficiently small torsion-free formal
subgroup. For $n$ larger than its $\ell$-primary exponent,
$\ell^nQ_n$ cannot be a nonzero $\ell$-primary torsion point.
Thus $P=0$, which eliminates the entire integral kernel.
\end{proof}

\begin{theorem}[Compact localization at exponent eleven]
\label{thm:eleven-compact-selmer}
The full Tate--Shafarevich group $\Sha(J/\mathbb Q)$ is finite,
$\operatorname{rank}J(\mathbb Q)=6$, and for every rational prime
$\ell$ the integral compact Selmer localization
\[
 \operatorname{Sel}(\mathbb Q,T_\ell J)
             \longrightarrow H^1(\mathbb Q_\ell,T_\ell J)
\]
is injective. At $\ell=3$ its source is a free $\mathbb Z_3$-module
of rank six.
\end{theorem}
\begin{proof}
The rank-zero factor $A$ has finite full $\Sha$ and rational
torsion $\mathbb Z/11\mathbb Z$ by Proposition~\ref{prop:eleven}.
The preceding proposition gives the corresponding assertions for
$B$, with rank one.  Finiteness of full $\Sha$ is invariant under
isogeny and finite products, so the integral isogeny lemma gives
finite $\Sha(J/\mathbb Q)$ and rank six.

For a number-field abelian variety $D$, the inverse-limit Kummer
sequence is
\[
 0\longrightarrow D(\mathbb Q)^{\wedge_3}
 \longrightarrow\operatorname{Sel}(\mathbb Q,T_3D)
 \longrightarrow T_3\Sha(D/\mathbb Q)\longrightarrow0.
\]
Finite $\Sha$ makes the last term zero.  Thus the Selmer group of
$A$ is zero and that of $B$ is $\mathbb Z_3$, with no three-primary
rational torsion.  The integral isogeny identifies the Selmer group
of $J$ with six copies of the latter group, respecting all local
Kummer conditions.

Lemma~\ref{lem:rank-one-compact-localization} applies to the
three rank-zero and six rank-one factors. It proves the claimed
integral localization for every $\ell$, including $\ell=11$,
where the isogeny degree need not be a unit.
\end{proof}

\begin{remark}[Compact and affine scope]
\label{rem:eleven-compact-scope}
This proves the compact localization condition for the base exponent
eleven.  It does not establish the corresponding condition for
$11\cdot3^r$ with $r\ge1$, nor the global cuspidality of all affine
pro-three Selmer sections at exponent eleven.  In the latter problem,
the distinct coordinate cusps have differences of order eleven, so
their compact three-adic Kummer classes coincide.  The good-place
argument that selects one cusp using a nonzero three-primary cusp
difference therefore cannot be imported from the exponents divisible
by three.  Separation of the affine cusp packets at three does not,
by itself, impose the needed global boundary support condition.
\end{remark}

\subsection{An exponent-thirteen Prym and its central value}
\label{subsec:thirteen}

The complementary rank-zero target at exponent thirteen is already
implicit in \citet[Proposition~2.1 and Corollary~2.2, p.~942]{Tzermias2004},
using the rank-one calculation of Gross--Rohrlich. The decomposition
over the cyclotomic field is due to
\citet[Theorem~6.2 and Corollary~6.3, pp.~89--90]{Lim1992}.
We give explicit quotient equations, the resulting Hecke constituents,
and a central-value certificate for this classical target. Its Selmer
input is proved separately in Section~\ref{subsec:lm-lz-rational}.

\begin{proposition}\label{prop:thirteenpartial}
The Fermat curve $X_{13}$ has a nonconstant rational morphism to an
abelian fourfold $P_{13}$ such that
\[
 P_{13}(\Q)\simeq\Z/13\Z,\qquad
 2.632<L(P_{13}/\Q,1)<2.636.
\]
Its conductor is $13^8$, its archimedean factor is
$(2\pi)^{-4s}\Gamma(s)^4$, and its functional-equation sign is $+1$.
\end{proposition}
\begin{proof}
We give the quotient, rational-rank argument, Hecke decomposition,
and central-value certificate explicitly.

\paragraph{The two curve quotients.}
There is a degree-thirteen map
\[
 X_{13}\longrightarrow\mathcal K,
 \qquad [X:Y:Z]\longmapsto[X^4Z:Y^4X:Z^4Y],
 \quad\mathcal K:\ A^4B+B^4C+C^4A=0.
\]
Substitution multiplies the Fermat equation by $(XYZ)^4$.
On the torus, the exponent matrix has determinant thirteen, and its
kernel is generated by $[X:Y:Z]\mapsto[\zeta X:\zeta^4Y:Z]$.
The invariant functions
$x=-X^{13}/Z^{13}$ and $y=XY^3/Z^4$ satisfy
$y^{13}=x(1-x)^3$. Thus $\mathcal K$ is the smooth proper cyclic
quotient $C_{1,3,9}$, of genus six.

The rational automorphism $c[A:B:C]=[C:A:B]$ has order three.
Its two fixed points are $[1:\zeta_3:\zeta_3^2]$ and its conjugate;
$[1:1:1]$ is not on $\mathcal K$. Riemann--Hurwitz gives genus two
for $B=\mathcal K/\langle c\rangle$.
For an explicit equation put
\[
 S=A+B+C,\quad T=AB+BC+CA,\quad U=ABC,\quad
 D=(A-B)(B-C)(C-A).
\]
Direct expansion gives
\[
 2(A^4B+B^4C+C^4A)
 =S^3T-3ST^2-S^2U+5TU-D(S^2-T).
\]
With $t=T/S^2$, $u=U/S^3$ and $d=D/S^3$, the equations are
\[
 d(1-t)=t-3t^2+(5t-1)u,\qquad
 d^2=t^2-4t^3-4u-27u^2+18tu.
\]
Eliminating $d$ gives
\[
 (13t^2-16t+7)u^2+(-12t^3+14t^2-7t+1)u+t^5=0.
\]
Its discriminant is $(1-t)^2f(t)$, where
\[
 B:\ v^2=f(t)=1-12t+52t^2-104t^3+104t^4-52t^5.
\]
The quotient coordinate is
\[
 v=\frac{2(13t^2-16t+7)u-12t^3+14t^2-7t+1}{1-t}.
\]
These invariants separate cyclic orbits generically: $S,T,U$
determine the unordered triple and $D$ distinguishes its two
cyclic orientations. Consequently this equation gives the quotient
function field, not just a map to a genus-two curve.
The discriminant of $f$ is $2^8 13^3$.

\paragraph{The finite rational target.}
Let $J=\Jac(\mathcal K)$ and
$P_{13}=\ker(\operatorname{Nm}:J\to\Jac(B))^0$.
Pullback followed by norm is multiplication by three, so
$J\sim_\Q\Jac(B)\times P_{13}$ and $\dim P_{13}=4$.
The exact counts
\[
 \begin{array}{c|ccc}
 q&\#B(\F_q)&\#B(\F_{q^2})&\#\Jac(B)(\F_q)\\\hline
 3&7&11&27\\
 5&6&26&26
 \end{array}
\]
use $\#\Jac(B)(\F_q)=(\#B(\F_q)^2+\#B(\F_{q^2}))/2-q$.
Prime-to-residue-characteristic torsion injectivity at both good
primes shows $\Jac(B)(\Q)_{\rm tors}=0$. The distinct rational
points $\infty$ and $(0,1)$ therefore give a non-torsion divisor
class. On the other hand,
\citet[\S4, table on p.~219]{GrossRohrlich1978} proves
$\operatorname{rank}J(\Q)=1$. It follows that
$\operatorname{rank}\Jac(B)(\Q)=1$ and
$\operatorname{rank}P_{13}(\Q)=0$.

Their Theorem~1.1 gives $J(\Q)_{\rm tors}=\Z/13\Z$.
This torsion is contained in $P_{13}$: its norm is zero because
$\Jac(B)(\Q)$ has no torsion, and
$3x=(1-c)(2+c)x$ whenever $(1+c+c^2)x=0$.
Since multiplication by three is invertible on a group of order
thirteen, $x$ belongs to the image of $1-c$, namely $P_{13}$.
Thus $P_{13}(\Q)=\Z/13\Z$.
Composing an Abel--Jacobi map of $\mathcal K$ with $1-c$ gives a
nonconstant map to $P_{13}$: otherwise $1-c$ would vanish on the
Jacobian generated by the curve, contrary to $\dim P_{13}=4$.

\paragraph{The Hecke constituents.}
Put $L=\Q(\zeta_{13})$, $H=\langle3\rangle\subset(\Z/13)^\times$,
and $K_0=L^H$, a quartic CM field. The automorphism
\[
 \tau[A:B:C]=[\zeta_{13}^4 A:\zeta_{13}B:C]
\]
satisfies $c\tau c^{-1}=\tau^3$ and gives the usual faithful
$L$-action on $J$. Over an algebraic closure of a coefficient field,
the twelve nontrivial $\tau$-eigenspaces of $V_\ell J$ are lines.
The automorphism $c$ groups them into four triples; the invariant
part $V_\ell\Jac(B)$ has one dimension in each triple.
The central field $K_0$ preserves this invariant part and acts
faithfully on it. Its endomorphisms are defined over $K_0$, because
$G_{K_0}$ acts on $L$ through $H$, which fixes $K_0$.
Thus $\Jac(B)$ has CM by $K_0$, defined over $K_0$.

Choose a cubic Dirichlet character $\chi$ modulo thirteen with
$\chi(2)=\zeta_3$; since $2^4=3$ modulo thirteen, its restriction
to $H$ sends $3$ to $\zeta_3$.
For $g\in G_{K_0}$ whose action on $L$ is $\zeta_{13}\mapsto
\zeta_{13}^{3^e}$, the operator $gc^{-e}$ preserves every
$\tau$-eigenline and commutes with $c$. It therefore acts by the
same scalar on the three lines in each orbit. Passing to the
$c$-eigenspaces shows that their characters differ by $\chi^j$,
for $j=0,1,2$.
To verify this comparison for the full Galois group, put
$D=\{1,5,8,12\}$, so $(\Z/13)^\times=H\times D$.
Choose a vector $e_d$ in the $\tau$-eigenline with eigenvalue
$\zeta_{13}^d$ for each $d\in D$; the vectors $c^ke_d$
then lie in the lines indexed by $9^kd$.
If $g(\zeta_{13})=\zeta_{13}^a$ and
$a^{-1}=9^{k(a)}d(a)$, then
$g(e_d)=u(g,d)c^{k(a)}e_{d(a)d}$ for a scalar $u(g,d)$.
Because $g$ commutes with $c$, the basis
$\sum_{j=0}^2\lambda^{-j}c^je_d$ of the $c$-eigenspace
with eigenvalue $\lambda$ gives its matrix as the matrix on the
invariant space multiplied by $\lambda^{k(a)}$.
For $\lambda=\zeta_3$, this multiplier is $\chi(a)$.
Thus, after extending coefficients,
\begin{equation}\label{eq:thirteen-twists}
 V_\ell J\simeq V_\ell\Jac(B)\oplus
 (V_\ell\Jac(B)\otimes\chi)\oplus
 (V_\ell\Jac(B)\otimes\chi^2),
\end{equation}
and $V_\ell P_{13}$ is the sum of the last two terms.

The CM theorem identifies the one-dimensional $G_{K_0}$ characters
with algebraic Hecke characters; see
\citet[Theorems~3.5 and 3.7, pp.~10--11]{ConradCM}.
Induction from one of the four eigenspaces gives
$V_\ell\Jac(B)$ over $\Q$. The published Hecke factorization
and completed functional equation also follow directly from
\citet[Theorem~4 and Remark~2, pp.~187--188]{Milne1972}:
the CM field $K_0$ is Galois-stable in the endomorphism algebra,
its field of definition is $K_0$, and its four embeddings form
one Galois orbit. Therefore each twist in
\eqref{eq:thirteen-twists} has a degree-four Hecke $L$-function,
with two factors $\Gamma(s)$. The two twists have complex-conjugate
Dirichlet coefficients and dual functional equations. Their root
numbers are complex conjugates of modulus one, so their product
has root number $+1$. In particular $L(P_{13},s)$ is entire.

\paragraph{Conductor and local factors.}
The displayed genus-two model has good reduction away from two
and thirteen. At two use the integral model
\[
 w^2+w=-3t+13t^2-26t^3+26t^4-13t^5,
 \qquad v=2w+1,
\]
whose reduction has the smooth odd-degree Artin--Schreier form.
At thirteen, take $\pi^8=13$ and set $t=\pi^{-2}u$,
$v=\pi^{-1}y$. The equation becomes
\[
 y^2=\pi^2-12u+4\pi^6u^2-8\pi^4u^3+8\pi^2u^4-4u^5.
\]
Its smooth special fiber is $y^2=u-4u^5$. Thus the reduction is
potentially good over a tame extension. A tame generator
$\pi\mapsto\zeta_8\pi$ acts on the special fiber by
$(u,y)\mapsto(\zeta_8^2u,\zeta_8y)$.
On its differentials $du/y,u\,du/y$ the eigenvalues are
$\zeta_8,\zeta_8^3$; on first cohomology the four eigenvalues are
the primitive eighth roots. This can equally be checked on the
same smooth equation in characteristic zero and specialized.
There are no inertia invariants, so the conductor exponent of
$\Jac(B)$ is four. Twisting by $\chi$ or $\chi^2$ multiplies
these eigenvalues by third roots; none becomes one. Each twist
still has conductor exponent four and no wild contribution.
Consequently $P_{13}$ has conductor $13^8$ and Euler factor one
at thirteen. With $B_0=13^4/(2\pi)^4$ its completed function is
\[
 \Lambda_P(s)=B_0^s\Gamma(s)^4L(P_{13},s),\qquad
 \Lambda_P(s)=\Lambda_P(2-s).
\]

The program computes Euler polynomials of $J$ from the Jacobi sums
for $(1,3,9)$, and those of $\Jac(B)$ by exact point counts;
polynomial division gives the factors of $P_{13}$.
For example, the factors at two and three are
\[
 \begin{split}
 P_2(T)&=1-4T^4+16T^8,\\
 P_3(T)&=1-3T+4T^2+3T^3-11T^4+9T^5
                       +36T^6-81T^7+81T^8.
 \end{split}
\]
For $J$, the degree-three extension counts at three are
$\#\mathcal K(\F_{27})=55$ and
$\#\mathcal K(\F_{729})=913$; they give
$1+9T^3+71T^6+243T^9+729T^{12}$.
At primes of even residue degree modulo thirteen, the standard
Jacobi identity of \citet[Lemma~1.1]{GrossRohrlich1978} supplies
the factor of $J$. At split primes the coefficients are integer
Newton polynomials in the thirteen Jacobi residue counts.
Only the exact initial coefficients required by the cutoff are
used or reported.

\paragraph{A central-value certificate with a universal tail.}
Write $L(P_{13},s)=\sum a_nn^{-s}$. The same Mellin argument gives
\[
 L(P_{13},1)=2\sum_{n\ge1}\frac{a_n}{n}Q_4(n/B_0),\qquad
 Q_4(x)\le\Gamma(k+1)^4x^{-k}.
\]
There are eight reciprocal roots of absolute value $\sqrt p$ at
every good prime. Hence $|a_n|\le d_8(n)\sqrt n$, where
$\sum d_8(n)n^{-\sigma}=\zeta(\sigma)^8$.
For $\sigma>1$ and $k+1/2>\sigma$, the omitted Dirichlet tail
after $N$ is therefore at most
\begin{equation}\label{eq:thirteen-tail}
 2\Gamma(k+1)^4B_0^kN^{\sigma-k-1/2}\zeta(\sigma)^8.
\end{equation}
This bound uses only the degree-eight Weil estimate.

The finite sum is evaluated on $z=2+it$, $0\le t\le20$, as
\[
 \frac2\pi\operatorname{Re}\int_0^{20}
 \frac{\Gamma(3+it)^4}{2+it}B_0^{2+it}
 \sum_{n\le N}\frac{a_n}{n^{3+it}}\,dt.
\]
Seventeen gamma recurrences bound its vertical tail by
\[
 \frac2\pi B_0^2\left(\sum_{n\le N}\frac{|a_n|}{n^3}\right)
 \frac{\Gamma(20)^4}{68\cdot20^{68}}.
\]
For $N=100000$, Arb arithmetic at $144$ bits gives the finite
integral $2.633619912614\pm9.68\cdot10^{-14}$ and vertical tail
less than $2.612\cdot10^{-20}$.
Equation~\eqref{eq:thirteen-tail}, with $k=8$ and $\sigma=3/2$,
is less than $0.001459$. These enclosures prove the asserted
rational interval and complete the proof.
\end{proof}

The finite rational group and central nonvanishing just proved do
not by themselves identify the divisible part of
$\Sha(\Q,P_{13})$. In the quartic-CM description the two constituents
over $\Q(\sqrt{13})$ are nontrivial cubic twists of a CM Hilbert
form. Their central characters are cubic, so the trivial-central-
character version of Nekov\'a\v r's Theorem~B cannot be applied
directly. Section~\ref{subsec:lm-lz-rational} supplies the needed
arithmetic input by a different Euler-system argument: it proves
$T_\ell\Sha(\Q,P_{13})=0$ for all but finitely many $\ell$ splitting
completely in one fixed finite Galois extension. This is sufficient
for Theorem~\ref{thm:split-prime-descent}. The map constructed above
generates $P_{13}$ by differences, so its image satisfies the
geometric hypothesis of that theorem.

Passing to the real cubic field cut out by $\chi$ does not supply
this hypothesis through Zhang's rank-one theorem. To check this
precisely, restrict the four primitive eighth-root inertia characters
computed above to $F_0=\Q(\sqrt{13})$: the characters become
$i,-i,i,-i$. Real multiplication and the polarization give
cyclotomic determinant on each two-dimensional Hilbert constituent,
so its two distinct inertia characters are $i$ and $-i$.
Frobenius at thirteen fixes both, since $13\equiv1\pmod4$.
Thus each local Weil representation is a sum of two characters,
with zero monodromy, and the local automorphic representation is
principal series of conductor exponent two. Restriction to the
further cubic field preserves this property, while the totally
real degree becomes six. All other finite places have good reduction.
Consequently even a simple-zero certificate for the untwisted form
would not satisfy the finite-place hypothesis of
\citet[Theorem~A and its preceding setup, p.~28]{Zhang2001},
which in even degree requires a conductor exponent one, or the
nonprincipal-place alternative of
\citet[Theorem~4.3.2, p.~588]{Zhang2010}.
Nekov\'a\v r's Theorems~B and~B$'$ require central nonvanishing;
they are not rank-one statements.

The other full cyclic-quotient factors do not give a verified
rank-zero substitute. At exponent thirteen their three types are
$(1,1,11)$, $(1,2,10)$ and $(1,3,9)$, with multiplicities three,
six and two. All three have sign $-1$ in
\citet[table on p.~219]{GrossRohrlich1978}; the first two have
Hecke-character conductor $(1-\zeta_{13})^2$.
Their Theorem~2.1, pp.~208--209, and the same table prove rank one
for the last two types, but give only rank at most one for the first.
The table's footnote invokes Birch--Swinnerton-Dyer to predict
equality in that remaining case. No unconditional positive-rank
claim for that hyperelliptic factor is inferred here.

\subsection{Nekov\'a\v r's third CM alternative and the cubic twists at thirteen}

Let $F$ be totally real and let $f$ be a parallel-weight-two CM
Hilbert eigenform.  Write $L^+$ and $L'$ for the real coefficient
field and its CM quadratic extension in Appendix B.6 of
\citet[Appendix~B.6]{Nekovar2012}.

\begin{lemma}[The split-coefficient condition is sufficient]
\label{lem:thirteen-split-coefficient}
Suppose that a rational prime $\ell$ splits completely in a
finite Galois extension containing $L'$.  Then the Galois-element
conditions (A1)--(A2) of that paper hold at every coefficient
prime above $\ell$.  In particular, after restriction to the
field trivializing the finite central character and the CM
quadratic extension, the two one-dimensional characters have
independent infinite images in the sense required by
Theorem B$'$(c2).
\end{lemma}
\begin{proof}
Proposition B.6.5(6), printed p.~664, applies because the real
coefficient completion is $\mathbb Q_\ell$ and its prime splits
in $L'/L^+$.  Proposition B.6.5(2) identifies the resulting
condition with the infinitude of $\psi_2(\ker\psi_1)$.
The proof on p.~665 is also explicit.  Dependence would force
powers of both characters to be powers of the cyclotomic
character.  Conjugacy exchanges them, so their exponents are
equal.  The cyclotomic exponent has odd weight; a potentially
crystalline character cannot then have the required half-integral
Hodge--Tate weight.  This proves the stated independence.
\end{proof}

This lemma removes character independence as an obstruction at
completely split coefficient primes.  It does not remove the
separate hypothesis $f\in S_2(\mathfrak n,1)$ in Theorems B and
B$'$ (printed p.~593).

For the exponent-thirteen factor, put
$L=\mathbb Q(\zeta_{13})$, $H=\langle3\rangle$, and
$K_0=L^H$.  Its maximal real subfield is
$F_0=\mathbb Q(\sqrt{13})$.
Let $B$ be the genus-two quotient used in the exponent-thirteen
calculation and let $\chi$ be the cubic Dirichlet character
with $\chi(2)=\zeta_3$.  The already established decomposition is
\[
 V_\ell P_{13}\simeq
 (V_\ell\operatorname{Jac}(B)\otimes\chi)
 \oplus
 (V_\ell\operatorname{Jac}(B)\otimes\chi^2).
\]
All finite coefficient extensions implicit in this display are
understood.  Over $F_0$, the real multiplication on
$\operatorname{Jac}(B)$ yields two-dimensional constituents
$W$ with determinant the cyclotomic character.  One may see this
directly from the polarization: the Rosati-fixed real
multiplication field acts self-adjointly, and the alternating
pairing on each constituent has cyclotomic similitude.
Consequently
\[
 \det(W\otimes\chi^j)=\chi_{\ell}\chi^{2j},
 \qquad j=1,2.
\]
The cubic character remains nontrivial on $G_{F_0}$, since its
real cubic field is disjoint from the quadratic field $F_0$.
Thus the corresponding Hilbert forms have nontrivial cubic
nebentypus.  Extending their coefficient fields, including
choosing completely split coefficient primes, does not change
this determinant character.

\begin{proposition}[Exact limit of this application]
\label{prop:thirteen-nebentypus-gate}
The split-coefficient lemma verifies the Galois-element
hypothesis for the exponent-thirteen cubic twists, but Theorems
B and B$'$ of Nekov\'a\v r (2012) do not apply directly to those
twists, because their nebentypus is nontrivial.
The elementary substitutions into Theorem A described below
do not repair this missing hypothesis.
\end{proposition}
\begin{proof}
The determinant calculation proves the first assertion.
Theorem A, printed p.~592, permits a nontrivial nebentypus
$\omega$, but requires an auxiliary CM quadratic extension
$K/F_0$ and a finite-order character $\theta$ with
\[
 \theta|_{\mathbb A_{F_0}^{\times}}\,\omega=1;
 \tag{*}
\]
its conclusion concerns the representation twisted by $\theta$
over $K$.  Taking $\theta=1$ leaves the desired representation
unchanged but violates (*) for the cubic twists.

More generally, start with the untwisted form $f_0$, replace it
by $f_0\otimes\chi^a$, and choose
$\theta=\chi^b\circ N_{K/F_0}$.  On the base ideles the norm is
squaring, so (*) is $\chi^{2(a+b)}=1$.  Since $\chi$ has order
three, this forces $a+b=0$ modulo three.  The resulting
representation over $K$ is untwisted, and hence is not either
of the two required cubic constituents.  Taking the auxiliary
CM field equal to the CM field of $f_0$ also violates the
explicit CM-field exclusion in Theorem A.

Finally, base change to the real cubic extension which kills
$\chi$ does not give a rank-zero untwisted problem: its cyclic
base-change $L$-function includes the untwisted
$\operatorname{Jac}(B)$ factor, and that variety already has
a non-torsion rational point.  In particular the entire
base-changed representation cannot have zero finite Selmer
group.  One would still have to isolate and control the
nontrivial cubic isotypic components; Theorem B applied to the
whole base change does not do this.
\end{proof}

This calculation concerns the precise scope of Nekov\'a\v r's
theorems, not a failure of the Tate--Shafarevich conclusion.
The following Euler-system argument treats the nontrivial-nebentypus
CM constituents directly and proves the needed vanishing for almost
all completely split rational primes.

\subsection{A rational quadratic-Hilbert consequence of the Euler-system machinery}
\label{subsec:lm-lz-rational}

In this subsection LZ denotes the February 2025 version of
\citet{LoefflerZerbes2025}, with theorem numbers and pages referring
to that version.  The distinction between its
stated large-image theorem and its formal intermediate results is
essential.  Its condition (BI), Definition 2.5.1, requires a product of
two special linear groups and fails for CM representations.  The next
proposition is instead a deduction from the constructions in Sections
3.4--3.5 and 6--8 and the abstract results in Sections 9--10.  Only a
rational Selmer conclusion is sought; the canonical integral period
normalization of Section 3.2.3 is not required.

This is an arithmetic input for the proper finite-descent application
at exponent thirteen. The regular-prime quotient-section classification
proved earlier does not depend on this separate source chain.
The following table fixes the versions and roles of its inputs.

\begin{center}
\begingroup
\small
\setlength{\tabcolsep}{3pt}
\renewcommand{\arraystretch}{1.12}
\begin{tabular}{@{}p{.17\linewidth}p{.17\linewidth}p{.30\linewidth}p{.30\linewidth}@{}}
\toprule
Source & Cited version & Results used & Role in the argument \\
\midrule
\citet{MazurRubin2004}
& Published memoir, 2004
& \S3.5, pp.~27--28: hypotheses (H.0)--(H.4)
& The lattice and auxiliary-twist hypotheses checked directly below. \\
\addlinespace
\citet{LoefflerRockwoodZerbes2024}
& Author version v3, 30 July 2024
& Theorems~2.6.2 and~7.1.1, pp.~4 and~15
& Finite generation and ordinary families of cohomology classes, used in the family Euler-system construction. \\
\addlinespace
\citet{LoefflerZerbes2025}
& Author version v2, 18 February 2025
& Theorems~3.4.4, 3.5.1 and~7.1.3; \S7.2; Theorem~8.1.1, p.~22
& The Hilbert family, its Yoshida lift and $p$-adic $L$-functions, and uniformly integral family Euler classes. \\
\addlinespace
\citet{LoefflerZerbes2026}
& Author version v4, 21 April 2026
& Theorem~A, p.~4; Theorem~6.6.4, p.~14
& Explicit reciprocity in the range $r_1-r_2\ge3$, followed here by continuation from a dense set of such weights. \\
\addlinespace
\citet{LoefflerZerbes2025}
& The same February 2025 version
& Theorem~8.2.1; Lemma~10.1.1; Propositions~10.1.4 and~10.2.2; Theorem~10.3.3; Corollary~9.6.2; proof of Theorem~12.3.1, p.~33
& Meromorphic comparison, leading terms and characteristic-zero Selmer control. These intermediate results supply the rational deduction; the stated (BI) theorem is not applied. \\
\bottomrule
\end{tabular}
\endgroup
\end{center}

\begin{proposition}[Rational quadratic-Hilbert criterion]
\label{prop:lm-rational-hilbert}
Let $F/\mathbb Q$ be real quadratic, and let $\pi$ be a cuspidal Hilbert
automorphic representation of weight $(2,2,0,0)$.  Suppose that its
finite central character factors through $N_{F/\mathbb Q}$ and that
$\pi^\sigma\not\simeq\pi$, where $\sigma$ is the nonidentity automorphism
of $F$.  Let $p\geq7$ split in $F$, be prime to the level, and be ordinary
at both places above $p$.  Write
\[
 V=\operatorname{Ind}_{G_F}^{G_{\mathbb Q}}\rho_{\pi,p}^{\vee}.
\]
Assume the following conditions for a stable lattice $T\subset V$.
The hypotheses (H.0)--(H.4) of
\citet[\S3.5, pp.~27--28]{MazurRubin2004} hold for
$T$ and for $T(\eta)$ whenever $\eta$ has $p$-power order and square-free
conductor prime to $p$ and the level and discriminant in this setup.
Assume also that $V$ is irreducible.
Then
\[
 L(\pi,1)\ne0\quad\Longrightarrow\quad H^1_f(F,\rho_{\pi,p}^{\vee})=0.
\]
\end{proposition}

\begin{proof}
We explain the deduction, including the hypotheses that are sometimes
hidden by citing the more restrictive main theorem.  Choose $S$ to
contain $p$ and all bad primes, and let $\mathcal R$ consist of
square-free products of primes $q\notin S$ with $q\equiv1\pmod p$.
The group $\Delta_m$ is the maximal $p$-power quotient of
$(\mathbb Z/m\mathbb Z)^\times$.  Thus the auxiliary characters in LZ
Section 9.1 are exactly of the type stipulated here.  They are not
arbitrary finite-order Dirichlet characters.

By LZ Theorem 3.4.4, there is an ordinary Hilbert family on a weight
polydisc through $\pi$.  Theorems 3.5.1 and 7.1.3 give its analytic
$p$-adic $L$-functions and twisted Yoshida lift.  Restrict to the
one-parameter family of Section 7.2, of weights
\[
 (2+4\mu,2+2\mu,-2\mu,-\mu).
\]
Let $\mathcal M_{\mathrm{et}}$ be the torsion-free pullback of the
Siegel Galois sheaf, a free rank-four module on this disc, and use
\[
 \mathcal M=\mathcal M_{\mathrm{et}}^{\vee}(\mathbf t_1)
\]
as in LZ Section 11.  Here $\mathbf t_1$ is the family cyclotomic
twist whose value at a classical point is $t_1$.  Proposition 7.2.4
identifies the semisimplification of the fiber of
$\mathcal M_{\mathrm{et}}$ with
$\operatorname{Ind}\rho_{\pi,p}(t_1)$; hence the fiber of
$\mathcal M$ at the central point, where $t_1=0$, has
semisimplification $V$.  The assumed irreducibility identifies the
representation itself.
The Hilbert family construction requires no smoothness of the entire
Siegel eigenvariety at this point.

LZ Theorem~8.1.1, p.~22, supplies a family of Euler-system classes
for this module. Its construction uses
\citet[Theorems~2.6.2 and~7.1.1, pp.~4 and~15]{LoefflerRockwoodZerbes2024},
which give finite generation and the ordinary family classes without
a large-image hypothesis.
Its last assertion places all the classes in a single finitely generated
module over the power-bounded functions on the disc.  This uniform
integrality assertion is prior to, and independent of, condition (BI).
There is no lattice mismatch in applying the hypotheses.  If one stable
lattice in an irreducible representation has absolutely irreducible
reduction, every stable lattice is homothetic to it: scale the second
lattice into the first but not into its maximal-ideal multiple.  Its
nonzero residual image is the whole irreducible reduction, and
Nakayama's lemma makes the inclusion an equality.  Thus the lattice
obtained from the family has the same Mazur--Rubin properties.  The
same argument applies after any of the allowed coefficient extensions
and auxiliary twists.
Fix the smoothing parameter $r=3$, and an integer $c>1$ prime to the
bad primes, with the finite central character taking value $1$ at $c$.
Enlarge $S$ by the prime factors of $c$ and use the corresponding
$\mathcal R$, so that every Euler-system index is prime to $c$.
For the sign containing the central critical value, Theorem 8.2.1
expresses the regulator of these classes as the product of the desired
family $p$-adic $L$-function, a smoothing factor, and an auxiliary
one-variable family $L$-value.  Here the completed explicit reciprocity
law is used in its current scope:
\citet[Theorem~A, p.~4; Theorem~6.6.4, p.~14 of the April 2026
author version]{LoefflerZerbes2026} assumes $r_1-r_2\geq3$.  For our family this difference
is $k_2-2=2\mu$.  The classical weights $\mu\geq2$ in the weight
disc form a Zariski-dense subset, satisfy this restriction, and also
have $r=3\leq2\mu$.  Thus the analytic-continuation argument proving
LZ Theorem~8.2.1 applies on that dense subset, including its
specialization at $\mu=0$; no small-weight reciprocity statement is
being assumed.  The auxiliary factor is not identically
zero: at sufficiently large classical $\mu$, its evaluation is at the
noncentral critical integer $-\mu+3$, and nonvanishing follows from
Proposition 2.3.7.  It may vanish at $\mu=0$, which is precisely why the
leading-term argument is needed.

The independence of the comparison ratio from the cyclotomic variable
can be checked explicitly in this continuation argument. On a fixed
cyclotomic component choose a nonzero integer $b$ divisible by $p-1$
and prime to $p$. At classical weight $\mu$, the allowed integral
cyclotomic exponents form $[-2\mu,-\mu-1]$. The pairs $(\mu,j)$ for
which both $j$ and $j+b$ lie in that interval are dense in the product
of the two discs: choose arbitrarily large positive $\mu$ in any
prescribed residue class, then choose $j$ in its growing interval in
any prescribed residue class. The classical comparison ratio depends
only on $\mu$, so its meromorphic continuations at $j+b$ and at $j$
agree on this dense set and hence identically. Integer multiples of
$b$ are dense in the additive $p$-adic parameter disc. Continuity,
or differentiation along translations tending to zero, therefore
makes the ratio independent of the cyclotomic variable. This is the
one-variable meromorphic comparison used in LZ Theorem~8.2.1; its
zeros or poles at the central weight are handled by the following
leading-term argument.

Let $X$ be a parameter at the central point.  Divide the entire family
of Euler classes by its largest common power $X^h$.  Lemma 10.1.1 and
Proposition 10.1.4 show that the specialized classes form a nonzero
Euler system in one fixed stable lattice.  In particular no separate
denominator depending on $m\in\mathcal R$ is introduced.  The two nested
ordinary local submodules have ranks two and three, as in Section 11.
Complex conjugation has two eigenvalues of each sign.  The local
vanishing conditions of Assumption 9.5.1 hold by Proposition 3.3.1,
also after the finite unramified extensions at $p$ supplied by $\Delta_m$.
One can equivalently check this by purity: the unramified part of each
ordinary one-dimensional character has a Frobenius eigenvalue of
nonzero Weil weight, so no power of it is $1$.  The same applies to
the dual local subquotients.

The precise hypotheses of Proposition 10.2.2 are now satisfied: they
are (H.0)--(H.4) for $T$ and its $\Delta_m$-twists, not (BI).  Thus the
leading-term Euler system has a nonzero regulator in its rank-one
local quotient.  Divide the meromorphic comparison isomorphism by
the auxiliary family $L$-value and apply Theorem 10.3.3.  It identifies
this regulator with the desired $p$-adic $L$-function times the
smoothing factor, up to a nonzero constant resulting from a basis of
the one-dimensional crystalline space.

At the central critical character the smoothing factor for $r=3$ is
\[
 (c^2-c^4)(c^2-c^{-2})\ne0.
\]
The ordinary interpolation factors are nonzero by Proposition 3.1.5.
Consequently $L(\pi,1)\ne0$ makes the regulator nonzero at that
character.  Corollary 9.6.2 with $k=0$, localized at this characteristic
zero height-one prime of the Iwasawa algebra, says that the ordinary
Selmer characteristic ideal is a unit there.  Multiplying the chosen
lattice or regulator by a power of $p$ does not affect this conclusion.
This is why the integral Betti lattice and the (BI)-dependent
normalization of Section 3.2.3 are unnecessary.

For completeness, let $\Lambda$ denote the cyclotomic Iwasawa algebra
on the component containing the central character, and let $\mathfrak q$
be its characteristic-zero height-one prime. Write
$\widetilde H^i_{\mathrm{Iw}}$ for the cohomology of the Selmer complex
with ordinary conditions at both places above $p$. The preceding
characteristic-ideal conclusion gives
$(\widetilde H^2_{\mathrm{Iw}})_{\mathfrak q}=0$.
By LZ Proposition~3.3.5, p.~11, its first cohomology injects into global
Iwasawa cohomology and is torsion-free; the Euler-characteristic rank
formula gives rank zero, so $\widetilde H^1_{\mathrm{Iw}}=0$ on this
component. The derived base-change argument in the proof of LZ
Theorem~12.3.1, p.~33, now gives
\[
0\longrightarrow
\widetilde H^1_{\mathrm{Iw}}\otimes_\Lambda\kappa(\mathfrak q)
\longrightarrow \widetilde H^1(\text{central fiber})
\longrightarrow
\operatorname{Tor}_1^\Lambda
 (\widetilde H^2_{\mathrm{Iw}},\kappa(\mathfrak q))
\longrightarrow0.
\]
Both end terms are zero. Global degree-zero and degree-three terms,
and the local exceptional terms, vanish rationally by irreducibility
and the local and dual invariant vanishing established above.
At this weight-two critical character the ordinary subspace is a
cyclotomic twist of an unramified line and its quotient is unramified; their nontrivial
Frobenius eigenvalues identify the ordinary local condition with the
Bloch--Kato condition. Thus the central-fiber group is the required
finite Selmer group, which vanishes. This is the formal base-change
deduction, rather than an application of the (BI) version of the stated
main theorem. The integral period normalization is unnecessary because
powers of $p$ are units at $\mathfrak q$.
\end{proof}

\subsection{The cubic twists of the cyclic-quartic CM surface}

Let $\mathcal K$ be the Klein quintic
\[
 A^4B+B^4C+C^4A=0,
\]
let $c$ be the rational order-three coordinate permutation, and let
$B=\mathcal K/\langle c\rangle$.  Thus $B$ has the genus-two model
\[
 v^2=1-12t+52t^2-104t^3+104t^4-52t^5.
\]
Put $A_B=\operatorname{Jac}(B)$, $L=\mathbb Q(\zeta_{13})$,
$H=\langle3\rangle\subset(\mathbb Z/13\mathbb Z)^\times$, and
$K=L^H$.  Then $K$ is cyclic quartic CM, with real quadratic subfield
$F=\mathbb Q(\sqrt{13})$.  Let $\chi$ be a nontrivial cubic character
modulo $13$.

\begin{lemma}[The CM algebra is defined over the quartic field]
\label{lem:lm-thirteen-cm-field}
There is a degree-four CM algebra
$K\hookrightarrow\operatorname{End}^0_K(A_B)$.  On the four CM lines
the action of $G_{\mathbb Q}$ is the natural transitive cyclic
permutation through $\operatorname{Gal}(K/\mathbb Q)$.
\end{lemma}
\begin{proof}
Let $\tau$ denote the order-thirteen diagonal automorphism of
$\mathcal K$ over $L$, with $c\tau c^{-1}=\tau^3$.  The operator
\[
 \theta=\tau+\tau^3+\tau^9
\]
commutes with $c$ and hence induces an element of
$\operatorname{End}^0_{\overline{\mathbb Q}}(A_B)$; explicitly use
pullback, this operator, and norm,
and divide by the degree $3$ of the quotient map.  The Galois
subgroup $H$ fixes this operator, since it permutes its three terms.
The induced endomorphism is consequently defined over $K$.

The quotient by $\tau$ is the projective line and its three branch
points are totally ramified.  Equivalently each nonidentity power of
$\tau$ has three fixed points, so its trace on first cohomology is
$-1$.  It follows that the nontrivial order-thirteen characters
occur once each in the Tate module of
$\operatorname{Jac}(\mathcal K)$.  Taking the
$c$-invariant part leaves one line for each of the four $H$-orbits.
The eigenvalues of $\theta$ are the Gaussian periods
\[
 \sum_{h\in H}\zeta_{13}^{ah},
 \qquad aH\in(\mathbb Z/13\mathbb Z)^\times/H.
\]
They are distinct: equality of two such sums would give a nonzero
rational polynomial of degree at most $12$, with constant coefficient
zero, divisible by $1+X+\cdots+X^{12}$.  They are therefore the four
conjugates of a generator of $K$.  This proves the asserted embedding.
Galois conjugation takes the period indexed by $aH$ to the period
indexed by $baH$, proving the permutation assertion.  Only this CM
subalgebra is needed below; no larger endomorphism-field assertion
is being assumed.
\end{proof}

\begin{lemma}[Explicit Mazur--Rubin conditions]
\label{lem:lm-thirteen-mr}
For all but finitely many primes $p\geq7$ splitting completely in
$K\mathbb Q(\zeta_3)$, each of the representations
\[
 V_p(A_B)\otimes\chi^a,\qquad a=1,2,
\]
and their duals and Tate twists satisfies (H.0)--(H.4).  The same is
true after a further twist by any $p$-power-order Dirichlet character
$\eta$ whose conductor is prime to $13p$.  In particular it is true
for the square-free auxiliary characters of
Proposition~\ref{prop:lm-rational-hilbert}.
\end{lemma}
\begin{proof}
Discard the finitely many primes dividing the order of the CM
endomorphism lattice, the polarization degree, or the bad reduction
conductor, and choose the four CM idempotents over $\mathbb Z_p$.
The Tate lattice is their direct sum, with each summand rank one.
Index them by $0,1,2,3$ cyclically, with complex conjugation adding
$2$.  Since the cyclic quartic field has no proper CM subfield, its
CM type is, after cyclic relabeling, $\{0,1\}$.

At a split good prime $p$, each CM line is a crystalline character
of $G_{\mathbb Q_p}$ with cyclotomic exponent $0$ or $1$. Restricted
to inertia it is therefore respectively $1$ or the cyclotomic
character $\varepsilon_p$; the remaining crystalline rank-one
factor is unramified.  The cyclotomic character maps inertia onto
$\mathbb Z_p^\times$.  The four Galois conjugates of this inertia
group thus supply, for arbitrary units $u_i$, the four diagonal
matrices
\[
 \begin{split}
 &\operatorname{diag}(u_1,u_1,1,1),\quad
 \operatorname{diag}(1,u_2,u_2,1),\\
 &\operatorname{diag}(1,1,u_3,u_3),\quad
 \operatorname{diag}(u_4,1,1,u_4).
 \end{split}
\]
They lie in $G_K$; both $\chi$ and every allowed $\eta$ are trivial
on them.  Their product is
\[
 \operatorname{diag}(u_1u_4,u_1u_2,u_2u_3,u_3u_4),
 \qquad \varepsilon_p= u_1u_2u_3u_4.
\]
In particular every integral torus element
\[
 \operatorname{diag}(a,b,c,d),\qquad ac=bd,
\]
is supplied: take $(u_1,u_2,u_3,u_4)=(1,b,c/b,a)$.
Its cyclotomic multiplier is $ac$.  This proves the required exact
torus-image assertion by local crystalline characters, without an
unverified index assertion for a global CM reciprocity image.

For (H.1), the four diagonal characters of this torus are pairwise
distinct modulo $p$, and its residual image has order prime to $p$.
Every invariant subspace over an algebraic closure of the residue
field must therefore be a sum of CM lines.  The cyclic Galois action
on these lines is transitive, so there is no proper such subspace.
This is absolute irreducibility.  Scalar twisting by $\chi^a$ does
not change it, and the residual character of $\eta$ is trivial.

For (H.2), choose $g\in I_{13}\cap G_K$ with $\chi(g)$ a primitive
cube root $\zeta$ in $\mathbb Z_p^\times$.  This is possible because
$13$ is totally ramified in $L$ and the subgroup $H$ fixing $K$
maps nontrivially under $\chi$.  The element $g$ has
$\varepsilon_p(g)=1$, and its action on the CM lines is diagonal.
Polarization implies that its two opposite products are both $1$.
Using the torus just constructed, multiply $g$ by a product of
conjugate $p$-inertia elements so that its untwisted Tate action becomes
\[
 D=\operatorname{diag}(\zeta^{-a},1,\zeta^{a},1).
\]
The adjusting product has cyclotomic multiplier $1$.  The resulting
element $\tau$ fixes $\mu_{p^\infty}$ and has $\eta(\tau)=1$;
on the $\chi^a$-twist its eigenvalues are
\[
 1,\ \zeta^a,\ \zeta^{2a},\ \zeta^a.
\]
Exactly one is $1$, and the other three differ from $1$ by units.
Thus $T/(\tau-1)T$ is free of rank one.  Duals have the same property,
and Tate twists do not alter the action of this $\tau$.

For (H.3), use the same inertia construction for $-I_4$, whose
cyclotomic multiplier is $1$.  Denote the resulting Galois element
by $z$.  It is trivial under $\chi$ and $\eta$, and its image is
central in the joint image of the representation and
$\varepsilon_p$.  It acts by $-1$ on the residual module and its
Cartier dual.  Sah's elementary cocycle identity consequently
annihilates the two required first cohomology groups by $2$.
Explicitly, for a continuous cocycle $f$ and a central $z$,
\[
 (z-1)f(g)=(g-1)f(z).
\]
Here $z-1=-2$ is invertible, so $f$ is the coboundary of
$-f(z)/2$.  Thus both groups vanish.  This also applies to the infinite
joint image by the identical continuous-cocycle calculation.
Hypothesis (H.0) is finite freeness and finite ramification; (H.4)
holds since $p\geq5$.  The argument proves every assertion for the
further $\eta$-twists as well.
\end{proof}

\begin{proposition}[The exponent-thirteen split-prime Selmer consequence]
\label{prop:lm-thirteen-split-prime}
Let $P_{13}$ be the rational fourfold complementary to $A_B$ in the
Klein-quintic Jacobian.  The decomposition in Proposition~\ref{prop:thirteenpartial} is
\[
 V_p(P_{13})\simeq
 (V_p(A_B)\otimes\chi)\oplus(V_p(A_B)\otimes\chi^2)
\]
after adjoining the cubic character values, and that proposition proves $L(P_{13},1)\ne0$.  For all but finitely many rational primes $p$ splitting completely in
$K\mathbb Q(\zeta_3)$,
\[
 H^1_f(\mathbb Q,V_p(P_{13}))=0,
 \qquad T_p\Sha(P_{13}/\mathbb Q)=0.
\]
\end{proposition}
\begin{proof}
The CM algebra over $K$ proved above gives algebraic Hecke characters
with the prescribed CM infinity types and the Tate-module realizations
used here: \citet[\S7, Theorem~10, pp.~511--512; Theorem~11(i),
p.~512]{SerreTate1968} constructs the algebraic idelic character and
identifies its $p$-adic realizations with the Tate characters; the
construction on pp.~513--514 gives the associated complex
Gr\"ossencharacters.
The primitive CM character over $K$ then induces over its real quadratic
subfield $F$ a weight-two cuspidal Hilbert representation $\pi_0$.
This is the quadratic induction theorem of
\citet[Proposition~12.1, p.~206 in the retypeset author version]{JacquetLanglands1970}.
Its norm-factorization obstruction to cuspidality is absent because
the two exponents of our CM character at every complex place differ
by $1$.  Its induced real Weil parameter is the weight-two discrete
series parameter; the relation $|a_w-b_w|=k_v-1$ is also recorded
in \citet[Proposition~B.2.7, pp.~654--655]{Nekovar2012}.
Fix the cohomological convention of LZ: $\rho_{\pi_0,p}$ is the
cohomological representation, so its dual induces to
$W=V_p(A_B)$.  Write $\chi$ also for the automorphic character
corresponding to the chosen Galois character in this convention;
explicitly its unramified local value is the value of the Galois
character on the geometric Frobenius used in LZ's Euler polynomial.
Thus automorphic twisting by this $\chi$ twists $\rho_{\pi_0,p}$
by the Galois character $\chi$.  Define
\[
 \pi_a=\pi_0\otimes\chi^{-a},\qquad
 \operatorname{Ind}_{G_F}^{G_{\mathbb Q}}\rho_{\pi_a,p}^{\vee}
      =W\otimes\chi^a,\qquad a=1,2.
\]
This fixes the normalization rather than choosing it after applying
the theorem.  Labeling Dirichlet characters by their values on
arithmetic instead of geometric Frobenius interchanges $a=1,2$
and does not change the pair.  Pairing opposite CM lines by the
polarization shows that its determinant on the Tate side is
cyclotomic.  The finite central character of $\pi_a$ is
$\chi^{-2a}\circ N_{F/\mathbb Q}$, which satisfies the
norm-factorization hypothesis of
Proposition~\ref{prop:lm-rational-hilbert}; it need not be trivial.
The two conjugate Hilbert constituents are distinct because the
four CM characters have distinct infinity types.  Their induction
to $G_{\mathbb Q}$ is irreducible by the same CM-line argument as
in Lemma~\ref{lem:lm-thirteen-mr}.

At the chosen split primes both Hilbert ordinary conditions hold:
the rank-one CM characters are crystalline, and each opposite pair
has one exponent $0$ and one exponent $1$.  All local invariants
required in the Euler-system argument vanish, also after the
unramified $p$-power auxiliary extensions at $p$: a Frobenius power
cannot turn a Weil eigenvalue of absolute value $p^{1/2}$ or
$p^{-1/2}$ into $1$.  Lemma~\ref{lem:lm-thirteen-mr} gives exactly
the auxiliary-twist conditions.  The product formula for the
$L$-function and $L(P_{13},1)\ne0$ imply the nonvanishing of each
of the two twisted central values.  Apply
Proposition~\ref{prop:lm-rational-hilbert} to the two Hilbert forms,
$\pi_1$ and $\pi_2$, with the displayed dual identification.
Shapiro's lemma then
gives finite-Selmer vanishing over $\mathbb Q$ for both summands.
Finite extension of the coefficient field is faithfully flat,
so the asserted vanishing for $V_p(P_{13})$ follows.

The Kummer exact sequence identifies the rational Bloch--Kato
Selmer group of an abelian variety with its pro-$p$ Selmer group
tensored with $\mathbb Q_p$.  Thus the Tate module of its
Tate--Shafarevich group becomes zero after tensoring with
$\mathbb Q_p$.  That Tate module is torsion-free over $\mathbb Z_p$,
so it is zero.  This proves the claim.  It does not claim finiteness
of the entire Tate--Shafarevich group.
\end{proof}



\subsection{The exponent-nineteen quotient and its arithmetic criterion}
\label{subsec:nineteen-arithmetic}

Let $C$ be the smooth projective curve with affine equation
\[
 y^{19}=x(1-x)^7,
\]
and put $J=\operatorname{Jac}(C)$.  Write
\[
 K=\mathbb Q(\zeta_{19}),\quad H=\langle7\rangle
 \subset\operatorname{Gal}(K/\mathbb Q),\quad
 K_0=K^H,\quad F=K^+,\quad F_0=K_0^+.
\]
Thus $[K:\mathbb Q]=18$, $[K_0:\mathbb Q]=6$,
$[F:\mathbb Q]=9$, and $[F_0:\mathbb Q]=3$.

\begin{proposition}[A sufficient simple-zero certificate]
\label{prop:lm-nineteen-rankone}
There are abelian varieties $A,P$ over $\mathbb Q$ of dimensions
three and six respectively, and a rational isogeny
$J\sim_{\mathbb Q} A\times P$, such that
\[
 \operatorname{rank} A(\mathbb Q)=1,
 \qquad \operatorname{rank} P(\mathbb Q)=0.
\]
If
\[
 L'(J/\mathbb Q,1)\ne0,
\]
then $\Sha(J/\mathbb Q)$ and $\Sha(P/\mathbb Q)$ are finite.
Consequently this analytic certificate supplies a nonconstant map
from the Fermat curve of exponent nineteen to a rank-zero abelian
variety with finite Tate--Shafarevich group.
\end{proposition}

\begin{proof}
The rational automorphism
\[
 c(x,y)=\left(\frac1{1-x},-\frac{y^7}{(1-x)^3}\right)
\]
has order three.  Indeed its second iterate is
\[
 c^2(x,y)=\left(\frac{x-1}{x},
                 -\frac{y^{11}}{x(1-x)^4}\right),
\]
and its third iterate is the identity, using $y^{19}=x(1-x)^7$.
The map on the $x$-line has the two fixed points
$x^2-x+1=0$.  Above either fixed point its action on the
nineteen-element geometric fiber has multiplier $7$ in
$\mathbb F_{19}$, hence exactly one fixed point.  It cyclically
permutes the three branch points, so there are no further fixed
points.  Riemann--Hurwitz gives genus three for
$B=C/\langle c\rangle$.  Put $A=\operatorname{Jac}(B)$ and
\[
 e=\frac{1+c+c^2}{3}\in\operatorname{End}^0_{\mathbb Q}(J).
\]
The image of $e$ is isogenous to $A$, while the image of $1-e$
defines the six-dimensional complementary factor $P$.

The curve $C$ over $\mathbb Q$ is isomorphic to the curve
$F_{7,1,-8}$ in \citet[Theorem~1.6(2), author p.~3]{McCallumTzermias2004},
by interchanging $x$ and $1-x$. Their theorem proves
$\operatorname{rank}J(\mathbb Q)=1$.
On the one-dimensional space $J(\mathbb Q)\otimes\mathbb Q$,
the order-three operator $c$ must act trivially, since
$\mathbb Q^\times$ has no nonidentity element of order three.
Thus the entire rank is in the image of $e$, proving the two
rank assertions.  This does not yet assert finiteness of any
full Tate--Shafarevich group.

We identify the Hilbert form needed for that finiteness assertion.
Let $\tau(x,y)=(x,\zeta_{19}y)$.  The relation
$c\tau c^{-1}=\tau^7$ gives an action over $K$ of the crossed
product algebra
\[
 K\rtimes H\simeq\operatorname{End}_{K_0}(K)
                  \simeq M_3(K_0)
\]
on $J$.  This algebra is simple, and its homomorphism to
$\operatorname{End}^0_K(J)$ is nonzero, so it is injective.
Here the generator of $H$ acts as $c$, with $c^3=1$; thus the
crossed product has the trivial factor set.
Complex conjugation on the field of definition takes
$\tau$ to $\tau^{-1}$ and fixes $c$.  Under the displayed
identification this is the conjugation on endomorphisms induced
by the involution on $K$.  Since $K=FK_0$, taking its fixed
subalgebra gives
\[
 \operatorname{End}_{K_0}(K)^{\operatorname{Gal}(K/F)}
       =\operatorname{End}_{F_0}(F)\simeq M_3(F_0)
       \subset\operatorname{End}^0_F(J).
\]
The idempotent $e$ is the rank-one trace projector in this
matrix algebra.  Matrix units therefore give an $F$-isogeny
\[
 J_F\sim_F A_F^3.
\]
This is an isogeny over the degree-nine real field, not an
unproved rational splitting into three copies.

For completeness, the CM types and Hecke coefficients also
identify the simple factor.  The CM type of the Jacobi-sum
character is
\[
 \{1,2,3,4,6,7,9,11,14\}\subset(\mathbb Z/19\mathbb Z)^\times.
\]
Its stabilizer is exactly $H$, by \citet[Lemma~1.6, pp.~206--207]{GrossRohrlich1978}.  Consequently the induced type on
$K_0$ is primitive and $A$ is absolutely simple with CM by
$K_0$.  The Jacobi sum indexed by $(1,7,11)$ is unchanged
when all three indices are multiplied by $7$, since this
cyclically permutes the triple.  Its values therefore lie in
$K_0$.  The real quadratic-induction traces lie in $F_0$.

More explicitly, after extending the $\ell$-adic coefficient field,
write $\psi_a$ for the CM character indexed by
$a\in(\mathbb Z/19\mathbb Z)^\times$ on $G_K$.
The Jacobi-sum symmetry gives $\psi_{7a}=\psi_a$.
The three representations
\[
 W_a=\operatorname{Ind}_{G_K}^{G_F}\psi_a,
 \qquad a\in(\mathbb Z/19\mathbb Z)^\times/(\{\pm1\}H),
\]
are pairwise nonisomorphic: their restrictions to $G_K$ are
$\psi_a\oplus\psi_{-a}$, and the stabilizer of the CM type
is exactly $H$.  Each occurs three times in $V_\ell(J_F)$.
On these multiplicity spaces the operator $c$ is the regular
order-three action, so $e$ selects one copy of each.  Therefore
\[
 V_\ell(A_F)\otimes\overline{\mathbb Q}_\ell
       =\bigoplus_a W_a.
\]
This describes the three real-multiplication components explicitly.

Let $f$ be the Hilbert form over $F$ obtained by quadratic
automorphic induction of the Jacobi-sum Hecke character over
$K$.  The character is not a norm character, because its
two infinity exponents at each complex place differ by one.
Thus it gives a cuspidal parallel-weight-two form.  In unitary
normalization its restriction to $\mathbb A_F^\times$ is
$\eta_{K/F}$, so the determinant formula for induction makes
its finite central character trivial.  These are the explicit
CM and induction constructions in \citet[\S\S2.1--2.2 and 3.4, author pp.~5--6 and 15]{ShuExpanded2021}; they can equally be obtained
from \citet[Theorems~10--11]{SerreTate1968} and
\citet[Proposition~12.1]{JacquetLanglands1970}.

The standard Jacobi-sum Euler factors give
\[
 L(f/F,s)=L(J/\mathbb Q,s).
\]
Every coefficient conjugate of $f$ has the same $L$-series:
conjugating its Hecke character is the same as reindexing
the primes of the abelian field $K$, which preserves norms.
The Hecke field is $F_0$.  Indeed it is contained in $F_0$
by the trace calculation, and a proper subfield would be
$\mathbb Q$.  In that case the existence statement immediately
before Zhang's Theorem~A supplies an attached elliptic curve
$E/F$, since $[F:\mathbb Q]$ is odd.  All three coefficient
conjugates would have the same two-dimensional system
$V_\ell(E)$, so the preceding representation identity and
the isogeny criterion of \citet[\S5, Corollary~2, p.~22]{Faltings1986} would give $A_F\sim_F E^3$,
contradicting the absolute simplicity of $A$.
It follows that $A_F$ is the abelian variety attached to
$f$, with
\[
 L(A/F,s)=\prod_{\iota:F_0\hookrightarrow\mathbb R}L(f^\iota/F,s)
          =L(J/\mathbb Q,s)^3.
\]
The equality here is compatible with the preceding isogeny;
in particular one does not apply a rank-one theorem to the
degree-fifty-four $L$-function of $A/F$ itself.

\citet[Theorem~A, p.~28; proof in \S7.2, p.~143]{Zhang2001}
applies to a parallel-weight-two Hilbert newform with trivial
central character if the totally real degree is odd or a
finite local conductor has exponent one.  It states that
analytic order at most one gives finite full
Tate--Shafarevich group for the associated abelian variety.
There is no CM exclusion in that theorem.  Here $[F:\mathbb Q]=9$
is odd, so no auxiliary finite-place condition is needed.
The root number of $L(J/\mathbb Q,s)$ is $-1$, as verified
below.  The assumed nonzero derivative therefore makes
$f$ have analytic order exactly one.  Zhang's theorem gives
$\Sha(A/F)$ finite.

Finiteness of the full Tate--Shafarevich group is invariant
under isogeny over a number field.  Hence the $F$-isogeny
$J_F\sim A_F^3$ gives $\Sha(J/F)$ finite.
The kernel of restriction
$\Sha(J/\mathbb Q)\to\Sha(J/F)$ lies in
\[
 H^1(\operatorname{Gal}(F/\mathbb Q),J(F)).
\]
This group is finite: $J(F)$ is finitely generated, the
Galois group is finite, and its positive-degree cohomology
is killed by the group order.  Thus $\Sha(J/\mathbb Q)$ is
finite.  The rational isogeny $J\sim A\times P$ gives
$\Sha(P/\mathbb Q)$ finite as well.

Finally, the quotient map from the exponent-nineteen Fermat
curve to $C$, followed by a rational-cusp Abel map and a
nonzero integral multiple of $1-e$, is nonconstant into $P$.
Indeed the Abel image of $C$ generates $J$, so its image
under a surjection onto the positive-dimensional factor
$P$ generates $P$ and cannot be constant.
\end{proof}

\paragraph{Exact analytic input to be certified.}
For the triple $(r,s,t)=(1,7,11)$, the residue
\[
 r^r s^s(t-19)^t\equiv68\pmod {19^2}
\]
has Fermat quotient zero modulo $19$.  Equivalently the
sum of inverse representatives of the displayed CM type
is $95\equiv0\pmod {19}$.  \citet[Proposition~3.12(1), author p.~16]{ShuExpanded2021} gives
Hecke conductor $\mathfrak P$, and therefore the induced
rational conductor is
\[
 N=|\operatorname{disc}K|\,N_{K/\mathbb Q}\mathfrak P
   =19^{17}\cdot19=19^{18}.
\]
\citet[Theorem~3.13, author pp.~17--18]{ShuExpanded2021} gives the tame global sign
$\left(\frac2{19}\right)=-1$.
The completed function is
\[
 \Lambda(s)=N^{s/2}(2\pi)^{-9s}\Gamma(s)^9L(J/\mathbb Q,s),
 \qquad \Lambda(s)=-\Lambda(2-s).
\]
Thus a rigorous interval excluding zero for
$L'(J/\mathbb Q,1)$, or equivalently for $\Lambda'(1)$,
is sufficient. Proposition~\ref{prop:nineteen-analytic-certificate}
proves such an interval using three degree-six factors.

\paragraph{The required character order.}
The subgroup $H$ of order three lies in the kernel of every
cubic Dirichlet character modulo nineteen, since that kernel
has order six.  Characters of order nine do restrict nontrivially to $H$
and give the factorization used in the next subsection. The arithmetic
criterion itself used the crossed-product algebra and odd-degree
real base change, independently of that analytic factorization.


\subsection{Exact degree-six Euler factors for the exponent-nineteen quotient}
\label{subsec:nineteen-euler}

Retain $C:y^{19}=x(1-x)^7$, its rational order-three automorphism
$c$, and $A=\operatorname{Jac}(C/\langle c\rangle)$ from
Proposition~\ref{prop:lm-nineteen-rankone}.  Put
$K=\mathbb Q(\zeta_{19})$, $H=\langle7\rangle$,
$K_0=K^H$, and let $\sigma$ be the automorphism induced by
$\zeta_{19}\mapsto\zeta_{19}^2$.  The cyclic group
$\operatorname{Gal}(K_0/\mathbb Q)$ has order six, and its
complex conjugation is $\sigma^3$.  Write $\lambda_0$ for the
unique prime over nineteen in $K_0$.

\begin{lemma}[An exact auxiliary class number]\label{lem:nineteen-class-one}
The field $K_0$ has class number one.
\end{lemma}
\begin{proof}
Put $t=\zeta_{19}+\zeta_{19}^7+\zeta_{19}^{11}$. Its minimal
polynomial is
\[
 f(T)=T^6+T^5+2T^4-8T^3-T^2+5T+7,
\]
and the conductor--discriminant formula gives
$\operatorname{disc}K_0=-19^5$. Its Minkowski bound is
\[
 \left(\frac4\pi\right)^3\frac{6!}{6^6}\sqrt{19^5}
 <\left(\frac{200}{157}\right)^3\frac5{324}\,19^2
       \frac{4359}{1000}<51.
\]
Here $\pi>157/50$ and $\sqrt{19}<4359/1000$. The splitting
orders in the cyclic quotient of $(\mathbb Z/19\mathbb Z)^\times$
show that the prime ideals of norm at most fifty are exactly the
six over seven, the six over eleven, and the unique prime over
nineteen. Define
\[
\begin{split}
 u_7&=-\frac{18t^5+30t^4+56t^3-81t^2-72t+42}{77},\\
 u_{11}&=1-u_7,\\
 u_{19}&=\frac{-9t^5-15t^4-28t^3+79t^2-41t-21}{77}.
\end{split}
\]
Reduction modulo $f$ computes monic integral characteristic
polynomials for these three elements and gives
$N(u_7)=7$, $N(u_{11})=11$, and $N(u_{19})=19$.
Thus they are integral and generate primes of the indicated norms.
Galois conjugation generates all primes above seven and eleven;
the prime above nineteen is unique. Every ideal class therefore has
a representative which is a product of principal prime ideals,
by the displayed Minkowski bound. This proves the assertion.
All reductions, norms, and the exhaustive list of small prime ideals
are verified by the exact arithmetic certificate accompanying the paper.
\end{proof}

\begin{proposition}[The degree-six Hecke factors]\label{prop:nineteen-euler}
The threefold $A$ has CM by $K_0$, with its CM algebra defined over
$K_0$. Its algebraic Hecke character has conductor $\lambda_0$ and
CM type $\{0,1,2\}$ in the cyclic order determined by $\sigma$,
up to conjugation. Its rational conductor is $19^6$ and its root
number is $-1$. For the order-nine Dirichlet character
$\chi(2)=\zeta_9$, one has
\[
 L(J,s)=L(A,s)L(A\otimes\chi,s)L(A\otimes\bar\chi,s),
\]
where all three factors have conductor $19^6$, gamma factor
$(2\pi)^{-3s}\Gamma(s)^3$, and local Euler factor one at nineteen.
The Euler polynomials at every other prime and the exact complex
root number are given by the formulas proved below.
\end{proposition}
\begin{proof}
\par\medskip\noindent\textbf{The CM character and its normalization.} 
The endomorphisms $eae$, for $a\in K$, identify the corner
$e(K\rtimes H)e$ with $K_0$.  For $a\in K_0$ the endomorphism
$ea$ is defined over $K_0$, since the rational projector $e$
commutes with $a$ and $H$ fixes $a$.  Thus all the CM algebra
of $A$ is defined over $K_0$.  Its rank-one CM character,
denoted $\theta$, has values in $K_0$ and restricts to the
Jacobi-sum character $j$ over $K$.  This follows either by
applying the CM Hecke-character construction to $A$, or by
restricting the explicit representation of $J$ and applying
the projector $e$.  In the cyclic order determined by
$\sigma$, its type is $\{0,1,2\}$ up to rotation or inverse.
The inverse type is a rotation of this same set, so this
choice has no effect on the rational Euler polynomial.

For a prime $\ell\ne19$, let $e_\ell$ be its order modulo
nineteen and let $f_\ell$ be its order in
$(\mathbb Z/19\mathbb Z)^\times/H$.  If $\mathfrak p$ is a
prime of $K_0$ over $\ell$ and $\mathfrak P$ a prime of $K$
above it, then
\[
 \beta=\theta(\mathfrak p)\in K_0,\qquad
 j(\mathfrak P)=\beta^{e_\ell/f_\ell}.
\]
The quotient $e_\ell/f_\ell$ equals one for $e_\ell=1,2$
and three for $e_\ell=3,6,9,18$.  A cube root in $K_0$ is
unique, since $\mu(K_0)=\{\pm1\}$.

The restriction identity also proves the conductor descent,
without assuming an unstated CM conductor formula.  At
nineteen the extension $K/K_0$ is tame of degree three,
so its two absolute Galois groups have the same wild
inertia.  \citet[Proposition~3.12(1), author p.~16]{ShuExpanded2021}
gives conductor one for $j$; in particular it is trivial
on wild inertia and nontrivial on inertia.  The same two
assertions therefore hold for $\theta$, whose conductor
over $K_0$ is exactly $\lambda_0$.  Away from nineteen it
is unramified, since $A$ is a factor of the Jacobian of
a curve with good reduction away from nineteen.

We use the Jacobi convention in \citet[\S2.1, author p.~5]{ShuExpanded2021}: $j(\mathfrak P)=-\sum\chi(x)\chi(1-x)^7$.
In this convention
\[
 j(\mathfrak P)\equiv1\pmod{\lambda_K^2},
 \qquad \lambda_K=(1-\zeta_{19}).
\]
For an elementary verification, expand each character value
to first order in $\zeta_{19}-1$.  Over the residue field of
cardinality $q\equiv1\pmod{19}$, the sums of the character
exponents of $x$ and of $1-x$, with $x\ne0,1$, are divisible
by nineteen.  The constant term is $-(q-2)\equiv1\pmod{19}$.
This proves the congruence.  Stickelberger's identity and
the absolute value $|j|=\sqrt q$ are given in \citet[equation~(2.1), author p.~5]{ShuExpanded2021},
on the same page.  These three properties determine $j$:
two candidates with the same ideal and all absolute values
have root-of-unity quotient, and no nonidentity root of
unity in $K$ is congruent to one modulo $\lambda_K^2$.

\par\medskip\noindent\textbf{Primes splitting in $K_0$.} 
Suppose $\ell\bmod19\in H=\{1,7,11\}$, so $f_\ell=1$.
Take an algebraic integer $\alpha\in K_0$ of norm $\ell$;
its principal ideal is one of the six primes above $\ell$.
Set
\[
 b=\alpha\sigma(\alpha)\sigma^2(\alpha).
\]
Then $b\bar b=N_{K_0/\mathbb Q}(\alpha)=\ell$.
Replacing $\alpha$ by another generator of the same prime
changes $b$ by $\pm1$: the corresponding product of a unit
has absolute value one in every embedding, and hence is
a root of unity.  Since $\ell^3\equiv1\pmod{19}$, its
residue satisfies $b\equiv\pm\ell^2\pmod{\lambda_0}$.
Choose the unique sign such that
\[
 \beta\equiv\ell^2\pmod{\lambda_0}.
\]
For $e_\ell=1$, this gives $\beta\equiv1\pmod{\lambda_0}$;
for $e_\ell=3$, it gives $\beta^3\equiv1\pmod{\lambda_0}$.
In $K$ the valuation of $\lambda_0$ is three, so either
congruence implies the required congruence modulo
$\lambda_K^2$.  To check the ideal multiplicities exactly,
the nine-element reflex type in $K$ is the inverse image
of a three-element type in $K_0$.  When $e_\ell=1$, each
prime of $K_0$ splits into three primes of $K$, and the
nine factors in \citet[equation~(2.1), author p.~5]{ShuExpanded2021} consist of all three
primes over each of the three reflex primes of $K_0$.
Thus $(j(\mathfrak P))=(\beta)\mathcal O_K$, up to a
cyclic conjugation of the type.  When $e_\ell=3$, each
prime of $K_0$ remains inert in $K$; the three members
of $H$ fix the prime and repeat each reflex prime three
times.  Thus $(j(\mathfrak P))=(\beta^3)\mathcal O_K$.
These are precisely the two claimed ideal identities.
The normalization just proved identifies the elements,
up to the harmless cyclic conjugation.  Consequently
\[
 P_{\ell,A}(T)=N_{K_0/\mathbb Q}(1-\beta T).
\]
This formula needs only a verified norm-$\ell$ generator,
not a conditional class-group computation.  Lemma~\ref{lem:nineteen-class-one} establishes the existence
of such generators for all these primes.

\par\medskip\noindent\textbf{Order-nine primes.} 
Suppose $e_\ell=9$, so $f_\ell=3$.  The decomposition group
in $K_0$ is $\langle\sigma^2\rangle$, and $\beta$ belongs
to its fixed field $k=\mathbb Q(\sqrt{-19})$.  For example,
this follows because $\beta^3=j(\mathfrak P)$ is fixed by
that group and the cube root in $K_0$ is unique.  The
type $\{0,1,2\}$ meets the two decomposition orbits with
multiplicities two and one.  Hence, if $\mathfrak l$ is
one of the primes of $k$ above $\ell$, then
\[
 (\beta)=\mathfrak l^2\bar{\mathfrak l}
          =\ell\mathfrak l
\]
after choosing one of the conjugate labels.  The imaginary
quadratic prime $\mathfrak l$ remains inert in $K_0/k$,
and its prime in $K_0$ remains inert in $K/K_0$.
More explicitly, the nine-element reflex type in $K$
meets the two decomposition orbits with multiplicities
six and three.  \citet[equation~(2.1), author p.~5]{ShuExpanded2021} therefore gives
$(j)=\mathfrak P^6\bar{\mathfrak P}^{3}$.
Since $j=\beta^3$, descending these exponents gives
$(\beta)=\mathfrak l^2\bar{\mathfrak l}$, as claimed.
The imaginary
quadratic field $k$ has class number one: its Minkowski
bound is $(2/\pi)\sqrt{19}<3$, and two is inert, so every
ideal class has a representative of norm one.
Choose $\alpha\in\mathcal O_k$ with norm $\ell$ and set
$\beta=\pm\ell\alpha$.  Select the sign by
\[
 \beta\equiv\ell^6\pmod{\lambda_0}.
\]
This is possible and unique because $\beta^2\equiv\ell^3$
and $\ell^9\equiv1\pmod{19}$.  Its cube is congruent to
one, and the same ideal/weight/congruence argument
identifies the correct Frobenius eigenvalue.  Thus
\[
 P_{\ell,A}(T)=1-\operatorname{Tr}_{k/\mathbb Q}(\beta)T^3
                   +\ell^3T^6.
\]
For example the formula gives $1+5T^3+125T^6$ at five
and $1+119T^3+4913T^6$ at seventeen.  Each can be computed
using a quadratic norm equation; no finite field of
cardinality $\ell^9$ is required.

\par\medskip\noindent\textbf{Even-order primes.} 
If $e_\ell$ is even, complex conjugation belongs to the
decomposition group.  The Jacobi sum is therefore real,
and its absolute values and normalization give
$j(\mathfrak P)=-\ell^{e_\ell/2}$.  Taking the unique
root when necessary gives $\beta=-\ell^{f_\ell/2}$ and
\[
 P_{\ell,A}(T)=
 \begin{cases}
 (1+\ell T^2)^3,&e_\ell=2\text{ or }6,\\
 1+\ell^3T^6,&e_\ell=18.
 \end{cases}
\]
These cases and the preceding two paragraphs exhaust all
primes other than nineteen.

\par\medskip\noindent\textbf{Three small completed functions.} 
Let $\chi$ be a primitive Dirichlet character of order nine
modulo nineteen.  It restricts nontrivially to $H$, whereas
a character of order three does not.  Induction in stages
and the extensions of $\theta|_{G_K}$ to $G_{K_0}$ give
\[
 V(J)\simeq V(A)\oplus V(A)\chi\oplus V(A)\chi^2.
\]
The representation $V(A)$ has the self-twist $\chi^3$,
since this character is trivial on $G_{K_0}$.  Thus the
last summand can also be written $V(A)\chi^{-1}$, and
\[
 L(J,s)=L(A,s)L(A\otimes\chi,s)L(A\otimes\bar\chi,s).
\]
The local twisted polynomials are simply
$P_{\ell,A}(\chi(\ell)T)$.  All three functions have
conductor $19^6$, gamma factor $(2\pi)^{-3s}\Gamma(s)^3$,
and local Euler factor one at nineteen.  To see the
conductor assertion, the Jacobi character over $K$ has
conductor exponent one at nineteen and a nontrivial
quadratic restriction to inertia there, while $K/\mathbb Q$
is tame.  Thus $V(J)$ is tame with no inertia invariants.
Its three dimension-six summands have no invariants
either, so each Artin conductor exponent is six.  The
order-nine twists are even and have the same infinity
type.  The root numbers of the two conjugate twists are
complex conjugates and have product one.  Since $J$ has
root number minus one, the root number of $A$ is minus one.
It is therefore sufficient to certify
\[
 L'(A,1)\ne0,\qquad L(A\otimes\chi,1)\ne0.
\]
Their product gives the desired nonzero derivative of
$L(J,s)$ at one.  The central nonvanishing and derivative are certified in
Proposition~\ref{prop:nineteen-analytic-certificate}.

\par\medskip\noindent\textbf{An explicit root number for the complex twist.} 
There is also an exact finite Gauss sum for the root number,
so it need not be estimated from a truncated functional
equation.  Normalize $\chi(2)=\zeta_9$, consistent with
the Euler twist $P_{\ell,A}(\chi(\ell)T)$.  With
$\zeta_{19}=\exp(2\pi i/19)$, the formula is
\[
 w(A\otimes\chi)=
 \frac{i}{\sqrt{19}}\sum_{u=1}^{18}
       \left(\frac{u}{19}\right)\chi(u)^6\zeta_{19}^{6u}.
\]
Here is a derivation including the local convention.
Use the positive finite additive character in \citet[\S3.1, equation~(3.4), author p.~9]{ShuExpanded2021} and its compatible
negative infinite character.  The complex local factors
are $i^{-1}$, as recorded in \citet[\S3.3, author p.~13]{ShuExpanded2021};
the three complex places therefore contribute $i$.
The trace additive character of $K_{0,\lambda_0}$ has
conductor $\lambda_0^{-5}$, since the degree-six local
extension is tame.  The multiplicative character has
conductor $\lambda_0$, so the integration shell in \citet[equation~(3.2), author p.~9]{ShuExpanded2021} is
$19^{-1}\mathcal O_{K_0,\lambda_0}^{\times}$.

Let $\varphi$ be the unitary character associated with
$\theta$.  On rational units $u\in\{1,\ldots,18\}$,
\[
 \varphi_{\lambda_0}(u)=\left(\frac{u}{19}\right).
\]
Indeed, for positive rational $u$ prime to nineteen,
restriction and ideal norm give
$\theta((u))^3=j((u))=u^9(u/19)$.
The last equality follows from the tame quadratic local
character in \citet[equation~(3.14), author p.~16]{ShuExpanded2021}, or
directly from the ideal, absolute value and Jacobi
congruence.  Taking the unique cube root in $K_0$ gives
$\theta((u))=u^3(u/19)$, and the principal-idele relation
gives the asserted local value.  The local component of the
Dirichlet twist is the inverse of its residue character,
and the residue norm from $K_0$ is $u\mapsto u^6$.
Thus the twisted local character on these units is
$(u/19)\chi(u)^{-6}$.  Its value at nineteen is one:
all other finite local factors on this principal idele
are one, and the infinite unitary factors at the positive
rational number nineteen are one.

On the integration shell, write $x=u/19$.  The residue
additive factor is $\zeta_{19}^{6u}$, since
$\operatorname{Tr}_{K_0/\mathbb Q}(u/19)=6u/19$.
The self-dual volume of the local integer ring is
$19^{-5/2}$.  Multiplication by $19^{-1}$ changes volume
by $19^6$, while $|x|^{-1/2}=19^{-3}$; dividing into
nineteen residue classes leaves the common factor
$19^{-1/2}$.  Substitution in equation~(3.2) therefore
gives precisely the displayed sum after multiplying by
the infinite factor $i$.
For the trivial twist, $(6/19)=1$ and the quadratic Gauss
sum is $i\sqrt{19}$, recovering $w(A)=-1$.  For the
inverse twist the root number is the complex conjugate,
as follows either from the functional equation or directly
from this formula using the oddness of the quadratic
character.  These checks also fix the sign convention.

\end{proof}

\begin{lemma}[The published canonical quartic]\label{lem:nineteen-canonical}
The genus-three quotient has the following plane-quartic
model, already published as $X_{15}$ by
\citet[author pp.~2 and 22; Table~1, case~15, p.~5]{KilicerEtAl2018}.
Its identification with this Fermat quotient is the known isomorphism of
\citet[Remark~7.2, p.~12]{Tafazolian2026}:
\[
\begin{split}
0={}&a^4-a^3b+2a^3d+2a^2bd+2a^2d^2-2ab^2d\\
 &+4abd^2-b^3d+3b^2d^2+2bd^3+d^4.
\end{split}
\]
\end{lemma}
\begin{proof}
The holomorphic differentials on $C$ are
\[
 \omega_h=(1-x)^{\lfloor7h/19\rfloor}\frac{dx}{y^h},\qquad
 h\in\{5,8,10,12,13,15,16,17,18\}.
\]
Their orders at the three branch points are respectively
$18-h$, $19b+18-7h$, and $8h-19b-20$, where $b=\lfloor7h/19\rfloor$;
all are nonnegative. Direct substitution gives
$c^*\omega_h=(-1)^{b-h}\omega_{7h\bmod19}$.
The three orbit sums form the invariant basis. Thus the invariant
holomorphic differentials have common
factor $dx/y^{18}$ and coefficient functions
\[
\begin{split}
 a&=(1-x)y^{13}+(1-x)^5y^2-(1-x)^6y,\\
 b&=(1-x)^2y^{10}+(1-x)^6+(1-x)^4y^6,\\
 d&=(1-x)^3y^8-(1-x)^4y^5+(1-x)^5y^3.
\end{split}
\]
Substitution and reduction by $y^{19}=x(1-x)^7$ prove
the quartic relation.  The three differentials are linearly independent, as their coefficient
functions have disjoint powers of $y$ of degree less than nineteen.
They form the full canonical system of the genus-three quotient.
Its image lies in the displayed smooth quartic, whose smoothness
is part of the cited published model. Since a plane quartic has degree
four and the canonical divisor of a genus-three curve has degree
four, the canonical map has degree one. This identifies the quotient
with the published curve. Exact substitution and rational linear
algebra also verify that its six quadratic canonical monomials are
independent. Small-prime plane counts at seven and eleven are eighteen
and eight, agreeing with traces $-10$ and $4$ from the CM formula.
No general assertion about reduction of this particular integral
quartic is required by the coefficient algorithm.
\end{proof}

\subsection{A certified simple central zero at exponent nineteen}
\label{subsec:nineteen-certificate}

We now certify the analytic hypothesis in the rank-one criterion.
Let $A$ be the genus-three quotient Jacobian, let $J$ be the
genus-nine Jacobi quotient Jacobian, and let $\chi$ be the order-nine
Dirichlet character modulo nineteen with $\chi(2)=\exp(2\pi i/9)$.
Write
\[
 L(A,s)=\sum_{n\geq1}a_n n^{-s},\qquad
 L_\chi(s)=L(A\otimes\chi,s),\qquad
 B=\left(\frac{19}{2\pi}\right)^3.
\]
The preceding Euler-factor calculation gives
$L(J,s)=L(A,s)L_\chi(s)L_{\bar\chi}(s)$.

\begin{proposition}[The central derivative is nonzero]
\label{prop:nineteen-analytic-certificate}
For the motivic normalization with center $s=1$,
\[
 2.4378<L'(A,1)<2.4383,\qquad
 4.087<|L_\chi(1)|^2<4.093,
\]
and
\[
 9.96<L'(J/\mathbb Q,1)<9.99.
\]
In particular $L(J/\mathbb Q,s)$ has a simple zero at one.
\end{proposition}

\begin{proof}
We give both truncation estimates, the finite exact input, and the
interval evaluation. Set $N=16384$. For $x>0$ define
\[
 H(x)=\frac1{2\pi i}\int_{\Re z=2}
          \frac{\Gamma(1+z)^3}{z}x^{-z}\,dz,
 \qquad
 D(x)=\frac1{2\pi i}\int_{\Re z=2}
          \frac{\Gamma(1+z)^3}{z^2}x^{-z}\,dz.
\]
If $X_1,X_2,X_3$ are independent exponential random variables of
mean one, Mellin inversion gives
\[
 H(x)=\Pr(X_1X_2X_3\geq x),\qquad
 D(x)=\mathbb E\!\left[\log_+\frac{X_1X_2X_3}{x}\right].
\]
Indeed $\mathbb E[(X_1X_2X_3)^z]=\Gamma(1+z)^3$;
the two inverse Mellin kernels are respectively the indicator of
$t>1$ and $\log_+t$. Thus, for every positive integer $k$,
\begin{equation}\label{eq:nineteen-kernel-bounds}
 0\leq H(x)\leq\Gamma(k+1)^3x^{-k},\qquad
 0\leq D(x)\leq\frac{\Gamma(k+1)^3}{k}x^{-k}.
\end{equation}
The first inequality is Markov's inequality. For the second, use
$\log_+t\leq t^k/k$ and take expectations.

Put $w=w(A\otimes\chi)$, the exact Gauss sum already calculated,
and
\[
 F_\chi=\sum_{n\geq1}\frac{a_n\chi(n)}n H(n/B).
\]
Contour shifting in the completed functional equations gives
\begin{equation}\label{eq:nineteen-central-sums}
 L_\chi(1)=F_\chi+w\overline{F_\chi},\qquad
 L'(A,1)=2\sum_{n\geq1}\frac{a_n}{n}D(n/B).
\end{equation}
For the first identity shift the contour of
$\Lambda_\chi(1+z)/z$, where
$\Lambda_\chi(s)=B^s\Gamma(s)^3L_\chi(s)$, and use
$\Lambda_\chi(s)=w\Lambda_{\bar\chi}(2-s)$.
The residue at zero is $\Lambda_\chi(1)$.
For the second use $\Lambda_A(1+z)/z^2$ and the sign
$w(A)=-1$; the residue is $\Lambda_A'(1)$.
Since $L(A,1)=0$ and $\Gamma(1)=1$, division by $B$
gives the stated derivative identity. The Hecke completed functions
are entire; their functional equations and gamma decay justify
the shifts. The kernels also give absolute convergence of the
resulting sums.

Let $b_n$ be the number of integral ideals of norm $n$ in the
degree-six CM field $K_0$. Its Hecke character has absolute value
$\sqrt{N\mathfrak a}$ on every ideal prime to nineteen. Grouping
the ideal series by norm gives
\[
 |a_n|\leq\sqrt n\,b_n,
 \qquad \sum b_n n^{-\sigma}=\zeta_{K_0}(\sigma)
 \quad(\sigma>1).
\]
Using \eqref{eq:nineteen-kernel-bounds}, the discarded tail of
$F_\chi$ is at most
\begin{equation}\label{eq:nineteen-coefficient-tail}
 E_N=\Gamma(k+1)^3B^k
       N^{\sigma-k-1/2}\zeta_{K_0}(\sigma).
\end{equation}
Here $k+1/2>\sigma$, and the inequality follows by extracting
$N^{\sigma-k-1/2}$ from each term with $n>N$.
The discarded derivative tail is at most $2E_N/k$.
We take $\sigma=11/10$ and $k=4$.

The zeta value in this bound is itself evaluated by an exact
Dirichlet factorization. If $\eta(2)=\exp(2\pi i/18)$ modulo
nineteen, then
\[
 \zeta_{K_0}(\sigma)=\zeta(\sigma)
     \prod_{j=1}^5 L(\eta^{3j},\sigma),
\]
and each nontrivial factor is the finite Hurwitz-zeta sum
\[
 L(\eta^k,\sigma)=19^{-\sigma}
   \sum_{j=0}^{17}\exp(2\pi i k j/18)
        \zeta\!\left(\sigma,\frac{2^j\bmod19}{19}\right).
\]
The factor $\zeta(\sigma)$ includes the prime above nineteen;
no Euler factor at that prime is omitted from the majorant.
Ball arithmetic gives
\[
 \zeta_{K_0}(11/10)<6.621273,
 \qquad E_N<0.000250861,
 \qquad 2E_N/4<0.000125431.
\]

For completeness, we also bound the infinite vertical contour
after the coefficient sum has been truncated. Put
$S_N=\sum_{n\leq N}|a_n|/n^3$ and $T=24$.
The recurrence for the gamma function and its Euler integral give
\[
 |\Gamma(3+it)|
   \leq\frac{\Gamma(3+m)}{|t|^m}\qquad(t\ne0).
\]
Consequently the discarded part of the half-sum contour for
$F_\chi$ is at most
\[
 V_H=\frac{B^2S_N\Gamma(3+m)^3}
             {\pi(3m)T^{3m}},
\]
while the discarded part of the derivative contour is at most
\[
 V_D=\frac{2B^2S_N\Gamma(3+m)^3}
             {\pi(3m+1)T^{3m+1}}.
\]
These estimates integrate $t^{-3m-1}$ and $t^{-3m-2}$,
respectively, on the two tails. With $m=16$ they give
$V_H<7.726\cdot10^{-19}$ and
$V_D<6.307\cdot10^{-20}$.

The finite input consists of all Euler factors needed through
$N$. Exactly 312 primes split completely in $K_0$ below this
cutoff. For each, an integral-basis vector of norm equal to the
prime is checked exactly; its three-conjugate product is checked
to have the correct Weil norm and residue sign, and its
characteristic polynomial supplies the Euler factor. This replay
uses only the number field and its maximal order, without a
class-group search or a class-number hypothesis. The remaining
residue-order cases use the explicit factors proved above and
the quadratic norm equation in $\mathbb Q(\sqrt{-19})$.
Inversion of these polynomials and multiplicativity give 468
nonzero coefficients through $N$. Each is an integer, and every
nonzero index is congruent to $1$, $7$ or $11$ modulo nineteen.

Here is the finite interval computation. For $z=2+it$ define
\[
 G(t)=B^z\Gamma(1+z)^3
       \sum_{n\leq N}\frac{a_n}{n^{1+z}}.
\]
The finite derivative is
$2\pi^{-1}\Re\int_0^T G(t)z^{-2}\,dt$.
For the complex half-sum, replace $a_n$ by $a_n\chi(n)$,
add its value at $t$ and at $-t$, and integrate with the
factor $(2\pi)^{-1}z^{-1}$ on the corresponding two terms.
This symmetric expression is meromorphic in the integration
variable and holomorphic near the real integration segments.
It does not apply complex conjugation to that variable
inside the interval integrator. The 24 unit intervals are
integrated at 128-bit precision by Arb, with requested absolute
and relative error tolerances $10^{-13}$. The returned enclosing
balls, rather than the requested tolerance alone, are used.
They give
\[
 \begin{split}
 D_N&\in[2.438030315780\mathbin\pm3.23\cdot10^{-13}],\\
 F_{\chi,N}&\in[0.9278931586\mathbin\pm3.40\cdot10^{-11}]\\
   &\hspace{12mm}+i[0.8339923732\mathbin\pm4.83\cdot10^{-11}].
 \end{split}
\]
The exact root number has enclosure
\[
 w\in[-0.9107472158164717387\mathbin\pm3\cdot10^{-20}]
       +i[0.4129642949245673588\mathbin\pm3\cdot10^{-20}].
\]
Adding the bounds $V_H+E_N$ to both coordinates of the
half-sum, adding $V_D+2E_N/4$ to the derivative, and then
using \eqref{eq:nineteen-central-sums} gives
\[
 \begin{split}
 L'(A,1)&\in[2.438\mathbin\pm0.000156],\\
 L_\chi(1)&\in[0.427\mathbin\pm0.000810]
                   +i[1.977\mathbin\pm0.000848],\\
 |L_\chi(1)|^2&\in[4.09\mathbin\pm0.00281],\\
 L'(A,1)|L_\chi(1)|^2&\in[9.97\mathbin\pm0.00891].
 \end{split}
\]
Finally, $L_{\bar\chi}(1)=\overline{L_\chi(1)}$ and
$L(A,1)=0$, so differentiating the factorization gives
$L'(J,1)=L'(A,1)|L_\chi(1)|^2$. The displayed enclosures
prove the inequalities and the simple zero.
\end{proof}

\begin{corollary}[The exponent-nineteen descent target]
\label{cor:nineteen-finite-sha}
The sixfold $P$ in Proposition~\ref{prop:lm-nineteen-rankone}
has rational rank zero and finite Tate--Shafarevich group. The groups
$\Sha(J/\Q)$ and $\Sha(A/\Q)$ are finite as well. Its generating
morphism $X_{19}\to P$ gives
$X_n(\A_\Q)_\bullet^{\mathrm{f\text{-}ab}}=X_n(\Q)$ for every
positive multiple $n$ of $19$.
\end{corollary}
\begin{proof}
Proposition~\ref{prop:nineteen-analytic-certificate} verifies the
analytic hypothesis of Proposition~\ref{prop:lm-nineteen-rankone}.
The latter gives full Tate--Shafarevich finiteness and a rational
rank-zero sixfold. Its proof constructs a map from $X_{19}$ whose
image generates $P$ by differences. The coordinate-power map
$X_n\to X_{19}$ therefore gives a nonconstant map to the same
rank-zero target with finite Tate--Shafarevich group.
\citet[Theorem~8.6, p.~385]{Stoll2007}, with its real-place erratum,
applies. This also proves Theorem~\ref{thm:nineteen-selmer-intro}.
\end{proof}


\begin{proof}[Finite-descent applications at exponents \(3,5,7,11,13,19\)]
First suppose $5\mid n$. The power map $X_n\to X_5$, followed by $f_Z$ and the Abel--Jacobi
embedding $C\to A$, is nonconstant. Stoll's Theorem 8.6 therefore gives
\eqref{eq:stoll} for $X_n$ by Proposition~\ref{prop:arithmetic}.
The inverse images under $f_Z$ of the three points in \eqref{eq:cpoints} are,
respectively, the divisors $Z=0$, $X=0$, and $Y=0$. Their inverse
images under the power map are precisely the coordinate divisor of
$X_n$. Thus every rational point of $X_n$ is a cusp.

For $3\mid n$, use the power map to the Fermat cubic $X_3$. Its
elliptic model is $E_{27}:v^2+v=u^3-7$, equivalently
$y^2=x^3-432$. The latter is the classical modular curve $X_0(27)$;
see \citet[\S10.1]{KrizLi2019}. Stoll's Corollary 8.8, together with
the rank-zero modular factor used in its proof and Corollary 6.3,
gives the required rank-zero target and finite-abelian-descent
equality. Since $X_0(27)$ has genus one, that nonzero modular factor
is the elliptic curve itself up to isogeny. Its rational group has
order three: good reduction gives orders $6$ at $5$ and $9$ at $7$,
while the three rational Fermat cusps already give three points.
Thus $X_3(\Q)$ consists of those cusps. The power map has no rational
point above a cusp except a coordinate cusp.

For $7\mid n$, use the known composite
\[
X_7\longrightarrow
\{a^3b+b^3c+c^3a=0\}\longrightarrow E_{49}\simeq X_0(49),
\qquad
[X:Y:Z]\longmapsto[X^3Z:Y^3X:Z^3Y].
\]
The second map is the quotient by cyclic coordinate permutation.
Its elliptic equation is $y^2=4x^3+21x^2+28x$;
\citet[\S2, equation (2.10), pp.~68--69; \S3.1--3.2, pp.~72--74]{Elkies1998}
gives these maps, the two rational points on $E_{49}$, and the
deduction that $X_7(\Q)$ consists precisely of its coordinate cusps.
Again Stoll's modular corollary gives a nonzero rank-zero factor,
necessarily the elliptic target up to isogeny, with vanishing
divisible part of $\Sha$. The composite $X_n\to X_7\to E_{49}$
therefore satisfies Stoll's Theorem 8.6, and rational points of
$X_n$ lie above the listed cusps of $X_7$.

For $11\mid n$, compose the power map $X_n\to X_{11}$ with
$f_{11}$ and the Abel--Jacobi embedding into $A_{11}$.
Proposition~\ref{prop:eleven} gives the required finite target and
finite $\Sha$. The three rational points of $C_{11}$ have preimages
respectively on $Z=0$, $X=0$, and $Y=0$, so the same argument forces
all rational points of $X_n$ to be coordinate cusps.

For $13\mid n$, compose the power map with the morphism
$X_{13}\to P_{13}$ of Proposition~\ref{prop:thirteenpartial}.
Its image generates this fourfold by differences, so it is nonconstant
and cannot be a torsion translate of an elliptic subvariety.
Proposition~\ref{prop:lm-thirteen-split-prime}, the finite rational
group computed in Proposition~\ref{prop:thirteenpartial}, and
Theorem~\ref{thm:split-prime-descent} give finite-abelian-descent
excellence. These assertions also prove
Theorem~\ref{thm:thirteen-selmer-intro}, since the power map is
surjective for every positive multiple of $13$.
The separate rational-point assertion
$X_{13}(\Q)\subset\{XYZ=0\}$ is
\citet[Theorem~0.5, p.~448]{Wiles1995}. The power map therefore forces
all rational points of $X_n$ to be coordinate cusps.

For $19\mid n$, use the generating map $X_{19}\to P$ of
Corollary~\ref{cor:nineteen-finite-sha}, preceded by the power map.
Its target has rational rank zero and finite Tate--Shafarevich group,
so \citet[Theorem~8.6, p.~385]{Stoll2007} gives
finite-abelian-descent excellence. The rational-point assertion
for $X_{19}$ is again \citet[Theorem~0.5, p.~448]{Wiles1995}, so the
rational points upstairs are coordinate cusps.

For odd $n$, the equation $t^n=-1$ has only the rational solution
$-1$, giving exactly the three listed centers; Theorem~\ref{thm:transfer}
applies. For even $n$, the real point set is empty.
\end{proof}

\subsection{The septic proper quotient}

The same CM arithmetic supplies the additional proper quotient needed
for the local statement at $3$.
Put $C_7:v^2+v=u^7$ and $A_7=\operatorname{Jac}(C_7)$.
\begin{proposition}\label{prop:septic}
The group $A_7(\mathbb Q)$ has order seven, and
$\Sha(\mathbb Q,A_7)$ is finite. In particular the $3$-adic Selmer group
of $T_3A_7$ is zero.
\end{proposition}
\begin{proof}
Use $(r,s,t,\delta)=(1,1,5,1)$ in Shu's Jacobi-sum description
\citep[\S\S2.1--2.2, pp.~5--6; Proposition 3.12(2), p.~16;
Theorem 3.13, pp.~17--18]{ShuExpanded2021}.
For $a=(-2)^5$, the integer $(a^6-1)/7$ is congruent to $3$ modulo
$7$. The Hecke character over $K=\mathbb Q(\zeta_7)$ therefore has
conductor $(1-\zeta_7)^2$. Its conductor over $\mathbb Q$ is $7^7$.
The finite root number at seven is $-\left(15/7\right)=-1$, while
its archimedean root number is $\left(-1/7\right)=-1$; all other
local signs are one. Thus the global sign is $+1$, and the completed
function is
\[
 \Lambda_7(s)=7^{7s/2}(2\pi)^{-3s}\Gamma(s)^3L(A_7,s),
 \qquad \Lambda_7(s)=\Lambda_7(2-s).
\]
These are algebraic Hecke-function inputs, not conclusions drawn from
numerical evaluations.

Here are the precise data and inequalities for central nonvanishing.
Write $B=7^{7/2}/(2\pi)^3$ and $L(A_7,s)=\sum a_nn^{-s}$.
The Mellin-integral derivation used for the degree-eleven quotient,
with three gamma factors, gives
\[
 L(A_7,1)=2\sum_{n\ge1}\frac{a_n}{n}Q_3(n/B),\qquad
 Q_3(x)=\frac1{2\pi i}\int_{(2)}
                   \frac{\Gamma(1+z)^3}{z}x^{-z}\,dz.
\]
Equivalently $Q_3(x)$ is the probability that the product of three
independent exponential random variables of mean one exceeds $x$.
For every positive integer $k$,
$Q_3(x)\le\Gamma(k+1)^3x^{-k}$ by Markov's inequality.
The Jacobi-sum Euler product gives
$|a_n|\le \sqrt n\,b_K(n)$, where
$\zeta_K(s)=\sum b_K(n)n^{-s}$.
Thus for $\sigma>1$ and $k+1/2>\sigma$ the omitted tail after $N$ is
at most
\[
 2\Gamma(k+1)^3 B^k N^{\sigma-k-1/2}\zeta_K(\sigma).
\]

All coefficients through $N=4096$ are computed exactly. At a prime
$p\ne7$ of even multiplicative order $f$ modulo seven, the polynomial
is $(1+p^{f/2}T^f)^{6/f}$, by
\citet[Lemma 1.1, p.~203]{GrossRohrlich1978}.
At residue degree three only $p=2,11$ require nonzero coefficients
below this cutoff. Direct counts over $\mathbb F_{p^3}$ give
\[
 P_2(T)=1-2T^3+8T^6,\qquad
 P_{11}(T)=1-44T^3+1331T^6.
\]
The first count is three points over $\mathbb F_8$; the second is
$1200$ points over $\mathbb F_{1331}$. At split primes the program
counts the seven character-value classes of the Jacobi sum, and
Newton identities give the required initial coefficients of its norm
polynomial. The factor at seven is one. Multiplicativity gives the
entire coefficient list; it has $124$ nonzero terms. For example,
$P_{29}(T)=1+T-27T^2+\cdots$.

For completeness, on the line $z=2+it$ the finite central sum is
\[
 \frac2\pi\operatorname{Re}\int_0^{20}
 \frac{\Gamma(3+it)^3}{2+it} B^{2+it}
 \sum_{n\le4096}\frac{a_n}{n^{3+it}}\,dt
\]
up to its vertical tail. Iterating the gamma recurrence seventeen
times and using the integral bound on the numerator gamma function
bounds that tail by
\[
 \frac2\pi B^2\left(\sum_{n\le4096}\frac{|a_n|}{n^3}\right)
 \frac{\Gamma(20)^3}{51\cdot20^{51}}.
\]
Arb arithmetic at $112$ bits encloses the displayed finite integral
by $0.97638934177248\pm5.84\cdot10^{-15}$ and the vertical tail by
$1.342\cdot10^{-16}$. With $k=4$ and $\sigma=11/10$, the Dirichlet
tail is less than $9.049\cdot10^{-6}$; the required zeta value is
$\zeta_K(11/10)=3.49720060785743449\ldots$, evaluated as the product
of the Riemann zeta factor and the five nonprincipal Dirichlet
$L$-functions of conductor seven. Consequently
\[
                 0.9763<L(A_7,1)<0.9765.
\]

Over the totally real field $F=\mathbb Q(\zeta_7+\zeta_7^{-1})$,
Shu's construction realizes the three two-dimensional coefficient
components of $A_7$ as parallel-weight-two Hilbert forms with trivial
central character: the restriction of the Hecke character to the real
idele group is the quadratic character of $K/F$, so automorphic
induction has central character equal to its square
\citep[\S3.4, p.~15]{ShuExpanded2021}. Each conjugate Jacobi Hecke function has the same
central value, since the conjugation reindexes the prime ideals of
$K$. The degree $[F:\mathbb Q]=3$ is odd, so
\citet[Theorem B and B-prime, p.~593; \S2.10.5, p.~641]{Nekovar2012}
applies without the alternative nonprincipal-local-component condition.
As in the preceding arithmetic arguments, it gives finite Mordell--Weil
group and finite whole Tate--Shafarevich group over $F$; restriction
then gives both finiteness assertions over $\mathbb Q$.
Finally \citet[Theorem 1.1, p.~202]{GrossRohrlich1978} identifies the
rational torsion as $\mathbb Z/7\mathbb Z$. Its exceptional septic
case does not occur: $1^3=1$ while $5^3=-1$ modulo seven.
Hence the rational group has order seven. Its $3$-adic completion
and the $3$-adic Tate module of its finite Tate--Shafarevich group both
vanish, proving the last assertion.
\end{proof}

\section{The local arithmetic source and a one-place criterion}
\label{sec:local-arithmetic-source}

We establish the arithmetic and local inputs for the unrestricted
pro-$3$ theorem. They also give a separate criterion that assumes
local geometry only at $3$ and a ramification condition on the
\emph{whole section cocycle}. That criterion applies to every
Galois-stable quotient retaining geometric abelianization.

Put $K=\Q(\zeta_3)$ and $\lambda=1-\zeta_3$. For a finite set $S$ of
rational primes containing $3$, let $K_S(3)$ be the maximal pro-$3$
extension of $K$ unramified outside the places above $S$, and set
\[
 P_S=\Gal(K_S(3)/K),\qquad
 \Gamma_S=\Gal(K_S(3)/\Q).
\]
The second extension is Galois because its defining ramification
condition is invariant under $\Gal(K/\Q)$. We shall use the condition
\begin{equation}
 H^1(P_S,\F_3)\longrightarrow H^1(K_\lambda,\F_3)
 \quad\hbox{is injective}.
 \label{eq:sourcecriterion}
\end{equation}

\begin{lemma}[Local source surjectivity]\label{lem:source}
If \eqref{eq:sourcecriterion} holds, the natural homomorphism
$G_{\Q_3}\to\Gamma_S$ is surjective. Condition
\eqref{eq:sourcecriterion} holds for
\[
 S=\{3\},\qquad S=\{3,7\},\qquad
 S=\{3,q\}\quad(q\text{ prime},\ q\equiv2\text{ or }5\pmod9).
\]
\end{lemma}
\begin{proof}
The image of $G_{K_\lambda}$ in $P_S$ is closed. If it were proper,
some finite $3$-group quotient would have proper image; a maximal
subgroup of that quotient has index three. The resulting nonzero
character $P_S\to\F_3$ would restrict to zero locally, contradicting
\eqref{eq:sourcecriterion}. Hence this local source surjects onto
$P_S$. The extension $K_\lambda/\Q_3$ has degree two, so the image of
$G_{\Q_3}$ also surjects onto $\Gamma_S/P_S$. It is all of $\Gamma_S$.

The imaginary quadratic field $K$ has class number one: its
Minkowski bound is $(2/\pi)\sqrt3<2$. Consequently Kummer theory
identifies $H^1(P_S,\F_3)$ with the group of $S$-units modulo cubes;
an element whose valuations outside $S$ are multiples of three can
have those valuations removed using principal ideals. Locally let
$U_i=1+\lambda^i\mathcal O_{K_\lambda}$. The groups of roots of unity,
the valuation, and the logarithm give
\[
 K_\lambda^\times/(K_\lambda^\times)^3
 =\langle\lambda\rangle\oplus\langle\zeta_3\rangle
   \oplus U_2/U_4.
\]
Here each of the first two factors has order three, and $U_2/U_4$
has dimension two over $\F_3$. Indeed $U_1=\mu_3\times U_2$, and
the logarithm identifies cubing on $U_2$ with multiplication by $3$,
whose $\lambda$-valuation is two. Thus $U_2^3=U_4$.

For $S=\{3\}$, the global basis is $\zeta_3,\lambda$, visibly
independent locally. If $q\equiv2$ or $5\pmod9$, then $q$ is inert
in $K$ and the global basis is $\zeta_3,\lambda,q$. Write
$c=(q+1)/3$, which is a $3$-adic unit. Since
$3=-\lambda^2/\zeta_3$,
\[
 -q=1-3c\equiv1+c\lambda^2+c\lambda^3\pmod{\lambda^4}.
\]
Its $U_2/U_4$ coordinate is nonzero, proving independence.
Finally, at $7$ put $\pi=2-\zeta_3=1+\lambda$; its norm is seven.
The global basis is $\zeta_3,\lambda,\pi,\bar\pi$.
Replacing the last element by $7=\pi\bar\pi$ preserves that basis.
The exact identity $\zeta_3\pi=1-\lambda^2$ and the congruence
\[
7=1-2\lambda^2/\zeta_3
 \equiv1+\lambda^2+\lambda^3\pmod{\lambda^4}
\]
give the two independent rows $(-1,0)$ and $(1,1)$ in $U_2/U_4$.
This proves \eqref{eq:sourcecriterion} in all the stated cases.
\end{proof}

\begin{lemma}[The compact cubic at three]\label{lem:cubiclocal}
For $E:v^2+v=u^3-7$, the group $E(\Q_3)$ is pro-$3$. Its
$3$-adic Kummer map is therefore injective. A global class in
$H^1(\Q,T_3E)$ that is Kummer at $3$ is the Kummer class of one of
the three points of $E(\Q)$.
\end{lemma}
\begin{proof}
The integral equation has discriminant $-3^9$, and is minimal since
its valuation is less than twelve. Its reduction has unique singular
point $(0,1)$ and three smooth $\F_3$-points, counting infinity.
Let $E_0$ be the inverse image of the smooth locus and $E_1$ the
formal subgroup. Thus $E_1$ is pro-$3$ and $E_0/E_1$ has order three.
The point $P=(3,4)$ has order three: its tangent has slope three,
and the addition formula gives $2P=(3,-5)=-P$.

We verify that every point with singular reduction belongs to
$P+E_0$ or $-P+E_0$. Such a point has integral coordinates
$u=3a,v=1+3b$. Its equation implies $b^2+b+1=3a^3$, so
$b=1+3w$. It follows that $a^3=1+3w+3w^2$. Cubes modulo nine
are $0,1,-1$, which forces $w\equiv0$ or $-1\pmod3$ and
$a\equiv1\pmod3$. Hence write
\[
 u=3+9z,\quad v=4+9w,\qquad
 w^2+w=3z+9z^2+9z^3.
\]
If $w\equiv0\pmod3$ and $z\ne0$, the slope in $Q-P$ is
$(w+1)/z$. If $z$ is a unit, the new $u$-coordinate is $1$ modulo
three and the reduction is smooth. If $z$ is not a unit, the new
coordinate has negative valuation and the sum reduces to infinity.
For $z=0$ this case gives $Q=P$. Thus $Q-P\in E_0$.
If $w\equiv-1\pmod3$, apply the same reasoning to $Q+P$,
whose slope is $w/z$; the case $z=0$ is $Q=-P$.
Consequently $E/E_0$ is generated by the nonzero class of $P$ and
has order three. The quotient $E/E_1$ has order nine, proving the
pro-$3$ assertion. The intersection of all $3^mE(\Q_3)$ is zero,
so the local Kummer map is injective.

The modular argument in the proof of Theorem~\ref{thm:regular-solvable-intro}
gives $E(\Q)\simeq\Z/3\Z$ and $T_3\Sha(\Q,E)=0$. At a finite
prime $v\ne3$, $E$ has good reduction and every class in
$H^1(\Q_v,T_3E)$ is Kummer. For example, the unramified cohomology
sequence has quotient $(T_3E(-1))^{\operatorname{Frob}_v}=0$ by
the Weil bounds, and its unramified term is the Kummer group of
the finite reduction. At the real place the cohomology is zero
because two is invertible on $T_3E$. Thus the hypothesized global
class lies in the inverse-limit $3$-Selmer group. The Kummer exact
sequence and $T_3\Sha=0$ identify that group with $E(\Q)$.
\end{proof}

\begin{lemma}[Action and change of tangent]\label{lem:action}
Let $n=3^r$, $r\geq1$, and put
$N=\pi_1((U_n)_{\bar\Q})^{(3)}$. Choose at a rational coordinate
cusp the tangent supplied by a coordinate parameter with leading
coefficient one. The resulting actual action on $N$ is unramified
outside $3$, and its restriction to $G_K$ has pro-$3$ image.
The relative cocycle between any two such standard rational
tangents is also unramified outside $3$.
\end{lemma}
\begin{proof}
The plane Fermat model and its coordinate boundary form a smooth
proper curve with an \emph{\'{e}tale} boundary over $\Z[1/3]$.
The chosen tangents extend as nonzero vectors there. At each
prime $v\ne3$, tame specialization identifies the pro-$3$
geometric covering category, including these tangential fiber
functors, with that of the geometric special fiber
\citep[Expos\'e XIII, \S2]{SGA1}. Inertia acts trivially on the
latter category and its fiber-functor isomorphisms. Therefore it
acts trivially on the actual based fundamental group and on every
path torsor between the stated tangents. This gives both
unramifiedness assertions, including the change-of-origin cocycle.
The integrality and nonvanishing of the tangent are essential in
this application of specialization.

There is also a direct pointed description using the triply punctured
line $B=\mathbb P^1\setminus\{0,-1,\infty\}$ with coordinate $s$.
The cover $U_n\to B$ is given by
\[
s=x^n,\qquad(-y)^n=1+s.
\]
If $F_0=\pi_1(B_{\bar\Q})^{(3)}$ with the unit tangent at zero,
then $F_0^{\ab}=\Z_3(1)^2$ and
\[
N_n=\ker(F_0\longrightarrow F_0^{\ab}/nF_0^{\ab}).
\]
This is an open characteristic subgroup. The equality with the full
pro-$3$ group of $U_n$ follows by taking the normal closure of a finite
$3$-cover under the $3$-group of deck transformations; its monodromy
is still a $3$-group. At $Q_0$ the explicit fiber functor has
$x=t$, $-y=(1+t^n)^{1/n}$ and $s=t^n$. Coefficientwise Galois action
on rational Puiseux series makes the arithmetic splittings compatible.
The other two standard cusps have the analogous coordinate description.
Thus the actual action on $N_n$ is the restriction of that on $F_0$,
and the inclusions $N_n\subset N_3\subset F_0$ are Galois stable.

On $F_0/\Phi(F_0)$ the action is scalar cyclotomic and is trivial
on $G_K$. The kernel of
$\operatorname{Aut}(F_0)\to\operatorname{Aut}(F_0/\Phi(F_0))$
is pro-$3$, by the finite $3$-group automorphism theorem and passage
to characteristic lower $3$-central quotients. Hence the actual
$G_K$-image on $F_0$, and on $N_n$, is pro-$3$.
For this punctured-line description the pointed unramification and
path-torsor statements above also follow directly from
\citet[Theorem 17.7 and Corollary 17.8, p.~51]{Wojtkowiak2009}.
\end{proof}

\Needspace{8\baselineskip}
\begin{theorem}[Restricted pro-$3$ Fermat sections]\label{thm:restricted3}
Let $n=3^r$, $r\geq1$, and fix one of the standard integral
tangents in Lemma~\ref{lem:action}. Let $M=N/J$, where $J$ is a
closed Galois-stable normal subgroup of $N$ contained in $[N,N]$;
the choice $J=1$ is permitted. Suppose $S$ is a finite set
containing $3$ and satisfies \eqref{eq:sourcecriterion}.

If a section of $M\rtimes G_\Q$ has its entire relative cocycle
trivial on inertia outside $S$ and its localization at $3$ is
point-or-cusp geometric, then it is globally cuspidal.
Within this ramification class its sections form exactly three
disjoint packets, one at each rational coordinate cusp, each
parametrized by
\begin{equation}
 H^1(G_{\Q,S},\Z_3(1))
    =\bigoplus_{q\in S}\Z_3\,\kappa_3(q).
 \label{eq:restrictedpackets}
\end{equation}
In particular this holds for all the sets $S$ in
Lemma~\ref{lem:source}, proving Theorem~\ref{thm:restrictedintro}.
No local condition away from $3$ is needed
for the implication to cuspidality.
\end{theorem}
\begin{proof}
Push the section to the compact cubic using $X_n\to X_3\simeq E$.
This is possible through $M$ since that map has abelian geometric
target and $J\subset[N,N]$. Lemma~\ref{lem:cubiclocal} identifies
the resulting global compact class with a point of $E(\Q)$.
These three points are its coordinate cusps. The local center at
$3$ maps to precisely that point, since the compact $3$-adic
Kummer map is injective. Its inverse image in $X_n$ is supported
on the coordinate boundary. Moreover, $\Q_3$ has no nontrivial
odd-order roots of unity: its roots of unity are $\{1,-1\}$.
Thus this center is one of the three rational cusps of $X_n$,
say $Q$. The local section must be cuspidal at $Q$.

Let $\alpha$ be the actual action and $z$ the section cocycle.
Over $G_K$ the graph homomorphism
$g\mapsto(z(g),\alpha(g))$ has pro-$3$ image by
Lemma~\ref{lem:action}, and is unramified outside $S$. It factors
through $P_S$ by that group's defining maximality. Consequently
both $\alpha$ and $z$ factor through $\Gamma_S$. Replace the
origin by the standard tangent at $Q$. Its change-of-origin
cocycle is unramified outside $3$ by Lemma~\ref{lem:action}, so
the new action and cocycle still factor through $\Gamma_S$.

Choose a single element of $M$ conjugating the local section at
$3$ into the cusp group at $Q$, and conjugate the global section
by that same element. Its new cocycle is $I_Q$-valued on
$G_{\Q_3}$, where $I_Q\simeq\Z_3(1)$ denotes the peripheral
subgroup in $M$. Lemma~\ref{lem:source} says this local source
surjects onto $\Gamma_S$. Every value of the global cocycle is
therefore already a value occurring on the local source, and
lies in $I_Q$. This is exactly a global cusp section at $Q$.

For completeness the peripheral subgroup embeds in $M$ because
it embeds already in $N^{\ab}$ by the single boundary relation.
Its parameters are unique: filling $Q$ gives the pro-$3$ analogue
of \eqref{eq:boundary}, and the Tate module of the resulting
generalized Jacobian has zero invariants globally and over
$\Q_3$. Thus the peripheral maps on $H^1$ are injective. Distinct
centers are separated by the compact cubic Kummer classes at $3$.
Since the standard tangents are unramified outside $3$, a global
cusp parameter is unramified outside $S$ exactly when its
valuations away from $S$ vanish. Kummer theory over
$\Z[S^{-1}]$, whose ideal class group is zero, gives
\eqref{eq:restrictedpackets}. Every such cusp section satisfies
all local point-or-cusp conditions. This proves the classification.
\end{proof}

The whole-cocycle ramification hypothesis is essential to this
one-place argument. The unrestricted Selmer theorem below instead
applies source surjectivity to successive additive central differences;
it permits cusp parameters with infinite support.

\section{Unrestricted pro-\texorpdfstring{$3$}{3} Selmer cuspidality}
\subsection{Compact cancellation and the unique abelian cusp parameter}
\label{subsec:lm-unrestricted-reduction}

Let $n=3^r$ with $r\ge1$, put $J_n=\operatorname{Jac}(X_n)$, and write
$N_n=\pi_1(U_{n,\overline\Q})^{\mathrm{pro}\text{-}3}$.
At a rational coordinate cusp $P$, let $t_P$ be the standard tangent
section with leading coefficient one.  The only compact arithmetic
input is the cubic statement in Lemma~\ref{lem:cubiclocal}:
$E=J_3$ has $E(\Q)=\Z/3\Z$, $T_3\Sha(E/\Q)=0$, and
$E(\Q_3)$ is pro-$3$.  The higher Jacobian $J_n$ requires no rank
or Tate--Shafarevich hypothesis.  Instead, arithmetic localization
will kill its compact difference from a cusp.

\begin{lemma}[The inherited arithmetic action and its localization]
\label{lem:lm-power-action}
For a standard integral tangent at any rational cusp of $U_n$,
the actual action of $G_{\mathbb Q}$ on $N_n$ factors through
\[
 \Gamma=\operatorname{Gal}(K_{\{3\}}(3)/\mathbb Q),
 \qquad K=\mathbb Q(\zeta_3).
\]
Its restriction to $G_K$ has pro-$3$ image, and the local map
$G_{\mathbb Q_3}\to\Gamma$ is surjective.
If $M$ is a finite free $\mathbb Z_3$-module with an action
obtained from this action by subquotients and tensor operations,
then an unramified-outside-$3$ class
$a\in H^1(\mathbb Q,M)$ with $a_3=0$ is zero.
\end{lemma}
\begin{proof}
Write $e=3^{r-1}$.  The coordinate-power cover
$U_n\to U_3$ is geometrically Galois with group $\mu_e^2$,
of $3$-power order.  Its geometric pro-$3$ group is therefore
an open subgroup $H=N_n$ of $N_3$, and its arithmetic
pro-$3$ group is the corresponding open arithmetic subgroup.
At the cusp $(0:-1:1)$ choose the local parameters
$u_n=X/Z$ and $u_3=X/Z$ on the respective curves; the cover
satisfies $u_3=u_n^e$.  Unit tangential basepoints are thus
compatible.  The same construction applies at the other
rational cusps by coordinate permutations.

The actual upstairs action is the restriction to $H$ of the
standard downstairs action, which has already been proved to
factor through $\Gamma$ and to have pro-$3$ image over $G_K$.
The same statements follow upstairs.  The surjection
$G_{\mathbb Q_3}\to\Gamma$ is the established arithmetic
source calculation for $K$ and the single prime over $3$.
Neither claim requires taking a derived quotient of the
open embedding.

For the cohomological assertion, choose a continuous cocycle
$z$ representing $a$.  Since the action on $M$ is unramified
outside $3$, unramifiedness of the class means that this
cocycle itself is zero on all arithmetic inertia groups
outside $3$.  Its graph over $G_K$ takes values in
$M\rtimes\rho(G_K)$, a pro-$3$ group.  Consequently both
the action and the cocycle factor through $\Gamma$.
The equality $a_3=0$ permits subtracting a single coboundary
with an element of $M$ so that the new cocycle is zero on
$G_{\mathbb Q_3}$.  It still factors through $\Gamma$.
Surjectivity of the local map makes it zero globally.
\end{proof}

This also proves the compact localization assertion of
Theorem~\ref{thm:unrestrictedthreefamiliesintro} for $n=3^r$.
Indeed, a class in $\Sel(T_3J_n)$ is unramified outside three:
at every other prime the Jacobian has good reduction and its local
completed Kummer classes are unramified. Apply
Lemma~\ref{lem:lm-power-action} to the quotient $T_3J_n$ of
$N_n^{\mathrm{ab}}$. A class whose localization at three is zero
is therefore zero globally.

\begin{lemma}[Inertia of a local geometric section]
\label{lem:lm-local-inertia}
Let $p\ne3$, and let $s_p$ be the image of a full local point or cusp
section in the geometrically pro-$3$ fundamental group.  Let $J_p$ be
arithmetic inertia and let $\tau_p:J_p\to\Z_3(1)$ be the tame
$3$-adic character.  Relative to an extending tangential origin, its
geometric cocycle has the following restriction, up to geometric conjugation:
\begin{enumerate}
\item If an ordinary local point reduces into the affine special fiber,
its restriction to $J_p$ is trivial.
\item If an ordinary local point reduces to a boundary point $d$, its
restriction has image in a conjugate of cusp inertia $I_d$ and is
$\tau_p\mapsto m\tau_p$ there, where
$m\in\Z_{>0}$ is its intersection multiplicity with the boundary.
\item At a local rational cusp $d$, with parameter
$\gamma\in H^1(\Q_p,\Z_3(1))$, its restriction has image in a conjugate
of $I_d$ and is $\tau_p\mapsto v_p(\gamma)\tau_p$ there, where
$v_p(\gamma)\in\Z_3$ is the completed Kummer valuation.
\end{enumerate}
\end{lemma}
\begin{proof}
The pair $(X_n,D_n)$ has good reduction at $p$, with finite \'etale disjoint
boundary.  Prime-to-$p$ specialization shows the first assertion
\citep[Expos\'e XIII, \S2]{SGA1}; the actual pointed action is
unramified by Lemma~\ref{lem:action}.
For the second, use a boundary parameter $t$ in the strict henselian
neighborhood of the reduction of $d$.  The prime-to-$p$ geometric
fundamental group of the punctured normal disc is tame cyclic; after
passing to its pro-$3$ quotient, its covers are generated by the
Kummer covers $t^{1/3^r}$.  Pulling such a cover back along the point
gives the Kummer class of $t(x)$, whose restriction to arithmetic inertia
is $v_p(t(x))\tau_p=m\tau_p$.  This computation on all finite $3$-group
covers gives the asserted map in the full pro-$3$ group.  Changing paths
conjugates the cusp inertia subgroup.  At a cusp, the same computation
uses its Kummer parameter $\gamma$, giving the third assertion.
\end{proof}

\begin{lemma}[Compact class and boundary center]
\label{lem:lm-power-compact}
Let $s$ be a Selmer section of an arithmetic quotient of
$\pi_1(U_n)$ which retains $N_n^{\mathrm{ab}}$.
Here every local section must be the image of a full point
section or a full rational-cusp section.
There is a unique rational coordinate cusp $P\in D_n(\mathbb Q)$
with the following properties:
\begin{enumerate}
\item Every full local representative at $3$ is cuspidal at $P$.
\item The global compact class of $s$ in
$H^1(\mathbb Q,T_3J_n)$ is the Kummer class of $P$.
\item At every $p\ne3$, any local cusp representative is at $P$,
and any ordinary local representative reducing to the boundary
reduces to $P$.
\end{enumerate}
\end{lemma}
\begin{proof}
Compose the abelian section with the compact map
$X_n\to X_3=E$.  Its global class is a compact $3$-primary
Selmer class of $E$, hence is the class of one of the three
rational cubic cusps.  Injectivity of the local Kummer map
$E(\mathbb Q_3)\to H^1(\mathbb Q_3,T_3E)$ identifies the
image of the actual local center with that cusp.
Its fiber on $X_n$ consists entirely of coordinate cusps.
Because $\mathbb Q_3$ has no nontrivial roots of unity of
odd order, exactly one point of this fiber is rational over
$\mathbb Q_3$, and it is the corresponding rational cusp
$P$ of $X_n$.  Thus the local center is $P$; in particular,
the full local representative is cuspidal.

Use the standard tangent at $P$ as global origin, and let
$a\in H^1(\mathbb Q,T_3J_n)$ be the compact class of $s$
relative to that origin.  The first assertion gives $a_3=0$.
For $p\ne3$, $X_n$ and $J_n$ have good reduction.  The
compact Kummer class of every local point is unramified:
the reduction kernel in $J_n(\mathbb Q_p)$ is pro-$p$ and
uniquely $3$-divisible, and the $3$-primary completion is
the $3$-primary part of the finite reduction group.
A cusp representative has the compact Kummer class of its
center, so the same assertion holds for it.  Thus $a$ is
unramified outside $3$.  Lemma~\ref{lem:lm-power-action}
applies to $T_3J_n$, a quotient of $N_n^{\mathrm{ab}}$,
and gives $a=0$.  This proves the second assertion without
any Selmer computation for $J_n$.

Every difference of coordinate cusps on $X_n$ is killed
by $n$.  Indeed, the tangent line at any coordinate cusp
$Q$ meets $X_n$ only at $Q$, with intersection multiplicity
$n$.  For two cusps $Q,P$, the ratio of their tangent-line
equations has divisor $n(Q-P)$.  Consequently $[Q-P]$
is $3$-primary torsion in $J_n$.

At $p\ne3$, a local cusp representative centered at
$Q\in D_n(\mathbb Q_p)$ has zero compact difference from
$P$.  The $3$-primary torsion of $J_n(\mathbb Q_p)$
injects into its $3$-primary completion, so $[Q-P]=0$.
The Abel--Jacobi embedding of the positive-genus curve
$X_n$ is injective, and therefore $Q=P$.

If an ordinary local point $x$ reduces to the boundary,
its reduction lifts uniquely to
$Q\in D_n(\mathbb Q_p)$ because $D_n$ is finite \'etale
over $\mathbb Z_p$.  The difference $[x-Q]$ belongs to
the formal reduction kernel of $J_n$, so its $3$-primary
Kummer class is zero.  Since the compact class of $x$
is that of $P$, the preceding argument again gives
$[Q-P]=0$ and $Q=P$.
The initial cubic image also proves uniqueness of $P$.
\end{proof}

\begin{theorem}[Abelian Selmer cuspidality for every power of three]
\label{thm:lm-power-abelian}
For every $r\ge1$, the Selmer sections with geometric group
$N_{3^r}^{\mathrm{ab}}$ are exactly the three disjoint
rational cusp packets, each parametrized by
$H^1(\mathbb Q,\mathbb Z_3(1))$.
The statement does not assume a global full pro-$3$ lift.
Moreover, for any quotient retaining this abelianization,
every full local representative at $3$ has the same cusp
parameter as this global cusp section.  At $p\ne3$ its
whole geometric inertia is, up to geometric conjugation,
\[
 \tau\longmapsto
 i_P\bigl(v_p(\beta)\tau_p(\tau)\bigr),
 \qquad \tau\in I_p,
\]
where $\beta$ is the global parameter and $\tau_p$ is
the tame $3$-adic character.
\end{theorem}
\begin{proof}
Take the cusp $P$ of Lemma~\ref{lem:lm-power-compact}
and use its standard tangent as origin.  There is an
exact sequence of $G_{\mathbb Q}$-modules
\[
 0\longrightarrow B_n\longrightarrow N_n^{\mathrm{ab}}
 \longrightarrow T_3J_n\longrightarrow0,
 \qquad
 B_n=\mathbb Z_3(1)[D_n(\overline{\mathbb Q})]
       /\mathbb Z_3(1)\sum_{Q\in D_n}e_Q.
\]
The compact class is zero.  The equality
$(T_3J_n)^{G_{\mathbb Q}}=0$ therefore gives a unique
preimage $b\in H^1(\mathbb Q,B_n)$ of the abelian class.
These fixed points vanish because the rational torsion
of an abelian variety over a number field is finite;
the identical fixed-point vanishing holds over each
$\mathbb Q_p$ because its torsion subgroup is finite.

Choose a second rational cusp $R\ne P$, and eliminate
$e_R$ from the boundary relation.  This splits
$B_n=I_P\oplus M_n$ equivariantly over $\mathbb Z_3$.
Write $b=(\beta,b_M)$.  At $3$, local cuspidality at $P$
and the injection
$H^1(\mathbb Q_3,B_n)\hookrightarrow
 H^1(\mathbb Q_3,N_n^{\mathrm{ab}})$ give $(b_M)_3=0$.
At $p\ne3$, the geometric local inertia calculation for
the smooth pair $(X_n,D_n)$ applies.  It is zero for
an ordinary point reducing into $U_n$; at a boundary
reduction or at a cusp it lies in the inertia of its
boundary center.  Lemma~\ref{lem:lm-power-compact}
identifies that center with $P$.  Hence $b_M$ is
unramified outside $3$.

The module $M_n$ is a quotient of a submodule of
$N_n^{\mathrm{ab}}$, so its action has the properties
of Lemma~\ref{lem:lm-power-action}.  The lemma gives
$b_M=0$.  Thus the abelian section is the cusp section
$t_{P,\beta}$.

The inclusion of $I_P$ induces injections in global
and local first cohomology: it is a direct summand of
$B_n$, and the following compact Tate module has no
fixed points.  Consequently the parameter is unique,
and the full local cusp parameter at $3$ equals
$\beta_3$.  Under the power cover its inertia parameter
is multiplied by $e=3^{r-1}$; this is consistent with
$u_3=u_n^e$.  Multiplication by $e$ is injective on
$H^1(\mathbb Q_3,\mathbb Z_3(1))
 \simeq\mathbb Z_3^2$, since $\mathbb Q_3$ has no
$3$-power roots of unity.

Finally the full local inertia formula is the tame
disc calculation, with its exponent identified by
the abelian class just obtained.  The uniqueness of
the cusp center was proved in the preceding lemma.
Every global cusp section is Selmer, giving the
asserted exact description of the Selmer set.
\end{proof}

\begin{corollary}[Filling the selected cusp]
\label{cor:lm-fill-full}
For a full geometric pro-$3$ Selmer section $s$ as in
Theorem~\ref{thm:lm-power-abelian}, let $W=U_n\cup\{P\}$.
The image of $s$ in the geometrically pro-$3$ fundamental group of
$W$ is the global point section at $P$.
\end{corollary}
\begin{proof}
Filling $P$ kills the closed normal subgroup generated by $I_P$.
Consequently the inertia formula in
Theorem~\ref{thm:lm-power-abelian} makes the entire projected
cocycle unramified outside $3$, including at primes where
$v_p(\beta)\ne0$.  The action is the quotient of the action defined
by the standard tangent $t_P$, whose image in $W$ is the point
section at $P$.  Thus both its action and its cocycle factor through
$\Gamma_{\{3\}}$, as before.  At $3$, filling the local cusp section
gives the point section at $P$, so the projected local cohomology
class is trivial.  Use its local geometric conjugator globally;
the resulting cocycle still factors through $\Gamma_{\{3\}}$ and
is zero on $G_{\Q_3}$.  Surjectivity of
$G_{\Q_3}\to\Gamma_{\{3\}}$ makes it zero globally.
\end{proof}

\begin{remark}[Infinite support and central differences]\label{rem:infinite-support}
Filling the selected cusp kills its entire ramification contribution.
Before filling, a local geometric conjugator can change the full
inertia subgroup even when abelianizations agree. The lower-central
argument below controls that conjugator at each graded level; it does
not claim that subtracting the cusp parameter makes the entire cocycle
unramified.

There is no general finite-support conclusion for $\beta$.
Indeed,
\[
 H^1(\Q,\Z_3(1))\simeq
 \varprojlim_m\Q^\times/(\Q^\times)^{3^m}
\]
identifies with exponent families $(b_p)_p\in\prod_p\Z_3$ such that,
for each $m$, only finitely many $b_p$ are nonzero modulo $3^m$.
The support itself may be infinite: if $p_1,p_2,\ldots$ enumerate the
primes other than $3$, take $b_{p_j}=3^j$ and $b_3=0$.
This is an actual global cusp parameter, whose cusp section is Selmer
and ramifies at every $p_j$.  Thus continuity and local geometry cannot
by themselves force the original section into a fixed finite
ramification set.

If the parameter's support is contained in one of the finite sets for
which the arithmetic source is dominated by $G_{\Q_3}$, the inertia
formula does place the whole section in that restricted theorem and
proves cuspidality.  The unrestricted conclusion is proved below by
applying the arithmetic-source argument to successive central differences.
\end{remark}
\subsection{Free groups and free solvable central factors}
\label{subsec:lm-central-structure}

Write $N=N_n$ and
\[
 \gamma_1N=N,\qquad \gamma_{m+1}N=\overline{[N,\gamma_mN]},
 \qquad L_m=\gamma_mN/\gamma_{m+1}N.
\]
These are the closed ordinary lower-central subgroups.  The following
facts provide both the local cusp normalizer and the torsion-freeness
needed in the derived quotients.

\begin{lemma}[Free pro-$3$ facts]
\label{lem:lm-free-central}
The group $N=N_n$ is free pro-$3$ of rank $n^2+1$.  Each $L_m$ is a finite free
$\Z_3$-module, and the graded Lie algebra $\bigoplus_mL_m$ is the free
$\Z_3$-Lie algebra on $N^{\rm ab}$.  In particular there is a natural
$G_\Q$-equivariant surjection
\[
 (N^{\rm ab})^{\otimes m}\longrightarrow L_m.
\]
Moreover $\bigcap_m\gamma_mN=1$.
For any cusp $P$, its inertia $I_P$ is generated by a member of a free
pro-$3$ basis of $N$.  If $0\ne b\in\Z_3$, then
\[
 N_N(I_P^{\,b})=I_P.
\]
Here $I_P^{\,b}$ denotes the image of multiplication by $b$ in $I_P$.
\end{lemma}
\begin{proof}
The geometric presentation has $2g(X_n)$ handle generators and $3n$
inertia generators, with the single surface relation.  Eliminating any
inertia generator other than the one at $P$ leaves a free basis of size
$2g(X_n)+3n-1=n^2+1$ containing a generator $a$ of $I_P$.

For a finitely generated free pro-$3$ group, the lower central factors
are finite free $\Z_3$-modules with the Witt ranks.  This is recorded by
\citet[\S3.5, p.~71, immediately after Theorem 3.5.1]{LubotzkySegal2002}.  The commutator maps give the
surjection from the free Lie algebra; equality of the Witt ranks makes
it an isomorphism in each degree.  Iterated brackets give the displayed
tensor surjection.  Every finite $3$-group is nilpotent, so any nonidentity
element of $N$ survives in a finite quotient in which some $\gamma_mN$
vanishes.  This proves the asserted intersection.

For completeness, the normalizer assertion follows from these same Lie
facts.  An element normalizing $\langle a^b\rangle$ acts on this
procyclic group by an element of $\Z_3^\times$.  Abelianization in $N$
shows that this element is $1$, since the image of $a$ is a basis vector
and $b\ne0$.  Thus the element centralizes $a^b$.
The centralizer of the basis vector $A$ in the free Lie algebra over
$\Q_3$ is $\Q_3A$: embed the free Lie algebra in the free associative
algebra, where a homogeneous polynomial commuting with $A$ is a scalar
multiple of a power of $A$; only the first power belongs to the Lie
algebra, since every Lie element is primitive for the standard
coproduct, whereas $A^m$ is not primitive for $m>1$ in characteristic zero.  If $h\in N$ centralizes $a^b$, its degree-one class must
therefore be a $\Z_3$-multiple of $A$.  Multiply $h$ by the corresponding
power of $a$ to arrange $h\in\gamma_2N$.  Its class in each successive
$L_m$, $m\ge2$, commutes with $bA$ and hence vanishes.  The intersection
just proved gives $h=1$ after this adjustment.  Thus the original
centralizer, and therefore the normalizer, is $I_P$.
\end{proof}

\begin{lemma}[Completion of nilpotent groups with free central factors]
\label{lem:lm-nilpotent-completion}
Let $G$ be a finitely generated nilpotent group whose ordinary
lower-central factors are free abelian.  Its pro-$p$ completion
$\widehat G_p$ satisfies
\[
 \gamma_i(\widehat G_p)/\gamma_{i+1}(\widehat G_p)
 \simeq\bigl(\gamma_iG/\gamma_{i+1}G\bigr)\otimes\mathbb Z_p.
\]
\end{lemma}
\begin{proof}
Choose bases of the successive free abelian central factors,
lift them to $G$, and order the resulting Malcev basis
$a_1,\ldots,a_h$ by central weight.  Every element has a
unique collected expression $a_1^{z_1}\cdots a_h^{z_h}$
with $z_i\in\mathbb Z$.  The finite nilpotent collection
process expresses multiplication and inversion by rational
polynomials in these coordinates (equivalently by finitely
many integer-valued binomial polynomials).

The same formulas extend continuously to $\mathbb Z_p^h$.
They take values there, because they do so on the dense set
of integer tuples, and the group identities extend by
continuity.  The resulting compact group $P$ has a central
series whose factors are the corresponding powers of
$\mathbb Z_p$, so $P$ is pro-$p$.  The closure of
$\gamma_iG$ consists of the tuples whose coordinates of
weight less than $i$ vanish.  Continuity of commutators and
density give $\gamma_iP=\overline{\gamma_iG}$.

Finally, every homomorphism $G\to B$ to a finite $p$-group
is continuous for this topology.  If $p^a$ is an exponent
of $B$, its kernel contains all collected expressions
whose coordinates are divisible by $p^a$.  It therefore
contains a neighborhood of the identity and extends to $P$.
Conversely every finite continuous quotient of $P$ is a
$p$-group.  Thus $P=\widehat G_p$, and the coordinate
description of the central factors proves the formula.
\end{proof}

\begin{lemma}[Free solvable pro-$p$ central factors]
\label{lem:lm-free-solvable-pro-p-grades}
If $F$ is a free pro-$p$ group of finite rank and $F^{[d]}$
is its $d$th closed derived subgroup, then every ordinary
lower-central factor of $F/F^{[d]}$ is a finite free
$\mathbb Z_p$-module.
\end{lemma}
\begin{proof}
Let $F_0$ be the discrete free group on the same generators
and put $S_d=F_0/F_0^{[d]}$.  \citet[Theorem~3.1]{Shmelkin1964} proves that
the lower-central factors of every free polynilpotent group
are free abelian; free solvable groups are the special case
in which every entry of the polycentral class row is $2$.
Thus each nilpotent truncation
$S_d/\gamma_{c+1}S_d$ has free abelian lower-central factors.

The universal properties of free groups and pro-$p$
completion identify
\[
 \bigl(S_d/\gamma_{c+1}S_d\bigr)^{\wedge_p}
 \simeq (F/F^{[d]})/\gamma_{c+1}(F/F^{[d]}).
\]
Indeed both sides have exactly the same homomorphisms to
finite $p$-groups of derived length at most $d$ and
nilpotency class at most $c$.  Apply
Lemma~\ref{lem:lm-nilpotent-completion} to the left side.
Taking $c\ge i$ gives the assertion for the $i$th factor.
\end{proof}
The original discrete group input is
\citet[\S3, Theorem~3.1, p.~101]{Shmelkin1964}; the free-solvable
specialization is specified on p.~91.  The preceding argument proves
the passage to pro-$p$ completion.

\subsection{Central induction in the retained quotient}
\label{subsec:lm-unrestricted-full}

The induction uses full local geometric representatives supplied by
the Selmer condition.  It never requires a global lift to the full group.
The local conjugator is controlled at each central level before the
next additive difference is passed through the arithmetic source.

\begin{lemma}[Fixed points in all central quotients]
\label{lem:lm-fixed-central}
Let $R$ be a $G_\Q$-stable quotient of $N_n$ retaining its
abelianization, and suppose $L_m(R)=\gamma_mR/\gamma_{m+1}R$ is
torsion-free for every $m$. Let $t=t_{P,\beta}$ be the image in this
quotient of any global cusp section of $U_n$, with arbitrary
$\beta\in H^1(\Q,\Z_3(1))$, possibly with infinite ramification support.
For every finite prime $p$ and every $m\ge1$,
\[
 L_m(R)^{G_{\Q_p}}=0,
 \qquad
 (R/\gamma_mR)^{t(G_{\Q_p})}=1.
\]
The action on $L_m(R)$ is independent of $\beta$: it is the action defined
by the standard extending tangent $t_P$.
\end{lemma}
\begin{proof}
Inner automorphisms of $N$ act trivially on every $L_m(R)$, so changing the
section from $t_P$ to $t$ leaves this action unchanged.
For $p\ne3$, the action on $N_n^{\rm ab}$ is unramified and fits into
\[
 0\longrightarrow B_n\longrightarrow N_n^{\rm ab}
 \longrightarrow T_3J_n\longrightarrow0.
\]
Arithmetic Frobenius eigenvalues on $T_3J_n\otimes\Q_3$ have complex
absolute value $p^{1/2}$, while those on $B_n\otimes\Q_3$ have absolute
value $p$: the latter are $p$ times roots of unity from the boundary
permutation.  The tensor surjection in
Lemma~\ref{lem:lm-free-central}, followed by the surjection onto
$L_m(R)$, implies that every Frobenius eigenvalue
on $L_m(R)\otimes\Q_3$ has absolute value $p^{w/2}$ for an integer
$m\le w\le2m$.  None is $1$.  Torsion-freeness of $L_m(R)$ therefore gives
$L_m(R)^{G_{\Q_p}}=0$.

At $3$ no new $3$-adic Hodge-theoretic input is required.  The actual
action defined by $t_P$ factors through
\[
 \Gamma_{\{3\}}=\operatorname{Gal}(K_{\{3\}}(3)/\Q),
 \qquad K=\Q(\zeta_3),
\]
because it is unramified outside $3$ and its restriction to $G_K$ has
pro-$3$ image.  The same is therefore true of its action on $L_m(R)$.
The surjection in Lemma~\ref{lem:source},
$G_{\Q_3}\twoheadrightarrow\Gamma_{\{3\}}$,
gives
\[
 L_m(R)^{G_{\Q_3}}=L_m(R)^{\Gamma_{\{3\}}}
 =L_m(R)^{G_\Q}\subseteq L_m(R)^{G_{\Q_2}}=0.
\]
Finally, filter $R/\gamma_mR$ by the images of its lower central
subgroups.  A fixed element has zero image in $L_1(R)$, then in $L_2(R)$,
and successively in all $L_j(R)$ with $j<m$.  It is therefore the identity.
This applies to the possibly ramified action defined by $t$, because
its actions on the graded factors are the actions just treated.
\end{proof}

\begin{lemma}[Nonzero local inertia forces a cusp]
\label{lem:lm-nonzero-inertia-cusp}
Let $p\ne3$, and put $N=N_n$. Suppose a local section $u$ has inertia
\[
 z_u(J_p)=g I_P^{\,b}g^{-1},\qquad 0\ne b\in\Z_3,
\]
relative to $t_P$.  Then $u$ is a cusp section at $P$ in the full
geometric pro-$3$ group.  If its abelian class equals that of
$t_{P,\beta}$, it is conjugate to $(t_{P,\beta})_p$.
\end{lemma}
\begin{proof}
Arithmetic inertia is normal in $G_{\Q_p}$, and its original tangent
action on $N$ is trivial.  The cocycle identity therefore gives
\[
 z_u(\sigma\tau\sigma^{-1})
 =z_u(\sigma)\,\sigma\bigl(z_u(\tau)\bigr)\,z_u(\sigma)^{-1}
 \quad(\sigma\in G_{\Q_p},\ \tau\in J_p).
\]
Thus the image of the local section normalizes the geometric subgroup
$z_u(J_p)=g I_P^{\,b}g^{-1}$.  Conjugate by $g^{-1}$.
The original tangent action preserves $I_P$ and all its nonzero closed
subgroups.  By Lemma~\ref{lem:lm-free-central}, the normalizer of
$I_P^{\,b}$ in $N\rtimes G_{\Q_p}$ is consequently
$I_P\rtimes G_{\Q_p}$.  Thus the conjugated section takes values in
the cusp decomposition group.  Its parameter belongs to
$H^1(\Q_p,I_P)$.  The injection
$H^1(\Q_p,I_P)\hookrightarrow H^1(\Q_p,N^{\rm ab})$
proved above identifies this parameter with $\beta_p$ whenever the
abelian classes agree.
\end{proof}

\begin{theorem}[Derived and central quotients for all powers of three]
\label{thm:lm-power-all-derived}
Let $n=3^r$ with $r\ge1$, and let $R$ be a
$G_{\mathbb Q}$-stable quotient of $N_n$ retaining
$N_n^{\mathrm{ab}}$.  Suppose every ordinary lower-central
factor $\gamma_mR/\gamma_{m+1}R$ is torsion-free over
$\mathbb Z_3$.  Every Selmer section of this arithmetic
quotient is globally cuspidal.

In particular, this holds for the full group $N_n$,
for every finite lower-central quotient
$N_n/\gamma_cN_n$ with $c\ge2$, and for every finite-derived
quotient $N_n/N_n^{[d]}$ with $d\ge1$, including the maximal
geometrically metabelian pro-$3$ quotient of $U_{3^r}$.
Each Selmer set consists of exactly three disjoint
packets parametrized by $H^1(\mathbb Q,\mathbb Z_3(1))$.
\end{theorem}
\begin{proof}
Theorem~\ref{thm:lm-power-abelian} supplies the unique
abelian cusp section $t=t_{P,\beta}$.  We spell out why
the central induction now takes place directly in $R$.

Lemma~\ref{lem:lm-fixed-central} supplies
$(R/\gamma_mR)^{t(G_{\mathbb Q_p})}=1$ at every finite prime,
and its proof identifies the action on each graded factor with
the original standard-tangent action.  The finite generation and
freeness of these factors follow from
Lemma~\ref{lem:lm-free-central} and the hypothesis on $R$.

At $3$ the full local representative is conjugate
to $t$.  At $p\ne3$ with $v_p(\beta)\ne0$, the
whole inertia of a full local representative is
a nonzero subgroup of a conjugate of $I_P$.
Lemma~\ref{lem:lm-nonzero-inertia-cusp} therefore makes the
full local representative cuspidal at $P$, with parameter
$\beta_p$ by the abelian parameter injection.  At $v_p(\beta)=0$, its
whole geometric inertia is zero.  These are
statements about the full \emph{local} representatives
supplied by the Selmer condition, and may then
be projected to $R$.

Assume that $s$ and $t$ agree modulo $\gamma_mR$.
Their next difference is an additive cocycle in
$L_m(R)=\gamma_mR/\gamma_{m+1}R$.
At a place where they are locally conjugate,
the image of a local conjugator in $R/\gamma_mR$
is fixed by $t$, hence is trivial.  The conjugator
therefore lies in $\gamma_mR$.
At $3$ this proves that the graded difference
has zero local cohomology class.  At $p\ne3$ with
nonzero $v_p(\beta)$, conjugation by the cusp inertia
is inner and hence trivial on $L_m(R)$, so the
graded difference is zero on arithmetic inertia.
At the other primes both whole inertias were
already zero.  The graded difference is thus
unramified outside $3$ and zero locally at $3$.
Lemma~\ref{lem:lm-power-action} makes it a global
coboundary.  Conjugating advances the agreement
to $\gamma_{m+1}R$.

Nonempty closed sets of conjugators at the
successive levels are nested in the compact
group $R$, so they have a common element.
The closed ordinary lower-central series of a
pro-$3$ group has trivial intersection, since
every finite $3$-group is nilpotent.  The common
conjugator consequently identifies the sections
in $R$ itself.

For the listed special cases, free pro-$3$
lower-central factors are finite free.  The
same is true of free solvable pro-$3$ groups
by the discrete free-polynilpotent theorem and
the nilpotent completion argument proved in
Lemma~\ref{lem:lm-free-solvable-pro-p-grades}.
Thus all of these quotients satisfy the
hypothesis. This proves the pure-power case of
Theorem~\ref{thm:unrestrictedthreefamiliesintro}.  The proof never
asserts injectivity of a derived quotient under the open embedding
$N_n\subseteq N_3$.
\end{proof}

\begin{remark}[Why the induction is stronger than cusp subtraction]\label{rem:central-induction}
The proof never asserts that subtracting the cusp parameter makes the
whole cocycle unramified.  Instead, at each central level the local
conjugator is forced into that level by the vanishing of fixed points.
Its effect on inertia then vanishes in the next central graded module.
Only this additive difference is passed through the arithmetic quotient
$\Gamma_{\{3\}}$.  This is the step that resolves the conjugator issue
without imposing finite support on the cusp parameter.
\end{remark}

\begin{remark}[Why the proof works directly upstairs]
\label{rem:lm-derived-cover-obstruction}
An open embedding of full pro-$3$ groups need not remain
injective after taking derived quotients.  Moreover, if
$H$ is an open subgroup of a free pro-$3$ group $F$,
the quotient $F/H'$ need not have torsion-free ordinary
central factors, even though it retains $F^{\mathrm{ab}}$.
For example, let $F$ be free on $x,y$ and let $H$ be the
kernel of $F\to\mathbb Z/3\mathbb Z$ taking $x$ to $1$
and $y$ to $0$.  In $F/H'$, the relation $[x^3,y]=1$
makes $3[x,y]=0$ in the second lower-central factor.
The image of $[x,y]$ there is nonzero: map $F$ onto the
Heisenberg group over $\mathbb F_3$, where the image of
$H$ is abelian and $[x,y]$ is a nontrivial central element.
The second lower-central factor is therefore
$\mathbb Z/3\mathbb Z$.

Consequently a full-group cover argument for $U_{3^r}\to U_3$
does not by itself prove the finite-derived assertion upstairs.
Theorem~\ref{thm:lm-power-all-derived} avoids this shortcut: it first
kills the compact class by arithmetic localization, then performs
the boundary and central arguments directly in the upstairs quotient.
\end{remark}

\section{Local cuspidality at three beyond the cubic tower}
\label{sec:localthree}

Two elementary maps to the same proper quotient give local information
that does not require an injection of geometric pro-$3$ groups.
Let $m\ge5$ be odd with $3\nmid m$, let
\[
 C_m:\ v^2+v=u^m,\qquad A_m=\operatorname{Jac}(C_m),
\]
where $C_m$ is the smooth proper model, and embed $C_m$ into $A_m$
using its unique point $\infty$ at infinity.
For a local field $L$, write $A_m(L)^{\wedge3}=\varprojlim_j
A_m(L)/3^jA_m(L)$.

\begin{lemma}[Four points in the zero Kummer locus]\label{lem:fourzero}
The points $Q\in C_m(\mathbb Q_3)$ whose classes $[Q-\infty]$ vanish
in $A_m(\mathbb Q_3)^{\wedge3}$ are exactly
\[
 \infty,\quad (0,0),\quad(0,-1),\quad
 W_m=(a_m,-1/2),\qquad a_m^m=-1/4.
\]
Here $a_m$ is the unique such element of $\mathbb Q_3$.
\end{lemma}
\begin{proof}
The hyperelliptic equation $(2v+1)^2=1+4u^m$ has good reduction at
$3$: its right-hand side is squarefree modulo $3$, and its odd degree
provides the smooth point at infinity. Since $m$ is odd, the map
$u\mapsto u^m$ is the identity on $\mathbb F_3$. Counting one $u$ for
each of the three values of $v$ gives
$\#C_m(\mathbb F_3)=4$, including infinity.

The compact $3$-adic Lie group $A_m(\mathbb Q_3)$ is an extension of
a finite group by an open subgroup isomorphic to a finite direct sum
of $\mathbb Z_3$. Its map to its pro-$3$ completion therefore has as
kernel precisely its finite prime-to-$3$ torsion subgroup. Good
reduction is injective on that subgroup, because the formal kernel
is pro-$3$ and multiplication by an integer prime to $3$ is invertible
there. The Abel--Jacobi map of the special fiber is injective. Thus at
most one zero-Kummer point can lie in any one residue disc, and there
are at most four in total.

The divisors of $v$ and $v+1$ are respectively
$m(0,0)-m\infty$ and $m(0,-1)-m\infty$. These two cusp classes have
order dividing $m$ and hence lie in the kernel. Raising a $3$-adic unit
to the odd power $m$, with $3\nmid m$, is an automorphism: it is an
automorphism on $\{\pm1\}$ and on $1+3\mathbb Z_3$ by the $3$-adic
logarithm. Hence $a_m$ exists uniquely. The point $W_m$ is a branch
point of the hyperelliptic map, so
$\operatorname{div}(u-a_m)=2W_m-2\infty$. Its class also lies in the
kernel. The four displayed points reduce respectively to infinity,
$(0,0)$, $(0,-1)$ and $(-1,1)$, which are distinct. They exhaust the
possible zero-Kummer points.
\end{proof}

\begin{proposition}[A vanishing proper Selmer group forces a local cusp]\label{prop:localtwomaps}
Suppose the $3$-adic Selmer group of $T_3A_m$ is zero; equivalently,
$A_m(\mathbb Q)^{\wedge3}=0$ and $T_3\Sha(\mathbb Q,A_m)=0$.
For every odd multiple $n$ of $m$, the localization at $3$ of any
point-or-cusp Selmer section of the geometrically pro-$3$ abelian
quotient of $U_n/\mathbb Q$ is cuspidal at a rational coordinate
cusp. The same conclusion holds for any geometric pro-$3$ quotient
retaining that abelianization, including the full pro-$3$ group.
\end{proposition}
\begin{proof}
Compose the coordinate-power morphism $X_n\to X_m$ with the two maps
\[
 f_Z(X:Y:Z)=\left(-XY/Z^2,X^m/Z^m\right),\qquad
 f_X(X:Y:Z)=\left(-YZ/X^2,Y^m/X^m\right)
\]
to $C_m$, and then with its Abel--Jacobi embedding. These extend to
morphisms of proper smooth curves. Rational coordinate cusps map to
$\infty,(0,0),(0,-1)$; their classes have order dividing $m$, so their
$3$-adic Kummer classes vanish. Thus projecting the global section
through either morphism gives an ordinary $3$-adic Selmer class of
$T_3A_m$ with respect to the indicated origin. Both classes are zero
by hypothesis. This uses only the geometric abelian quotient of the
global section and its full local point-or-cusp realizations.

If a local realization at $3$ had center in $U_n(\mathbb Q_3)$, its
image $(X:Y:Z)\in U_m(\mathbb Q_3)$ would have all coordinates nonzero.
The preceding lemma applies to both its images on $C_m$. None can be
$\infty,(0,0),(0,-1)$ because the inverse images of these points are
coordinate divisors. Consequently both images must be $W_m$.
The first equation says $X^m/Z^m=-1/2$. The Fermat equation gives
$Y^m/Z^m=-1/2$ as well, hence $Y^m/X^m=1$. The second image would
instead require $Y^m/X^m=-1/2$, a contradiction.

Therefore every local realization at $3$ is a coordinate cusp.
The only roots of unity in $\mathbb Q_3$ are $\{\pm1\}$, so for odd
$n$ the only $n$th root of unity there is $1$. The three local
coordinate cusps are exactly the three rational coordinate cusps.
This proves the stated local conclusion. No conclusion about global
conjugacy follows merely from this argument.
\end{proof}

\begin{proof}[Proof of Theorem~\ref{thm:localthreeintro}]
If $n$ is divisible by $m\in\{5,7,11\}$, the arithmetic
Propositions~\ref{prop:arithmetic}, \ref{prop:septic}, and
\ref{prop:eleven}, respectively, make the proper $3$-adic Selmer
group zero: the rational Mordell--Weil group has order $m$, and the
Tate--Shafarevich group is finite. Apply Proposition~\ref{prop:localtwomaps}.
If $3\mid n$, the coordinate-power map $X_n\to X_3=E$
sends the compact $3$-adic class of any Selmer section to one of
$E(\mathbb Q)=\{O,(3,4),(3,-5)\}$, by Lemma~\ref{lem:cubiclocal}.
The local Kummer map at $3$ is injective, so every local center maps
to that rational cusp. The inverse image of a coordinate cusp under
the power map consists entirely of coordinate cusps. Since
$\mathbb Q_3$ has no nontrivial odd-order roots of unity, the local
center is again one of the three rational coordinate cusps upstairs.
These projections factor through geometric abelianization, proving
the assertion for every quotient in the statement.
\end{proof}


\subsection{Local separation of rational cusp packets}
\label{subsec:lm-affine-packet-separation}

\begin{proposition}[The affine abelianization separates the rational cusps]
\label{prop:lm-affine-ab-cusp-separation}
For every odd $n\ge3$, the three rational cusp packets of
$U_n/\mathbb Q_3$ have pairwise disjoint images in the
geometrically pro-$3$ abelian quotient.
Consequently the three global rational cusp packets are
also pairwise disjoint in that quotient.
If $3\nmid n$, this assertion uses the affine
abelianization: the compact Jacobian Kummer classes of
the rational cusps are all equal.
\end{proposition}
\begin{proof}
If $3\mid n$, compose with the compact cubic map.
The three rational cubic cusp points have distinct
$3$-primary Kummer classes over $\mathbb Q_3$, because
$E(\mathbb Q_3)$ is pro-$3$ and its Kummer map is
injective.  Thus their inverse-image rational cusp
packets are distinct.

Suppose $3\nmid n$.  Choose a prime $q\ge5$ dividing $n$.
The coordinate-power morphism $U_n\to U_q$ carries
distinct rational coordinate cusps to distinct rational
coordinate cusps.  It is enough to separate the packets
on $U_q$; no injectivity of its map on pro-$3$
fundamental groups is required.

By coordinate permutations it is enough to compare
\[
 P=(0:-1:1),\qquad R=(-1:0:1).
\]
Put $L=\mathbb Q_3(\zeta_q)$, with $\zeta_q$ primitive,
and consider the rational unit on $U_{q,L}$
\[
 \varphi=\frac{Y+\zeta_q Z}{Y+\zeta_q^2 Z}.
\]
Its divisor on the compact Fermat curve is
$q(Q_{\zeta_q}-Q_{\zeta_q^2})$, where
$Q_\zeta=(0:-\zeta:1)$.  In particular $\varphi$
is regular and nonzero at both $P$ and $R$, with
\[
 \varphi(P)=\frac1{1+\zeta_q},
 \qquad
 \varphi(R)=\zeta_q^{-1}.
\]
For any cusp section centered at $P$ or $R$, the
map induced by $\varphi$ to $\mathbb G_m$ is the
Kummer section of the corresponding displayed
value.  The tangent parameter has no effect,
because $\varphi$ has order zero at that cusp.

Equality of two abelian cusp sections at $P$ and
$R$, after restriction to $L$, would therefore give
\[
 \left[\frac{\zeta_q}{1+\zeta_q}\right]=0
 \quad\text{in}\quad H^1(L,\mathbb Z_3(1)).
\]
For a finite extension $L/\mathbb Q_3$, the kernel
of $L^\times\to (L^\times)^{\wedge_3}$ is exactly
the finite subgroup of roots of unity of order
prime to $3$.  The displayed algebraic number is
not a root of unity: choose a complex embedding
taking $\zeta_q$ to $\exp(2\pi i/q)$, under which
its absolute value is
\[
 \frac1{2\cos(\pi/q)}<1
 \qquad(q\ge5).
\]
This is a contradiction.  Coordinate permutations
separate each pair of rational cusps.

Finally, the difference of any two coordinate
cusps on $X_n$ is killed by $n$, using the ratio
of their total-flex tangent lines.  When $3\nmid n$,
these differences have zero $3$-primary Kummer
class.  This verifies the asserted distinction
between compact and affine abelianization.
Global equality would imply local equality at
$3$, so the global packets are disjoint as well.
\end{proof}

\begin{remark}[Scope of local packet separation]\label{rem:affine-packet-scope}
The proposition identifies a unique rational cusp
from a known local abelian cusp class at $3$.
It does not imply that a global abelian Selmer
section is globally cuspidal.  For mixed exponents, the nonzero relative
localization kernel in Proposition~\ref{prop:mixed-kernel} remains,
and membership in this kernel does not itself
give a local geometric representative at the
other places.  Both distinctions are necessary:
compact ambiguity of the cusp center is repaired
by the proposition, whereas the global Selmer
obstruction is not.
\end{remark}


\section{Arbitrary sections and their local obstruction}
\label{sec:arbitrary}

The distinction between arbitrary sections and Selmer sections already
appears when ramification is allowed only at $3$. The construction below
uses Shafarevich's arithmetic freeness theorem and the free-source method
of Hoshi. We give the equivariant construction over $\Q$ explicitly.
The exchanged arithmetic basis also appears in
\citet[Theorem 12, p.~12]{Medvedovsky2016}; the elementary argument is
included to make the descent under complex conjugation transparent.

Put $P=P_{\{3\}}$ and $\Gamma=\Gamma_{\{3\}}$ in the notation of
Section~\ref{sec:local-arithmetic-source}. Complex conjugation supplies an involution $c$, so
\[
\Gamma=P\rtimes\langle c\rangle.
\]
The field $K=\Q(\zeta_3)$ contains $\mu_3$, has class number one, and
has just one prime above $3$. Thus
\citet[\S4, example 1 after Theorem 5, p.~82]{Shafarevich1963}
gives relation rank zero and generator rank
$[K:\Q]-\operatorname{rank}\mathcal O_K^\times=2$.
Consequently $P$ is free pro-$3$ of rank two; see also
\citet[Lemma 3.2]{Hoshi2010}.

\begin{lemma}[An exchanged free basis]\label{lem:exchanged}
There is a free pro-$3$ basis $g,h$ of $P$ with $h=cgc$.
Thus $c$ exchanges the two generators exactly.
\end{lemma}
\begin{proof}
The extensions $K(\zeta_9)$ and $K(\sqrt[3]{3})$ are independent
cyclic cubic extensions unramified outside $3$. Their Kummer classes
are $\zeta_3$ and $3$: the first has zero valuation, whereas the
second has $\lambda$-valuation two, nonzero modulo three. They
therefore give the two independent quotients of $P/\Phi(P)$.
Complex conjugation acts trivially on the cyclotomic quotient,
since $\Gal(\Q(\zeta_9)/\Q)$ is abelian, and by inversion on the
other quotient, whose extension over $\Q$ has group $S_3$.

Choose $g\in P$ whose Frattini vector has both eigencoordinates
nonzero, say $(a,b)$. The vector of $h=cgc$ is $(a,-b)$.
Their determinant is $-2ab\ne0$ in $\F_3$, so Burnside's basis
theorem says $g,h$ generate $P$. A generating pair of a rank-two
free pro-$3$ group is a free basis: the resulting surjective
endomorphism is invertible by the Hopfian property of finitely
generated profinite groups. Since $c^2=1$, it also sends $h$ to $g$.
\end{proof}

\begin{lemma}[Direct full-group cocycles]\label{lem:directcocycles}
For every pro-$3$ group $M$ with a continuous $\Gamma$-action and
every $a\in M$, there is a continuous $M$-cocycle $z_a$ with
\[
z_a(c)=1,\qquad z_a(g)=a,\qquad z_a(h)=c(a).
\]
\end{lemma}
\begin{proof}
The universal property of the free pro-$3$ basis gives a continuous
homomorphism
\[
P\longrightarrow M\rtimes P,\qquad
g\longmapsto(a,g),\quad h\longmapsto(c(a),h).
\]
The target is pro-$3$ and its projection to $P$ is the identity
on the free generators, hence everywhere. Conjugation by $(1,c)$
exchanges the two specified images. The homomorphism is therefore
$c$-equivariant and extends to a section
$\Gamma\to M\rtimes\Gamma$ sending $c$ to $(1,c)$. Its relative
cocycle is the asserted $z_a$. This constructs an actual full-group
cocycle, with no passage from graded lifts to an inverse limit.
\end{proof}

Let $E$ be the proper cubic and put $T=T_3E$, with the action supplied
by a standard cusp tangent.
\begin{lemma}[The compact cohomology]\label{lem:compactrank}
There is an isomorphism of compact $\Z_3$-modules
\[
H^1(\Gamma,T)=T/(g-1)T^{c=1}.
\]
This group is finitely generated and has rank one.
\end{lemma}
\begin{proof}
Since multiplication by two is invertible on $T$, every cocycle is
cohomologous to one whose value at $c$ is zero. Such a normalized
cocycle is determined by its arbitrary value $x$ at $g$, with value
$c(x)$ at $h$. Conversely every $x\in T$ occurs by
Lemma~\ref{lem:directcocycles}. A coboundary preserves the
normalization exactly when it comes from $y\in T^{c=1}$, and its
value at $g$ is $(g-1)y$, proving the formula.

The determinant of complex conjugation on $T$ is $-1$, so its
$+1$ eigenspace has rank one. The map
$(g-1):T^{c=1}\to T$ is injective: a vector in its kernel is fixed
by $c$ and $g$, hence by $h=cgc$ and all of $\Gamma$. Such a vector
is a compatible system of rational $3$-power division points and
is zero because $E(\Q)$ is finite. Thus the displayed quotient
has rank $2-1=1$. No assertion about its finite torsion is needed.
\end{proof}

\begin{proof}[Proof of Theorem~\ref{thm:arbitraryintro}]
Let $N_n=\pi_1((U_n)_{\bar\Q})^{(3)}$ for $n=3^r$.
Lemma~\ref{lem:action} identifies $N_n$ with a $G_\Q$-stable open
subgroup of $N_3$. The stability uses the compatible coordinate
tangents: the local parameter downstairs is the
$3^{r-1}$-st power of the parameter upstairs with leading
coefficient one, so the tangential arithmetic splittings commute
with the covering map. Therefore the action on this actual open
subgroup factors through the same $\Gamma$ as the cubic action.

The image $T'$ of $N_n$ under the compact projection
$N_3\twoheadrightarrow T$ has finite index in $T$. Indeed the
image of an open subgroup under a surjective continuous
homomorphism of profinite groups has finite index. In particular
$T'\otimes\Q_3=T\otimes\Q_3$. By Lemma~\ref{lem:compactrank}
choose $x\in T'$ whose class $\alpha$ in $H^1(\Gamma,T)$ has
infinite order, and choose $a\in N_n$ above $x$.

For each $t\in\Z_3$, use Lemma~\ref{lem:directcocycles} with
$a^{1+3t}$ in place of $a$, and inflate its cocycle along
$G_\Q\twoheadrightarrow\Gamma$. The exponent is well defined
in the procyclic closure of $a$. Its compact class is
\[
\alpha_t=(1+3t)\alpha.
\]
These classes have infinite order and are pairwise distinct.
Inflation $H^1(\Gamma,T)\to H^1(\Q,T)$ is injective: the same
vector exhibiting an inflated class as a coboundary would exhibit
the original class as a coboundary. Therefore the resulting
sections are pairwise geometrically nonconjugate. Their whole
cocycles are unramified outside $3$, since they factor through
$\Gamma$. There are continuum many parameters $t$.

A rational cusp of $X_n$ maps to a rational cusp of $E$. Its
compact class is thus one of the three torsion Kummer classes
in $E(\Q)$, independently of its inertia parameter. Each
$\alpha_t$ is nontorsion, so none of the constructed sections
is globally cuspidal. The compact projection factors through
$N_n^{\ab}$, so this distinction, and pairwise nonconjugacy,
persist in every Galois-stable geometric quotient retaining that
abelianization.

Finally, local point-or-cusp geometry at $3$ would make
$\alpha_t$ a local Kummer class there. Lemma~\ref{lem:cubiclocal}
would then make it a global Kummer class in the finite group
$E(\Q)\simeq\Z/3\Z$, contradicting its infinite order.
Every constructed section therefore fails the Selmer condition
at $3$.
\end{proof}

The choices of an arithmetic basis and of $a$ are noncanonical, but
once made they give the displayed family of full cocycles by a fixed
formula on generators. This construction supplies abstract
noncuspidal sections and proves their local obstruction. It makes
no claim of a noncuspidal Selmer section.

\subsection{A family with zero compact abelianization}

The preceding construction detects noncuspidality through a nontorsion
compact class. A different family shows that the obstruction can be
entirely in the boundary part of geometric homology.

\begin{lemma}[Lifting over a source of $p$-cohomological dimension one]
\label{lem:cdlift}
Let $\Gamma_0$ be profinite with $\operatorname{cd}_p\Gamma_0\leq1$,
and let $N_0$ be a finitely generated pro-$p$ group with a continuous
$\Gamma_0$-action. Every continuous cocycle
$\Gamma_0\to N_0^{\ab}$ lifts to a continuous $N_0$-valued cocycle.
\end{lemma}
\begin{proof}
Pull $N_0\rtimes\Gamma_0\to N_0^{\ab}\rtimes\Gamma_0$ back along
the section specified by the cocycle. The resulting extension
$E\to\Gamma_0$ has kernel $D=[N_0,N_0]$.
Choose a nested cofinal sequence $V_i$ of characteristic open subgroups
of $N_0$. Then $E_i=E/(D\cap V_i)$ has finite $p$-group kernel and
$E=\varprojlim_iE_i$.

Every extension of $\Gamma_0$ by a finite $p$-group $L$ splits.
To see this, induct on $|L|$. The nontrivial characteristic elementary
abelian subgroup $Z=\Omega_1(Z(L))$ is normal in the extension.
First split the quotient extension with kernel $L/Z$ by induction,
then pull back along its section. The remaining extension has kernel
$Z$, and its obstruction lies in $H^2(\Gamma_0,Z)=0$.
The action on $Z$ is the one induced by this extension. Continuous
cohomology applies: an extension with finite kernel has local
continuous sections, by choosing an open subgroup meeting its kernel
trivially, and hence a continuous set-theoretic section on finitely
many clopen translates.

If a section into $E_i$ has been chosen, pull $E_{i+1}\to E_i$ back
along it. The same finite-$p$-kernel argument supplies a compatible
section into $E_{i+1}$. Taking this sequence's inverse limit gives a
continuous section $\Gamma_0\to E$. Its $N_0$-component lifts the
specified abelian cocycle exactly.
\end{proof}

\begin{proposition}\label{prop:boundaryfamily}
For every $n=3^r$, $r\geq1$, there are continuum many pairwise
geometrically nonconjugate full pro-$3$ sections of $U_n/\Q$ whose
whole cocycles are unramified outside $3$ and whose classes in
$H^1(\Q,T_3\Jac(X_n))$ are zero. Each fails point-or-cusp geometry
at $3$. All these assertions persist after passing to any
Galois-stable geometric quotient retaining $N_n^{\ab}$.
\end{proposition}
\begin{proof}
Use $x=X/Z$, $y=Y/Z$ and the standard tangent $t=x$ at
$Q_0=[0:-1:1]$. In its Puiseux fiber functor,
\[
x=t,\qquad y=-(1+t^n)^{1/n}.
\]
The leading constants of $x,y$ are $1,-1$; both have compatible
rational $3$-power roots. The reference section therefore has zero
Kummer classes under these two invertible functions.

Set $M=N_n^{\ab}$. The rational cusps $Q_0$ and
$Q_1=[-1:0:1]$ supply Galois-equivariant maps
$j_0,j_1:\Z_3(1)\to M$ from their inertia groups.
The map induced by the functions $x,y$ satisfies
\[
(q_x,q_y)j_0(z)=(z,0),\qquad
(q_x,q_y)j_1(z)=(0,z),
\]
since these are their orders of zero at the two cusps.
All boundary classes have zero image in $T_3\Jac(X_n)$.

Let $\Gamma=\Gamma_{\{3\}}$ be the arithmetic source above.
The free subgroup $P$ has index two, so
$\operatorname{cd}_3\Gamma=1$: the Hochschild--Serre spectral sequence
and invertibility of two on every $3$-primary module reduce the
assertion to $\operatorname{cd}_3P=1$.
The Kummer cocycle $\kappa_3(3)$ factors through $\Gamma$.
Indeed the fields obtained by adjoining compatible $3$-power roots
of $3$ and roots of unity are unramified outside $3$, and over
$K=\Q(\zeta_3)$ their Galois closures are pro-$3$.
For every $b\in3+9\Z_3$ the abelian cocycle
\[
a_b=(j_0+bj_1)\kappa_3(3)
\]
therefore lives on $\Gamma$. Lemma~\ref{lem:cdlift} lifts it to an
actual $N_n$-cocycle on $\Gamma$. Inflate to $G_\Q$.
The entire cocycle is unramified outside $3$ because this is true
of $G_\Q\to\Gamma$, and its compact abelianization is zero.

Localize at $3$ and apply the valuation to the two coordinate Kummer
classes. The resulting vector is $(1,b)$, so different parameters
give different geometric conjugacy classes.
If a point $(x,y)\in U_n(\Q_3)$ gave this vector, then
$v_3(x)=1$. The equation $y^n=-1-x^n$ would imply $v_3(y)=0$,
contradicting $b\in3+9\Z_3$.
The only rational cusps over $\Q_3$ are $Q_0,Q_1$ and
$Q_\infty=[1:-1:0]$, since $\Q_3$ has no nontrivial
$3$-power roots of unity. Their local coordinate orders give,
respectively, the three lines of possible cusp valuation vectors
\[
(z,0),\qquad(0,z),\qquad(-z,-z)\qquad(z\in\Z_3).
\]
The vector $(1,b)$ is on none of them: $b$ is nonzero and is not one.
Thus the sections are neither point geometric nor cusp geometric
at $3$. This also proves global noncuspidality. The coordinate maps
and all these valuation distinctions factor through $M$, so they
remain valid in every stated quotient.
\end{proof}

The two constructions exhibit different failures of local geometry.
The first is visible in the proper elliptic Tate module; the second
vanishes in compact homology and is detected by two boundary Kummer
coordinates. Neither construction gives a Selmer counterexample.



\section{Mixed-exponent classifications and compact images}
\label{sec:mixed-exponents}

The power-of-three theorem does not transfer to arbitrary odd
multiples by the same open-subgroup argument. We first describe the
module left by transfer alone and a nonzero localization kernel in it.
The later arithmetic arguments in this section resolve the families
$5\cdot3^r$ and $7\cdot3^r$ of Theorem~\ref{thm:lm-power-all-derived}.

\begin{proposition}[Prime-to-three transfer]
\label{prop:mixed-transfer}
Let $n=dm$ be odd, where $d=3^r\ge3$ and $(m,3)=1$, and let
$f:U_n\to U_d$ be the coordinate-power map. Put
$N_j=\pi_1((U_j)_{\bar\Q})^{(3)}$ and $A_j=N_j^{\ab}$.
Then $f_*:N_n\to N_d$ is surjective, and there is a
$G_\Q$-equivariant decomposition
\[
A_n\simeq A_d\oplus M,\qquad
M=\ker(f_*:A_n\to A_d),\qquad
\operatorname{rank}_{\Z_3}M=n^2-d^2.
\]
For a Selmer section of the full geometric pro-$3$ group of $U_n$,
subtract the compatible global cusp section whose pushdown is
given by Theorem~\ref{thm:lm-power-all-derived}. Its abelian difference
belongs to
\[
\ker\bigl(H^1(\Q,M)\longrightarrow H^1(\Q_3,M)\bigr).
\]
No vanishing of that kernel follows from transfer.
\end{proposition}
\begin{proof}
The geometric cover has degree $m^2$. A subgroup of index prime
to three maps onto every finite $3$-group quotient of the ambient
group: its image has index dividing both that index and a power
of three. Thus $N_n\to N_d$ is onto. The geometric group $N_j$
is free pro-$3$ of rank
$2g(X_j)+3j-1=j^2+1$. Homological transfer is Galois equivariant
and satisfies $f_*f^!=m^2\operatorname{id}$. Dividing by the
unit $m^2\in\Z_3$ gives the asserted splitting and rank.

Theorem~\ref{thm:localthreeintro} identifies the center at three
with the compatible rational cusp. Choose compatible parameters
$t_d=t_n^m$. Their inertia map is multiplication by the unit $m$.
Thus the global cusp parameter downstairs lifts uniquely upstairs.
Its localization agrees with the given local parameter, by the
local injectivity of peripheral Kummer cohomology. As in the proof
of Theorem~\ref{thm:restricted3}, this injectivity follows by
filling the cusp: the boundary Tate modules have no local
invariants, and the compact Tate module has none because rational
torsion on an abelian variety over $\Q_3$ is finite.

The abelian difference has zero pushdown and zero restriction
at three. The split decomposition puts it in the displayed
kernel. Transfer cannot remove $M$, on which the pushdown is zero.
\end{proof}

\begin{proposition}[A nonzero relative localization kernel]
\label{prop:mixed-kernel}
If $m>1$ in Proposition~\ref{prop:mixed-transfer}, then
\[
\ker\bigl(H^1(\Q,M)\longrightarrow H^1(\Q_3,M)\bigr)\ne0.
\]
This is a cohomological assertion. The constructed class is not
asserted to be Selmer or to lift to a nonabelian section.
\end{proposition}
\begin{proof}
Choose a prime $q\mid m$; it is odd. Write
$P_\zeta=(0:-\zeta:1)$, and let $e_\zeta$ be its inertia
generator in the boundary Tate module of $A_n$. The vector
\[
w=\sum_{\substack{\zeta^q=1\\\zeta\ne1}}e_\zeta-(q-1)e_1
\]
defines a Galois-equivariant map $j:\Z_3(1)\to A_n$.
These cusps lie over one rational cusp of $U_d$, and their
inertia generators all map with index $m$. The coefficients
sum to zero, so $j$ takes values in $M$.

The function
\[
u=\frac{Y^q+Z^q}{(Y+Z)^q}
\]
is a unit on $U_n$. Each line $Y+\zeta Z=0$ meets $X_n$
only at $P_\zeta$, with multiplicity $n$. Consequently
\[
\operatorname{div}(u)=n\left(
\sum_{\zeta^q=1,\,\zeta\ne1}P_\zeta-(q-1)P_1\right).
\]
The geometric Kummer map of $u$, restricted to $M$, gives
$h:M\to\Z_3(1)$ with
\[
h\circ j=n\bigl((q-1)+(q-1)^2\bigr)=nq(q-1).
\]

Let $a=\log_3(-2)$ and $b=\log_3(-5)$, using the injective
logarithm on $1+3\Z_3$. Both are nonzero. The global class
\[
\kappa=b[2]-a[5]\in H^1(\Q,\Z_3(1))
\]
has completed valuation $b\ne0$ at two, so has infinite order.
At three both terms have zero valuation and their logarithms
cancel; the order-two signs have zero pro-$3$ Kummer class.
Thus $\kappa|_{\Q_3}=0$. It follows that $j_*(\kappa)$ is in
the asserted kernel, and it is nonzero because
$h_*j_*(\kappa)=nq(q-1)\kappa\ne0$.
\end{proof}

\subsection{A smaller abelian gate for two-prime mixed exponents}
\label{subsec:lm-mixed-compact-gate}

Let $q\ne3$ be an odd prime, and let
\[
 n=3^r q^s,\qquad r,s\ge1,\qquad S=\{3,q\}.
\]
The boundary-localization range will be determined exactly below. It
includes every $3q^s$ and, for arbitrary $r$, all primes
$q\equiv2,4,5,7\pmod9$.
As before, write $X_n$ for the compact Fermat curve,
$D_n$ for its coordinate cusps, $U_n=X_n\setminus D_n$,
$J_n=\operatorname{Jac}(X_n)$, and
$N_n=\pi_1(U_{n,\overline{\mathbb Q}})^{\mathrm{pro}\text{-}3}$.
The statements below retain an explicit compact Selmer hypothesis.
It and the higher graded conditions are established for $n=5\cdot3^r$
in Theorem~\ref{thm:lm-five-power-selmer} and for $n=7\cdot3^r$
by Corollary~\ref{cor:lm-septic-class-gate} and
Proposition~\ref{prop:septic-three-class}. For other exponents
they remain explicit hypotheses.

\begin{lemma}[Exact order of a cusp difference]
\label{lem:lm-exact-cusp-order}
For every $n\ge3$, the difference of two distinct coordinate
cusps on $X_n$ has exact order $n$ in $J_n(\overline{\mathbb Q})$.
\end{lemma}
This is the classical cusp-order assertion recorded by
\citet[p.~191, immediately before Theorem~A]{Coleman1986}.

\begin{lemma}[The abelian unit localization input]
\label{lem:lm-abelian-leopoldt}
If $F/\mathbb Q$ is a finite abelian extension, then
\[
 \mathcal O_F^\times\widehat{\ }_3
 \longrightarrow
 \prod_{v\mid3}(F_v^\times)^{\wedge_3}
\]
is injective.
\end{lemma}
\begin{proof}
This is the abelian case of Leopoldt's theorem.
\citet[pp.~192--193, Remark~1 after Theorem~$1_p$]{Brumer1968}
records the nonvanishing regulator for real abelian fields;
the unit-index assertion on p.~195 gives the CM extension.
It follows from the Baker--Brumer theorem.  We recall the application
to make its precise strength clear.  Modulo roots of unity,
the rational representation of
$G=\operatorname{Gal}(F/\mathbb Q)$ on global units contains
each relevant nontrivial even character with multiplicity
one, by Dirichlet's unit theorem.  The diagonal $3$-adic
logarithm map to $\prod_{v\mid3}F_v$ is $G$-equivariant.
A nonzero kernel after tensoring with $\mathbb Q_3$ would,
after scalar extension, contain a full character line.
That line is defined using algebraic character values,
so it would give a nonzero algebraic-coefficient relation
among the logarithms of multiplicatively independent
global units, at any fixed embedding into
$\overline{\mathbb Q}_3$.

Those logarithms are linearly independent over $\mathbb Q$:
an integral relation would make a product of independent
units a local root of unity, hence a global root of unity.
The Baker--Brumer theorem makes them independent over the
algebraic closure of $\mathbb Q$ as well, a contradiction.
This proves injectivity on the free part after tensoring
with $\mathbb Q_3$.  Any integral kernel is therefore
torsion.  The remaining finite $3$-power roots of unity
inject into every local multiplicative $3$-completion,
so the integral kernel is also zero.

The transcendence input is the Baker--Brumer theorem in
\citet[Theorem, p.~11-01; proof, pp.~11-02--11-08]{Rousseau1969}.
\end{proof}

\begin{lemma}[A single prime in the real subfield]
\label{lem:root-real-q-units}
Let $F/\mathbb Q$ be a finite abelian extension, let $F^+$ be its maximal
real subfield, and let $q\ne3$ be a rational prime having exactly one prime
in $F^+$.  Let $S_F$ consist of the primes above $3$ and $q$.  Then
\[
 H^1(G_{F,S_F},\mathbb Z_3(1))
 \longrightarrow\prod_{v\mid3}H^1(F_v,\mathbb Z_3(1))
\]
is injective.
\end{lemma}
\begin{proof}
Integral Kummer theory identifies the source with the $3$-completion
of the finitely generated $S_F$-unit group.  The inverse-limit ideal-class
term vanishes because the ideal class group is finite.  A class in the
kernel has zero valuation at every prime above $3$.  Exactness of
completion for finitely generated abelian groups places it in the
completed group $E_q(F)=\mathcal O_F[1/q]^\times\widehat{\ }_3$.

If $F$ is totally real, the unique prime $\mathfrak q\mid q$ has a
valuation $a\in\mathbb Z_3$.  The norm to $\mathbb Q$, followed by
localization at $3$ and the logarithm, gives
$f_{\mathfrak q}a\log_3(q)=0$.  Since $\log_3(q)\ne0$, $a=0$.
The class is then an ordinary unit class and abelian Leopoldt makes it
zero, including its integral torsion part.

Otherwise $F$ is a CM field.  Write $c$ for complex conjugation.
The integral projectors $(1+c)/2$ and $(1-c)/2$ split $E_q(F)$, since
$2$ is a unit in $\mathbb Z_3$.  Norm and inclusion identify the plus
part with $E_q(F^+)$: norm composed with inclusion is multiplication by two.
The localization kernel on this part is zero by the preceding real-field
argument, applied to $F^+$.  Here norm at $3$ uses all the places above
each place of $F^+$, as does the localization map in the statement.

It remains to consider the minus part.  Modulo roots of unity the minus
part of the ordinary unit group has rank zero.  Indeed, for every
ordinary unit $u$ the element $u/\bar u$ is a root of unity, by
Kronecker's theorem.  The valuation lattice at $q$ has minus rank at
most one: the unique prime of $F^+$ either stays a single prime in $F$
or splits into one conjugate pair.  The ideal class group is finite,
so the valuation image has finite index in the permitted lattice.
Consequently the minus part of $E_q(F)$ modulo finite torsion has
$\mathbb Z_3$-rank zero or one.

In the rank-one case choose a nonzero algebraic anti-invariant
$q$-unit after tensoring with $\mathbb Q$.  Concretely, if
$(a)=\mathfrak q^h$, then $u=a/\bar a$ has nonzero divisor
$h(\mathfrak q-\bar{\mathfrak q})$ and spans the required rational
line.  At every place above $3$, it is a local unit and its $3$-adic
logarithm is nonzero: a local unit with logarithm zero is a root of
unity, and an embedding of $F$ cannot turn this algebraic non-root
of unity into a root of unity.  Localization is therefore injective
on that line after tensoring with $\mathbb Q_3$.  Its integral kernel
can only be torsion.  The finite global $3$-power roots of unity inject
into local multiplicative $3$-completion.  This kills the remaining
kernel and proves the lemma.
\end{proof}

\begin{proposition}[Exact range of the two-prime boundary argument]
\label{prop:root-real-boundary-range}
Let $n=3^r q^s$, where $r,s\ge1$ and $q\ne3$ is an odd prime.
Choose distinct rational cusps $P,R$ and split the boundary Tate module
as $B_n=I_P\oplus M_n$ by eliminating $e_R$ from the diagonal boundary
relation. For $S=\{3,q\}$, the map
\[
 H^1(G_{\mathbb Q,S},M_n)\longrightarrow H^1(\mathbb Q_3,M_n)
\]
is injective if and only if
\[
 r=1\quad\hbox{or}\quad q\equiv2,4,5,7\pmod9.
 \tag{R}\label{eq:root-real-prime-condition}
\]
\end{proposition}
\begin{proof}
The boundary consists of three copies of $\mu_n$. Removing the two
rational basis vectors gives an integral equivariant decomposition
\[
 M_n\simeq\mathbb Z_3(1)\oplus
 \bigoplus_{\substack{d\mid n\\d>1}}
 \left(\operatorname{Ind}_{G_{F_d}}^{G_{\mathbb Q}}
       \mathbb Z_3(1)\right)^{\oplus3},
 \qquad F_d=\mathbb Q(\zeta_d).
\]
No division by a Galois-group order is involved. Shapiro's lemma reduces
the question to $S$-unit localization in these fields and in $\mathbb Q$;
the local target is the product over all places above three.
Write $d=3^a q^b$.  The decomposition subgroup at $q$ is
\[
 \langle q\bmod3^a\rangle\times(\mathbb Z/q^b\mathbb Z)^\times.
\]
Its image in $\operatorname{Gal}(F_d^+/\mathbb Q)$ is the whole group
exactly when $q$ and $-1$ generate $(\mathbb Z/3^a\mathbb Z)^\times$.
For $a\le1$ this always holds.  For every $2\le a\le r$ it holds
exactly when $v_3(q^2-1)=1$, equivalently when $q\equiv2,4,5,7\pmod9$.
Under (R), every $F_d^+$ has a unique prime above $q$, so
Lemma~\ref{lem:root-real-q-units} proves the desired injectivity.

Conversely suppose $r\ge2$ and $q\equiv1$ or $8\pmod9$.
Set $F=\mathbb Q(\zeta_{3^r})$ and $E=F^+$, of degree $g$.
There are $t\ge3$ primes of $E$ above $q$: its decomposition group
has index divisible by three in the cyclic real cyclotomic group.
The $q$-unit group of $E$ has rank $g-1+t$, whereas the product of
local unit groups at all its places above $3$ has $\mathbb Z_3$-rank $g$.
All $q$-units have zero $3$-valuations.  Thus the localization kernel
on their $3$-completion has rank at least $t-1>0$.
Restriction from $E$ to $F$ is injective on integral first cohomology,
because corestriction composed with restriction is multiplication by two.
The resulting nonzero kernel in $F$ occurs in the induced summands of
$M_n$ corresponding to $d=3^r$.  Hence the displayed map is not
injective.  This is a cohomological obstruction; it does not produce
a Selmer section with the additional local geometric conditions.
\end{proof}

\begin{theorem}[A compact Selmer criterion for mixed abelian cuspidality]
\label{thm:root-real-mixed-criterion}
For $n=3^r q^s$ as above satisfying
\eqref{eq:root-real-prime-condition}, suppose
\[
 \ker\!\left(
 \operatorname{Sel}(T_3J_n)
 \longrightarrow H^1(\mathbb Q_3,T_3J_n)
 \right)=0.
 \tag{$\mathrm C_n$}
\]
Then every Selmer section in the geometrically
pro-$3$ abelian quotient of $U_n/\mathbb Q$ is
globally cuspidal.  Its center and parameter are
unique.  Thus the Selmer set consists of the
three rational cusp packets, each parametrized
by $H^1(\mathbb Q,\mathbb Z_3(1))$.

A sufficient arithmetic condition for
$(\mathrm C_n)$ is rank $J_n(\mathbb Q)=0$
and $T_3\Sha(J_n/\mathbb Q)=0$.
The conjunction of these sufficient conditions is not asserted here
for any mixed-exponent Jacobian.
\end{theorem}
\begin{proof}
The compact cubic detector identifies the local
center at $3$ with one of the three rational
coordinate cusps $P$ of $X_n$: this uses only
the morphism $X_n\to X_3$ and the established
cubic Selmer input, not injectivity of a map
between geometric pro-$3$ groups.

Use a standard tangent at $P$ as origin.
The compact class $a$ of the section belongs
to $\operatorname{Sel}(T_3J_n)$ and has
$a_3=0$.  Hypothesis $(\mathrm C_n)$ gives
$a=0$ globally.

At every $p\notin S$, the pair $(X_n,D_n)$
has good reduction.  If a local representative
is a cusp $Q$, or an ordinary point reducing
to a cusp $Q$, its compact Kummer class is
that of $Q$.  In the ordinary case this uses
the pro-$p$, uniquely $3$-divisible formal
reduction kernel and the unique lift of the
finite \'etale boundary point.  Since this
class equals that of $P$, the $3$-primary
Kummer class of $Q-P$ is zero.
If $Q\ne P$, Lemma~\ref{lem:lm-exact-cusp-order}
makes its order exactly $n$, divisible by $3$;
its nonzero $3$-primary torsion component
cannot vanish in local $3$-completion.
Thus $Q=P$.

The zero compact class has a unique lift
$b\in H^1(\mathbb Q,B_n)$, because
$(T_3J_n)^{G_{\mathbb Q}}=0$.  Write
$b=(\beta,b_M)$ using the splitting of
Proposition~\ref{prop:root-real-boundary-range}.
The exact local disc inertia calculation
shows that $b_M$ is unramified at every
$p\notin S$: nonzero geometric inertia
can occur only along $I_P$.
At $3$, local cuspidality at $P$ and
$(T_3J_n)^{G_{\mathbb Q_3}}=0$ give
$(b_M)_3=0$.
No unramifiedness at the bad prime $q$
is claimed or needed.
Proposition~\ref{prop:root-real-boundary-range} gives
$b_M=0$ globally.

Hence the section is the cusp section
$t_{P,\beta}$.  The direct summand
$I_P\subset B_n$ and the vanishing of
compact Tate-module fixed points make
its parameter unique.  The cubic
local detector makes its center unique.
Every global cusp section is Selmer.

Finally, under the stated sufficient
arithmetic condition,
$\operatorname{Sel}(T_3J_n)$ is the finite
group $J_n(\mathbb Q)[3^\infty]$.
This group injects into local
$3$-completion at $3$, so
$(\mathrm C_n)$ follows.
\end{proof}

\begin{proposition}[The additional gates for nonabelian quotients]
\label{prop:lm-real-mixed-higher-criterion}
Retain \eqref{eq:root-real-prime-condition} and $(\mathrm C_n)$,
and let $R$ be a
$G_{\mathbb Q}$-stable quotient of $N_n$
retaining $N_n^{\mathrm{ab}}$, with finite
free ordinary lower-central factors
$L_m(R)=\gamma_mR/\gamma_{m+1}R$.
For every $m\ge2$ assume
\[
 L_m(R)^{G_{\mathbb Q_3}}=0,\qquad
 \ker\!\left(
 H^1(G_{\mathbb Q,S},L_m(R))
 \longrightarrow H^1(\mathbb Q_3,L_m(R))
 \right)=0.
 \tag{$\mathrm G_m$}
\]
Then every Selmer section in this quotient
is globally cuspidal.  For a finite
lower-central quotient only its nonzero
grades need be checked.  For a finite-derived
quotient all its nonzero grades are involved.
These additional cohomological conditions
are not proved by the $S$-unit argument.
\end{proposition}
\begin{proof}
The preceding theorem supplies the global
abelian cusp section $t=t_{P,\beta}$.
At $3$ the full local representative is
cuspidal at $P$, with parameter $\beta_3$,
by the same abelian parameter injection.
At $p\notin S$ the compact-center argument
and local inertia formula apply.  If
$v_p(\beta)\ne0$, the full free pro-$3$
normalizer argument makes the local
representative a cusp section conjugate
to $t_p$.  If $v_p(\beta)=0$, both
whole geometric inertias are zero.

For good $p\notin S$, positive Frobenius
weights give $L_m(R)^{G_{\mathbb Q_p}}=0$.
At $3$, this is automatic for $L_1$,
which is an extension of the boundary
Tate module by $T_3J_n$, and is explicitly
assumed for the higher grades.
Thus the quotient $R/\gamma_mR$ has no
fixed points under the action defined
by $t$ at the places needed here.
Inner twisting by $t$ leaves the graded
actions unchanged.

The central induction in the proof of Theorem~\ref{thm:lm-power-all-derived} therefore
makes the next difference cocycle
unramified outside $S$ and locally
a coboundary at $3$.  The original
tangent action on the grade is also
unramified outside $S$, so this
cocycle belongs to
$H^1(G_{\mathbb Q,S},L_m(R))$.
The second part of $(\mathrm G_m)$
makes it a global coboundary.
No condition is imposed on its
inertia at $q$, which is retained
inside the allowed set $S$.
The induction and compactness of
nested conjugator sets identify
the two sections in $R$.
\end{proof}

\begin{remark}[The reduction and its limits]\label{rem:real-mixed-criterion-limits}
This proves that, for the stated two-prime
mixed exponents, compact Selmer localization
injectivity is sufficient for the complete
\emph{abelian} classification.
The boundary obstruction has been removed
by a proved integral localization theorem,
after the other local geometric conditions
restrict its support to $S$.
This does not claim that the unrestricted
relative one-place cohomology kernel is zero;
the previously exhibited nonzero classes
in that kernel need not have the support
and geometry established in this argument.

The criterion is sufficient, not asserted
necessary.  In particular, failure of
$(\mathrm C_n)$ would not construct a
noncuspidal Selmer section.
Nor does $(\mathrm C_n)$ alone prove the
full or metabelian classification: the
higher-grade conditions above remain.
This proves Theorem~\ref{thm:realconditionalmixedintro}.
The proof does not transfer across a
prime-to-$3$ cover by assuming an
arithmetic splitting or a derived-quotient
embedding.
\end{remark}


\subsection{An arithmetic source for the exponents \texorpdfstring{$5\cdot3^r$}{5 times powers of three}}
\label{subsec:lm-fifteen-source}

Let $A=\operatorname{Jac}(C)$, where $C$ is the genus-two curve
$v^2+v=u^5$, and put
\[
 F=\mathbb Q(\zeta_5),\qquad K=\mathbb Q(A[3]),\qquad S=\{3,5\}.
\]
The following calculation controls the full arithmetic action, not only
one rational isogeny factor of a higher Fermat Jacobian.

\begin{proposition}[The residual field]
\label{prop:lm-a5-residual}
The field $K$ contains $F$ and $\mathbb Q(\zeta_3)$, and
\[
 [K:\mathbb Q]=80,\qquad
 \operatorname{Gal}(K/F)\simeq C_{20}.
\]
There is exactly one prime of $K$ above $3$ and exactly one above $5$.
Their ramification indices and residue degrees are respectively
$(20,4)$ and $(40,2)$.  If $H\simeq C_5$ is the unique subgroup of
order five in $\operatorname{Gal}(K/F)$ and $K_0=K^H$, then $K_0/\mathbb Q$
is Galois of degree $16$, and
\begin{equation}
 |D_K|=3^{76}5^{94},\qquad |D_{K_0}|=3^{12}5^{14}.
 \label{eq:lm-residual-discriminants}
\end{equation}
The ramification index of $K_0$ above $5$ is $8$.
\end{proposition}
\begin{proof}
The automorphism $(u,v)\mapsto(\zeta_5u,v)$ gives the full CM ring
$\mathbb Z[\zeta_5]$ on $A_F$. Its CM type is primitive, since the
cyclic quartic field $F$ has no imaginary quadratic subfield; hence
$A$ is absolutely simple.  Since $3$ is inert in $F$,
$A[3]$ is a one-dimensional vector space over
$\mathbb Z[\zeta_5]/3\simeq\mathbb F_{81}$.
If a Galois element acts trivially on $A[3]$, its conjugation on
the reduction of $\zeta_5$ is trivial.  The element $\zeta_5$
has exact order five in $\mathbb F_{81}^{\times}$; hence the
element fixes $F$.  This proves $F\subset K$.
The Weil pairing gives $\mathbb Q(\zeta_3)\subset K$.

The principal polarization has Rosati involution equal to
complex conjugation on $F$.  On $\mathbb F_{81}$ this is
$a\mapsto a^9$.  Consequently the image of $G_F$ is contained in
\[
 \{a\in\mathbb F_{81}^{\times}:a^{10}\in\mathbb F_3^{\times}\}
 =\mu_{20}.
\]
This is a direct upper bound; no general index assertion for CM
representations is needed.

The curve has good reduction at $3$.  The differential basis
$du/(2v+1),u\,du/(2v+1)$ gives CM type $\{1,2\}$ on $F$,
whose reflex type is $\{1,3\}$.  The local CM reciprocity
formula identifies the inertia action at $3$ with the inverse
reflex norm on $\mathcal O_{F_3}^{\times}$.
Since Frobenius acts on $\zeta_5$ by exponent $3$, its reduction is
\[
 \mathbb F_{81}^{\times}\longrightarrow\mathbb F_{81}^{\times},
 \qquad a\longmapsto a^{-4}.
\]
Its image has order $80/\gcd(80,4)=20$.  Thus the global
image over $F$ is exactly $C_{20}$, and the local inertia at
$3$ already fills it.  The field $F_3/\mathbb Q_3$ is
unramified of degree four.  It follows that the decomposition
group at $3$ is the entire group of order $80$, with
ramification index $20$ and residue degree four.
The reciprocity formula used here is stated in
\citet[Theorem~5.1, pp.~483--484]{Lombardo2017}; the
finite unit factor is trivial at a place of good reduction.

For the other ramified prime, use the isomorphism
\[
 x=4u,\qquad y=16(2v+1),\qquad y^2=x^5+256.
\]
\citet[Proposition~2.2(4)--(5), author p.~4]{StollYang2003}
give no conductor at $2$ and, up to inverse and coefficient
conjugation, the following character on
units at $\lambda=1-\zeta_5$:
\[
 \alpha\longmapsto
 \left(\frac{\alpha}{\lambda}\right)_2
 \rho(\alpha)^{q(256)},\qquad
 q(256)=\frac{256^4-1}{5}\equiv4\pmod5.
\]
Here $\rho$ has exact order five and conductor $\lambda^2$;
the quadratic factor has conductor $\lambda$.
Both factors retain their orders modulo $3$.
Thus inertia in $K/F$ at $5$ is $C_{10}$, with wild part
$C_5$ and tame part $C_2$.  The inertia index over
$\mathbb Q_5$ is therefore $4\cdot10=40$.
Because $K$ contains $\zeta_3$, its residue degree at $5$
is at least two.  The product already equals $[K:\mathbb Q]$,
so the prime above $5$ is unique and its residue degree is two.
There is no ramification elsewhere, by good reduction of
$A$ away from $5$ and the prime-to-residue-characteristic
torsion criterion.

Apply the conductor--discriminant formula to the abelian
extension $K/F$ of degree twenty.  At $3$, all nineteen
nontrivial characters have conductor exponent one; the
norm of the prime $(3)$ of $F$ is $3^4$.
At $5$, sixteen characters are nontrivial on the wild
$C_5$ and have conductor exponent two, two further
characters are nontrivial only on the tame $C_2$ and
have exponent one, and one nontrivial character is
unramified.  Thus
\[
 |D_K|=(5^3)^{20}\,3^{4\cdot19}\,5^{16\cdot2+2}
       =3^{76}5^{94}.
\]
The subgroup $H$ is characteristic in $C_{20}$ and hence
normal in $\operatorname{Gal}(K/\mathbb Q)$.
The extension $K_0/F$ has group $C_4$.
Its three nontrivial characters have exponent one at
$3$; at $5$ two have exponent one and one is unramified.
Consequently
\[
 |D_{K_0}|=(5^3)^4\,3^{4\cdot3}\,5^2=3^{12}5^{14}.
\]
Its ramification index at $5$ is $4\cdot2=8$.
\end{proof}

\begin{lemma}[The exact residual group and ramification]
\label{lem:lm-quintic-residual-group}
With the notation of Proposition~\ref{prop:lm-a5-residual},
\[
 \operatorname{Gal}(K/\mathbb Q)=
 \langle\tau,\varphi\mid\tau^{20}=1,\quad
 \varphi\tau\varphi^{-1}=\tau^3,\quad \varphi^4=\tau^{10}\rangle.
\]
Thus its extension of $C_{20}$ by $C_4$ is nonsplit: every lift of the
indicated generator of $C_4$ has order eight. At five the inertia group
is $C_5\rtimes C_8$, with the indicated generator of $C_8$ acting by
third powers on $C_5$. Its lower ramification groups are
\[
 |G_0|=40,\qquad G_1=G_2=C_5,\qquad G_j=1\quad(j\ge3).
\]
\end{lemma}
\begin{proof}
Put $k=\mathbb F_{81}$ and identify $A[3]$ with $k$.
The Rosati adjoint relation writes the alternating Weil pairing as
\[
 \langle x,y\rangle=\operatorname{Tr}_{k/\mathbb F_3}(cxy^9),
 \qquad c\ne0,\quad c^9=-c.
\]
Indeed every linear functional on $k$ is a trace functional, and the
adjoint of a scalar $a$ is $a^9$; alternation gives the condition on $c$.
Every residual Galois element is semilinear, of the form
$x\mapsto ax^{3^i}$. A lift of the third-power generator of
$\operatorname{Gal}(\mathbb Q(\zeta_5)/\mathbb Q)$ has $i=1$.
For its multiplier $m\in\mathbb F_3^\times$, the trace pairing gives
\[
 (ca^{10})^{27}=mc,\qquad a^{10}=mc^2.
\]
Consequently its fourth power is the scalar
\[
 a^{1+3+9+27}=a^{40}=(mc^2)^4=c^8=-1.
\]
For a generator $\tau$ of the scalar subgroup $C_{20}$, this scalar
is $\tau^{10}$, and conjugation acts on $\tau$ by its third power.
The displayed presentation has at most eighty normal forms
$\tau^a\varphi^i$, $0\le a<20$, $0\le i<4$. The already proved degree
is eighty, so it is the exact presentation. Multiplying a lift by a
scalar does not change its fourth power, since $1+3+9+27=40$.

At five the mod-three cyclotomic multiplier is unramified. The known
inertia group of order forty is therefore its kernel. It surjects
onto $C_4$, since the fifth cyclotomic field is totally ramified at
five. A lift in inertia of the indicated generator has order eight,
and together with its normal wild subgroup $C_5$ it generates
$C_5\rtimes C_8$ of order forty.

Finally the unique prime above five has residue degree two. The
exponent ninety-four in the established discriminant therefore makes
the local different exponent forty-seven. All positive lower groups
are subgroups of the wild group $C_5$. If exactly $t$ of them are
nontrivial, the different formula is
$47=(40-1)+t(5-1)$, whence $t=2$. This proves the asserted filtration.
\end{proof}

\begin{theorem}[Local arithmetic source dominance]
\label{thm:lm-fifteen-source}
Let $K_S(3)$ be the maximal pro-$3$ extension of $K$
unramified outside the primes above $3$ and $5$, and put
\[
 \Gamma=\operatorname{Gal}(K_S(3)/\mathbb Q).
\]
For the unique prime $\mathfrak p\mid3$ of $K$, the maps
\[
 G_{K_{\mathfrak p}}\longrightarrow
 \operatorname{Gal}(K_S(3)/K),\qquad
 G_{\mathbb Q_3}\longrightarrow\Gamma
\]
are surjective.
\end{theorem}
\begin{proof}
Set $G=\operatorname{Gal}(K/\mathbb Q)$ and
\[
 V=\ker\bigl(H^1(G_{K,S},\mathbb F_3)
       \longrightarrow H^1(K_{\mathfrak p},\mathbb F_3)\bigr).
\]
The unique prime above $3$ makes $V$ a $G$-module.  This
module is finite dimensional, and is semisimple because
$3\nmid|G|=80$.

First suppose $V^H\ne0$, where $H=C_5$ is as above.
Inflation--restriction for the degree-five extension
$K/K_0$ identifies $H^1(G_{K_0,S},\mathbb F_3)$ with
$H^1(G_{K,S},\mathbb F_3)^H$.  The same restriction is
injective locally, since $[K_{\mathfrak p}:(K_0)_{\mathfrak p_0}]=5$.
Thus a nonzero element gives a cyclic cubic extension
$L_0/K_0$ unramified outside $S$, and split completely
at the prime above $3$.
Only the primes above $5$ can ramify in $L_0/K_0$;
this ramification is tame and has index at most three.
The tame different formula, and $e_5(K_0/\mathbb Q)=8$,
give
\[
 \operatorname{rd}(L_0)
 \le 3^{12/16}5^{14/16}\,5^{(1-1/3)/8}
 =3^{3/4}5^{23/24}<10.659.
\]
But $[L_0:\mathbb Q]=48$, whereas Odlyzko's unconditional
bound at degree at least $48$ is $\operatorname{rd}>13.736$.
This is a contradiction, so $V^H=0$.

If $V\ne0$, choose an irreducible $G$-submodule $W\subset V$.
Every nontrivial irreducible $\mathbb F_3[C_5]$-module has
dimension four: $\Phi_5(T)$ is irreducible over $\mathbb F_3$
because $3$ has order four modulo five.
Since $W^H=0$, we have $\dim W\ge4$.
The common kernel of the characters in $W$ cuts out an
elementary abelian extension $L/K$ of degree
$3^{\dim W}\ge81$.  It is Galois over $\mathbb Q$ because
$W$ is $G$-stable.  It is split at $3$ and unramified
outside $S$.

At $5$, tame inertia in an elementary abelian $3$-extension
is cyclic of exponent three, so its order is at most three.
Using $e_5(K/\mathbb Q)=40$ and
\eqref{eq:lm-residual-discriminants} gives
\[
 \operatorname{rd}(L)
 \le 3^{76/80}5^{94/80}\,5^{(1-1/3)/40}
 =3^{19/20}5^{143/120}<19.329.
\]
On the other hand $[L:\mathbb Q]\ge80\cdot81=6480$.
Odlyzko's unconditional bound at every degree at least
$4800$ is $\operatorname{rd}>21.843$.
This contradiction proves $V=0$.

For precision, these are the degree-$48$ and degree-$4800$ rows
of \citet[Table~2]{Odlyzko1976}. The accompanying description,
paragraphs (i)--(ii), specifies that the complex-column
bound applies to every number field of degree at least
the stated threshold.  No generalized Riemann hypothesis
or numerical class-group assumption enters the argument.

Let $P=\operatorname{Gal}(K_S(3)/K)$ and let $D$ be the
closed image of $G_{K_{\mathfrak p}}$ in $P$.
If $D$ were proper, a finite $3$-group quotient of $P$
would have proper image of $D$; that image lies in an
index-three maximal subgroup.  The resulting nonzero
character $P\to\mathbb F_3$ would vanish on $D$,
contradicting $V=0$.  Thus $D=P$.
Finally $G_{\mathbb Q_3}$ already surjects onto
$\operatorname{Gal}(K/\mathbb Q)$ by
Proposition~\ref{prop:lm-a5-residual}, and its intersection
with $P$ contains all of $P$ by what was just proved.
Hence it surjects onto $\Gamma$.
\end{proof}

\begin{lemma}[The full geometric action over the residual field]
\label{lem:lm-five-power-action}
Let $n=5\cdot3^r$, $r\ge0$, and let $N_n$ be the geometric
pro-$3$ fundamental group of the affine Fermat curve $U_n$.
With a standard unit tangent at any rational coordinate cusp,
the actual $G_{\mathbb Q}$-action on $N_n$ is unramified
outside $S$, factors through $\Gamma$, and has pro-$3$
image on $G_K$.
For every finite free $\mathbb Z_3$-module $M$ obtained
from this action by tensor operations and subquotients,
\[
 H^1(G_{\mathbb Q,S},M)\longrightarrow H^1(\mathbb Q_3,M)
 \quad\text{is injective}.
\]
\end{lemma}
\begin{proof}
First consider $n=5$.  The established three cyclic maps
give a $\mathbb Q$-isogeny $J_5\sim A^3$.
The image of $G_K$ on $T_3A$ lies in the kernel of
$\operatorname{GL}(T_3A)\to\operatorname{GL}(A[3])$,
which is pro-$3$.  The rational Tate-module isomorphism
coming from the isogeny therefore identifies the image
on $V_3J_5$ with a pro-$3$ group.  Being pro-$3$ is a
topological property and is unaffected by changing a
stable lattice.  In particular the image on $J_5[3]$
is a finite $3$-group, so all its composition factors
over $G_K$ are trivial.  This argument does not require
the isogeny degree to be prime to three.

The boundary of $X_5$ is three copies of $\mu_5$.
Because $K$ contains both $\mu_5$ and $\mu_3$, all
composition factors of the boundary module modulo
$3$ are also trivial over $G_K$.
The exact boundary--compact sequence for $N_5^{\mathrm{ab}}$
therefore shows that $G_K$ has a finite $3$-group image
on $N_5/\Phi(N_5)$: a representation all of whose
composition factors are trivial can be put in upper
unipotent form.
The kernel of
$\operatorname{Aut}(N_5)\to\operatorname{Aut}(N_5/\Phi(N_5))$
is pro-$3$ for a finitely generated free pro-$3$ group.
Consequently the actual image on $N_5$ is pro-$3$.

For $n=5\cdot3^r$, the coordinate-power map $U_n\to U_5$
has geometrically Galois group $\mu_{3^r}^2$.
Thus $N_n$ identifies with an open subgroup of $N_5$.
At the rational cusp $(0:-1:1)$ its local parameters
satisfy $u_5=u_n^{3^r}$, so the unit tangents are
compatible and the upstairs action is the restriction
of the downstairs action.  The subgroup is stable for
this action because the map and the compatible tangents
are defined over $\mathbb Q$.
The same statement holds at the other rational cusps
by coordinate permutation.  This is a genuine
$3$-power cover; no prime-to-$3$ cover is treated as an
injection of pro-$3$ groups.

The smooth pair and the chosen extending tangent give
unramified actual action at every prime outside $S$.
Its restriction to $G_K$ has pro-$3$ image, hence
factors through $\operatorname{Gal}(K_S(3)/K)$; the
whole action consequently factors through $\Gamma$.
These properties pass to $M$.

Let $z$ represent a class in $H^1(G_{\mathbb Q,S},M)$.
Over $G_K$ its graph takes values in the pro-$3$ group
$M\rtimes\rho(G_K)$, so both the action and the cocycle
factor through $\Gamma$.
If the local class is zero, subtract one coboundary
so that $z$ is identically zero on $G_{\mathbb Q_3}$.
Theorem~\ref{thm:lm-fifteen-source} makes this cocycle
zero on $\Gamma$, proving injectivity.
\end{proof}

\begin{theorem}[Unconditional mixed exponents $5\cdot3^r$]
\label{thm:lm-five-power-selmer}
For every $r\ge1$ and $n=5\cdot3^r$, the compact
localization criterion $(\mathrm C_n)$ holds:
\[
 \ker\!\left(\operatorname{Sel}(T_3J_n)
       \longrightarrow H^1(\mathbb Q_3,T_3J_n)\right)=0.
\]
Every Selmer section of the geometric pro-$3$ abelian
quotient of $U_n/\mathbb Q$ is globally cuspidal.
The same holds for the full geometric pro-$3$ group,
every finite lower-central quotient, and every finite-derived
quotient (in particular the metabelian quotient).
More generally it holds for any arithmetic-stable quotient
retaining $N_n^{\mathrm{ab}}$ whose ordinary lower-central
factors are torsion-free over $\mathbb Z_3$.
In each case the Selmer set consists of exactly the
three disjoint rational cusp packets parametrized by
$H^1(\mathbb Q,\mathbb Z_3(1))$.
As throughout, a local Selmer representative means the
image of a full point section or a full rational-cusp section;
no global full lift is assumed for a quotient section.
\end{theorem}
\begin{proof}
A compact Selmer class is unramified outside $S$:
there $J_n$ has good reduction and its local Kummer
classes are unramified.  Apply
Lemma~\ref{lem:lm-five-power-action} to $T_3J_n$, a
quotient of $N_n^{\mathrm{ab}}$, to obtain $(\mathrm C_n)$.
The mixed compact criterion then gives the abelian
conclusion.  To spell out its two essential geometric
steps, the morphism $X_n\to X_3$ and the established
cubic compact Selmer input force the local center at
$3$ to be one rational coordinate cusp $P$.
At any $p\notin S$, equality of compact Kummer classes
forces every boundary reduction or cusp center to be
$P$: the difference of distinct coordinate cusps has
exact order $n$, whose $3$-primary part is nonzero.
This is where $r\ge1$ is used.
The boundary component transverse to $I_P$ is therefore
unramified outside $S$ and zero at $3$, so it vanishes
by the same cohomological injection.
The surviving component gives the unique global
cuspidal parameter $\beta$.

For the nonabelian conclusion write
$L_m=\gamma_mR/\gamma_{m+1}R$.
These finite free grades are subquotients of tensor powers of
$N_n^{\mathrm{ab}}$, so their actions factor through $\Gamma$ and
Lemma~\ref{lem:lm-five-power-action} makes their $S$-ramified first
cohomology localize injectively at three. At a good prime, the
compact and boundary factors have strictly positive Frobenius weights,
as do their tensor subquotients. Hence $L_m$ has no global invariants.
Surjectivity of $G_{\mathbb Q_3}\to\Gamma$ gives
$L_m^{G_{\mathbb Q_3}}=0$ as well. We have verified every condition
$(\mathrm G_m)$ in Proposition~\ref{prop:lm-real-mixed-higher-criterion}.
That proposition therefore gives the claimed classification directly
in $R$, using full local representatives and no global lift.
The induction allows arbitrary ramification at five and arbitrary,
possibly infinite, support of the cusp parameter. Only each successive
graded difference is required to be unramified outside $S$.
The free-group and completion lemmas give torsion-free grades for
all the specific quotients listed in the statement.
This proves the $5\cdot3^r$ case of
Theorem~\ref{thm:unrestrictedthreefamiliesintro}.
\end{proof}

\begin{remark}[Scope]\label{rem:lm-five-power-scope}
This proves a mixed-exponent family, including $n=15$,
without computing the Mordell--Weil rank or Tate--Shafarevich
group of $J_{15}$.  It does not assert the same result
for $n=5$, for $5^s3^r$ with $s>1$, or for other
prime-to-three factors.  For $n=5$, the residual/source
calculation remains true but distinct compact cusp
differences have prime-to-three order, so the indicated
boundary-center step is unavailable.  For $s>1$, the
map to $U_5$ has a prime-to-three part and does not
give the open subgroup used in the action lemma.
\end{remark}

\subsection{A class-number certificate for the quintic residual field}

The next calculation proves three-primary class-group vanishing
directly and supplies a second proof of source dominance.
It uses the small-prime term in the unconditional explicit formula.

\begin{proposition}\label{prop:quintic-three-class}
For $K=\Q(\Jac(v^2+v=u^5)[3])$, one has
$\operatorname{Cl}(K)[3]=0$.
\end{proposition}
\begin{proof}
If the three-primary class group were nonzero, class field theory
would give an everywhere-unramified cyclic cubic extension $L/K$.
It would be totally complex of degree $240$, with
\[
 \operatorname{rd}(L)=\operatorname{rd}(K)=3^{76/80}5^{94/80}.
\]
The unique prime $\mathfrak p\mid3$ of $K$ has norm $81$, and
$(3)=\mathfrak p^{20}$. Its ideal class has order dividing twenty,
so its Artin image in $\operatorname{Gal}(L/K)$ is trivial.
Consequently $L$ has three distinct primes of norm $81$.

Define an even function $H$, zero outside $[-2,2]$, by
\[
 H(t)=\frac{(2-t)(1+\tfrac12\cos\pi t)}3
       +\frac{\sin\pi t}{2\pi}\quad(0\le t\le2).
\]
For $g(x)=1+\cos(\pi x)$ on $[-1,1]$ and $g(x)=0$ elsewhere,
direct integration gives
\[
 H(t)=\frac13\int_{\mathbb R}g(x)g(x-t)\,dx.
\]
Thus $H\ge0$, $H(0)=1$, $H\le1$ by Cauchy--Schwarz, and its
Fourier transform is nonnegative. Set $b=51/5$ and
\[
 F(x)=\frac{H(x/b)}{\cosh(x/2)},\quad
 I=\int_0^\infty\frac{1-F(x)}{2\sinh(x/2)}\,dx,
 \quad \log B=\log(8\pi)+\gamma-I.
\]
The explicit formula and unconditional positivity criterion of
\citet[Eqs.~(2.3)--(2.4), p.~122]{Odlyzko1990} give, for any
totally complex field of degree $d$,
\[
 \log\operatorname{rd}\ge\log B+\frac{f-E}{d},\qquad
 E=4\int_0^\infty F(x)\cosh(x/2)\,dx,
\]
where
\[
 f=2\sum_{\mathfrak l}\sum_{m\ge1}
 \frac{\log N\mathfrak l}{(N\mathfrak l)^{m/2}}
 F(m\log N\mathfrak l).
\]
The autocorrelation verifies the positivity condition throughout
the critical strip; no hypothesis on the zeros is imposed.
Since $\int_0^2H(t)\,dt=2/3$, we have $E=8b/3=136/5$.
Every prime contribution is nonnegative. Retaining only the first
powers of the three known norm-$81$ primes gives
\[
 f\ge f_{81}:=
 \frac{6\log81}{9\cosh(\log81/2)}H(\log81/b).
\]

Here are the complete endpoint and tail specifications of the
certified integral. Take $\epsilon=1/16384$. Integrate the displayed
meromorphic expression for $(1-F(x))/(2\sinh(x/2))$ on
$[\epsilon,2b]$ by complex-ball quadrature at $192$ bits, with
absolute and relative tolerances $2^{-120}$, using the certified
Arb integration algorithm described by \citet{Johansson2018}.
Beyond $2b$, its integral is exactly $-\log\tanh(b/2)$.
At the origin, $H'(0)=0$ and
$|H''(t)|\le\|g'\|_2^2/3=\pi^2/3<6$.
For $0\le x\le\epsilon$ this gives
\[
 1-H(x/b)\le3x^2/b^2,
 \qquad 1-\operatorname{sech}(x/2)\le x^2/8.
\]
As $b\ge2$ and $2\sinh(x/2)\ge x$, the integrand lies between
zero and $x$. Its omitted endpoint integral is therefore at most
$\epsilon^2/2$. Subtracting this upper bound from the expression
for $\log B$ gives a certified lower bound $\log B_-$.

Arb ball arithmetic, with exact rational inputs and thresholds,
certifies
\begin{gather*}
 B_->\frac{21046}{1000},\qquad f_{81}>\frac{472}{1000},\\
 \exp\!\left(\log B_-+\frac{f_{81}-136/5}{240}\right)
       >\frac{18828}{1000},\qquad
 3^{76/80}5^{94/80}<\frac{18818}{1000}.
\end{gather*}
For orientation, the two last enclosures are respectively
$18.8285116161758712974\ldots$ and
$18.8171976695261105380\ldots$; the proof uses the rational
inequalities. The explicit-formula lower bound for $L$ exceeds
its established root discriminant, a contradiction.
\end{proof}

\begin{corollary}[A ray-class proof of source dominance]
\label{cor:quintic-ray-source}
Proposition~\ref{prop:quintic-three-class} implies the two
surjections in Theorem~\ref{thm:lm-fifteen-source}.
\end{corollary}
\begin{proof}
By the Frattini criterion it is enough to exclude a nontrivial cyclic
cubic extension of $K$, unramified outside $3,5$, and split at its
prime above three. Such an extension is unramified at three and
has at most tame conductor exponent one at the unique prime
$\mathfrak q\mid5$. Its group is a quotient of the ray class group
modulo $\mathfrak q$. The residue group is $\mathbb F_{25}^{\times}$,
whose three-primary part has order three. The global unit $\zeta_3$
reduces to an element of that exact order, so the ray-class exact
sequence and Proposition~\ref{prop:quintic-three-class} leave no
three-primary quotient. The local decomposition at three already
surjects onto $\operatorname{Gal}(K/\Q)$ by
Proposition~\ref{prop:lm-a5-residual}; the full source surjection follows.
\end{proof}


\subsection{A class-group gate for the septic residual source}

\begin{lemma}[Split-at-three ray classes]
\label{lem:lm-septic-ray-gate}
Let $K/\mathbb Q$ be Galois and contain $\mu_3$.
Suppose it has one prime $\mathfrak p$ above $3$, with
$(3)=\mathfrak p^e$ and $3\nmid e$, and exactly two primes
$\mathfrak q_1,\mathfrak q_2$ above $7$, each of residue
degree two.  Put $S=\{v:v\mid21\}$.
Then
\[
 \ker\bigl(H^1(G_{K,S},\mathbb F_3)
       \longrightarrow H^1(K_{\mathfrak p},\mathbb F_3)\bigr)
 \simeq\operatorname{Hom}(\operatorname{Cl}(K),\mathbb F_3).
\]
In particular, the local group at $3$ surjects onto the
maximal pro-$3$ extension unramified outside $S$ if and
only if $\operatorname{Cl}(K)[3]=0$.
\end{lemma}
\begin{proof}
A character in the kernel cuts out a cyclic cubic extension
unramified at $3$, split at $\mathfrak p$, and unramified
outside the two primes over $7$.  Its possible ramification
there is tame, so its conductor divides
$\mathfrak m=\mathfrak q_1\mathfrak q_2$.
Class field theory identifies these characters with the
characters of the ray class group modulo $\mathfrak m$,
quotiented by the class of $\mathfrak p$.
The relevant exact sequence is
\[
 \mathcal O_K[1/3]^\times\longrightarrow
 (\mathcal O_K/\mathfrak m)^\times\longrightarrow
 C_{\mathfrak m}^{\{\mathfrak p\}}\longrightarrow
 \operatorname{Cl}(\mathcal O_K[1/3])\longrightarrow0.
\]
Each residue field is $\mathbb F_{49}$, whose multiplicative
group has three-primary part of order three.  The two
primes restrict to the distinct primes of
$\mathbb Q(\zeta_3)$ above $7$: otherwise the Galois action
on primes and the existence of both restrictions would
give more than two primes in $K$.
Choose the common generator $2$ of the order-three subgroup
in $\mathbb F_7^\times$ inside each residue field.
The three-primary projection of $\zeta_3$ has coordinate
vectors $(1,2)$ or $(2,1)$; that of the rational
$S$-unit $3$ has equal nonzero coordinates, because
$3^{16}=4$ in $\mathbb F_7$.
These two vectors are linearly independent over $\mathbb F_3$.
Thus the images of $\mu_3$ and $3$ generate the entire
three-primary residue group.
The three-primary part of $C_{\mathfrak m}^{\{\mathfrak p\}}$
therefore equals that of
$\operatorname{Cl}(\mathcal O_K[1/3])$.
Finally $[\mathfrak p]$ has order dividing $e$, which is
prime to three.  Removing $\mathfrak p$ does not change
the three-primary class group, proving the isomorphism.
The last assertion is the Frattini criterion for a
pro-$3$ group, as in the quintic residual-source proof.
\end{proof}

\begin{proposition}[Residual data and the class-group reduction]
\label{prop:lm-septic-residual-gate}
Let $A_7=\operatorname{Jac}(v^2+v=u^7)$,
$F=\mathbb Q(\zeta_7)$, and $K=\mathbb Q(A_7[3])$.
Then $K$ satisfies the hypotheses of
Lemma~\ref{lem:lm-septic-ray-gate}, with $e=56$.
More precisely,
\[
 [K:\mathbb Q]=336,\quad\operatorname{Gal}(K/F)=C_{56},
 \quad |D_K|=3^{330}7^{380}.
\]
At $3$ there is one prime with $(e,f)=(56,6)$.
At $7$ there are two primes, each with $(e,f)=(84,2)$.
If $H=C_7\subset C_{56}$ and $K_0=K^H$, then
\[
 [K_0:\mathbb Q]=48,\qquad |D_{K_0}|=3^{42}7^{44},
 \qquad\operatorname{Cl}(K_0)[3]=0.
\]
The $H$-invariant part of the localization kernel for $K$
is zero.  Its possible remaining part has
$\mathbb F_3$-dimension divisible by six; its vanishing
is exactly the condition
$\operatorname{Cl}(K)[3]=0$.
\end{proposition}
\begin{proof}
The same CM endomorphism argument as for $A_5$ gives
$F,\mathbb Q(\zeta_3)\subset K$.
The prime $3$ is inert in $F$, and the residual module is
one dimensional over $\mathbb F_{729}$.
Polarization bounds the $G_F$-image by
$\{a:a^{28}\in\mathbb F_3^\times\}=C_{56}$.
The differential type is $\{1,2,3\}$, with reflex
$\{1,4,5\}$; local CM reciprocity at the good prime
$3$ gives the power $-(1+3^4+3^5)=-325$.
Since $\gcd(325,728)=13$, its image has order $56$.
Thus local inertia fills the global $G_F$-image, and
the local decomposition group at $3$ has order
$56\cdot6=336$, the whole global group.

\citet[Proposition~3.12(2), author p.~16]{ShuExpanded2021}, applied
to $(r,s,t,\delta)=(1,1,5,1)$, gives conductor $\lambda_7^2$.
His Equation~(3.14) on that page expresses the local
unit character as the product of a seventh-power
Hilbert-symbol character and a quadratic character of
conductor one.  The exact conductor two makes the
order-seven factor nontrivial.  Thus the unit image is
$\mu_{14}$, retaining its order modulo three, and
$e_7(K/F)=14$.

For the residue degree, take $\lambda=1-\zeta_7$.
The global reciprocity formula on the principal idele
$\lambda$ shows that its local character at $7$, modulo
inertia, is the reflex product
\[
 \alpha=(1-\zeta_7)(1-\zeta_7^4)(1-\zeta_7^5)
       \quad\text{in }\mathbb F_{729}^{\times},
\]
up to inversion and coefficient conjugation.
Indeed all other finite unit factors are trivial by
good reduction, while the local-three factor is the
inverse reflex norm.
In the field
$\mathbb F_3[z]/(z^6+z^5+z^4+z^3+z^2+z+1)$, direct
reduction gives
\[
 \alpha^2=-z^3,\qquad \alpha^{14}=-1,\qquad
 \alpha^{28}=1.
\]
Hence its class modulo $\mu_{14}$ has order two.
Therefore $f_7=2$, and the number of primes is
$336/(84\cdot2)=2$.
In particular the prime above $7$ is \emph{not} unique.

The conductor--discriminant calculation has fifty-five
nontrivial characters of conductor exponent one at
$3$, whose prime norm in $F$ is $3^6$.
At $7$, forty-eight characters have conductor exponent
two, four have exponent one, and three nontrivial
characters are unramified.  As $|D_F|=7^5$, this gives
\[
 |D_K|=7^{5\cdot56}\,3^{6\cdot55}\,7^{2\cdot48+4}
       =3^{330}7^{380}.
\]
The corresponding quotient $K_0/F$ has degree eight,
inertia of order two at $7$, and gives
$|D_{K_0}|=3^{42}7^{44}$.
The stronger class-number-one assertion in
Corollary~\ref{cor:septic-auxiliary-class-one} below proves
the required three-class vanishing, using only this
discriminant and the explicit-formula test function used above.

The field $K_0$ has the same two residue-degree-two
primes above $7$ and one prime of ramification index
eight above $3$.  Lemma~\ref{lem:lm-septic-ray-gate}
therefore makes its entire split-at-three localization
kernel zero.  Inflation--restriction through the
degree-seven extension $K/K_0$, together with local
restriction of degree seven, identifies this with the
$H$-invariant kernel for $K$.
Finally $\Phi_7$ is irreducible over $\mathbb F_3$,
since $3$ has order six modulo seven.  All nontrivial
$\mathbb F_3[C_7]$-constituents have dimension six.
The earlier ray-class lemma identifies their possible
sum with the remaining three-class-group gate.
\end{proof}

\begin{corollary}[The auxiliary class number]\label{cor:septic-auxiliary-class-one}
For the degree-$48$ field $K_0$ above, one has $h(K_0)=1$.
\end{corollary}
\begin{proof}
If its class number were greater than one, its Hilbert class
field would be a totally complex field of degree $d\ge96$
with the same root discriminant
$3^{7/8}7^{11/12}$. Use the kernel of
Proposition~\ref{prop:quintic-three-class} with $b=51/5$,
and discard every nonnegative prime contribution. The
unconditional explicit formula gives
\[
 \operatorname{rd}\ge B\exp\!\left(-\frac{136}{5d}\right),
 \qquad d\ge96.
\]
The endpoint and tail bounds for $B$ are exactly those already
specified there. Re-evaluating that integral with $192$-bit
ball arithmetic gives the rational certified inequalities
\[
 B\exp(-17/60)>\frac{1585}{100},
 \qquad 3^{7/8}7^{11/12}<\frac{1557}{100}.
\]
The lower bound increases with $d$. The incompatible bounds
exclude the Hilbert class field and prove the assertion.
The calculation does not compute the class group of $K$ or that of
the degree-$56$ field $M$ introduced below.
\end{proof}

\begin{remark}[Why the smaller field is used]
\label{rem:septic-smaller-field}
The root discriminant of $K$ is approximately $26.5685$, larger
than the unconditional asymptotic complex-field discriminant threshold.
The quintic discriminant argument therefore does not exclude a
nontrivial $C_7$-constituent by its dimension alone. We will first
control the whole geometric action, identify the field, and reduce
its three-class group to a degree-$56$ field. Principal prime ideals
in that smaller field then supply the stronger explicit-formula bound
needed to prove vanishing.
\end{remark}


\subsection{The whole septic geometric action}

Put $F=\mathbb Q(\zeta_7)$ and $F_0=\mathbb Q(\sqrt{-7})$,
and write $H=\operatorname{Gal}(F/F_0)=\{1,2,4\}$.
Let
\[
 A=\operatorname{Jac}(v^2+v=u^7),\qquad
 B=\operatorname{Jac}(y^7=x(1-x)^2),
\]
and let $E$ be the elliptic curve
$y^2+xy=x^3-x^2-2x-1$ of conductor $49$ and CM by
$\mathcal O_{F_0}$.  Its discriminant is $-7^3$ and
its $j$-invariant is $-3375$.
The substitution $(x_0,y_0)=(x-2,2y+x)$ gives
$y_0^2=4x_0^3+21x_0^2+28x_0$, the classical quotient in
\citet[Eq.~(2.10), pp.~68--69]{Elkies1998}.
All character comparisons below are made over $F$.
No rational splitting $B\sim E^3$ is asserted.

\begin{lemma}[Two distinct norm operations]
\label{lem:lm-septic-cm-comparison}
Let $\eta_A,\eta_B$ be the Jacobi Hecke characters over
$F$ for the triples $(1,1,5)$ and $(1,2,4)$, and let
$\eta_E$ be the CM Hecke character over $F_0$, choosing
the one of its two embeddings compatible with $H$.
Define
\[
 \eta_E^F(\mathfrak a)=
 \eta_E(N_{F/F_0}^{\mathrm{ideal}}\mathfrak a),\qquad
 \nu(\mathfrak a)=N_{F/\mathbb Q}\mathfrak a.
\]
Then, on ideals prime to $7$,
\[
 \eta_B=\eta_E^F,\qquad
 N_{F/F_0}^{\mathrm{coefficient}}(\eta_A)
       =\nu\,\eta_E^F.
 \tag{$\dagger$}\label{eq:lm-septic-character-identities}
\]
The superscripts distinguish the norm of the input ideal
from the norm of a character value.
\end{lemma}
\begin{proof}
The class number of $F$ is one.  Indeed, Minkowski's
bound is less than five, while primes above $2$ have
norm eight and primes above $3$ have norm $3^6$.
There is no nontrivial ideal of norm at most four.

For $A$, the reflex infinity type is
$R=\{1,4,5\}$; for $B$, it is $H=\{1,2,4\}$.
The latter is exactly the infinity type of $\eta_E^F$.
The CM character of $E/F_0$ has conductor the prime
above seven: the induction conductor formula is
$49=7\cdot7$.
Its local unit character is the nontrivial quadratic
residue character.  Pullback through the degree-three
totally ramified extension $F/F_0$ preserves that
quadratic unit character.

For $B$, \citet[Proposition~3.12(1), author p.~16]{ShuExpanded2021} applies because
\[
 1^1 2^2(4-7)^4=324,\qquad
 \frac{324^6-1}{7}\equiv0\pmod7.
\]
Thus $\eta_B$ has conductor exponent one at seven.
Equation~(3.14) expresses its unit character as a
seventh-power Hilbert-symbol character times the
quadratic character of conductor one.
The order-seven factor is trivial on units in the
conductor-one case.  Hence its unit character is
the same quadratic character as that of $\eta_E^F$.
Both characters are unramified away from seven.
Their quotient therefore has trivial infinity type
and is unramified at every finite place.  It is a
finite ideal-class character, so class number one
makes it trivial.  This proves the first identity.

For the second identity, the coefficient norm
multiplies the three coefficient conjugates indexed
by $H$.  Of the elements of $R$, two lie in $H$ and
one in its other coset.  Its infinity type is
therefore twice $H$ plus once $-H$.
Dividing by the ideal norm character $\nu$, whose
infinity type is once each of the six embeddings,
leaves exactly the type $H$ of $\eta_E^F$.

At seven, the finite unit character of $\eta_A$
has image $\mu_{14}$.  Coefficient norm to $F_0$
kills $\mu_7$, since $1+2+4=7$, and sends $-1$
to $-1$.  It therefore leaves precisely the same
quadratic unit character.  The character $\nu$ is
unramified on units.  Away from seven all characters
are unramified.  The quotient of the two sides of
the second identity is again an everywhere-unramified
finite ideal-class character of $F$, hence trivial.
\end{proof}

\begin{lemma}[The integral septic isogeny]\label{lem:septic-seven-isogeny}
The five cyclic quotient maps of $X_7$ induce homomorphisms
$\Phi:J_7\to A^3\times B^2$ and $\Psi:A^3\times B^2\to J_7$
such that $\Phi\Psi=[7]$ and $\Psi\Phi=[7]$. Consequently
\[
 T_3J_7\simeq(T_3A)^{\oplus3}\oplus(T_3B)^{\oplus2}
\]
integrally and $G_\Q$-equivariantly.
\end{lemma}
\begin{proof}
For $a=1,\ldots,5$, use the quotient curve
$C_a:y^7=x(1-x)^a$ and the rational morphism determined by
\[
 x=-X^7/Z^7,\qquad y=(-1)^{a+1}XY^a/Z^{a+1}.
\]
Its geometric deck subgroup in $\mu_7^2$ is
$H_a=\{(\zeta_7^u,\zeta_7^v):u+av=0\}$.
The thirty characters occurring in $H^1(X_7)$ are the pairs
$(i,j)$ with $i,j,i+j\ne0$ modulo seven. They partition into
the five lines $j=ai$ just listed. The triples for $a=1,3,5$
are equivalent to $(1,1,5)$ by coordinate permutations and
scaling all exponents by a unit modulo seven; those for
$a=2,4$ are equivalent to $(1,2,4)$. The resulting monomial
changes of cyclic coordinates are defined over $\Q$, so
the product of their Jacobians is $A^3\times B^2$.

Let $f_a$ denote these quotient maps. The trace identity gives
$f_a^*f_{a,*}=\sum_{h\in H_a}h$ on $J_7$ and
$f_{a,*}f_a^*=[7]$ on $\Jac(C_a)$. On a character line the
first operator is seven if the line is fixed by $H_a$ and
zero otherwise. Thus $f_{a,*}f_b^*=0$ for $a\ne b$, and
$\sum_a f_a^*f_{a,*}=[7]$. These identities on a rational
Tate module are identities of homomorphisms of abelian varieties,
by faithfulness of that realization. Taking the product of
the pushforwards and the sum of the pullbacks proves both
composition formulas. Since seven is a unit in $\Z_3$, they
give the asserted integral Tate-module isomorphism.
\end{proof}

\begin{proposition}[All compact residual factors are controlled]
\label{prop:lm-whole-septic-action}
Let $K=\mathbb Q(A[3])$.  Then $E[3]$ is rational over
$K$, and the image of $G_K$ on the entire rational
Tate module of the degree-seven Fermat Jacobian is
pro-$3$.  Consequently, with a standard unit tangent,
the actual action of $G_K$ on the entire geometric
pro-$3$ group $N_7$ has pro-$3$ image.
The same is true for $N_{7\cdot3^r}$, $r\ge0$.
All these actions are unramified outside $\{3,7\}$.
\end{proposition}
\begin{proof}
On $G_F$, the second identity in
\eqref{eq:lm-septic-character-identities} gives the
$3$-adic character identity
\[
 \rho_{E,3}=
 \chi_3^{-1}N_{F_3/(F_0)_3}^{\mathrm{coefficient}}(\rho_{A,3}),
\]
up to the harmless choice of the conjugate CM embedding.
The field $K$ contains $F$ and $\mu_3$.
On $G_K$, both $\rho_{A,3}$ and $\chi_3$ are congruent
to one modulo three, so the same is true of
$\rho_{E,3}$.  Thus $E[3]\subset E(K)$.
The images on $T_3A$ and $T_3E$ over $G_K$ are pro-$3$.

Lemma~\ref{lem:septic-seven-isogeny} gives the decomposition
of $V_3J_7$ into three copies of $V_3A$ and two of $V_3B$.
After restriction to $G_F$, the first character
identity in \eqref{eq:lm-septic-character-identities}
identifies $V_3B$ with three copies of $V_3E$:
the six coefficient embeddings yield three copies
of each of the two $F_0$-embeddings.
Hence $V_3J_7|_{G_K}$ has pro-$3$ image.
This topological property does not depend on the
choice of a stable integral lattice.

The boundary of the Fermat curve consists of three
copies of $\mu_7$.  Since $K$ contains both $\mu_7$
and $\mu_3$, the boundary modulo three is trivial
over $G_K$.  The compact and boundary terms show
that the image on $N_7/\Phi(N_7)$ is a three-group.
The automorphism kernel over the Frattini quotient
is pro-$3$, proving the assertion for the full $N_7$.

Finally $U_{7\cdot3^r}\to U_7$ is a genuine geometric
three-power cover.  Compatible unit tangents,
$u_7=u_{7\cdot3^r}^{3^r}$, identify the upstairs group
with an arithmetic-stable open subgroup, so the
action statement passes upstairs.  Good logarithmic
reduction of the pair and the extending tangent
give unramified actual action outside $\{3,7\}$.
\end{proof}

\begin{corollary}[A finite sufficient gate for $7\cdot3^r$]
\label{cor:lm-septic-class-gate}
If $\operatorname{Cl}(\mathbb Q(A_7[3]))[3]=0$, then
for every $r\ge1$ all full, abelian, finite lower-central,
and finite-derived geometric pro-$3$ Selmer sections
of $U_{7\cdot3^r}/\mathbb Q$ are globally cuspidal,
with the three rational cusp packets as before.
The same conclusion holds for an arithmetic-stable
quotient retaining abelianization and having finite
free ordinary lower-central factors. Moreover the compact
integral Selmer localization at three is injective for each such $n$.
\end{corollary}
\begin{proof}
The ray-class lemma makes the local map at three
surjective onto the maximal $\{3,7\}$-ramified
pro-$3$ arithmetic source over $K$; the residual
decomposition group already equals
$\operatorname{Gal}(K/\mathbb Q)$.
The preceding proposition places every compact,
boundary, and central coefficient action in that
source.  The cocycle-graph argument therefore gives
the required integral cohomology localization
injections.  The cubic compact detector, exact
three-primary order of cusp differences, and the
same central induction as for $5\cdot3^r$ apply.
More explicitly, the coefficient-graph homomorphism proves
$H^1(G_{\Q,\{3,7\}},L)\hookrightarrow H^1(\Q_3,L)$
for each finite free coefficient module just described, exactly as
in Lemma~\ref{lem:lm-five-power-action}. Compact Selmer classes
are unramified outside $\{3,7\}$ by good reduction, proving
the compact assertion. Positive good-prime weights give zero
global invariants in every lower-central factor; source surjectivity
gives the same at three. Thus all hypotheses of
Proposition~\ref{prop:lm-real-mixed-higher-criterion} hold with $q=7$.
The class-group hypothesis is verified below in
Proposition~\ref{prop:septic-three-class}, using the explicit field
calculation and degree-$56$ reduction. This corollary then supplies
the septic case of Theorem~\ref{thm:unrestrictedthreefamiliesintro}.
\end{proof}

\subsection{An explicit description of the septic residual field}

Let $A=\operatorname{Jac}(v^2+v=u^7)$,
$F=\mathbb Q(\zeta_7)$, $\lambda=1-\zeta_7$, and
$E:y^2+xy=x^3-x^2-2x-1$.
Retain the residual field and character calculations of the
preceding septic lemmas: $K=\mathbb Q(A[3])$ contains $F$,
$\operatorname{Gal}(K/F)=C_{56}$, its conductor divides
$(3)\lambda^2$, and the coefficient-norm identity controls
$E[3]$.

\begin{proposition}\label{prop:lm-explicit-septic-field}
There is an equality of subfields of an algebraic closure of
$\mathbb Q$,
\[
 K=\mathbb Q\bigl(\zeta_7,\sqrt[7]{3},E[3]\bigr).
\]
Its degree is $336$.  The intermediate field fixed by the
normal subgroup of order seven is
$F(E[3])$, of degree $48$ over $\mathbb Q$.
\end{proposition}
\begin{proof}
On the residual coefficient field $\mathbb F_{3^6}$, the
coefficient norm to $\mathbb F_{3^2}$ is $a\mapsto a^{91}$.
Polarization identifies the mod-three cyclotomic character
on $C_{56}$ with $a\mapsto a^{28}$.  The character identity
for $E$ therefore gives $a\mapsto a^{91-28}=a^7$ on
$C_{56}$.  Thus $F(E[3])/F$ is precisely its degree-eight
subextension.

The seven-primary part of the ray class group of $F$ for
the modulus $(3)\lambda^2$ has order at most seven.
Indeed, $F$ has class number one.  The seven-primary parts
of $(\mathcal O_F/(3))^\times$ and
$(\mathcal O_F/\lambda^2)^\times$ are each cyclic of
order seven.  The global group $\mu_7$ maps nontrivially
to both factors: reduction at three preserves order seven,
and $\zeta_7=1-\lambda$ is nontrivial modulo $\lambda^2$.
The quotient of these two local factors by its diagonal
image has order seven.  Further global units can only
make the ray quotient smaller.  The degree-seven
subextension of $K/F$ shows that it is not smaller.
There is consequently a unique cyclic extension of degree
seven with conductor dividing this modulus.

We verify directly that $F(\sqrt[7]{3})$ has this conductor.
It has degree seven, since the valuation of $3$ at the
unique prime over three is one.  It is cyclic because
$\mu_7\subset F$.  At three it is totally and tamely
ramified of degree seven, and hence has conductor exponent
one.  It is unramified at primes other than three and seven
by the Kummer equation.

At seven put $\alpha=\sqrt[7]{3}$ and $\beta=\alpha-3$.
The polynomial
\[
 (T+3)^7-3
\]
is Eisenstein over $\mathbb Q_7$: all its nonleading
coefficients are divisible by seven and
$v_7(3^7-3)=1$.  Thus
$L_7=\mathbb Q_7(\alpha)$ is totally ramified of degree
seven with $\mathcal O_{L_7}=\mathbb Z_7[\beta]$.
Its different has exponent seven, since the derivative
is $7\alpha^6$ and $\alpha$ is a unit.

Let $F_7=\mathbb Q_7(\zeta_7)$ and
$M_7=F_7L_7$.  The coprime totally ramified degrees six
and seven give $[M_7:\mathbb Q_7]=42$ and
$e(M_7/L_7)=6$.  This last extension is tame, with
different exponent five.  Transitivity of the different
therefore gives
\[
 d(M_7/\mathbb Q_7)=5+6\cdot7=47,
 \qquad
 d(M_7/F_7)=47-7\cdot5=12.
\]
For a cyclic degree-seven extension, all six nontrivial
characters have the same conductor and the
conductor--discriminant formula gives their common
conductor exponent $12/6=2$.
This proves the required local conductor at $\lambda$.

The uniqueness from the ray group now identifies
$F(\sqrt[7]{3})$ with the degree-seven subextension of
$K/F$.  Its degree is coprime to that of $F(E[3])/F$,
so their compositum is $K$.  The degree assertions
follow from $6\cdot7\cdot8=336$.
\end{proof}

\begin{corollary}[The split residual group]\label{cor:septic-residual-group}
The full residual Galois group is
\[
 \operatorname{Gal}(K/\Q)\simeq C_{56}\rtimes_3 C_6,
\]
where a generator of $C_6$ acts on $C_{56}$ by cubing.
In particular the extension by $\operatorname{Gal}(F/\Q)$ splits.
\end{corollary}
\begin{proof}
Put $T=\Q(E[3])$ and $F_0=\Q(\sqrt{-7})$.
The CM endomorphism action on $E[3]$ shows that $F_0\subset T$.
At the good prime three, which is inert in $F_0$, the local
CM reciprocity formula gives the full scalar inertia
$\mathbb F_9^\times=C_8$: the rank-one CM reflex norm is the
identity. Hence $\operatorname{Gal}(T/F_0)=C_8$.
Complex conjugation has order two and acts on this scalar
group by cubing, so $\operatorname{Gal}(T/\Q)=C_8\rtimes_3 C_2$.
These uses of good-prime reciprocity are the same
all-units formula as in \citet[Theorem~5.1, pp.~483--484]{Lombardo2017}.

The field $F(\sqrt[7]{3})$ has group $C_7\rtimes_3 C_6$:
the lift sending $\zeta_7$ to $\zeta_7^3$ and fixing the
chosen real seventh root has order six. Its intersection
with $T$ is precisely $F_0$, since both fields contain $F_0$
and their degrees are $42$ and $16$. Their compositum is
$K$ by Proposition~\ref{prop:lm-explicit-septic-field}.
Its Galois group is therefore the fiber product of these
two split groups over $C_2=\operatorname{Gal}(F_0/\Q)$.
Pairing the indicated order-six lift with complex conjugation
in the second factor gives an order-six lift in the fiber
product. It acts by cubing on both $C_7$ and $C_8$,
and therefore on $C_7\times C_8=C_{56}$. This proves the claim.
\end{proof}

\paragraph{An explicit polynomial specification.}
The three-division polynomial of $E$ is
\[
 f(X)=3X^4-3X^3-12X^2-12X-1.
\]
For $Z=2y+x$, put $g(X)=4X^3-3X^2-8X-4$.
An exact resultant computation gives
\[
 R(Z)=\operatorname{Res}_X(f(X),Z^2-g(X))
 =27Z^8-1323Z^6+6174Z^4-117649.
\]
The preceding subresultant is
\[
 (-36Z^4+1323Z^2-4116)X
       +9Z^4+588Z^2+1029.
\]
Its $X$-coefficient is coprime to $R(Z)$, so it recovers
$x$ from each root $z$, and then $y=(z-x)/2$.
It follows that $\mathbb Q(E[3])$ is the splitting
field of $R$.  Thus $K$ is explicitly the splitting
field of
\[
 (T^7-3)\bigl(27T^8-1323T^6+6174T^4-117649\bigr).
\]
The condition $\operatorname{Cl}(K)[3]=0$ will follow from the
degree-$56$ reduction and the principal-prime certificate below.
The defining polynomials alone do not establish class-group vanishing.

\subsection{An equivalent septic gate in degree fifty-six}

Let $E:y^2+xy=x^3-x^2-2x-1$, let
$\alpha=\sqrt[7]{3}$ be real, and put
\[
 K=\mathbb Q(\zeta_7,\alpha,E[3]),\qquad
 L=\mathbb Q(\alpha,E[3]).
\]
Let $c$ be complex conjugation in this fixed embedding and
write $K'=K^{\langle c\rangle}$ and
$M=L^{\langle c\rangle}$.
We shall verify $[L:\mathbb Q]=112$ and
$D_L=3^{110}7^{124}$ as part of the argument.
The preceding field calculations give $[K:\mathbb Q]=336$ and
$\operatorname{Cl}(K_0)[3]=0$ for the degree-$48$
field $K_0=K^{C_7}=\mathbb Q(\zeta_7,E[3])$.

\begin{proposition}\label{prop:lm-septic-degree56}
Let $z$ be either real root of
\[
 R(Z)=27Z^8-1323Z^6+6174Z^4-117649.
\]
Then
\[
 M=\mathbb Q(z,\alpha),\qquad [M:\mathbb Q]=56,
 \qquad \operatorname{sign}(M)=(2,27),
 \qquad D_M=-3^{55}7^{62}.
\]
Moreover
\[
 \operatorname{Cl}(K)[3]=0
 \quad\Longleftrightarrow\quad
 \operatorname{Cl}(M)[3]=0.
\]
Consequently this explicit degree-$56$ class-group condition
is an equivalent finite gate for the preceding septic
arithmetic-source argument.  Its vanishing is proved below.
\end{proposition}
\begin{proof}
Put $T=\Q(E[3])$ and $F_0=\Q(\sqrt{-7})$.
The CM action embeds $\mathcal O_{F_0}/3=\mathbb F_9$ in
$\operatorname{End}(E[3])$, making $E[3]$ a rank-one module.
The nontrivial automorphism of $F_0$ acts nontrivially on this
endomorphism field, so $F_0\subset T$. Thus
$\operatorname{Gal}(T/F_0)$ is a subgroup of $\mathbb F_9^\times$.
At the good prime three the CM type has one embedding; its
reflex norm on local units is the identity. The local CM reciprocity
formula of \citet[Theorem~5.1]{Lombardo2017} therefore gives the full
scalar group $\mathbb F_9^\times$ as inertia. Consequently
$[T:\Q]=16$, there is one prime over three with $(e,f)=(8,2)$,
and the scalar subgroup acts transitively on the eight nonzero
three-torsion points.

The valuation of $3$ in $T$ is eight, prime to seven, so
$T(\alpha)/T$ has degree seven. This proves $[L:\Q]=112$
and $[K:L]=3$. The latter extension is cyclic and unramified at
three, since the inertia group of $K/\Q$ there has order $56$.
At seven its cubic group lies in decomposition (which has index
two) and then in inertia (whose residue quotient has order two).
Thus both seven-primes are tamely ramified of degree three.
Their residue degrees in $L$ are two, so their contributions to
$N_{L/\Q}\mathfrak d_{K/L}$, the relative discriminant norm, multiply to $7^8$.
The discriminant tower formula gives
\[
 3^{330}7^{380}=D_L^3\,7^8,
 \qquad D_L=3^{110}7^{124}.
\]

First we compare $K'$ and $M$.  The cyclic subgroup
$\operatorname{Gal}(K/L)$ has order three; it fixes
$\alpha$ and $E[3]$ and acts on $\zeta_7$ by a power
of $\zeta_7\mapsto\zeta_7^2$.  It commutes with $c$,
as is checked on these three generators.  Thus $K'/M$
is cyclic of degree three.

The two primes above seven in $K$ are interchanged
by $c$: they are already distinguished by the two
embeddings of $\mu_3$, which lies in $\mathbb Q(E[3])$.
Their contractions to $L$ are likewise interchanged.
Consequently $K'$ and $M$ each have exactly one prime
above seven.  The respective local ramification and
residue indices are $(84,2)$ and $(28,2)$, so that
this prime is totally and tamely ramified in $K'/M$.
At three the cyclic cubic extension is unramified:
inertia in $K/\mathbb Q$ has order $56$, prime to
three.  There is no ramification at other finite
primes, and a cyclic cubic extension cannot ramify
at infinity.  Hence precisely one prime ramifies
in $K'/M$.

Chevalley's ambiguous class number formula, as stated
by \citet[Theorem~1, author p.~1]{Lemmermeyer2013}, now reads
\[
 \bigl|\operatorname{Cl}(K')^{\operatorname{Gal}(K'/M)}\bigr|
 =\frac{h_M}{
 [\mathcal O_M^\times:
  \mathcal O_M^\times\cap N_{K'/M}(K')^\times]}.
\]
The denominator is a power of three, since the cube of every base
unit is a norm. If $3\nmid h_M$,
integrality forces it to have three-adic valuation zero,
so the fixed class group has no three-primary part.
A nonzero finite abelian three-group with an action of
$C_3$ has nonzero fixed elements: on its subgroup
killed by three, $(\sigma-1)^3=0$.  Therefore
$\operatorname{Cl}(K')[3]=0$.

Conversely, the norm from $\operatorname{Cl}(K')$ to
$\operatorname{Cl}(M)$ is surjective.  Indeed, the
cyclic cubic extension ramifies at seven, and hence
intersects the Hilbert class field of $M$ trivially.
Base change of that Hilbert class field and Artin
reciprocity identify the resulting surjection with
the ideal-class norm.  This proves
\[
 \operatorname{Cl}(K')[3]=0
 \quad\Longleftrightarrow\quad
 \operatorname{Cl}(M)[3]=0.
\]

We next compare $K$ and $K'$.  Since two is invertible
on finite three-groups, extension and norm identify
the three-primary class group of $K'$ with the
$c$-invariants of that of $K$.
If $\operatorname{Cl}(K')[3]=0$, then $c$ acts as
$-1$ on $V=\operatorname{Cl}(K)[3]$.
For $g$ in the normal subgroup $C_7$ of
$\operatorname{Gal}(K/\mathbb Q)$, the explicit Kummer
description gives $cgc^{-1}=g^{-1}$.  Since $c$
acts as a scalar on $V$, the actions of $g$ and
$g^{-1}$ on $V$ are equal.  Their common action has
order dividing both two and seven and is therefore
trivial.  Thus $V=V^{C_7}$.

Prime-to-three extension and norm identify these
$C_7$-invariants with $\operatorname{Cl}(K_0)[3]$,
which has already been proved zero.  Hence $V=0$.
The converse follows immediately from the quadratic
norm and extension maps.  Combining this with the preceding equivalence
proves the claimed equivalence.

For the field description, the degree-$16$ Galois
group of $\mathbb Q(E[3])$ acts transitively on its
eight nonzero three-torsion points.  A nonzero real
point is fixed by $c$, so its degree-eight point
field is exactly $\mathbb Q(E[3])^{\langle c\rangle}$.
The previous division-polynomial and subresultant
calculation shows that $z=2y+x$ generates this point
field and has polynomial $R$.
Adjoining the real root $\alpha$ gives degree
$8\cdot7=56$, because the two degrees are coprime.
The resulting field is contained in $M$ and has
its full degree, so they coincide.
There are exactly two nonzero real three-torsion
points and the ordinate coordinate is injective
on these eight points.  Thus the point field has
two real embeddings.  The pure degree-seven field
has one real embedding, giving signature $(2,27)$
for their compositum.

Finally $L/M$ is unramified at finite places.
At seven its two primes are interchanged by $c$.
At three $c$ acts nontrivially on the unramified
quadratic subfield $\mathbb Q_3(\sqrt{-7})$, and
so is outside inertia.  All other finite primes
are already unramified.  The discriminant tower
formula therefore gives $D_L=D_M^2$.
The signature supplies the negative sign, proving
$D_M=-3^{55}7^{62}$.
\end{proof}

\subsection{An unconditional three-class bound for the degree-fifty-six field}

We now remove the finite class-group hypothesis by exhibiting enough
principal prime ideals to exclude an unramified cyclic cubic extension.
The method is the explicit-formula class-number argument: principal
primes split in an unramified abelian extension, and their contributions
raise the lower bound for its discriminant.

\begin{proposition}\label{prop:septic-three-class}
For the degree-$56$ field $M$ of
Proposition~\ref{prop:lm-septic-degree56},
\[
 \operatorname{Cl}(M)[3]=0.
\]
\end{proposition}
\begin{proof}
Suppose otherwise. Class field theory gives an unramified cyclic
cubic extension $T/M$. Since the extension has odd degree, its real
places remain real. Hence
\[
 [T:\mathbb Q]=168,\qquad (r_1(T),r_2(T))=(6,81),\qquad
 \operatorname{rd}(T)=\operatorname{rd}(M)
   =3^{55/56}7^{62/56}<25.366.
 \tag{$*$}\label{eq:septic-hilbert-rd}
\]

Here is the finite algebraic input to the discriminant calculation.
The field can be constructed from $\alpha^7=3$ and
$w=3z$, whose monic polynomial is
\[
 w^8-441w^6+18522w^4-28588707=0.
\]
The exact maximal-order computation verifies its degree, signature
and discriminant. For $592$ distinct prime ideals of $M$ of norm
at most $5000$, the accompanying certificate gives an algebraic
integer generating each ideal. It verifies the equality of the
two ideal lattices in Hermite normal form and the exact norm of
the generator. There are $204$ distinct norms, with largest norm
$4973$. The replay verifies irreducibility of the degree-eight
and degree-$56$ polynomials, the two subfield embeddings, the
maximal order, the residue degrees and primality of each ideal,
and absence of repeated ideals. In particular the input includes one ideal of norm $3$,
two of norm $5$, and three of norm $25$.
The verification uses only arithmetic in the certified maximal
order and the explicit generators; it does not assume a computed
class group or the generalized Riemann hypothesis.
Write $c_q$ for the number of these verified ideals of norm $q$.
Each splits completely in $T$, so $T$ has at least $3c_q$ primes
of norm $q$.

For completeness we specify the analytic calculation, including
the error bounds. Put $g(t)=1+\cos(\pi t)$ for $|t|\le1$ and
$g(t)=0$ otherwise. Its normalized autocorrelation is the even
function $H$, supported on $[-2,2]$, whose expression for
$0\le t\le2$ is
\[
 H(t)=\frac{(2-t)(1+\tfrac12\cos\pi t)}3
       +\frac{\sin\pi t}{2\pi}.
\]
It has $H(0)=1$, $0\le H\le1$, and nonnegative Fourier
transform. Set $b=14$ and
\[
 F(x)=\frac{H(x/b)}{\cosh(x/2)}.
\]
This is an unconditional admissible kernel in the explicit formula \citep[Eqs.~(2.3)--(2.4), p.~122]{Odlyzko1990}.
Define
\[
 I=\int_0^\infty\frac{1-F(x)}{2\sinh(x/2)}\,dx,
 \qquad
 A=\int_0^{2b}\frac{H(x/b)}{2\cosh^2(x/2)}\,dx.
\]
The real-place term is $A$, since
$\int_0^\infty(2\cosh(x/2))^{-1}\,dx=\pi/2$.
Also $\int_0^2H(t)\,dt=2/3$, so the pole term is
$4\int_0^\infty F(x)\cosh(x/2)\,dx=8b/3$.
Discarding the nonnegative zero contributions and all unused
prime ideals in the explicit formula therefore gives
\[
\begin{split}
 \log\operatorname{rd}(T)\ge{}&
 \log(8\pi)+\gamma-I+\frac A{28}-\frac29\\
 &+\frac1{56}\sum_q c_q
   \sum_{\substack{m\ge1\\m\log q\le28}}
       \frac{4\log q\,H(m\log q/14)}{q^m+1}.
\end{split}
\tag{$**$}\label{eq:septic-explicit-lower}
\]
The factor $1/56$ is $3/168$: every verified prime of $M$
gives three primes in $T$. This calculation does not require
$M/\mathbb Q$ or $T/\mathbb Q$ to be Galois.

The finite sum in~\eqref{eq:septic-explicit-lower}, before division by $56$, has the certified
value
\[
 19.51514734319693587857082907544961
 \quad\text{with error less than }10^{-31}.
\]
The integrals are enclosed using $192$-bit ball arithmetic.
For $\varepsilon=1/16384$, the integral defining $I$ from
$\varepsilon$ to $28$ is evaluated by rigorous complex-ball
quadrature. Its tail is exactly $-\log\tanh7$.
The omitted integral from $0$ to $\varepsilon$ lies between
$0$ and $\varepsilon^2/2$: indeed $H'(0)=0$, $|H''|\le6$,
and $1-\operatorname{sech}(x/2)\le x^2/8$ imply that its
integrand is at most $x$. Using this upper endpoint bound for
$I$ gives
\[
 I<0.726948814142097,
 \qquad A>0.973804608359249.
\]
Substitution in~\eqref{eq:septic-explicit-lower}, with outward rounding throughout, yields
\[
 \operatorname{rd}(T)>25.4185>25.366,
\]
contradicting~\eqref{eq:septic-hilbert-rd}. Thus no such cubic extension exists, and
$\operatorname{Cl}(M)[3]=0$.
\end{proof}

A second evaluation with $b=27/2$ and $256$-bit arithmetic gives
the compatible lower bound $25.4115>25.366$, using the same exact
principal ideals. The proof establishes the required three-class
vanishing and does not need the exact class number of $M$.

By Proposition~\ref{prop:lm-septic-degree56}, the residual field
$K$ also has $\operatorname{Cl}(K)[3]=0$. Applying
Corollary~\ref{cor:lm-septic-class-gate} completes the septic
case, and hence all three families, of
Theorem~\ref{thm:unrestrictedthreefamiliesintro}.

\subsection{Compact localization at the base mixed exponents}
\label{subsec:compact-base-cases}
\begin{corollary}[The base quintic and septic arithmetic sources]
\label{cor:compact-base-source-five-seven}
For $a\in\{5,7\}$ and $S_a=\{3,a\}$, the map
\[
 H^1(G_{\Q,S_a},T_3J_a)\longrightarrow H^1(\Q_3,T_3J_a)
\]
is injective integrally. In particular the compact localization
$\Sel(T_3J_a)\to H^1(\Q_3,T_3J_a)$ is injective.
Together with the pure-power and mixed-family source
arguments and Theorem~\ref{thm:eleven-compact-selmer}, this proves the abelian localization assertions of
Theorem~\ref{thm:compact-localization-sections-intro}.
\end{corollary}
\begin{proof}
The quintic action and local-source theorem were proved first on
$U_5$ itself, before passing through the three-power covers. The
same is true at seven, where Proposition~\ref{prop:lm-whole-septic-action}
explicitly includes $r=0$. The verified class-group certificate
supplies its arithmetic-source surjectivity. Thus the cocycle-graph
argument applies to the compact coefficient modules $T_3J_5$ and
$T_3J_7$ directly. More explicitly, a cocycle on $G_{\Q,S_a}$
and its coefficient action define a pro-three graph over the
residual field. Both factor through the established arithmetic
source. Its image from $G_{\Q_3}$ is the entire source; after
subtracting the same integral local coboundary globally, the
cocycle is zero. This proves the full cohomology injection,
including at seven where the residual-field degree is divisible
by three and cannot be used for prime-to-three averaging.

A compact Selmer class is unramified outside $\{3,5\}$ or
$\{3,7\}$, respectively: at any other prime the abelian variety
has good reduction and the geometric Kummer class is unramified.
The integral cohomology localization injections from the source
theorems therefore kill any such class that is zero at three.
This argument uses neither the cubic local detector nor separation
of cusp differences by their three-primary order. Those additional
steps are used only in the affine classification, whose mixed
families still require $r\ge1$.
\end{proof}

\subsection{The full compact pro-three image}
\begin{proposition}[Nonabelian localization from a surjective arithmetic source]\label{prop:nonabelian-source-localization}
Let $S$ contain $3$, let $K/\mathbb Q$ be finite Galois and unramified
outside $S$, and put $\Gamma=\operatorname{Gal}(K_S(3)/\mathbb Q)$.
Suppose that the image of $G_{\mathbb Q_3}$ in $\Gamma$ is $\Gamma$.
Let $M$ be a finitely generated pro-$3$ group with continuous actual
$G_{\mathbb Q,S}$-action factoring through $\Gamma$, whose restriction
to $G_K$ has pro-$3$ image in $\operatorname{Aut}(M)$.
Then
\[
 H^1(G_{\mathbb Q,S},M)\longrightarrow H^1(\mathbb Q_3,M)
\]
is injective as a map of nonabelian cohomology sets.
\end{proposition}
\begin{proof}
For every continuous cocycle $c$, its graph on $G_{K,S}$ takes values
in the pro-$3$ group $M\rtimes\rho(G_{K,S})$; hence it factors through
$\operatorname{Gal}(K_S(3)/K)$. Thus $c$ itself factors through
$\Gamma$. If two cocycles are locally cohomologous, conjugate one by
the same element of $M$ globally to make their restrictions to
$G_{\mathbb Q_3}$ equal. Surjectivity onto $\Gamma$ makes them equal
everywhere.
\end{proof}

\begin{corollary}[The full compact image at the base quintic and septic]\label{cor:compact-full-base-five-seven}
Let $a\in\{5,7\}$. Every point-or-cusp Selmer section of the full
geometrically pro-$3$ arithmetic fundamental group of $U_a$ has
compact image globally conjugate to the proper point section of the
rational coordinate cusp determined by its localization at $3$.
More generally, for any stable quotient $R$ of the affine geometric
group retaining its abelianization, the same holds in the compact
quotient of $R$, obtained by killing all boundary inertia images.
\end{corollary}
\begin{proof}
Put $I=\overline{\langle\text{boundary inertia}\rangle}^{\mathrm{normal}}$
in $N_a$ and $K_R=\ker(N_a\to R)$. The compact quotient of $R$ is
$N_a/(K_RI)$, a quotient of $N_a/I$, not a presumed map from $R$
to the full proper group. The product $K_RI$ is closed because the
groups are compact.
Use the local-three conclusion to choose the rational coordinate cusp
$P$. Keep a fixed standard rational cusp $O$ as the reference section.
Let $c$ be the compact cocycle of the given section and $b$ that of
the proper point $P$, both relative to $O$. The compact geometric
group, and every quotient of it, inherit the established whole action
based at $O$ and surjective arithmetic source for
$S=\{3,a\}$. At $p\notin S$, smooth proper good reduction makes the
actual inertia action trivial and every proper local point section
unramified, relative to a rational proper basepoint extending over the
integers. The precise specialization input is
\citet[Expos\'e X, Theorem~3.8 and Corollary~3.9, corrected p.~217]{SGA1}.
Comparing the pointed arithmetic sequences gives the fibre product with
$G_{\Q_p}\to G_{\F_p}$, so the statement concerns the actual action. The compact image of a local cusp section is also such a
proper point section. Triviality on inertia is unchanged by geometric
conjugacy when inertia acts trivially. Therefore the entire compact
relative cocycle $c$ factors through $G_{\mathbb Q,S}$. The same
proper-point argument applies to $b$. Their restrictions at $3$ are
cohomologous, by the local cusp conclusion. Apply Proposition~\ref{prop:nonabelian-source-localization}. This proves the compact-section assertion of Theorem~\ref{thm:compact-localization-sections-intro}.
\end{proof}

\begin{remark}\label{rem:compact-cusp-frobenius-fibre}
At a good prime $p\ne3$, an affine local representative ramifying at a
coordinate cusp $Q$ has the proper Frobenius section of $Q$.
The corollary therefore restricts such $Q$ to the fibre of the proper
geometrically pro-$3$ point map containing $P$.
It does not assert that this fibre is a singleton. In particular,
coordinate-cusp differences are prime-to-$3$ torsion, so their
proper abelian pro-$3$ classes coincide. A finite-field separation
statement or a boundary-support argument remains necessary for the
affine conclusion. The partial separation proved in
Theorem~\ref{thm:finitefield-boundary-intro} leaves precisely the
same-component alternative at split fields.
\end{remark}

\subsection{The remaining finite-field cusp fibre}

Let \(a\in\{5,7\}\), let \(X_a\) be the proper Fermat curve, and let
\(D_a\) be its coordinate boundary. Suppose a full geometrically
pro-\(3\) point-or-cusp Selmer section has compact image equal to that
of a rational coordinate cusp \(P\). For \(p\notin\{3,a\}\), put
\[
 E_{a,P}(p)=
 \left\{Q\in D_a(\mathbf F_p):
 [s_Q]=[s_P]\text{ in }
 \operatorname{Sec}\bigl(\pi_1^{\mathrm{geom},(3)}
             (X_a/\mathbf F_p)\bigr)\right\}.
\]
Section classes here are taken modulo the geometric pro-\(3\) group;
the arithmetic quotient remains \(G_{\mathbf F_p}\).
Smooth proper specialization shows that every centre of nontrivial
local affine inertia belongs to \(E_{a,P}(p)\). Indeed, a local cusp,
or a local point reducing to a cusp \(Q\), has compact specialization
\(s_Q\); a point with interior reduction is unramified at \(p\).
For a section defined only in a geometric quotient \(R\), this
conclusion uses its compact quotient \(R^{\mathrm{cp}}\), and the
corresponding fibre can be larger. No full global lift is implicit.

General point injectivity for a single geometric pro-\(\ell\) quotient
over a finite field is false. Hoshi's Remark~10(i), which records an
example of Tamagawa, gives distinct points with the same full
geometrically pro-\(\ell\) section class
\cite[pp.~742--743]{FCHoshi2009}. After a finite constant extension,
the arithmetic group has the form
\[
 \Pi_X\simeq A\times \prod_{q\ne\ell}\mathbf Z_q,\qquad
 1\longrightarrow\Delta_X\longrightarrow A
       \longrightarrow\mathbf Z_\ell\longrightarrow1.
\]
For a further constant extension \(L/K\) of degree prime to \(\ell\),
new \(L\)-points conjugate under \(\operatorname{Gal}(L/K)\) have the
same section class: conjugation by a constant-field lift differs from
conjugation inside \(\Pi_{X_L}\) by a central factor. Since \(G_L\) is
abelian, conjugacy by \(\Pi_{X_L}\) is equivalent to geometric conjugacy
for sections. This example does not identify the points as Fermat
coordinate cusps, does not involve a prescribed rational cusp \(P\),
and does not compute \(E_{a,P}(p)\) over a prime field.

The almost-injectivity theorem of Saidi--Tamagawa requires a cofinite
set \(\Sigma\) of retained geometric primes
\cite[Proposition~1.8(vi)--(vii), Corollary~1.11(iii)]
 {FCSaidiTamagawa2009}.
It cannot be applied with \(\Sigma=\{3\}\). Its affine inertia
statements also do not restore point separation after boundary inertia
is killed in the proper quotient.

Even a proof that \(E_{a,P}(p)=\{P\}\) on a density-one set of primes
would not by itself prove finite ramification support for the
transverse boundary class. It would only make that class unramified on
the density-one set; the complement may be infinite. It would not make
the local class zero, since an interior reduction need not belong to
the \(P\)-cusp packet. The unrestricted cusp parameter itself is
allowed to have infinitely many nonzero valuations, so the needed
support assertion concerns the transverse component.

\begin{lemma}[Density-one Kummer localization]\label{lem:density-one-kummer-localization}
Let \(F\) be a number field. Localization on any density-one set of its
finite places, in Dirichlet density, is injective on \(H^1(F,\mathbf Z_3(1))\). The analogous
statement over \(\mathbf Q\) holds for finite direct sums of induced
Tate modules
\(\operatorname{Ind}_{G_F}^{G_{\mathbf Q}}\mathbf Z_3(1)\).
\end{lemma}
\begin{proof}
For \(C=\mu_{3^m}\) and \(E=F(C)\), the cyclic group
\(G=\operatorname{Gal}(E/F)\subset(\mathbf Z/3^m)^\times\) satisfies
\(H^1(G,C)=0\). If a generator \(u\) is nontrivial modulo \(3\), then
\(u-1\) is invertible. Otherwise, writing
\(v_3(u-1)=s<m\), the order is \(3^{m-s}\), the norm has valuation
\(m-s\), and its kernel is \(3^sC=(u-1)C\). The trivial group is
immediate. Restriction to \(E\) is therefore injective.

A nonzero finite-level class consequently defines a nontrivial
Kummer extension over \(E\). A nonidentity element of its Kummer
subgroup, viewed in the Galois splitting field over \(F\), is detected
by Frobenius at a positive-density set of primes splitting in \(E\).
At such a prime the coefficient action is trivial and the local
class is nonzero. This contradicts density-one local vanishing.
The inverse systems of finite zeroth cohomology groups satisfy the Mittag--Leffler condition, so passage through all \(m\) proves the first assertion. Shapiro's lemma
proves the second: the sum of $(N\mathfrak p)^{-s}$ over primes of residue degree at least
two remains bounded as $s\to1^+$, while the excluded degree-one prime
sum is at most $[F:\mathbf Q]$ times the excluded rational prime sum.
Thus lifting a rational Dirichlet-density-one set gives a
Dirichlet-density-one set of places of $F$.
\end{proof}

Thus density-one \emph{local vanishing} of the transverse boundary
class would suffice, whereas density-one separation of coordinate
cusps supplies only an inertia restriction. The next section proves
singleton boundary fibres when \(p\not\equiv1\pmod a\) and separates
different coordinate components at every good prime. The
same-component alternative at split primes, the required global
vanishing of the transverse class, and the affine base-exponent
classification remain separate questions.

\section{Proper pro-three boundary fibres of Fermat curves}
\label{sec:finitefield3-boundary-reconstruction}

\subsection{Exact scope and results}
\label{subsec:ff3-scope}

We determine the possible identifications among rational coordinate
boundary points in the proper geometric pro-three group.
Let \(n\ge3\), let \(k=\mathbb F_Q\) have characteristic prime to \(3n\),
and let
\[
 C_n:\ X^n+Y^n+Z^n=0.
\]
This smooth proper curve has genus \((n-1)(n-2)/2\ge1\). Put
\[
 D_X=C_n\cap\{X=0\},\quad D_Y=C_n\cap\{Y=0\},\quad
 D_Z=C_n\cap\{Z=0\},\qquad D=D_X\sqcup D_Y\sqcup D_Z.
\]
Fix a geometric base point and write
\[
 N=\pi_1(C_{n,\bar k})^{\mathrm{pro}\text{-}3},\qquad
 \Pi=\pi_1(C_n)/\ker\bigl(\pi_1(C_{n,\bar k})\longrightarrow N\bigr).
\]
Thus
\begin{equation}
 1\longrightarrow N\longrightarrow\Pi\longrightarrow G_k
 \longrightarrow1,\qquad G_k\simeq\widehat{\mathbb Z}.
 \label{eq:ff3-1-1}
\end{equation}
The arithmetic quotient is full \(G_k\). Let \(\kappa(P)\) be the
proper point-section class of \(P\in C_n(k)\), modulo conjugacy by \(N\).
All fibres below are intersected with \(D(k)\); no assertion excludes
an ordinary proper point from the full fibre of a boundary point.

\begin{theorem}[Separation of coordinate components]
\label{thm:ff3-different-components}
If \(P,Q\in D(k)\) lie on different coordinate components, then
\begin{equation}
 \kappa(P)\ne\kappa(Q).
 \label{eq:ff3-1-2}
\end{equation}
This holds for both odd and even \(n\).
\end{theorem}

\begin{theorem}[Uniform boundary blocks and coefficient-prime separation]
\label{thm:ff3-dichotomy}
Suppose \(D(k)\ne\varnothing\), and put \(d=\gcd(n,Q-1)\).
Each component has \(d\) rational points. There is one divisor \(h\mid d\)
such that the fibres on each component are the cosets of the unique
subgroup \(\mu_h(k)\subset\mu_d(k)\) under its regular diagonal action.
There are exactly \(3d/h\) boundary section classes.
Moreover,
\[
 3\nmid h,\qquad 3\mid n\ \Longrightarrow\ h=1.
\]
In particular, the boundary map is injective for every exponent
divisible by three, at every allowed finite field.

For a prime exponent \(a\ge5\), if \(Q\not\equiv1\pmod a\), the
boundary consists of three points with distinct classes. If
\(Q\equiv1\pmod a\), either all \(3a\) classes are distinct or the
three whole coordinate components are exactly the fibres.
\end{theorem}

For odd \(n\), the point \((0:-1:1)\) shows that \(D(k)\ne\varnothing\).
For even \(n\), nonemptiness is equivalent to
\(-1\in(k^\times)^n\), or to \((Q-1)/d\) being even.
When this condition fails, all three rational boundary components
are empty.

\begin{corollary}[Specified primes with singleton boundary fibres]
\label{cor:ff3-density}
Let \(n\ge3\) be odd. If \(3\mid n\), the boundary map is injective
at every rational prime \(q\nmid3n\). If \(3\nmid n\), the primes
\(q\nmid3n\) satisfying \(\gcd(n,q-1)=1\) give singleton boundary
fibres. This specified sufficient set has natural density
\begin{equation}
 \prod_{p\mid n}\frac{p-2}{p-1}.
 \label{eq:ff3-1-3}
\end{equation}
In particular the specified nonsplit sets for \(n=5\) and \(n=7\)
have densities \(3/4\) and \(5/6\), respectively.
\end{corollary}

The proof first uses Pál's order-three action theorem to separate
different components. A tangent-line calculation then proves the
stronger injectivity when three divides the exponent. The remaining
same-component alternative for prime exponents different from three
is not decided; in particular no density-one separation or affine
global classification at the base exponents five and seven follows.

\subsection{Geometric centralizers over a finite field}
\label{subsec:ff3-centralizers}

We need uniqueness, not just existence, of invariant geometric paths. The exact published input is \citet[Lemma~1.7, p.~144]{FCSaidiTamagawa2009}. Its standing assumptions on p.~140 allow any retained prime set containing a prime different from the residue characteristic, and allow a proper hyperbolic curve. Thus \(\Sigma=\{3\}\) is allowed when the genus is at least two. The direct proper-curve proof below also covers genus one. The additional cofinite-\(\Sigma\) assumptions occurring later in their paper are not present in this lemma.

\begin{lemma}[Trivial geometric centralizers]
\label{lem:ff3-centralizer}

Let \(C/k\) be a smooth proper geometrically connected curve of positive genus over a finite field of characteristic different from \(\ell\). Put \[
 N_C=\pi_1(C_{\bar k})^{\mathrm{pro}\text{-}\ell},
\] and retain full \(G_k\) in the corresponding arithmetic quotient. For any arithmetic section \(s\), \begin{equation}
Z_{N_C}(s(G_k))=1.
\label{eq:ff3-2-1}
\end{equation}

\end{lemma}
\begin{proof}

In genus at least two, apply Saïdi--Tamagawa Lemma 1.7 to \(Z=s(G_k)\). Projection is injective on \(Z\) and its image is all of \(G_k\), hence open. The lemma injects the commensurator of \(Z\), and therefore its centralizer, into \(G_k\). Its intersection with the geometric kernel is consequently trivial.

For every positive genus, the following proper, singleton-prime argument proves the assertion directly. Use \(s\) to identify the arithmetic quotient with \(N_C\rtimes G_k\), and suppose a nonidentity \(z\in N_C\) is fixed by the \(G_k\)-action. Choose an open characteristic subgroup \(M\subset N_C\) with \(z\notin M\), and put \[
 H=\overline{\langle M,z\rangle}.
\] This is an open \(G_k\)-stable subgroup. The image of \(z\) in \(H/M\) is nonzero, and \(H/M\) is cyclic; hence the image of \(z\) in \(H^{\mathrm{ab}}\) is nonzero.

The open subgroup \(H\rtimes G_k\) corresponds to a finite étale geometrically connected proper cover \(Y\to C\) over \(k\). The geometric pro-\(\ell\) topology induced on its geometric subgroup is the full topology: an open subgroup of a pro-\(\ell\) group has this property. At the original fundamental-group level, one proves it by first passing to the normal core of the covering subgroup, then taking the intersection of the finitely many conjugates of the kernel of a finite \(\ell\)-quotient. The resulting quotient is an extension of finite \(\ell\)-groups. Thus \[
 H=\pi_1(Y_{\bar k})^{\mathrm{pro}\text{-}\ell},
 \qquad H^{\mathrm{ab}}=T_\ell\operatorname{Jac}(Y).
\] Since \(Y\) is proper, \(\operatorname{Jac}(Y)(k)\) is finite. Its rational \(\ell\)-primary torsion group is finite and therefore admits no nonzero compatible Tate-module system, so \[
 (T_\ell\operatorname{Jac}(Y))^{G_k}=0.
\] But the nonzero image of \(z\) is fixed by \(G_k\), a contradiction. This proves \eqref{eq:ff3-2-1}.

\end{proof}

This proper proof avoids any need to infer integral fixed-point vanishing from a rational vector-space statement without checking torsion. The cited original lemma proves the more general assertion by the same open-subgroup device and Frobenius weights.

\subsection{Arithmetic equality as an invariant-path statement}
\label{subsec:ff3-paths}

Take \(b\in C(k)\), and use its point section to fix the actual \(G_k\)-action on \[
 N_C=\pi_1(C_{\bar k},b)^{\mathrm{pro}\text{-}\ell}.
\] For geometric points \(x,y\), let \(\mathcal P_\ell(x,y)\) be the torsor of pro-\(\ell\) étale paths between their fibre functors.

If \(x,y\) are rational, this torsor has its natural \(G_k\)-action. Choosing a path identifies its descent cocycle with the difference between the two point sections. Therefore \begin{equation}
\kappa(x)=\kappa(y)
 \quad\Longleftrightarrow\quad
 \mathcal P_\ell(x,y)^{G_k}\ne\varnothing.
\label{eq:ff3-3-1}
\end{equation} This is the usual torsor description of nonabelian \(H^1\), with no change of the arithmetic quotient. Two invariant paths differ by a \(G_k\)-fixed element of the appropriate geometric fundamental group. Lemma~\ref{lem:ff3-centralizer} makes that element the identity. Thus, in our setting, an invariant path is unique whenever it exists.

Now let \(A\) be a finite group of \(k\)-automorphisms of \(C\). Define a geometric action extension \begin{equation}
1\longrightarrow N_C\longrightarrow E_A
 \overset{\mathrm{pr}}\longrightarrow A\longrightarrow1.
\label{eq:ff3-3-2}
\end{equation} One concrete construction uses pairs \((a,p)\), with \(a\in A\) and \(p\in\mathcal P_\ell(b,a(b))\). Multiplication concatenates \(p\) with the image under \(a\) of the second path. This construction is independent, up to the usual inner equivalence, of the base-point choices.

The quotient \(A\) in \eqref{eq:ff3-3-2} is retained in its entirety. Only the geometric kernel has been made pro-\(\ell\). In particular, for \(A=S_3\) and \(\ell=3\), \(E_A\) is \textbf{not} being defined as the maximal pro-three quotient of an orbifold group; doing that would lose its quotient of order two.

Because the automorphisms are defined over \(k\) and \(b\) is rational, \(G_k\) acts on \(E_A\), preserves its group law, and acts trivially on its quotient \(A\).

\begin{lemma}[Coherent lifting of an invariant point class]
\label{lem:ff3-invariant-action}

Suppose \(C/k\) satisfies Lemma~\ref{lem:ff3-centralizer}. If \(\kappa(b)\) is fixed by the action of every element of \(A\), then \eqref{eq:ff3-3-2} has a section \[
 A\longrightarrow E_A
\] whose image is pointwise fixed by \(G_k\).

\end{lemma}
\begin{proof}

For each \(a\in A\), invariance says \(\kappa(a(b))=\kappa(b)\). By \eqref{eq:ff3-3-1}, the fibre of \(E_A\to A\) above \(a\) contains a \(G_k\)-fixed element. Its uniqueness follows from Lemma~\ref{lem:ff3-centralizer}. Therefore \[
 E_A^{G_k}\longrightarrow A
\] is bijective. It is a homomorphism because \(G_k\) acts on \(E_A\) by automorphisms. Its inverse is the required section.

\end{proof}

The uniqueness step is essential. Choosing unrelated paths for each automorphism would not by itself prove a finite-group splitting.

\subsection{Pál’s theorem and its admissible application}
\label{subsec:ff3-pal}

Let \(C/\mathbb C\) be a smooth projective connected curve of positive genus, and let a cyclic group \(B\) of prime order \(\ell\) act nontrivially. \citet[Theorem~5.2, author pp.~12--13]{FCPal2014} states that the fixed points of \(B\) correspond bijectively to conjugacy classes of sections of the pro-\(\Sigma\) completed action extension, when \(\ell\in\Sigma\). For \(\Sigma=\{\ell\}\) this is the full pro-\(\ell\) action extension.

Pál's definition retains the finite quotient only when every prime dividing its order belongs to \(\Sigma\). Our only use of this theorem is \begin{equation}
B=C_3,\qquad \Sigma=\{3\}.
\label{eq:ff3-4-1}
\end{equation} We do not apply it to \(S_3\) or to a group of order \(5\) or \(7\).

For \(B=C_3\), the path extension \(E_B\) of Section~\ref{subsec:ff3-paths} is the same extension as in Pál's theorem. Its kernel is pro-three and its quotient has order three, so the total group is pro-three. Equivalently, Pál's exact completion sequence (Definition 2.3) identifies the pro-three completion of the topological action extension with this geometric path extension.

There is an elementary naturality that we will need. If \(\tau\) normalizes \(B\), it acts on section classes of \(E_B\) by conjugating the finite quotient and transporting paths. The fixed-point-to-section bijection is equivariant: \begin{equation}
\tau\cdot[s_x]=[s_{\tau x}].
\label{eq:ff3-4-2}
\end{equation} Indeed, in Pál's Definition 2.1, the subgroup associated with \(x\) is the subgroup of lifts fixing a chosen point above \(x\) in the universal cover. Applying a normalizing automorphism sends it to the subgroup fixing the transported point over \(\tau x\). Changing the chosen lift only conjugates by the geometric group.

\paragraph{Specialization to finite characteristic}

The Fermat model and its coordinate permutations are defined over \(\mathbb Z[1/n]\). Adjoin a primitive cube root \(\omega\), and use a strict henselian discrete valuation ring above the residue characteristic, with \(3n\) invertible. The Fermat curve is smooth and proper over this ring. When \(3\nmid n\), the two points \begin{equation}
f_+=(1:\omega:\omega^2),\qquad
 f_-=(1:\omega^2:\omega)
\label{eq:ff3-4-3}
\end{equation} are sections over it.

Smooth proper specialization identifies the geometric pro-three fundamental groups and their path torsors, compatibly with morphisms and composition. The exact original references are \citet[Exposé~X, Section~2, pp.~206--207; Theorem~3.8 and Corollary~3.9, p.~217 (original marginal pp.~283--284)]{SGA1}. The generic geometric fibre can be compared with the complex curve by characteristic-zero base change and Riemann existence \citep[Exposé~XII, Theorem~5.1]{SGA1}.

The needed naturality follows from the specialization construction on finite étale covers and their fibre functors: pullback by each extending coordinate automorphism commutes with restriction to the geometric generic and special fibres, and evaluation at an extending point commutes with those restriction functors. Consequently the fixed-point splitting classes and the normalizer action in \eqref{eq:ff3-4-2} specialize together.

Only a \textbf{geometric} comparison is asserted. The arithmetic Galois groups in characteristic zero and finite characteristic are not identified. Nor is a specialization theorem for a possibly wild \(S_3\)-quotient stack required.

This also covers residue characteristic \(2\). The geometric coefficient prime and the order of \(C_3\) are still invertible; the transposition is used only as an extending automorphism, not as a tame group quotient of order two.

\subsection{The coordinate permutation extension does not split}
\label{subsec:ff3-s3}

We apply the finite-action theorem to the cyclic subgroup of coordinate
permutations. On \(C_n\), put
\[
 \rho(X:Y:Z)=(Y:Z:X),\qquad \tau(X:Y:Z)=(X:Z:Y).
\]
They act faithfully and generate \(S_3\), with
\(\rho^3=\tau^2=1\) and \(\tau\rho\tau=\rho^{-1}\).
A projective fixed point of \(\rho\) must be one of its three eigenlines,
\[
 (1:1:1),\quad (1:\omega:\omega^2),\quad(1:\omega^2:\omega).
\]
The first has Fermat value \(3\ne0\). If \(3\mid n\), the other two
also have Fermat value three, so the cyclic fixed locus is empty.
If \(3\nmid n\), precisely those two lie on the curve, and \(\tau\)
interchanges them.

\begin{lemma}[The coordinate action extension does not split]
\label{lem:ff3-s3-nonsplit}
The relative geometric action extension
\begin{equation}
 1\longrightarrow N\longrightarrow E_{S_3}
 \longrightarrow S_3\longrightarrow1
 \label{eq:ff3-5-1}
\end{equation}
has no splitting.
\end{lemma}
\begin{proof}
Suppose \(t:S_3\to E_{S_3}\) were a splitting, and restrict it to
\(B=\langle\rho\rangle\), obtaining \(u:B\to E_B\).
If \(3\mid n\), Pál's theorem and the specialization just established
identify its section classes with an empty fixed locus, a contradiction.

If \(3\nmid n\), there are exactly two possible classes, those of
the two cyclic fixed points. The transposition interchanges them.
On the other hand,
\[
 t(\tau)\,u(\tau^{-1}g\tau)\,t(\tau)^{-1}=u(g)\qquad(g\in B).
\]
Changing the chosen lift of \(\tau\) changes this action only by
geometric conjugacy. Thus the class of \(u\) would be fixed by the
transposition, again a contradiction.
\end{proof}

Only the cyclic order-three subgroup is an input to Pál's theorem.
The finite quotient \(S_3\) in \eqref{eq:ff3-5-1} has not been
replaced by a pro-three completion.

\subsection{Separation over both odd and even exponents}
\label{subsec:ff3-canonical}

Twisted transpositions compare a pair of components without assuming
that the exponent is odd. After permuting coordinates, write
\[
 P=(0:u:1),\qquad Q=(v:0:1),\qquad u^n=v^n=-1,
\]
with \(u,v\in k^\times\). Let
\[
 A=\operatorname{diag}(1,u^2,1),\qquad
 B=\operatorname{diag}(v^2,1,1),\qquad \sigma=(XZ).
\]
These diagonal maps preserve \(C_n\), since \(u^{2n}=v^{2n}=1\).
They satisfy
\[
 A(Q)=Q,\quad A\tau(P)=(0:u^2:u)=P,\qquad
 B(P)=P,\quad B\sigma(Q)=(v^2:0:v)=Q
\]
as projective points.
If \(\kappa(P)=\kappa(Q)=c\), both \(A\) and \(A\tau\) preserve \(c\),
so \(\tau\) preserves \(c\). Likewise \(\sigma\) preserves \(c\).
The two transpositions generate \(S_3\).
Lemma~\ref{lem:ff3-invariant-action} would split
\eqref{eq:ff3-5-1}, contradicting
Lemma~\ref{lem:ff3-s3-nonsplit}. This proves
Theorem~\ref{thm:ff3-different-components}. For odd \(n\), in particular,
\begin{equation}
 \kappa(0:-1:1),\quad\kappa(-1:0:1),\quad\kappa(1:-1:0)
 \quad\hbox{are pairwise distinct}.
 \label{eq:ff3-6-1}
\end{equation}
The automorphisms used are defined over \(k\). The cyclic fixed points
used to obstruct a geometric splitting may be considered over \(\bar k\);
no cube root of unity in \(k\) is required.

\subsection{Uniform cyclic blocks}
\label{subsec:ff3-components}

The diagonal symmetries determine the form of every possible fibre.
If \(u^n=-1\) has one solution in \(k\), all its solutions are
\(u\zeta\) for \(\zeta\in\mu_n(k)=\mu_d(k)\), where \(d=\gcd(n,Q-1)\).
Thus each coordinate component is a regular orbit of this cyclic group.
The relation
\[
 P\sim Q\quad\Longleftrightarrow\quad\kappa(P)=\kappa(Q)
\]
is invariant under every \(k\)-automorphism. The elements of the cyclic
group carrying a fixed point to an equivalent point form a subgroup:
if \(tP\sim P\) and \(uP\sim P\), invariance gives \(tu^{-1}P\sim P\).
The fibres are its cosets. A cyclic group has a unique subgroup of
each order, so their common size is a divisor \(h\mid d\).

Coordinate permutations identify these relations on the three
components, possibly inverting the cyclic coordinate. They preserve
the subgroup order, so the same \(h\) works on all three.
Theorem~\ref{thm:ff3-different-components} excludes any identification
between components. There are exactly \(3d/h\) classes.

If \(n=a\ge5\) is prime, then \(d=1\) or \(a\).
In the first case the three boundary points are distinct; in the
second, primality leaves \(h=1\) or \(a\), giving the asserted
singleton or whole-component dichotomy. We prove the stronger
conclusion \(h=1\) when \(3\mid n\) next.

\subsection{Boundary divisor classes and the coefficient prime}
\label{subsec:ff3-abelian}

An explicit pole calculation distinguishes the exponents divisible
by three from the remaining ones.

\begin{lemma}[Exact order of a coordinate-point difference]
\label{lem:ff3-exact-cusp-order}
Over an algebraically closed field of characteristic not dividing
\(n\ge3\), the divisor class \(P-Q\) of any two distinct coordinate
points on \(C_n\) has exact order \(n\) in its Jacobian.
\end{lemma}
\begin{proof}
At a coordinate point the tangent line intersects the curve in
\(n\) times that point. Therefore the ratio of the tangent lines at
\(P,Q\) has divisor \(n(P-Q)\).

It remains to exclude a proper divisor of \(n\) as the order.
After permuting coordinates, write \(P=(0:u:1)\), with \(u^n=-1\).
The tangent is \(L=Y-uZ\), and its zero divisor on the curve is \(nP\).
On \(C_n\setminus\{P\}\), use affine coordinates
\[
 x=X/L,\qquad z=Z/L.
\]
Their coordinate ring is
\[
 \bar k[x,z]\big/\bigl(x^n+(1+uz)^n+z^n\bigr).
\]
The polynomial in \(z\) has degree \(n-1\), since \(u^n=-1\)
and its leading coefficient is \(n u^{n-1}\ne0\).
At \(P\), \(X/Z\) is a uniformizer, \(L\) vanishes to order \(n\),
and \(Z\) is a unit. Hence \(x\) has pole order \(n-1\), and \(z\)
has pole order \(n\).

Every function regular away from \(P\) has a unique expression
\[
 \sum_{i=0}^{n-1} a_i(z)x^i,\qquad a_i(z)\in\bar k[z],
\]
because the equation is monic of degree \(n\) in \(x\).
The monomials \(z^j x^i\), \(j\ge0\), \(0\le i<n\), have pole orders
\(jn+i(n-1)\). These orders are pairwise distinct: reduction modulo
\(n\) determines \(i\), and then determines \(j\).
Thus leading poles cannot cancel, and every nonconstant function
whose only possible pole is \(P\) has pole order at least \(n-1\).

If the order of \(P-Q\) were a proper divisor \(m\) of \(n\),
a function with divisor \(m(Q-P)\) would have its only pole at
\(P\), of order \(m\le n/2<n-1\). This is impossible.
\end{proof}

Suppose now \(P,Q\in D(k)\) are distinct and \(3\mid n\).
If their proper pro-three point sections were equal, their images
in the geometric abelianization \(T_3\operatorname{Jac}(C_n)\)
would have the same Kummer class. The finite Kummer exact sequences
would then give
\[
 [P-Q]\in\bigcap_{j\ge1}3^j\operatorname{Jac}(C_n)(k).
\]
The group \(\operatorname{Jac}(C_n)(k)\) is finite; that intersection
is its prime-to-three subgroup. This contradicts the exact order
\(n\) in Lemma~\ref{lem:ff3-exact-cusp-order}.
Consequently every boundary fibre is a singleton when \(3\mid n\).
If \(3\nmid n\), then \(3\nmid d\), so the block size is automatically
prime to three. This proves all of Theorem~\ref{thm:ff3-dichotomy}.

When \(3\nmid n\), the same divisor calculation shows that every
boundary difference has zero three-primary Kummer image. The
original Jacobian therefore does not separate these boundary points.
The component-separation theorem in that range uses the full
nonabelian proper group. Distinguishing affine cusp inertia alone
would not prove separation after that inertia is filled.

For the density assertion, if \(3\mid n\), injectivity holds at every
good prime just proved. Otherwise, for odd \(n\), the condition
\(\gcd(n,q-1)=1\) gives precisely the three rational coordinate
points. For each distinct prime \(p\mid n\), it excludes the identity
class among the \(p-1\) invertible residues modulo \(p\).
Chinese remaindering and the cyclotomic form of Chebotarev give
\eqref{eq:ff3-1-3}; this is natural density, hence also Dirichlet density
\citep[Definition~8.30, Theorem~8.31, and Example~8.34,
pp.~147--148]{Milne2020ANT}.
This proves Corollary~\ref{cor:ff3-density} and
Theorem~\ref{thm:finitefield-boundary-intro}.

Each proved inequality of full point-section classes is visible in
some finite characteristic pro-three geometric quotient. Otherwise
the closed sets of conjugators giving equality in finite quotients
would have the finite-intersection property, and compactness would
give equality in the full group. No uniform small quotient is asserted.

\subsection{A sufficient automorphism criterion for collisions}
\label{subsec:ff3-automorphisms}

The following criterion is useful for testing potential constructions. It is a sufficient condition only.

\begin{lemma}[The Frattini automorphism kernel]
\label{lem:ff3-frattini}

For a finitely generated pro-\(\ell\) group \(M\), \[
 \ker\bigl(\operatorname{Aut}(M)\to
            \operatorname{Aut}(M/\Phi(M))\bigr)
\] is pro-\(\ell\), where \(\Phi(M)=\overline{M^\ell[M,M]}\). The corresponding kernel in the outer automorphism group is also pro-\(\ell\).

\end{lemma}
\begin{proof}

For a finite \(\ell\)-group \(P\), choose an ordered basis of \(P/\Phi(P)\), of dimension \(d\). It has \(|\Phi(P)|^d\) ordered lifts to \(P\). Every lift generates \(P\), by Burnside's basis theorem. An automorphism trivial on \(P/\Phi(P)\) acts on the set of lifts, and the stabilizer of any lift is trivial: an automorphism fixing a generating tuple is the identity. The order of the automorphism kernel therefore divides the \(\ell\)-power \(|\Phi(P)|^d\), so it is an \(\ell\)-group.

Pass to a cofinal family of characteristic finite quotients of \(M\). The inverse-limit kernel is pro-\(\ell\). Inner automorphisms act trivially on the Frattini quotient, so the outer kernel is a quotient of this pro-\(\ell\) kernel.

\end{proof}

This is a statement about the kernel. The full automorphism group is not asserted to be pro-\(\ell\).

\begin{lemma}[A residual-trivial automorphism criterion]
\label{lem:ff3-trivial-residual}

Let \(C/k\) be a geometrically connected proper curve with \(b\in C(k)\) and finitely generated geometric pro-\(\ell\) group \(M\). Let \(\alpha\in\operatorname{Aut}_k(C)\) have finite order \(r\) prime to \(\ell\). If \(\alpha\) acts trivially on \(M/\Phi(M)\), then \begin{equation}
\kappa(\alpha b)=\kappa(b).
\label{eq:ff3-9-1}
\end{equation}

\end{lemma}
\begin{proof}

The outer action of \(\alpha\) has order dividing \(r\) and lies in the pro-\(\ell\) kernel of Lemma~\ref{lem:ff3-frattini}. It is therefore trivial.

Form the geometric action extension \(1\to M\to E_{\langle\alpha\rangle}\to C_r\to1\). Since every element of \(C_r\) acts on \(M\) by an inner automorphism, the centralizer \[
 L=Z_{E_{\langle\alpha\rangle}}(M)
\] surjects onto \(C_r\), with kernel \(Z(M)\). This is a central extension of \(C_r\) by the abelian pro-\(\ell\) group \(Z(M)\).

Choose a lift \(e\in L\) of a generator of \(C_r\). Then \(e^r\in Z(M)\). Multiplication by \(r\) is an automorphism of an abelian pro-\(\ell\) group. There is a unique \(z\in Z(M)\) with \(z^r=e^{-r}\), so \(t=ze\) has order exactly \(r\) and maps to that generator. The subgroup \(\langle t\rangle\) is the unique complement to \(Z(M)\) in \(L\): any other lift of the same generator differs by an element of \(Z(M)\), and its \(r\)-th power forces that element to be the identity.

The \(G_k\)-action preserves \(L\), its kernel, and this unique complement. It acts trivially on the quotient \(C_r\) because \(\alpha\) is defined over \(k\). Hence it fixes the complement pointwise. In particular it fixes a path over \(\alpha\), and \eqref{eq:ff3-3-1} gives \eqref{eq:ff3-9-1}.

\end{proof}

No centre-freeness of \(M\) is required in this proof. The central kernel is retained and the complement is constructed explicitly.

\paragraph{Failure of this criterion for the Fermat diagonal rotations}

For the rest of this section specialize \(n=a\in\{5,7\}\), in
characteristic different from \(3,a\). Thus \(N\) is the proper
geometric pro-three group of \(C_a\). These are the
base exponents whose affine boundary problem remains open.


The diagonal group \[
 A_a=(\mu_a)^2\subset\operatorname{Aut}(C_{a,\bar k})
\] acts faithfully on \(H_1(C_{a,\bar k},\mathbb F_3)=N/\Phi(N)\).

Indeed, every nonidentity element \(\gamma\in A_a\) has order \(a\). The quotient map \(C_a\to C_a/\langle\gamma\rangle\) is tame, since the residue characteristic is not \(a\). Riemann--Hurwitz gives \[
 g(C_a/\langle\gamma\rangle)<g(C_a).
\] The invariant subspace of rational three-adic \(H_1\) has dimension twice the quotient genus, so \(\gamma\) acts nontrivially on \(H_1(C_a,\mathbb Q_3)\).

If its action modulo three were trivial, its integral matrix would lie in \(\ker(\operatorname{GL}_{2g}(\mathbb Z_3)\to
       \operatorname{GL}_{2g}(\mathbb F_3))\), a pro-three group. A nontrivial element of order \(a\) cannot lie there. Therefore its residual action is nontrivial. Every nonidentity element of \(A_a\) has order \(a\), proving faithfulness.

The criterion of Lemma~\ref{lem:ff3-trivial-residual} therefore does not construct the same-component boundary collision sought here. Failure of a sufficient criterion is not proof that such a collision cannot occur.

\subsection{Hoshi’s general counterexample mechanism and the boundary field}
\label{subsec:ff3-hoshi}

\citet[Remark~10(i), pp.~742--743]{FCHoshi2009} explains a general source of collisions in proper geometrically pro-\(\ell\) point maps over finite fields. One first makes the outer Galois action pro-\(\ell\) by a finite extension of constants; one then uses a further extension of degree prime to \(\ell\) that supplies new points. Conjugate points in that extension yield coincident section classes in the stated setting.

Hoshi writes conjugacy by the arithmetic group. For sections of the full procyclic \(G_k\), this is equivalent to geometric conjugacy: a conjugator with arithmetic image \(g\) can be multiplied on the right by \(s(g)^{-1}\), which centralizes \(s(G_k)\), to obtain a conjugator in the geometric kernel.

This is a genuine warning against unrestricted singleton-prime point injectivity. It does not specify the Fermat boundary or the fibre above a rational coordinate point.

For the present boundary, the mechanism has a concrete obstruction.

\begin{proposition}[The boundary field in the constant-extension construction]
\label{prop:ff3-hoshi-boundary}

Let \(k'=\mathbb F_{Q^t}\) be a finite constant-field extension over which the outer \(G_{k'}\)-action on \(N\) has pro-three image. Then \[
 \operatorname{ord}_a(Q^t)
\] is a power of three. Any further extension of \(k'\) of degree prime to three adds no points to the coordinate boundary \(D(k')\).

\end{proposition}
\begin{proof}

The residual action of \(G_{k'}\) on \(N/\Phi(N)\) is a finite three-group. Thus its Frobenius matrix has three-power order.

The faithful diagonal subgroup \(A_a\) from Section~\ref{subsec:ff3-automorphisms} is normalized by Frobenius. On \(A_a\), conjugation is the \(Q^t\)-power action on its root-of-unity coordinates (up to replacing Frobenius by its inverse, which does not change the order). Faithfulness on the residual module therefore identifies this conjugation with conjugation by the residual Frobenius matrix. Its order must be a power of three. The order of scalar \(Q^t\) on \((\mathbb Z/a\mathbb Z)^2\) is exactly \(\operatorname{ord}_a(Q^t)\), proving the first assertion.

The field of all coordinate points over \(k'\) is \(k'(\mu_a)\), of degree \(\operatorname{ord}_a(Q^t)\). For \(a=5\), this degree divides four, hence must already be one. For \(a=7\), it divides six, hence is one or three. More generally its three-power part cannot be removed by an extension of degree prime to three: \[
 \operatorname{ord}_a(Q^{tu})
 =\frac{\operatorname{ord}_a(Q^t)}
        {\gcd(\operatorname{ord}_a(Q^t),u)}
 =\operatorname{ord}_a(Q^t)
 \qquad(3\nmid u).
\] Since \(a\) is prime, a coordinate component has either its single canonical rational point or all \(a\) points. Thus the further extension supplies no new coordinate boundary points.

\end{proof}

This excludes this particular use of Hoshi's two-stage construction for producing new Fermat boundary points. It does not decide the same-component alternative over a field where all boundary points are already rational, nor does it exclude other collision constructions.

\subsection{A depth-three limitation at quintic primes}
\label{subsec:ff3-quintic-depth-three}

The full pro-three boundary theorem cannot in general be detected
in the first three ordinary lower-central grades.  The following
result gives a specified positive-density family and an explicit prime.

\begin{proposition}[Constant quintic point map through class three]
\label{prop:ff3-quintic-depth-three}
Let $q\ne3,5$ be a rational prime with $q\equiv1\pmod5$, let
\[
 C/\mathbb F_q:\ X^5+Y^5+Z^5=0,\qquad
 C_A:\ v^2+v=u^5,\qquad A=\operatorname{Jac}(C_A),
\]
and suppose arithmetic Frobenius on $A[3]$ has order twenty.
Use $b=(0:-1:1)$ as base point and put
\[
 N=\pi_1(C_{\overline{\mathbb F}_q},b)^{\mathrm{pro}\text{-}3}.
\]
With closed ordinary lower-central terms, continuous nonabelian
cohomology satisfies
\[
 H^1(G_{\mathbb F_q},N/\gamma_4N)=\{1\}.
\]
Thus every section of this arithmetic quotient is geometrically
conjugate to the point section of $b$.  In particular, the proper
point-section map is constant in this quotient and in every finite
Galois-stable quotient of $N$ of nilpotency class at most three.

The set of rational primes satisfying these conditions has natural
density $1/10$, and it contains $q=11$.
\end{proposition}

\begin{proof}
The exact residual calculation for the quintic CM Jacobian gives
\[
 E=\mathbb Q(\zeta_5)\subset K=\mathbb Q(A[3]),\qquad
 [K:\mathbb Q]=80,\qquad
 \operatorname{Gal}(K/E)=C_{20}\subset\mathbb F_{81}^{\times};
\]
see Lemma~\ref{lem:lm-quintic-residual-group}.
Here $A[3]$ has rank one over
$\mathbb Z[\zeta_5]/3=\mathbb F_{81}$.
The three degree-five quotient maps
\[
 \begin{split}
 f_Z&=(-XY/Z^2,X^5/Z^5),\\
 f_X&=(-YZ/X^2,Y^5/X^5),\\
 f_Y&=(-ZX/Y^2,Z^5/Y^5)
 \end{split}
\]
give the integral isomorphism
\[
 T_3\operatorname{Jac}(X_5)\simeq(T_3A)^3.
\]
Indeed their joint pushforward $\Phi$ and summed pullback $\Psi$
satisfy $\Phi\Psi=[5]$ and $\Psi\Phi=[5]$, by the three disjoint
diagonal-character summands.  Hence $5^{-1}\Psi$ is an integral
three-adic inverse.  Good reduction at $q$ transports this identity
to the finite-field Frobenius representation.

Because $q$ splits in $E$, its residual Frobenius is multiplication
by an element $t\in\mathbb F_{81}^{\times}$ of order twenty.
The eigenvalues over $\overline{\mathbb F}_3$ on $A[3]$ are
\[
 t,\quad t^3,\quad t^9,\quad t^{27}=t^7.
\]
Each occurs three times on
$\overline V=N^{\mathrm{ab}}/3$.
The action is semisimple since twenty is prime to three.
None of the exponents $1,3,7,9$ is zero modulo twenty;
every sum of two lies between two and eighteen;
every sum of three is odd.
Consequently $1$ is not an eigenvalue of Frobenius on
\[
 \overline V^{\otimes m},\qquad 1\le m\le3.
\]
Repeated tensor factors are included.  Replacing $t$ by another
generator, or Frobenius by its inverse, does not change this assertion.

Put $L_m=\gamma_mN/\gamma_{m+1}N$.
Iterated commutators give a continuous equivariant surjection
\[
 (N^{\mathrm{ab}})^{\otimes m}\longrightarrow L_m;
\]
the grades are finitely generated compact $\mathbb Z_3$-modules.
The commutator description and its compatibility with the Galois
action are also described in
\citet[Introduction, author pp.~1--2]{DavisPriesWickelgren2024}.
It follows that $F-1$ is invertible modulo three on $L_m$ for
$m\le3$. Nakayama's lemma makes it surjective integrally, and a
surjective endomorphism of a finitely generated $\mathbb Z_3$-module
is injective. Thus
\[
 H^0(G_{\mathbb F_q},L_m)=H^1(G_{\mathbb F_q},L_m)=0
 \qquad(1\le m\le3).
\]
For $H^1$, explicitly choose a coboundary with the prescribed
cocycle value at the topological Frobenius generator; equality on
the full procyclic group follows by continuity.  The arithmetic
quotient remains full $G_{\mathbb F_q}$.

Induct through the central extensions
\[
 1\longrightarrow L_m\longrightarrow N/\gamma_{m+1}N
 \longrightarrow N/\gamma_mN\longrightarrow1.
\]
The abelian case follows from the displayed vanishing.  In a
successive stage, conjugate a cocycle so that its image in the
preceding quotient is trivial, lifting one conjugating element
through the surjection.  The remaining cocycle takes values in
$L_m$ and is a coboundary.  This proves the asserted singleton
cohomology set through $m=3$.  Equivalently, an invariant path lifts
at each central step because its fibre is an $L_m$-torsor with
zero $H^1$; the zero $H^0$ also makes these low-depth paths unique.
No section in the full geometric group is assumed to exist above
an arbitrary quotient section.

The residual elements defining the prime set are the eight generators
of the normal subgroup $C_{20}$ in the group of order eighty.
This is a conjugacy-stable subset, so Chebotarev gives natural
density $8/80=1/10$.

For an explicit prime, exact counting on $v^2+v=u^5$ gives
\[
 \#C_A(\mathbb F_{11})=13,\qquad
 \#C_A(\mathbb F_{121})=103.
\]
One can use $\mathbb F_{121}=\mathbb F_{11}[s]/(s^2-2)$,
where $2$ is a nonsquare, and count the square roots of
$1+4u^5$.  The Frobenius polynomial is therefore
\[
 T^4+T^3-9T^2+11T+121.
\]
Modulo three it is $T^4+T^3+2T+1$.
Its derivative is $(T-1)^3$ and its value at one is two, so it is
square-free.  Its companion matrix has order twenty by exact
matrix multiplication over $\mathbb F_3$.  Thus $q=11$ meets
the hypothesis.  These finite calculations are recorded in the
reproducibility certificate.
\end{proof}

At $q=11$, the same isogeny gives $\#C(\mathbb F_{11})=15$,
so all proper rational points are the fifteen coordinate points.
Their equality in $N/\gamma_4N$ is compatible with
Theorem~\ref{thm:ff3-different-components}: the latter concerns
the full geometric pro-three group and separates different
coordinate components. The proposition makes no assertion of a
full collision within one component. Already in tensor degree four,
$1+1+9+9\equiv0\pmod{20}$, so the eigenvalue exclusion stops.
An actual integral higher-path class remains to be determined;
this limitation of low-depth quotients does not settle the affine
base-five or base-seven Selmer-section problem.

\subsection{Two obstructions to diagonal-symmetry shortcuts}
\label{subsec:ff3-diagonal-shortcuts}

The next two observations explain why coprime complement conjugacy
does not settle the remaining full same-component boundary question.
They concern the entire Fermat geometric group, not just one
hyperelliptic quotient.

\begin{proposition}[A nontrivial diagonal twist cannot remove the residual obstruction]
\label{prop:ff3-diagonal-twist-obstruction}
Let $a\in\{5,7\}$, let $k=\mathbb F_Q$ have characteristic prime
to $3a$, and assume $\mu_a\subset k$.  On
$C_a:X^a+Y^a+Z^a=0$, let
\[
 \delta(X:Y:Z)=(\zeta^rX:\zeta^sY:Z),\qquad
 r,s\in\mathbb Z/a\mathbb Z,
\]
where $\zeta$ is primitive.  Let $N$ be the geometric pro-three
fundamental group and let $F$ denote arithmetic Frobenius.
If the action of $\delta^{-1}F$ on $N/\Phi(N)$ has three-power
order, then $r=s=0$.
In particular, for a nonidentity $\delta$, the closed subgroup of
$\operatorname{Out}(N)$ generated by its twisted Frobenius cannot
be pro-three.
\end{proposition}

\begin{proof}
Here $\delta$ and $F$ are initially outer automorphisms of $N$;
their actions on $N/\Phi(N)$ are canonical. They commute as outer
automorphisms because $\delta$ is defined over $k$. Thus the
displayed twisted action is defined without choosing a
$\delta$-fixed base point.

Put $J_a=\operatorname{Jac}(v^2+v=u^a)$, and write $\beta$ for
the automorphism induced by $(u,v)\mapsto(\zeta u,v)$.
The three degree-$a$ morphisms to this same curve are
\[
 \begin{split}
 f_Z&=(-XY/Z^2,X^a/Z^a),\\
 f_X&=(-YZ/X^2,Y^a/X^a),\\
 f_Y&=(-ZX/Y^2,Z^a/Y^a).
 \end{split}
\]
They satisfy $u^a=v(v+1)$ and extend to the proper smooth models.
Each is the quotient by its order-$a$ diagonal deck subgroup.
The induced pushforward on three-torsion is surjective: its
composition with pullback is multiplication by $a$, which is
invertible on three-torsion.

The three $u$-coordinate weights of $\delta$ are
\[
 w_Z=r+s,\qquad w_X=s-2r,\qquad w_Y=r-2s
 \pmod a.
\]
Consequently these quotient maps intertwine the action of
$\delta$ with $\beta^{w_i}$ and the action of $F$ with the same
Frobenius operator $B$ on $J_a[3]$.

The quotient of $v^2+v=u^a$ by $\beta$ is the projective $v$-line.
The $\beta$-invariants in its rational first homology therefore
vanish. Since $a$ is invertible in $\mathbb Z_3$, invariants are
an integral direct summand and commute with reduction modulo three.
Moreover $\dim_{\mathbb F_3}J_a[3]=a-1$, and $\Phi_a(T)$ is
irreducible over $\mathbb F_3$: the order of $3$ modulo $a$ is
$a-1$ for both $a=5$ and $a=7$. Thus $J_a[3]$ has rank one over
\[
 \mathbb F_3[\beta]\simeq\mathbb F_{3^{a-1}}.
\]
Frobenius commutes with $\beta$ because $\zeta\in k$, so $B$ is
a nonzero scalar in this field.

The three quotient actions of $\delta^{-1}F$ are therefore the
scalars $B\beta^{-w_i}$. Under the hypothesis each has
three-power order. Its order also divides $3^{a-1}-1$, which is
prime to three, so each is the identity. Hence
\[
 \beta^{w_Z}=\beta^{w_X}=\beta^{w_Y}.
\]
As $\beta$ has exact order $a$, this gives $3r=3s=0$ modulo $a$,
and therefore $r=s=0$. The final assertion follows because the
finite residual image of a pro-three group is a three-group.
\end{proof}

This excludes a nontrivial diagonal twist as the first step in the
pro-three-outer-action constant-extension construction. It does not
exclude other constructions of collisions, or decide the fibres
over a field where all boundary points are already rational.

\begin{proposition}[The geometric diagonal centralizer has a nonzero quotient]
\label{prop:ff3-diagonal-centralizer}
Let $a\in\{5,7\}$ and let $C_a$ be the proper Fermat curve over
an algebraically closed field of characteristic prime to $3a$.
Let $N$ be its geometric pro-three fundamental group and
$H=(\mu_a)^2$ its diagonal automorphism group. Form the relative
geometric action extension
\[
 1\longrightarrow N\longrightarrow E_H\longrightarrow H
 \longrightarrow1,
\]
retaining the full finite quotient $H$. Choose any complement
$S\subset E_H$, mapping isomorphically onto $H$, and let it act
on $N$ by conjugation. Such a complement exists by coprime
Schur--Zassenhaus. Put $L_2=\gamma_2N/\gamma_3N$ and
$g=(a-1)(a-2)/2$. Then
\[
 Z_N(S)\subset\gamma_2N,\qquad
 \frac{Z_N(S)}{Z_N(S)\cap\gamma_3N}
       \simeq L_2^H\simeq\mathbb Z_3^{g-1}.
\]
In particular this quotient has rank five for $a=5$ and rank
fourteen for $a=7$.
\end{proposition}

\begin{proof}
The choice of complement supplies an actual action on $N$.
Its action on each ordinary lower-central grade is the canonical
geometric $H$-action: changing lifts by inner automorphisms has
no effect on a grade. No arithmetic Galois invariance of $S$ is
assumed or asserted.

Write $V=N^{\mathrm{ab}}$. After an unramified coefficient
extension of $\mathbb Z_3$ containing $\mu_a$, the characters of
$H$ on $V$ are
\[
 (i,j)\in(\mathbb Z/a\mathbb Z)^2,
 \qquad i\ne0,\quad j\ne0,\quad i+j\ne0,
\]
each with multiplicity one. In characteristic zero, the
differentials $x^{i-1}y^{j-a}\,dx$, with $i,j\ge1$ and
$i+j\le a-1$, have characters $(i,j)$; adjoining their conjugates
gives precisely the displayed set. Smooth proper
prime-to-characteristic specialization gives the same geometric
representation in the stated characteristics. The degree-two
surface description used below is also given in
\citet[author pp.~2--4]{DavisPriesWickelgren2024}.
In particular $V^H=0$, so every element of $Z_N(S)$ lies in
$\gamma_2N$.

The degree-two surface presentation gives
\[
 L_2\simeq\bigwedge^2V/\mathbb Z_3\omega,
\]
where $\omega$ is the primitive symplectic relation. A trivial
diagonal character in the exterior square pairs $(i,j)$ with
$(-i,-j)$. There are exactly $g$ such pairs. The relation is
invariant and primitive, with a unit coefficient in a symplectic
basis. Since $|H|=a^2$ is invertible in $\mathbb Z_3$, the
averaging projector makes invariants exact. Therefore taking
invariants removes exactly one primitive rank-one summand,
and $L_2^H$ is finite free of rank $g-1$. The unramified
coefficient extension is faithfully flat, so this rank and
freeness assertion descends to $\mathbb Z_3$.

It remains to check that these graded fixed elements actually
lift to the geometric centralizer. For
$\overline x\in L_2^H$, choose $x\in\gamma_2N$.
The elements
\[
 c_h=x^{-1}h(x)\in\gamma_3N\qquad(h\in H)
\]
form a continuous nonabelian cocycle for the action furnished by
$S$. Its class is trivial because $H$ is finite of order prime
to three and $\gamma_3N$ is pro-three. Explicitly, finite
Schur--Zassenhaus conjugacy kills it in every finite
$H$-stable quotient; the corresponding nonempty closed sets of
conjugating elements have the finite-intersection property, so
compactness kills it in $\gamma_3N$ itself. Hence there is
$u\in\gamma_3N$ with $c_h=u\,h(u)^{-1}$ for every $h$.
Then $xu$ is fixed by $H$ and lifts $\overline x$.
Thus $Z_N(S)\to L_2^H$ is surjective, with kernel
$Z_N(S)\cap\gamma_3N$, proving the stated quotient.
\end{proof}

The absence of diagonal invariants in the original Jacobian does
not make this geometric centralizer trivial. Coprime conjugacy
of complements therefore does not provide the Frobenius
compatibility required for a rational invariant path. Neither
proposition proves a full same-component collision or full
separation, and neither supplies the missing affine
base-five or base-seven Selmer-section classification.


\subsection{Exact remaining gate}
\label{subsec:ff3-gate}

For \(Q\equiv1\pmod a\), choose a primitive \(\zeta\in k\), put \[
 b=(0:1:-1),\qquad
 \delta(X:Y:Z)=(X:Y:\zeta Z).
\] The remaining question is exactly whether \begin{equation}
\mathcal P_3(b,\delta b)^{G_k}
\label{eq:ff3-11-1}
\end{equation} is empty. By \eqref{eq:ff3-3-1}, nonemptiness is equivalent to \(\kappa(b)=\kappa(\delta b)\). Since \(a\) is prime, that one equality implies the collapse of the whole coordinate component, and then of each component by permutation symmetry. Theorem~\ref{thm:ff3-different-components} still separates the three resulting classes.

One can express the same gate as a nonabelian Lang equation. Choose a path \(p\) and use the right \(N\)-torsor convention \(g(p)=p\,c_g\). Then \(c_g\) is a continuous cocycle. For arithmetic Frobenius \(F\), a path \(pn\) is fixed if and only if \begin{equation}
c_F F(n)=n,\qquad\text{equivalently}\qquad
 c_F=nF(n)^{-1}.
\label{eq:ff3-11-2}
\end{equation} The value at a topological generator determines invariance under \(G_k\), by continuity. The abelianization of this equation is always soluble for the boundary difference, as Section~\ref{subsec:ff3-abelian} explains. Full integral pro-three solvability is the unresolved issue.

No specific solution or obstruction to \eqref{eq:ff3-11-1} at split rational primes has been obtained in this proof. Consequently the following claims remain unproved here:

\begin{itemize}
\item
  that another boundary point lies over a chosen rational cusp for even one specified split rational prime;
\item
  that collisions occur at a set of positive Dirichlet density;
\item
  that singleton boundary fibres occur at a density-one set;
\item
  that the affine pro-three base-five or base-seven Selmer section classification follows from this partial result.
\end{itemize}

The cofinite-prime almost-injectivity theorems do not fill this gap. Saïdi--Tamagawa Proposition 1.8(vi)--(vii) and Corollary 1.11(iii), printed pp.~144--148, impose a cofinite retained prime set. The set \(\{3\}\) is not cofinite. Their earlier Lemma 1.7, used in Section~\ref{subsec:ff3-centralizers}, is valid in the singleton setting and supplies only the centralizer input stated there.

\section{Auxiliary local calculations}
\subsection{A limitation of the auxiliary level-twenty-seven Kummer condition}
\label{subsec:auxiliary-level27-kummer}

One possible auxiliary test uses a rational three-torsion family
and the rank-zero curve \(X_0(27)\). The following explicit
calculation shows that its integral three-adic Kummer condition
fails to distinguish the boundary at infinity from either rational
modular cusp class, on a positive-density set of good primes.
The conclusion concerns this auxiliary condition; it supplies no
global Selmer counterexample and does not decide the full compact
Fermat point-map fibre discussed above.

Write
\[
 C=X_0(27):\quad y^2+y=x^3-7,\qquad
 h=-3+\frac{y+5}{x},\qquad
 \lambda=\left(\frac{h}{h+9}\right)^3.
\]
The model and the map \(h:C\to X_0(9)\) are given in
\cite[proof of Proposition 24, p.~18 of the 2020 author version]{BourdonGillRouseWatson2024}.
Their displayed \(j\)-map, followed by direct substitution, gives
\[
 j=\frac{(h+9)^3(h^3+243h^2+2187h+6561)^3}
          {h^9(h^2+9h+27)}
   =\frac{27(9-8\lambda)^3}{\lambda^3(1-\lambda)}.
\]
Thus this is the relevant level-\(27\) cover for the family
\[
 \mathcal E_\lambda:\quad
 Y^2+XY+\frac{\lambda}{27}Y=X^3.
\]
This rational \(3\)-torsion family is recalled in
\cite[\S2.3, (2.3.3) and Lemma 2.3.5, p.~7]{FreiHonigsVoight2025}.
Its marked point \((0,0)\) fixes a vector modulo three, while the
geometric determinant is one. Thus the residual geometric image
is contained in the unipotent group of order three. The higher
congruence kernel is pro-\(3\), so the entire geometric Tate image
is pro-\(3\). This makes the family available on the geometric
pro-\(3\) quotient, but does not make its auxiliary Kummer condition
separate the boundary points.
Let \(O\) be the origin of \(C\), and put \(T=(3,4)\).
These are rational modular cusps, with \(\lambda(O)=1\) and
\(\lambda(T)=0\). The third rational point \((3,-5)\) is
noncuspidal: it has \(\lambda=-1/8\) and \(j=-12288000\).
The distinction between the two rational modular cusps and the
rational noncuspidal self-point is recorded by
\citet[\S6.2, p.~17 of the 17 March 2008 author version]{WuthrichSelfPoints2009}.
Put \(D_\infty=\lambda^{-1}(\infty)\).
For \(q\ne3\), write
\[
 \kappa_q:C(\Q_q)\longrightarrow H^1(\Q_q,T_3C)
\]
for the integral \(3\)-adic Kummer map based at \(O\).

\begin{lemma}[Failure of auxiliary Kummer separation]
\label{lem:auxiliary-level27-kummer-failure}
For each \(A\in\{O,T\}\), there is a positive-density set of rational
primes \(q\ne3\) for which some \(R_q\in D_\infty(\Q_q)\) satisfies
\[
 \kappa_q(R_q)=\kappa_q(A).
\]
The equality also holds on a \(q\)-adic neighborhood of \(R_q\) in
\(C(\Q_q)\). Explicit witnessing primes are \(q=17\) for \(A=O\)
and \(q=71\) for \(A=T\).
\end{lemma}

\begin{proof}
The discriminant of \(C\) is \(-3^9\), so \(C\) has good reduction at
every \(q\ne3\). Its reduction kernel is a pro-\(q\) group on which
multiplication by \(3\) is an automorphism. Consequently
\[
 C(\Q_q)^{\wedge3}\simeq C(\F_q)^{\wedge3},
\]
and equality of the \(3\)-primary reduction classes is equivalent to
equality under \(\kappa_q\).

Substitution of \(y=x(h+3)-5\) in the equation of \(C\) gives
\[
 F(x,h)=x^3-(h+3)^2x^2+(9h+27)x-27=0.
\]
The fibre \(D_\infty\) is \(h=-9\); hence its points have
\begin{equation}\label{eq:auxiliary-level27-fibre}
 f(x)=x^3-36x^2-54x-27=0,\qquad y=-6x-5.
\end{equation}
Here \(\operatorname{disc}(f)=-3^{13}\).
Modulo \(17\), equation~\eqref{eq:auxiliary-level27-fibre} has the point
\[
 \overline R=(12,8).
\]
It is nonzero \(2\)-torsion, since \(2\cdot8+1=0\) in \(\F_{17}\).
Also \(f'(12)=7\pmod {17}\). Hensel's lemma therefore supplies a
point \(R_{17}\in D_\infty(\Q_{17})\) reducing to \(\overline R\).
Its \(3\)-primary reduction class is zero, proving the claim for
\(A=O\) at \(17\).

Modulo \(71\), the point \(\overline R=(69,7)\) belongs to~\eqref{eq:auxiliary-level27-fibre}, and
\(f'(69)=31\pmod {71}\). With \(\overline D=\overline R-\overline T\),
the group law gives
\[
 \overline D=(36,69),\quad
 2\overline D=(39,55),\quad
 4\overline D=(38,35),\quad
 8\overline D=O.
\]
Thus the Hensel lift \(R_{71}\in D_\infty(\Q_{71})\) has the same
\(3\)-adic Kummer class as \(T\). Constancy on the corresponding
residue neighborhoods follows from the reduction comparison.

We give the finite-extension argument for positive density explicitly.
Fix one of these pairs \((q_0,A)\), and choose the algebraic point
\(R\) of~\eqref{eq:auxiliary-level27-fibre} with the specified embedding into \(\Q_{q_0}\).
Set \(K=\Q(R)\). Since \(f\) is monic with discriminant a power of
\(3\), \(K\) is unramified away from \(3\).
The point \(R-A\) is divisible by \(9\) in \(C(\Q_{q_0})\):
its reduction has order dividing \(8\), and multiplication by \(9\)
is an automorphism on the reduction kernel. Choose an algebraic
point \(B\), embedded in \(\Q_{q_0}\), such that
\begin{equation}\label{eq:auxiliary-level27-division}
 [9]B=R-A.
\end{equation}
Let \(L/\Q\) be a finite Galois extension containing
\(K\), the coordinates of \(B\), \(C[9]\), and \(\mu_{27}\).
It may be chosen unramified away from \(3\). Indeed, \(C\) is an
abelian scheme away from \(3\), every \(K\)-point extends there by
properness, and multiplication by \(9\) is finite \'etale there.
Adjoining the finite \'etale fibre above \(R-A\), taking its Galois
closure, and adjoining \(\mu_{27}\) introduce no ramification
away from \(3\).

Choose a prime of \(L\) above \(q_0\), and let \(\sigma\) be its
arithmetic Frobenius. Since the chosen \(R\) and \(B\) are
\(\Q_{q_0}\)-rational, \(\sigma\) fixes both. Its action on
\(\mu_{27}\) is the \(17\)-th power: both \(17\) and \(71\) are
\(17\pmod {27}\).
Chebotarev's theorem
\cite[Chapter I, \S2.2, pp.~I-7--I-8]{Serre1989}
gives a positive-density set of primes whose Frobenius class in
\(\operatorname{Gal}(L/\Q)\) is the class of \(\sigma\).
For each such \(q\), choose a prime above \(q\) with Frobenius
exactly \(\sigma\). The extension is unramified, so its decomposition
group is generated by \(\sigma\). Therefore \(R\) and \(B\) give
points in \(C(\Q_q)\), and~\eqref{eq:auxiliary-level27-division} holds there. Moreover
\begin{equation}\label{eq:auxiliary-level27-congruence}
 q\equiv17\pmod {27}.
\end{equation}

For \(q\equiv2\pmod3\), cubing is a bijection on \(\F_q\).
For each \(y\in\F_q\), the equation \(x^3=y^2+y+7\) has exactly one
solution, whence \(\#C(\F_q)=q+1\).
By~\eqref{eq:auxiliary-level27-congruence}, the \(3\)-primary subgroup has order exactly \(9\).
Equation~\eqref{eq:auxiliary-level27-division} consequently makes the \(3\)-primary reduction class
of \(R-A\) zero. The good-reduction comparison proves the required
Kummer equality for every prime in this positive-density set.
\end{proof}

\begin{remark}[Local Fermat points and the remaining scope]
\label{rem:auxiliary-level27-local-fermat}
These collisions persist on actual local points of the level cover of
\(U_a\), for \(a=5\) or \(7\).
At any of the preceding primes \(q\ne a\), take \(w=q^{-m}\) with
\(m\ge1\), and set
\[
 \lambda=w^{3a},\qquad
 h=\frac{9w^a}{1-w^a}.
\]
As \(m\) increases, \(h\) approaches \(-9\).
Since \(f'\) is a unit at the chosen residue root, the implicit
function theorem gives a point of \(C(\Q_q)\) near \(R_q\) with
this value of \(h\). Its Kummer class remains that of \(A\).
At the same time,
\[
 x_F=w^3,\qquad
 y_F=-w^3(1-w^{-3a})^{1/a}
\]
are defined over \(\Q_q\) by Hensel's lemma and satisfy
\(x_F^a+y_F^a=1\), hence give the points
\([x_F:y_F:-1]\) of the all-plus model of \(U_a\).
They reduce to the coordinate cusp \((1:-1:0)\), where
\(\lambda=\infty\).
The modular point supplies the cyclic level-\(27\) subgroup for
\(\mathcal E_\lambda\); away from the exceptional \(j\)-values,
quadratic twisting does not change the field of definition of such
a subgroup.

Thus the auxiliary \(X_0(27)\) Kummer condition does not separate
this boundary from either rational cusp class, even after deleting
a finite set of primes. This is a calculation about that condition
only. It does not construct a global Selmer section, and it does
not identify the full compact pro-$3$ section of \(X_a\) at this
boundary with the prescribed rational cusp. The latter finite-field
fibre condition remains an additional, unresolved requirement.
\end{remark}


\subsection{Further obstructions from proper quotients}

\begin{proposition}[No good-outside-three abelian quotient]
\label{prop:lm-no-good-outside-three-quotient}
If $3\nmid n$, the Jacobian of $X_n$ has no positive-dimensional
abelian quotient over $\mathbb Q$ with good reduction outside
$3$.  In particular, there is no nonconstant morphism over
$\mathbb Q$ from $X_n$ to a genus-one curve whose Jacobian has
good reduction outside $3$.
\end{proposition}
\begin{proof}
The Fermat curve and its Jacobian have good reduction at $3$.
Good reduction passes to abelian quotients by the
N\'eron--Ogg--Shafarevich criterion.  A quotient as in the
statement would therefore have good reduction everywhere
over $\mathbb Q$, contradicting Fontaine's theorem that no
positive-dimensional abelian variety over $\mathbb Q$ has
everywhere good reduction.  A nonconstant map to a
genus-one curve induces a surjection onto its Jacobian.
\end{proof}

For this input see \citet[Corollary, p.~517]{Fontaine1985}.

\subsection{A local obstruction to lifting rational points of the three-punctured line}
\label{subsec:lm-local-metabelian-separation}

\begin{proposition}[Local point/cusp separation]
\label{prop:lm-local-metabelian-separation}
Let $U_1=\mathbb P^1\setminus\{0,1,\infty\}$ over
$\mathbb Q_3$.  No ordinary point section of $U_1(\mathbb Q_3)$
is cuspidal in the geometrically pro-$3$ metabelian quotient.
The three rational cusp packets in that quotient are pairwise
disjoint.
\end{proposition}
\begin{proof}
Fix $a\in\mathbb Q_3\setminus\{0,1\}$ and consider
\[
 V_a:\quad a x^9+(1-a)y^9=1,\qquad xy\ne0,
 \qquad t=a x^9.
\]
This is a finite \'etale cover of $U_1$, geometrically Galois
with deck group $(\mathbb Z/9\mathbb Z)^2$.  It has the point
$Q=(1,1)$ above $a$, and its compactification $\overline V_a$
is geometrically a Fermat curve of degree $9$.

Put $N=\pi_1(U_{1,\overline{\mathbb Q}_3})^{\mathrm{pro}\text{-}3}$.
The geometric group $H$ of $V_a$ embeds as an open subgroup of
$N$.  Since $N/H$ is abelian, the closed derived groups satisfy
\[
 N'\subset H,\qquad N''\subset H'.
\]
Let $\mathcal H$ be the arithmetic subgroup corresponding to
$V_a$, using $Q$ to lift the section of $a$.

Suppose that the section of $a$ is cuspidal modulo $N''$.
After lifting its conjugator to $N$, its full point section
and a conjugate of a full cusp section differ elementwise by
$N''$.  Because $N''\subset\mathcal H$, the conjugated cusp
section is itself contained in $\mathcal H$.  The intersection
of its cusp decomposition group with $\mathcal H$ is a cusp
decomposition group of the cover and maps onto
$G_{\mathbb Q_3}$.  Hence its boundary center
$P\in\overline V_a\setminus V_a$ is rational over
$\mathbb Q_3$.

The two sections agree after passing from $H$ to $H^{\mathrm{ab}}$,
since $N''\subset H'$.  After compactification their Kummer
classes in $H^1(\mathbb Q_3,T_3\operatorname{Jac}(\overline V_a))$
agree.  The local completion kernel calculation in the proof of Lemma~\ref{lem:fourzero} shows that
$[Q-P]$ is torsion of order prime to $3$.  \citet[Theorem~A, p.~191]{Coleman1986} applies to the geometric degree-nine Fermat curve,
since $9=(19-1)/2$ and $2\le8$.  It forces a point
torsion-equivalent to a coordinate cusp to be a coordinate
cusp.  This contradicts the fact that $Q=(1,1)$ is interior.

For the second assertion, use the two Kummer coordinates
$t$ and $1-t$ to identify the geometric abelianization with
$\mathbb Z_3(1)^2$.  The three cusp packets there are the
lines with directions $(1,0)$, $(0,1)$ and $(-1,-1)$.
Here constants $-1$ have trivial pro-$3$ Kummer class.
Two different packets intersect only in the zero class.
If two full cusp sections agreed modulo $N''$, both cusp
parameters would therefore be zero.  They lift to the
standard cover $x^3+y^3=1$.  As above, equality modulo
$N''$ persists in the geometric abelianization of this
cover, because its geometric deck group is abelian.
Their compact images are two different rational coordinate
cusps of the smooth cubic.  Their difference is a nonzero
point of order $3$ on its elliptic Jacobian.  Such a point
is not in the prime-to-$3$ torsion kernel of local completion,
a contradiction.
\end{proof}

\begin{corollary}
\label{cor:lm-no-rational-point-lift}
Suppose that every geometrically pro-$3$ Selmer section of
$U_n/\mathbb Q$ is locally cuspidal at $3$, as in
Theorem~\ref{thm:localthreeintro}.  Its pushdown under the
coordinate map $U_n\to U_1$ cannot be the section of an
ordinary rational point of $U_1$, even if this equality is
requested only in the geometrically pro-$3$ metabelian quotient.
In particular, the rational point $t=2$ gives no such Selmer lift.
\end{corollary}
\begin{proof}
Localizing a putative equality at $3$ would make the ordinary
point section cuspidal, contrary to
Proposition~\ref{prop:lm-local-metabelian-separation}.
\end{proof}


\section{Comparison with Chabauty--Kim and single-prime quotients}

\subsection{What Theorem E supplies}
\citet[Theorem E and \S5.7]{BKL2026} assumes a proper curve of genus at
least two with a rational point and a Jacobian quotient of dimension
at least two, rank zero, and finite $\Sha$. It proves containment of
the Chabauty--Kim loci in one finite closed subscheme at every prime
of good reduction. Our $A$ is a dimension-two quotient of
$\Jac(X_5)$, so its hypotheses hold. Their proof uses the zero
$V_pA$-Selmer group and Manin--Mumford; their subsequent finite-descent
argument recovers the proper Selmer section conclusion already
covered by Stoll's theorem.

There are three distinct statements here. The explicit isogeny is a
geometric calculation. Its rank-zero and $\Sha$ consequence needs the
separate arithmetic theorem and nonvanishing. The assertion about
affine prosolvable or metabelian cusp packets additionally needs
Theorem~\ref{thm:transfer}; it is not a literal assertion of Theorem E.
Indeed one rank-zero factor already suffices for Stoll, so finiteness
for the entire quintic Jacobian is stronger arithmetic information
than is needed for the cusp conclusion. The finite subscheme in
Theorem E is not asserted to consist solely of rational cusps.
The genus-five factor $A_{11}$ also satisfies Theorem E, while the
whole Fermat Jacobian at exponent eleven has rank six.

\subsection{Why geometrically pro-\texorpdfstring{$3$}{3} is a different problem}
For a proper abelian variety over a nonarchimedean completion $k_v$,
the full Tate-module Kummer
map is injective on its compact group of local points. A single
$\ell$-adic Kummer map generally kills prime-to-$\ell$ torsion and,
when $v\nmid\ell$, an open pro-residue-characteristic subgroup.
Thus neither compact-center identification nor the full finite
abelian descent input above survives unchanged after retaining only
$\Delta^{(3)}$. Theorem~\ref{thm:unrestrictedthreefamiliesintro} therefore needs its separate
single-prime argument; it does not follow by projecting
Theorem~\ref{thm:transfer}.

Likewise, constructing noncuspidal global sections of a geometric
pro-$3$ quotient does not disprove its Selmer assertion. One must check
the stated local condition at every place, including $3$.
\citet[Theorem A; Theorem 3.5 and Corollary 3.6]{Hoshi2010} constructs
nongeometric pro-$p$ sections for suitable covers of the triply
punctured line over finite extensions of cyclotomic fields. His
Fermat examples and the lifting of compact Jacobian classes concern
abstract sections; they do not establish an everywhere locally
geometric counterexample for $U_3/\Q$. A family failing local
geometry at the prime retained by the quotient cannot supply such
a counterexample.

\citet[Theorem~A, pp.~2--3; Theorem~2.4 and Corollary~2.5,
pp.~9--11]{HoshiTsujimura2021} give explicit nongeometric pro-$p$
sections of punctured projective lines. Their hypotheses include
an odd regular $p$, a base field containing $\mu_p$, and prescribed
$p$-unit Kummer coordinates in the cyclotomic extension unramified
outside $p$. The theorem lifts those coordinates and gives criteria
excluding both ordinary points and cusps; it imposes no Selmer
condition. The two families over $\Q$ in Section~\ref{sec:arbitrary}
use a free arithmetic source with an explicit descent under complex
conjugation. Both fail local geometry at $3$. Neither that
construction nor the cited preprint supplies a noncuspidal Selmer
section of $U_{11}/\Q$.

\begin{proposition}[Exponent one]\label{prop:one}
For $n=1$, the all-cuspidal Selmer assertion is false for every
geometric quotient considered in this paper, including the full
geometrically pro-$3$ group.
\end{proposition}
\begin{proof}
The curve $U_1$ is $\mathbb P^1\setminus\{0,1,\infty\}$ with
coordinate $t=X/(-Z)$. Its point $t=3$, or $[3:-2:-1]$, supplies
a global ordinary point section and hence a Selmer section. The
two invertible functions $t,1-t$ send its geometric abelian class
to $(\kappa(3),\kappa(-2))$. A cusp at zero has second coordinate
zero; a cusp at one has first coordinate zero; a cusp at infinity
has difference of coordinates $\kappa(-1)$. Valuations at two and
three exclude all three possibilities. The same calculation in
$3$-adic Kummer groups remains valid, with $\kappa_3(-1)=0$.
Thus the point section is noncuspidal even after abelianization,
and therefore in every stated quotient retaining that abelianization.
\end{proof}

\begin{proposition}[Even plus exponents]\label{prop:even-plus}
For every even $n\geq2$, $X_n(\mathbb R)=\varnothing$. Consequently
every arithmetic quotient of $\pi_1(U_n)$ has empty point-or-cusp
Selmer section set under the all-place convention above.
\end{proposition}
\begin{proof}
Three even real powers are nonnegative, and their sum is zero only
when all three coordinates are zero. This is not a projective point.
Neither an ordinary real point nor a real cusp is therefore available
to satisfy the local condition. This does not assert emptiness of
the abstract section set of a geometric pro-$3$ quotient.
\end{proof}

\subsection{Unresolved boundaries}
The unrestricted geometrically pro-$3$ Selmer assertion for odd
exponents outside $3^r$, $5\cdot3^r$, and $7\cdot3^r$, $r\ge1$, is not decided here.
Compact localization at the base exponent eleven is proved in
Theorem~\ref{thm:eleven-compact-selmer}; the corresponding higher
$11\cdot3^r$ compact groups and the base exponent's affine boundary
step remain untreated. Compact localization also holds at base
exponents five and seven, but this alone does not identify all their
affine Selmer sections with global cusp sections. At these two base
exponents, Theorem~\ref{thm:finitefield-boundary-intro} confines a
possible finite-field boundary collision to one coordinate component
over a field of size \(1\) modulo the exponent. The remaining question
is the existence of a Frobenius-fixed pro-three path from a rational
cusp to its nontrivial diagonal translate, equivalently the Lang
equation \eqref{eq:ff3-11-2}. Even a density-one boundary separation
statement would need an additional argument to give global vanishing
of the transverse affine boundary class.
Theorem~\ref{thm:realconditionalmixedintro} gives an explicit sufficient
criterion for the stated two-prime mixed exponents, with its compact
and higher-grade hypotheses retained. The
unrestricted single-prime result also leaves quotients whose ordinary
lower-central factors have torsion in need of further input. The
finite-derived and prosolvable classification leaves open odd
exponents all of whose prime divisors are irregular. Its first
unresolved prime is \(37\). The finite local-centre bounds in
Theorem~\ref{thm:finite-local-centres-intro} reduce that local
question but do not prove global cuspidality. In particular, the residual geometric
$\mathrm{SL}_2(\F_\ell)$ representation used in the modular Fermat
argument for $\ell\geq5$ is not a representation of a prosolvable
geometric quotient: its image is a nontrivial perfect finite group.
Passing to an abelian geometric deck cover does not remove that
obstruction: its kernel contains the geometric commutator subgroup,
whose image is the whole perfect residual group. One cannot transplant
that representation to a geometrically metabelian or pro-$3$ group.
Separate retained-quotient and local
arguments are required.

\subsection{A limitation of the single-place arithmetic-source method}
\label{subsec:source-dimension-obstruction}

\begin{proposition}[A dimension obstruction to a single local arithmetic source]
\label{prop:pro3-source-dimension-obstruction}
Let $K/\mathbb Q$ be a finite Galois extension containing $\mu_3$,
let $d=[K:\mathbb Q]$, and let $t$ be the number of its primes above
three.  Let $S$ be a finite set of finite primes of $K$, containing
all primes above three and stable under $\operatorname{Gal}(K/\mathbb Q)$,
and write $s=|S|-t$.  For any $v\mid3$ one has
\[
 \dim_{\mathbb F_3}\ker\left(
 H^1(G_{K,S},\mathbb F_3)\longrightarrow H^1(K_v,\mathbb F_3)
 \right)
 \ge \frac{d(t-2)}{2t}+t+s-2.
\]
In particular, the homomorphism from $G_{K_v}$ to the maximal
pro-three quotient of $G_{K,S}$ cannot be surjective if $t\ge3$,
or if $t=2$ and $s>0$.

For every such $K$ containing $\mathbb Q(\zeta_{11})$, the second
conclusion applies when $S$ contains all primes above both three and
eleven.  For every such $K$ containing $\mathbb Q(\zeta_{13})$,
the first conclusion applies even when $S$ consists only of the
primes above three.
\end{proposition}
\begin{proof}
Since $\mu_3\subset K$, the constant coefficient module $\mathbb F_3$
can be identified with $\mu_3$.  Kummer theory embeds
$\mathcal O_{K,S}^{\times}/(\mathcal O_{K,S}^{\times})^3$
in $H^1(G_{K,S},\mu_3)$.  The field $K$ is totally imaginary,
with $r_2=d/2$ complex places.  Dirichlet's unit theorem gives
$S$-unit rank $r_2-1+|S|$; the three-primary roots of unity
contribute one further dimension modulo cubes.  Hence
\[
 \dim H^1(G_{K,S},\mathbb F_3)\ge d/2+t+s.
\]
No class-number assumption is used.

Local Kummer theory identifies the target with $K_v^\times/K_v^{\times3}$.
The valuation supplies one dimension.  A sufficiently small principal
unit group is the additive group of a finite-index $\mathbb Z_3$
lattice in $K_v$, and the three-primary roots of unity supply one
additional dimension.  Thus the target has dimension
$[K_v:\mathbb Q_3]+2=d/t+2$; the equality of local degrees uses
Galoisness.  Subtracting gives the claimed lower bound.  Surjectivity
on the pro-three arithmetic source would make restriction injective
on its homomorphisms to $\mathbb F_3$, contradicting a positive bound.

There are two primes above three in $\mathbb Q(\zeta_{11})$, since
$\operatorname{ord}_{11}(3)=5$.  Every extension containing that
field has at least two such primes, while the allowed primes above
eleven give $s>0$.  Similarly
$\operatorname{ord}_{13}(3)=3$, so $\mathbb Q(\zeta_{13})$ has
four primes above three, and every containing extension has
$t\ge4$.  This proves the final assertions.
\end{proof}

\begin{remark}[What the obstruction does not say]
\label{rem:pro3-source-dimension-scope}
This is an obstruction to the particular universal arithmetic-source
argument, not to injectivity for a specified geometric coefficient
module, and not a construction of a noncuspidal Selmer section.
A global $G_{\mathbb Q}$-equivariant cocycle that vanishes at a chosen
place also has conjugate local restrictions constrained at the other
places above three.  Joint localization or an isotypic cohomology
calculation can therefore be stronger than the single-place argument
tested here.  Compact geometric Kummer classes may also be unramified
at the exponent prime after potential good reduction, allowing a
smaller set $S$.  When $t=2$ and $s=0$, the displayed dimension bound
is zero and supplies no obstruction.
\end{remark}

\subsection{Why conjugate local generators do not suffice}

Nonzero first-cohomology classes with locally varying coboundary parameters
are standard local--global cohomology phenomena; see
\citet[\S2 and Examples~3.4--3.6, pp.~321 and 327--328]{DvornicichZannier2001}.
The relative profinite formalism is developed by
\citet[\S2.1, Proposition~2.4]{WilkesRelative2019}, and its extension and
projectivity interpretation by \citet{WilkesRelative2025}.
The following explicit example verifies the additional generation and
invariant hypotheses which arise in the present source argument.

\begin{proposition}\label{prop:conjugate-source-not-enough}
There exist a profinite group \(\Gamma=P\rtimes\Delta\), with \(P\) a
finitely generated free pro-three group and \(|\Delta|\) prime to three,
a closed subgroup \(D\), and a finite free continuous
\(\mathbf Z_3[\Gamma]\)-module \(T\), such that:
\begin{enumerate}
\item the image of \(P\) on \(T\) is pro-three;
\item the conjugates of \(D\) generate \(\Gamma\);
\item \(T^D=0\);
\item \(H^1(\Gamma,T)\longrightarrow H^1(D,T)\) is not injective,
even after tensoring with \(\mathbf Q_3\).
\end{enumerate}
The same nonzero class restricts to zero on every conjugate of \(D\).
\end{proposition}
\begin{proof}
Index five symbols by \(i\in\mathbf Z/5\mathbf Z\). Let \(P\) be free
pro-three on \(x_i\), and let the dihedral group \(\Delta=D_{10}\) act
by its usual permutations of the symbols. Put
\[
 T=\left\{(a_i)\in\mathbf Z_3^5:\sum_i a_i=0\right\}.
\]
Let \(\Delta\) permute the coordinates. The homomorphism
\(\epsilon:P\to\mathbf Z_3\), defined by \(\epsilon(x_i)=1\), is
\(\Delta\)-invariant. Let \(g\in P\) act by \(4^{\epsilon(g)}\).
This gives a continuous action of \(\Gamma=P\rtimes\Delta\) on \(T\);
its \(P\)-image is pro-three.

Let \(\Delta_i\) be the reflection subgroup fixing \(i\), and put
\[
 D_i=\langle x_i\rangle\rtimes\Delta_i,\qquad
 b_i=5e_i-\sum_j e_j .
\]
Then \(b_i\in T^{\Delta_i}\), and \(\sigma b_i=b_{\sigma i}\).
Assigning \(x_i\) to \((3b_i,4)\) in the pro-three semidirect product
\(T\rtimes\overline{\langle4\rangle}\) gives, by freeness, a continuous
cocycle \(c:P\to T\) with \(c(x_i)=3b_i\). It is
\(\Delta\)-equivariant, so it extends to \(\Gamma\) with \(c|_\Delta=0\).

On \(D_i\), the cocycle is the coboundary of \(b_i\): it takes the value
\((4-1)b_i=3b_i\) on \(x_i\) and zero on \(\Delta_i\).
If it were a global coboundary with parameter \(b\in T\), its zero
restriction to \(\Delta\) would imply \(b\in T^\Delta=0\), a
contradiction. The same argument applies over \(\mathbf Q_3\).

Multiplication by \(4-1=3\) is injective on \(T\), so \(T^{D_i}=0\).
The subgroups \(D_i\) generate \(P\) through the \(x_i\), and generate
\(\Delta\) through the five reflections. Hence they generate \(\Gamma\).
They are conjugate under \(\Delta\). Take \(D=D_0\).
Restriction to every further \(\Gamma\)-conjugate is also zero, by
the conjugacy compatibility of restriction in cohomology.
\end{proof}

\begin{remark}\label{rem:conjugate-source-scope}
This is an abstract continuous group-cohomology counterexample, not an
arithmetic Selmer-section counterexample. It shows that equivariance,
generation by conjugate local subgroups, and zero local invariants do
not justify replacing a single surjective local source by its normal
closure. The local coboundary parameters still require compatibility.
\end{remark}

\begin{lemma}[A module-specific sufficient condition]
\label{lem:source-local-parameter}
Let \(\Gamma=P\rtimes\Delta\), with \(P\) pro-three and
\(\Delta\) finite of order prime to three. Let
\(D=P_0\rtimes\Delta_0\) be a closed subgroup, where
\(\Delta_0\subset\Delta\), and suppose the \(\Delta\)-conjugates of
\(P_0\) generate \(P\). If \(T\) is finite free over \(\mathbf Z_3\)
and \(T^{P_0}=0\), there is a natural injection
\[
 \ker\!\left(H^1(\Gamma,T)\to H^1(D,T)\right)
       \hookrightarrow T^{\Delta_0}/T^\Delta.
\]
In particular equality of the two invariant lattices implies
localization injectivity.
\end{lemma}
\begin{proof}
Normalize the cocycle to vanish on \(\Delta\), using
\(H^1(\Delta,T)=0\). A local coboundary then has a unique parameter
\(b\), since \(T^{P_0}=0\), and \(b\in T^{\Delta_0}\).
Changing the normalized cocycle changes \(b\) by an element of
\(T^\Delta\). If its class in the displayed quotient is zero,
subtract that global coboundary. The cocycle now vanishes on
\(\Delta\) and \(P_0\), hence on every \(\Delta\)-conjugate of \(P_0\)
and therefore on \(\Gamma\).
\end{proof}
\subsection{The relation and character conditions for joint localization}

The preceding counterexample leaves a precise compatibility condition.
It can be expressed either through relations among local source groups
or, for abelian coefficient actions, through selected character components
of an arithmetic Galois module.

\begin{proposition}[The exact local-parameter condition]
\label{prop:source-exact-relation-condition}
Let \(\Gamma=P\rtimes\Delta\), where \(P\) is pro-three and
\(\Delta\) is finite of order prime to three. Let
\(D=P_0\rtimes\Delta_0\), with \(\Delta_0\subset\Delta\), be a
closed subgroup. Suppose the \(\Delta\)-conjugates of \(P_0\)
generate \(P\). Let \(T\) be a finite free continuous
\(\mathbf Z_3[\Gamma]\)-module with \(T^{P_0}=0\).

Form the free pro-three product \(F\) of copies of \(P_0\) indexed
by \(\Delta/\Delta_0\), with the action induced by the conjugation
action of \(\Delta_0\) on \(P_0\). The conjugation maps give a
surjection \(F\to P\). Put \(R=\ker(F\to P)\).
For \(b\in T^{\Delta_0}\), let \(c_b\) be the cocycle on \(F\)
whose restriction to the factor indexed by \(\sigma\Delta_0\)
is the coboundary with parameter \(\sigma b\).
Then
\begin{equation}
 \ker\!\left(H^1(\Gamma,T)\longrightarrow H^1(D,T)\right)
 \simeq
 \frac{\{b\in T^{\Delta_0}:c_b|_R=0\}}{T^\Delta}.
 \label{eq:source-exact-relation-condition}
\end{equation}
\end{proposition}
\begin{proof}
The definition of \(c_b\) is independent of the chosen coset
representative because \(b\) is fixed by \(\Delta_0\).
The action of \(F\) on \(T\) factors through \(P\), whose image
is pro-three. The universal property of the free pro-three product,
applied to the cocycle graph in \(T\rtimes\rho(P)\), gives a
continuous cocycle \(c_b\). It is \(\Delta\)-equivariant.
Since \(R\) acts trivially on \(T\), its restriction to \(R\)
is a homomorphism, and it descends to \(P\) precisely when this
homomorphism is zero. In that case it extends to \(\Gamma\)
by taking value zero on \(\Delta\), and it is a coboundary
on \(D\).

Conversely, a class in the displayed kernel has a representative
vanishing on \(\Delta\), because \(|\Delta|\) is a unit in
\(\mathbf Z_3\). Its unique local coboundary parameter lies in
\(T^{\Delta_0}\); uniqueness uses \(T^{P_0}=0\).
Equivariance determines its restrictions to every factor of \(F\).
The existence of the original cocycle on \(P\) is exactly
\(c_b|_R=0\). Two normalized representatives differ by a
coboundary with parameter in \(T^\Delta\). This proves
\eqref{eq:source-exact-relation-condition}.
\end{proof}

\paragraph{An arithmetic module testing only the required characters.}
The compact kernel of localization at the places over the coefficient
prime is a familiar fine-Selmer object; compare
\citet[Eq.~(1.2), author p.~2]{WuthrichFineSelmer2007}.
The following formulation gives a direct character condition and
does not assert vanishing of that kernel.

\begin{proposition}[An abelian-character localization condition]
\label{prop:cm-character-localization-condition}
Let \(E/\mathbf Q\) be finite Galois. Let \(S\) be a finite
set of rational primes containing three and the ramification
of \(E/\mathbf Q\) and of a finite-dimensional continuous
\(\mathbf Q_3\)-representation \(V\).
Let \(E_S\) be the maximal extension of \(E\) unramified outside
the primes over \(S\). Suppose \(E_\infty\subset E_S\) is Galois
over \(\mathbf Q\), is abelian over \(E\), and trivializes
the action of \(G_E\) on \(V\). Put
\[
 H=\operatorname{Gal}(E_\infty/E),\qquad
 \Delta=\operatorname{Gal}(E/\mathbf Q).
\]
Suppose that, over a finite coefficient extension
\(\Omega/\mathbf Q_3\), the \(H\)-representation
\(V_\Omega=V\otimes_{\mathbf Q_3}\Omega\) is a direct sum
of characters, each nontrivial on \(H\) and on every local
decomposition image \(H_v\) for \(v\mid3\).

Let \(X_S^{[3]}(E_\infty)\) be the maximal abelian pro-three
quotient of \(\operatorname{Gal}(E_S/E_\infty)\) in which
the decomposition groups at every prime above three vanish.
Thus it classifies abelian pro-three extensions unramified
outside \(S\) and split at every three-adic prime.
There is a natural isomorphism
\begin{equation}
\begin{split}
 &\ker\!\left(H^1(G_{\mathbf Q,S},V)
            \longrightarrow H^1(\mathbf Q_3,V)\right)
                \otimes_{\mathbf Q_3}\Omega\\
 &\qquad\simeq
 \operatorname{Hom}^{\mathrm{cts}}_H
       \bigl(X_S^{[3]}(E_\infty),V_\Omega\bigr)^\Delta.
\end{split}
\label{eq:cm-character-localization-condition}
\end{equation}
In particular only the character components which occur in
\(V_\Omega\) need to vanish.
\end{proposition}
\begin{proof}
Let \(\chi\) be one of the characters. Choose \(h\in H\)
with \(\chi(h)\ne1\). The inner action of \(h\) on continuous
group cohomology is trivial. Since \(H\) is abelian, the same
action is multiplication by \(\chi(h)\) on
\(H^j(H,\Omega(\chi))\). The scalar \(\chi(h)-1\) is invertible,
so these cohomology groups vanish for all \(j\).
Inflation--restriction therefore identifies
\[
 H^1(G_{E,S},\Omega(\chi))
 \simeq
 \operatorname{Hom}^{\mathrm{cts}}_H
 \bigl(\operatorname{Gal}(E_S/E_\infty)^{\mathrm{ab},(3)},
       \Omega(\chi)\bigr).
\]
Continuous additive images in the coefficient line are pro-three,
which explains the indicated abelian pro-three quotient.

The same argument applies locally because
\(\chi|_{H_v}\ne1\). Consequently vanishing of a local class
is equivalent to the corresponding homomorphism killing the
decomposition subgroup over that place. Summing over the
characters gives the joint-localization kernel over \(E\).

Finite-group cohomology with coefficients in \(\Omega\) vanishes
in positive degrees by averaging. Globally, inflation--restriction
through \(E/\mathbf Q\) identifies the rational cohomology
with the \(\Delta\)-invariants over \(E\). Locally, its analogous
form identifies \(H^1(\mathbf Q_3,V_\Omega)\) with the
\(\Delta\)-invariants of the product over \(v\mid3\).
Taking kernels proves \eqref{eq:cm-character-localization-condition}.
The action of \(\Delta\) on the right is well defined: changing
a lift in \(\operatorname{Gal}(E_\infty/\mathbf Q)\) changes
both source and target by the same \(H\)-action.
\end{proof}

\begin{lemma}[Integral recovery from rational injectivity]
\label{lem:localization-integral-recovery}
Let \(T\) be a stable finite free \(\mathbf Z_3\)-lattice in \(V\),
and assume \(V^{G_{\mathbf Q_3}}=0\).
If
\[
 H^1(G_{\mathbf Q,S},V)\longrightarrow H^1(\mathbf Q_3,V)
\]
is injective, then so is the corresponding map with coefficients
in \(T\).
\end{lemma}
\begin{proof}
A class in the integral kernel is zero after passing to \(V\).
The exact sequence \(0\to T\to V\to V/T\to0\) makes it the
boundary of an element of \((V/T)^{G_{\mathbf Q,S}}\).
The local boundary map from \((V/T)^{G_{\mathbf Q_3}}\) is
injective because \(V^{G_{\mathbf Q_3}}=0\).
The global invariant is also a local invariant, and its local
boundary is zero; it is therefore zero. This proves the claim.
\end{proof}

\begin{lemma}[Local nontriviality for the eleven coefficient characters]
\label{lem:eleven-coefficient-local-nontriviality}
Let \(N_{11}\) be the geometric pro-three group of \(U_{11}\),
and put \(E=\mathbf Q(\zeta_{11})\). Over a finite splitting
coefficient field, every character in a positive ordinary
lower-central grade of \(N_{11}\) is nontrivial on
\(G_{E_v}\) for each \(v\mid3\). The same holds for every
nonzero character in a stable quotient of such a grade.
In particular these rational coefficient modules have zero
local invariants at three.
\end{lemma}
\begin{proof}
The decomposition \(J_{11}\sim A_{11}^3\times B^6\), with
\(B=\operatorname{Jac}(y^{11}=x(1-x)^2)\), reduces the
compact degree-one calculation to two CM types. The
holomorphic types are
\[
 H_A=\{1,2,3,4,5\},\qquad H_B=\{1,2,3,6,7\};
\]
equivalently the reflex types are
\[
 R_A=\{1,6,4,3,9\},\qquad R_B=\{1,6,4,2,8\}.
\]
These follow from the Fermat quotient type formula
\citep[\S2.1, author p.~5]{ShuExpanded2021}.
The two Frobenius-three orbits are
\[
 C=\langle3\rangle=\{1,3,9,5,4\},\qquad -C.
\]
The intersections of \(R_A\) with these orbits have sizes
four and one; those of \(R_B\) have sizes two and three.
Multiplication by a coefficient automorphism either preserves
or exchanges the two orbits, so every conjugate type meets
both.

At the good prime three, local CM reciprocity on an open
group of principal units is the inverse reflex norm
\citep[Theorem~5.1, pp.~483--484]{Lombardo2017}.
Every compact character therefore has, at every \(v\mid3\),
a nonempty set of algebraic unit exponents, all of the same
sign. The boundary characters are cyclotomic and have the
same sign. A character in a positive tensor power is a
product with nonnegative multiplicities, so its unit exponent
vector cannot be zero. Indeed, after taking the three-adic
logarithm on a sufficiently small unit group, the distinct
embeddings are linearly independent and a nonzero exponent
vector gives a nonzero linear map.

Ordinary lower-central grades are generated by brackets in
degree one, and their actions are subquotients of these tensor
powers. Coordinate-cusp differences are killed by eleven, so the
generalized-Jacobian extension splits on three-adic Tate modules,
as explained in Remark~\ref{rem:boundary-finite-invariant-obstruction}.
Over the splitting coefficient field the action is therefore
diagonal, so subquotients retain only characters from that
list. This proves the claim.
\end{proof}

\begin{remark}[A degree-two boundary component]
\label{rem:boundary-finite-invariant-obstruction}
For \(a=11\) or \(13\), put \(E=\mathbf Q(\zeta_a)\).
The boundary module of \(U_a\) has the integral decomposition
\[
 B_a\simeq\mathbf Z_3(1)^2\oplus
     \left(\operatorname{Ind}_{G_E}^{G_{\mathbf Q}}
                       \mathbf Z_3(1)\right)^3.
\]
There are three rational boundary points and three primitive
cyclotomic orbits; eliminating one rational basis vector proves
the decomposition.
Coordinate-cusp differences are killed by \(a\), so the
generalized-Jacobian extension splits on three-adic Tate modules.
Hence \(B_a\) is a summand of the geometric abelian module.
For the free geometric pro-three group, its second lower-central
grade is the exterior square of that module.

Choose a nontrivial character \(\eta\) of
\((\mathbf Z/a)^\times/\langle3\rangle\), using a coefficient
extension when necessary. Bracketing a rational boundary line
with the \(\eta\)-line in an induced summand gives
\(\Omega(2)\otimes\eta\) in the second grade.
For eleven one can take the odd quadratic character; for thirteen
the displayed quotient is \(C_4\).
The finite character \(\eta\) is trivial on the local group at
three and nontrivial globally. The square of the finite
mod-three cyclotomic character is trivial.
Thus the associated finite coefficient representation has more
local invariants than global invariants already in degree two.
An invariant-equality criterion alone cannot handle all these
grades. This does not establish a nonzero localization kernel:
the relation condition or the arithmetic character condition
must still be tested.
\end{remark}

\begin{remark}[Finite ramification and the base-exponent problem]
\label{rem:joint-localization-ramification-scope}
The character condition concerns classes in
\(H^1(G_{\mathbf Q,S},V)\). It does not force an arbitrary
affine Selmer section to have its whole relative cocycle
unramified outside \(S\).
At a prime exponent different from three, compact
three-adic Kummer classes do not detect differences between
coordinate cusps, since those differences have order prime to
three. Compact localization and one locally cuspidal
representative therefore do not by themselves restrict boundary
ramification to one global cusp direction.

If the global abelian class is already that of one cusp,
the point-or-cusp local condition can control inertia in the
subsequent central induction. A nonzero boundary inertia vector
forces the same cusp direction; a zero vector makes the whole
inertia trivial. This is a separate step, preceding application
of a finite-ramification cohomological condition.
The propositions above consequently supply formal
module-specific conditions, not a new base-eleven or
base-thirteen pro-three cuspidality theorem.
\end{remark}


\Needspace{10\baselineskip}
\subsection{The obstruction to the rank-zero CM route at exponent seventeen}

Let $F_{17}$ be the smooth projective Fermat curve
$X^{17}+Y^{17}=Z^{17}$, rationally isomorphic to $X_{17}$ by
changing the sign of $Z$, and put
\[
 C_s:\quad v^{17}=u(1-u)^s,\qquad
 A_s=\operatorname{Jac}(C_s),\qquad 1\leq s\leq15.
\]
The genus of each $C_s$ is eight.  The hyperelliptic quotient $C_1$
is isomorphic over $\mathbb Q$ to $y^2+y=x^{17}$.

\begin{proposition}[The degree-eight Hilbert route has zero central values]
\label{prop:lm-seventeen-obstruction}
There is a $\mathbb Q$-isogeny
\[
 \operatorname{Jac}(F_{17})\sim_{\mathbb Q}
 A_1^3\times A_2^6\times A_3^6.
\]
Each $A_i$ is absolutely simple and has CM by
$K=\mathbb Q(\zeta_{17})$.  Its associated weight-two Hilbert
representation over $F=K^+$ has trivial finite central character and
is supercuspidal at the unique place of $F$ above $17$; its conductor
at that place has exponent three.  Thus the nonprincipal-series
alternative in Nekov\'a\v r's 2012 Theorem~B applies to these forms.
Nevertheless,
\[
 L(A_1/\mathbb Q,1)=L(A_2/\mathbb Q,1)
                  =L(A_3/\mathbb Q,1)=0.
\]
In particular, none of these factors supplies the nonzero central
value required by that rank-zero theorem.
\end{proposition}

\begin{proof}
The quotient maps $(x,y)\mapsto(x^{17},xy^s)$ give the rational
isogeny to $\prod_{s=1}^{15}A_s$; see \citet[\S2, pp.~207--208]{GrossRohrlich1978}.  Permuting the triple $(1,s,16-s)$ and
rescaling modulo $17$ gives exactly the following orbits:
\[
 \begin{array}{c|c|c}
 \text{representative}&s\text{-values}&\text{multiplicity}\\ \hline
 (1,1,15)&1,8,15&3\\
 (1,2,14)&2,5,7,9,11,14&6\\
 (1,3,13)&3,4,6,10,12,13&6.
 \end{array}
\]
These equivalences are rational birational equivalences of the
quotient curves.  The displayed product follows.
The CM simplicity assertion follows from \citet[Lemma~1.6 and the following remark, pp.~206--207]{GrossRohrlich1978}: their only
imprimitive case requires a nonidentity cube root of unity modulo
the exponent, and $3\nmid16$.  Equivalently, the cyclic field
$K$ has no proper CM subfield: every nonidentity subgroup of its
cyclic group of order sixteen contains complex conjugation.

Here is an explicit calculation of the local condition.  For
$(r,s,t)=(1,i,16-i)$, set
\[
 z=r^r s^s(t-17)^t\in(\mathbb Z/17^2\mathbb Z)^\times,
 \qquad d=\frac{z^{16}-1}{17}\pmod {17}.
\]
The three values of $(z,d)$, in the order above, are
\[
 (178,8),\qquad(76,13),\qquad(215,11).
\]
In particular, all three $d$ are nonzero.  \citet[Proposition~3.12(2), author p.~16]{ShuExpanded2021} therefore gives conductor
$\mathfrak P^2$ for the Jacobi-sum Hecke character at the prime
$\mathfrak P$ of $K$ over $17$.  Its restriction to
$U^1/U^2\simeq\mathbb F_{17}$ is nontrivial.  With
$\lambda=\zeta_{17}-\zeta_{17}^{-1}$, complex conjugation takes
$\lambda$ to $-\lambda$, and hence acts by $-1$ on $U^1/U^2$.
The character and its conjugate thus have inverse, distinct,
order-seventeen restrictions.  Their induction from $K_{\mathfrak P}$
to $F_{\mathfrak p}$ is an irreducible two-dimensional Weil
representation, giving a supercuspidal local representation.
The tame quadratic extension has different exponent one and residue
degree one; the induction formula for the Artin conductor gives
$1+2=3$.

The central character is also explicit.  In unitary normalization
the Jacobi-sum character restricts to the quadratic character
$\eta_{K/F}$ on $\mathbb A_F^\times$; see \citet[\S3.4, author p.~15]{ShuExpanded2021}.  The determinant character of quadratic automorphic
induction is this restriction times $\eta_{K/F}$, hence is trivial.
The CM infinity exponents differ by one at every complex place,
so quadratic induction gives parallel weight two.  Thus the even
degree $[F:\mathbb Q]=8$ creates no obstruction to
\citet[Theorem~B(b), p.~593]{Nekovar2012}.

For $\delta=1$, \citet[Theorem~3.13, author pp.~17--18]{ShuExpanded2021} gives
\[
 w(A_i)=-\left(\frac{rst\,d}{17}\right).
\]
The archimedean factor is $(-1/17)=1$ and the other finite local
factors are one.  Since the three residues $rst\,d$ are respectively
$1,7,4$ modulo $17$, the signs are $-1,+1,-1$.  This also agrees
with the original table of \citet[p.~219]{GrossRohrlich1978}.  The functional
equation forces the first and third central values to vanish.

For the remaining factor, \citet[Theorem~2.1, p.~208]{GrossRohrlich1978}
gives a rational point of infinite order in $A_2(\mathbb Q)$.
If its central value were nonzero, the Hilbert form just described
would satisfy every hypothesis of Nekov\'a\v r Theorem~B(b).
That theorem would give zero rational Bloch--Kato Selmer group
for each coefficient realization.  But its representation induces
to $V_\ell(A_2)$ over $\mathbb Q$, by the Jacobi-sum construction;
Shapiro's lemma would therefore give
$H^1_f(\mathbb Q,V_\ell(A_2))=0$.
The Kummer class of the rational point of infinite order contradicts
this conclusion.  Hence the middle central value vanishes as well.
Its root number is positive, so its analytic order of vanishing is
at least two.
\end{proof}

\begin{corollary}[No rank-zero abelian quotient at seventeen]
\label{cor:seventeen-no-rankzero-quotient}
Every positive-dimensional abelian quotient over $\Q$ of
$\Jac(F_{17})$ has positive Mordell--Weil rank over $\Q$.
\end{corollary}
\begin{proof}
\citet[Concluding remark~2, p.~665]{McCallum1988} states
$\mathbb Z[\zeta_{17}]$-rank one over $K$ for the two wild-split
classes at exponent seventeen, namely the
classes of $(1,1,15)$ and $(1,3,13)$; the remaining class is the
wild-nonsplit class $(1,2,14)$. The residues $-2rst\,d$ are
$15,3,9$ modulo $17$, respectively, identifying those two split classes;
also $17\nmid B_8B_{10}$. This gives rational rank one:
the vector space $M=A_i(K)\otimes_{\mathbb Z}\mathbb Q$ is
one-dimensional over $K$, and $\operatorname{Gal}(K/\mathbb Q)$
acts semilinearly through its natural action on the CM
endomorphisms.  Galois descent (equivalently Hilbert~90 in dimension
one) gives $\dim_{\mathbb Q}M^{\operatorname{Gal}(K/\mathbb Q)}=1$;
these invariants equal $A_i(\mathbb Q)\otimes\mathbb Q$, since
taking invariants commutes with tensoring by $\mathbb Q$ for a
finite group.  Together with Gross--Rohrlich's
point of infinite order on $A_2$, this shows that every absolutely
simple factor in the displayed isogeny has positive rational rank.
Consequently no positive-dimensional abelian quotient of
$\operatorname{Jac}(F_{17})$ has rank zero over $\mathbb Q$.
Indeed a simple isogeny factor of such a quotient must occur in
the displayed decomposition, and rank is additive and isogeny
invariant.  This final rank statement uses McCallum's published
conclusion; the central-value obstruction above does not need it.
\end{proof}

\paragraph{Scope of the obstruction.}
The calculation verifies the finite-place hypothesis rather than
leaving it as an unresolved arithmetic condition.  The failed input
is precisely central nonvanishing: all three values are zero.
It rules out a rank-zero abelian target obtained from the Fermat
Jacobian and the present rank-zero Selmer theorem.  It proves no
noncuspidal section and no failure of a section conjecture.  A
rank-one or rank-two descent calculation together with a curve-level
local intersection argument would be a different route; none is
asserted here.

The section classification at exponent $17$ is nevertheless given
by Theorem~\ref{thm:regular-solvable-intro}; its proof uses the
regular-prime local argument rather than a rank-zero quotient.

\subsection{A complementary factor and the remaining arithmetic question at thirty-one}
\label{subsec:thirtyone-gate}

The next calculation concerns one cyclic quotient of the Fermat curve.
It identifies its only possible rational rank-zero factor, but does not
establish a new section classification. Let $C/\Q$ be the smooth proper
model of
\[
 y^{31}=x(1-x)^5,
\]
put $J=\Jac(C)$, and define
\[
 c(x,y)=\left(\frac1{1-x},-\frac{y^5}{1-x}\right),\quad
 B=C/\langle c\rangle,\quad A=\Jac(B),\quad
 e=\frac{1+c+c^2}{3}.
\]
Let $P$ be the image of $1-e$ up to rational isogeny; equivalently,
one may take the connected image of $2-c-c^2$. Put
\[
 K=\Q(\zeta_{31}),\quad H=\langle5\rangle=\{1,5,25\},\quad
 M=K^H,\quad F_0=M^+.
\]

\begin{proposition}[The rational factors and their rank bounds]
\label{prop:thirtyone-rational-factors}
There is a rational isogeny $J\sim_\Q A\times P$, where
$\dim A=5$, $\dim P=10$, and
\[
 \operatorname{End}^0_\Q(J)=\Q\times\Q(\zeta_3),\qquad
 \operatorname{End}^0_\Q(A)=\Q,\qquad
 \operatorname{End}^0_\Q(P)=\Q(\zeta_3).
\]
Both factors are $\Q$-simple, and
\[
 \operatorname{rank}A(\Q)\ge1,\qquad
 \operatorname{rank}J(\Q)\le4,\qquad
 \operatorname{rank}P(\Q)\in\{0,2\}.
\]
Thus $P$ is the only possible positive-dimensional rational
rank-zero isogeny factor of this particular $J$.
\end{proposition}
\begin{proof}
The degree-thirty-one map to the $x$-line has three fully ramified
points, so Riemann--Hurwitz gives $g(C)=15$.
Direct substitution gives
\[
 c^2(x,y)=\left(\frac{x-1}{x},-\frac{y^{25}}{x(1-x)^4}\right),
 \qquad c^3=1.
\]
The fixed $x$-coordinates satisfy $x^2-x+1=0$.
For a primitive sixth root $\rho$ with $\rho+\rho^{-1}=1$,
the point $(\rho,\rho^2)$ is fixed. Multiplication of its
$y$-coordinate by $\zeta_{31}^i$ changes the fixed-point condition
to $5i=i$, so this is the unique fixed point in that fiber.
The conjugate fiber contributes one more, and the three branch
points are cyclically permuted. Hence
$28=3(2g(B)-2)+4$, giving $g(B)=5$ and the stated dimensions.
Norm and pullback give the displayed rational isogeny.

The Jacobi-sum CM type on $K$ is
\[
 S=\{1,2,3,4,5,7,8,9,10,13,14,15,19,20,25\};
\]
it consists of the residues $b$ with $b+(5b\bmod31)<31$.
Its multiplicative stabilizer is exactly $H$. In the cyclic
quotient of order ten, with generator induced by $3\bmod31$,
the type is
\[
 T=\{0,1,2,4,8\}\subset\mathbb Z/10\mathbb Z,
\]
whose translation stabilizer is trivial. The CM decomposition
of \citet[Lemma~1.6, pp.~206--207]{GrossRohrlich1978} therefore
gives $J_{\overline\Q}\sim D^3$, where $D$ is absolutely simple
of dimension five with CM by $M$. In particular
$\operatorname{End}^0_{\overline\Q}(J)=M_3(M)$.
The automorphism $c$ and the cyclotomic endomorphisms satisfy
$cac^{-1}=\sigma_5(a)$. Their crossed product $K\rtimes H$ is
$\operatorname{End}_M(K)=M_3(M)$, so these endomorphisms generate
the whole geometric algebra. Galois acts on the unique expansion
$\sum_{j=0}^2a_jc^j$ through its usual action on $a_j\in K$,
and fixes $c$. The rational fixed algebra is therefore
$\Q[c]=\Q\times\Q(\zeta_3)$. Its two central factors give the
claimed endomorphism algebras and rational simplicity.
The same argument over $F_0$ gives $F_0[c]$; in particular $A$
is absolutely simple and $A_{F_0}$ has real multiplication by $F_0$.

\citet[Theorem~2.1, pp.~208--209]{GrossRohrlich1978} gives a
rational divisor class of infinite order on this cyclic quotient:
\[
 D_0=[(\rho,\rho^2)+(\rho^{-1},\rho^{-2})-2\infty].
\]
The exclusions in that theorem are $s=1,15,29$, so $s=5$ is
permitted. The two displayed points are fixed by $c$.
Differences of branch points are killed by $31$, and $c$
permutes those branch points; hence $cD_0-D_0$ is torsion.
For the quotient map $\pi:C\to B$,
$\pi^*\pi_*D_0=(1+c+c^2)D_0$ is $3D_0$ plus torsion.
Thus $\pi_*D_0\in A(\Q)$ has infinite order.

For the upper bound use the equivalent triple $(5,1,-6)$.
The numerators of $B_2,B_4,\ldots,B_{28}$ have respective
nonzero residues modulo $31$
\[
 1,30,1,30,5,22,7,10,2,12,29,28,17,28.
\]
Kummer's criterion makes $31$ regular. If
$q(u)=(u^{30}-1)/31\pmod{31}$, direct integer arithmetic gives
\[
 q(5^5(-6)^{-6})=0,\qquad B_{28}=14\pmod{31},\qquad
 q(5^5(-6)^{-6})^3-30B_{28}=14\pmod{31}.
\]
The quotient is therefore tame. The required odd unit eigenspace
in \citet[Theorem~1.2 and the following criterion]{McCallumTzermias2004}
vanishes because $31\nmid B_{14}$.
\citet[Theorems~4.16--4.17 and Corollary~4.18, author
pp.~35--36]{ShuExpanded2021} now gives
$\dim_{\mathbb F_{31}}\Sel_\lambda(J/K)=7$.
Here $\lambda=\zeta_{31}-1$; replacing it by
$\zeta_{31}-\zeta_{31}^{-1}$ changes it by a unit.
The tame parameter is zero, for which that source explicitly uses
$(0/31)=-1$ in its dimension formula.

Write $R=\operatorname{rank}J(\Q)$. Semilinear Galois descent
on the $K$-vector space $J(K)\otimes\Q$ gives dimension $R$
over $K$. Its finite torsion contributes one dimension modulo
$\lambda$, since $J[\lambda]$ has order $31$ and is rational
over $K$. Thus the Selmer exact sequence gives
\[
 7=R+1+\dim_{\mathbb F_{31}}\Sha(J/K)[\lambda].
\]
\citet[Theorems~1.2 and 1.4; Proposition~5.1]{McCallumTzermias2004}
proves that $\Sha[\lambda^3]$ is free over
$\mathcal O_K/(\lambda^3)$ and that its pairing with
$\Sha[\lambda]$ is nonzero. The latter group cannot have
dimension one: if $a$ generated the former, $\lambda^2a$
would generate the latter, whereas the Rosati adjoint satisfies
$\widehat\lambda^2\equiv\lambda^2\pmod{\lambda^3}$.
Functoriality and alternation would then give
\[
 [a,\lambda^2a]=[\widehat\lambda^2a,a]
              =[\lambda^2a,a]=-[a,\lambda^2a],
\]
forcing the pairing to vanish at the odd prime $31$.
Consequently $R\le4$. Since $A(\Q)$ has positive rank,
$\operatorname{rank}P(\Q)\le3$; its $\Q(\zeta_3)$-action
makes that rank even, proving the assertion.
\end{proof}

\begin{proposition}[The exact CM class numbers]
\label{prop:thirtyone-cm-class-numbers}
The fields $F_0$ and $M$ have discriminants $31^4$ and $-31^9$,
and class numbers
\[
 h(F_0)=1,\qquad h(M)=3.
\]
The type norm for $T=\{0,1,2,4,8\}$ acts trivially on
$\operatorname{Cl}(M)$.
\end{proposition}
\begin{proof}
The conductor--discriminant formula gives the displayed
discriminants, since every nontrivial character has conductor $31$.
For $R_0=\{\pm1,\pm5,\pm25\}$, put
$\eta=\sum_{a\in R_0}\zeta_{31}^a$.
In the basis of the five Gaussian periods indexed by
$3^jR_0$, cyclic convolution gives the multiplication matrix
\[
 \begin{pmatrix}
 -4&0&0&2&1\\ -6&1&2&1&2\\ -6&2&2&1&1\\
 -4&1&1&0&2\\ -5&2&1&2&0
 \end{pmatrix}.
\]
Its characteristic polynomial is
$t^5+t^4-12t^3-21t^2+t+5$, so $N(\eta)=-5$.
The Minkowski bound is
$5!\,31^2/5^5=115320/3125<37$.
The only prime ideals below that bound lie over $2,5,31$:
the first has norm $2^5=32$ and is the principal ideal $(2)$;
the five primes over $5$ are generated by the conjugates of
$\eta$; the prime over $31$ is generated by
$N_{K/F_0}(1-\zeta_{31})$. This exhausts the bounded ideals,
and proves $h(F_0)=1$.

The roots of unity in $M$ are $\{\pm1\}$. For a unit $u$,
Kronecker's theorem makes $u/\bar u$ a root of unity.
If it were $-1$, then $M=F_0(u)$ with $u^2$ a negative unit.
That quadratic extension would have discriminant dividing $4$,
contrary to ramification at $31$. Thus every unit is real and
the Hasse unit index is one. Let $z$ be a primitive tenth root
of unity and use the character taking $3\bmod31$ to $z$.
Reduction modulo $\Phi_{10}(z)$ gives the exact identity
\[
 \sum_{j=0}^{29}(3^j\bmod31)z^j=-62z.
\]
Hence the four primitive odd order-ten characters have
Bernoulli product $16$. The remaining odd quadratic character has
\[
 B_{1,\chi}=\frac1{31}\sum_{j=0}^{29}(3^j\bmod31)(-1)^j=-3.
\]
The analytic class-number formula, applied to the Dirichlet
factorization of $\zeta_M/\zeta_{F_0}$, therefore gives
\[
 \frac{h(M)}{h(F_0)}=\frac{2}{2^5}
       \prod_{\chi\ \mathrm{odd}}(-B_{1,\chi})
       =\frac{3\cdot16}{16}=3.
\]
Finally complex conjugation acts by inversion on $\operatorname{Cl}(M)$,
since $h(F_0)=1$. A generator $\sigma$ of its cyclic Galois
group of order ten consequently acts as $-1$ on this cyclic
class group of order three, because $\sigma^5$ is conjugation.
The type has four even and one odd exponent, so its norm acts
by $4-1=3$, and is zero on the class group.
\end{proof}

\begin{proposition}[The central vanishing and the cubic-character restriction]
\label{prop:thirtyone-central-obstruction}
The whole-Jacobian function has root number $+1$ and
\[
 L(J/\Q,1)=L'(J/\Q,1)=L(A/\Q,1)=0.
\]
If $\chi$ is a cubic Dirichlet character modulo $31$, then,
over a coefficient field containing its values,
\[
 V(J)=V(A)\oplus V(A)\otimes\chi\oplus V(A)\otimes\chi^2.
\]
The two complementary weight-two Hilbert constituents over $F_0$
have nontrivial cubic central characters. Neither the rank-zero
assertion for $P$ nor finiteness of its Tate--Shafarevich group
is established by this calculation.
\end{proposition}
\begin{proof}
The restriction of $\chi$ to $H$ is nontrivial, since
$5=3^{20}\pmod{31}$. If $\varphi$ is the algebraic CM Hecke
character of $M$ associated with $A$, and
$\nu=\chi\circ N_{M/\Q}$, then $\nu$ cuts out $K/M$.
The three characters of this cyclic cubic extension give
\[
 L(J/\Q,s)=L_M(\varphi,s)L_M(\varphi\nu,s)
                            L_M(\varphi\nu^2,s),\qquad
 L(A/\Q,s)=L_M(\varphi,s).
\]
Indeed the three-dimensional multiplicity space over $G_M$
has these three characters, and $c$-invariants select the
untwisted constituent. Inducing to $G_\Q$ proves the asserted
representation identity.

Quadratic induction from $M/F_0$ gives a weight-two Hilbert
form $\pi$. Its normalized central character is trivial:
real multiplication on $A_{F_0}$ and its polarization give
cyclotomic determinant on every two-dimensional constituent.
As in Proposition~\ref{prop:lm-nineteen-rankone},
the CM construction, the attached-variety statement and Faltings'
isogeny theorem identify the abelian variety to which
\citet[Theorem~A, p.~28]{Zhang2001} applies.
The real degrees are $[F_0:\Q]=5$ and $[K^+:\Q]=15$, both
odd. A nonzero central value for $A/\Q$, respectively $J/\Q$,
would therefore force rank zero over these real fields.
The rational infinite-order points already proved rule this
out. The tame root-number formula
of \citet[Theorem~3.13 and the tame case in \S5]{ShuExpanded2021}
gives $W(J/\Q)=(2/31)=+1$.
Its functional equation then makes the central order even,
so $L'(J/\Q,1)=0$ as well.

Put $\chi_{F_0}=\chi\circ N_{F_0/\Q}$.
The two complementary forms $\pi\otimes\chi_{F_0}$ and
$\pi\otimes\chi_{F_0}^2$ have central characters
$\chi_{F_0}^2$ and $\chi_{F_0}$, respectively.
They are nontrivial because $[F_0:\Q]=5$ and the character
cuts out a cubic extension. The trivial-central-character
hypothesis in Zhang's theorem therefore excludes these direct
applications. Adjoining the real cubic field kills the twists,
but makes $P$ isogenous to two copies of the positive-rank $A$
after base change; it does not produce a rank-zero application.
\end{proof}

For comparison with another possible arithmetic input, use the
September 2025 author version of \citet{Hsieh2014}, whose
Theorem~2 and Corollaries~2--3 are on pp.~3--4.
Its inclusion $(F_{\psi,\Sigma})\subset(L_{\psi,\Sigma})$
is a lower-bound divisibility for the Selmer module, so it does
not bound that module above from a nonzero central value.
Corollary~2 supplies equality for an anticyclotomic finite branch,
but that condition fails here already at $31$.
Indeed $\varphi^c\varphi$ is the norm character and
$\nu^c=\nu$, so
$(\varphi\nu)^c(\varphi\nu)=\operatorname{Norm}\,\nu^2$.
For any permitted coefficient prime different from $31$ and $3$,
the cyclotomic character is unramified at $31$, whereas $\nu^2$
has nontrivial tame cubic inertia. This obstruction survives
passage to the finite branch and dual or Tate twists.
Corollary~3 instead requires a CM type induced from an imaginary
quadratic field. Although $M$ contains $\Q(\sqrt{-31})$, the
primitive type $T$ contains both parities and is not so induced.
These are failures of the specified hypotheses, not a claim that
no other arithmetic theorem could apply.

The nonvanishing of $L(P/\Q,1)$ remains unproved here.
Even such a certificate would still require an applicable Selmer
upper-bound or descent theorem for the cubic complementary
constituents. The present rank bound leaves
$\operatorname{rank}P(\Q)=0$ or $2$, and no exponent-thirty-one
proper finite-descent conclusion follows from this calculation.
The all-prime quotient section classification at $31$ follows
separately from Theorem~\ref{thm:regular-solvable-intro}.

\subsection{The cyclic quotient at thirty-seven and its remaining arithmetic question}
\label{subsec:thirtyseven-gate}

The order-three symmetry of a cyclic Fermat quotient at thirty-seven
produces a rational sixfold, but that sixfold has positive rank.
The complementary twelvefold is the remaining rank-zero candidate
within this particular cyclic quotient.  The following calculation
does not establish a new section classification.

Let $C/\mathbb Q$ be the smooth proper model of
\[
 y^{37}=x(1-x)^{10},\qquad J=\operatorname{Jac}(C),
\]
and put
\[
 c(x,y)=\left(\frac1{1-x},\frac{y^{10}}{(1-x)^3}\right),\quad
 B=C/\langle c\rangle,\quad A=\operatorname{Jac}(B),\quad
 P=\operatorname{im}(2-c-c^2)\subset J.
\]
Let $R_0,R_1,R_\infty$ be the unique points above $0,1,\infty$.
Write
\[
 K=\mathbb Q(\zeta_{37}),\quad H=\langle10\rangle=\{1,10,26\},
 \quad M=K^H,\quad F_0=M^+.
\]

\begin{proposition}[The rational factors at thirty-seven]
\label{prop:thirtyseven-factors}
There is a rational isogeny $J\sim_{\mathbb Q}A\times P$, with
$\dim A=6$, $\dim P=12$, and
\[
 \operatorname{End}_{\mathbb Q}^0(J)
 =\mathbb Q\times\mathbb Q(\zeta_3),\qquad
 \operatorname{End}_{\mathbb Q}^0(A)=\mathbb Q,\qquad
 \operatorname{End}_{\mathbb Q}^0(P)=\mathbb Q(\zeta_3).
\]
Both factors are rationally simple, and
\[
 \operatorname{rank}A(\mathbb Q)\geq1,\qquad
 \operatorname{rank}P(\mathbb Q)\in2\mathbb Z_{\geq0},\qquad
 P(\mathbb Q)_{\mathrm{tors}}\simeq\mathbb Z/37\mathbb Z.
\]
Thus $P$ is the only possible positive-dimensional rational rank-zero
isogeny factor of this $J$.
\end{proposition}
\begin{proof}
The degree-thirty-seven map to the $x$-line is fully ramified at
$0,1,\infty$, so $g(C)=18$.  The identity
$10^2+10+1=3\cdot37$ gives, by direct substitution,
\[
 c^2(x,y)=\left(\frac{x-1}{x},
       -\frac{y^{26}}{x(1-x)^7}\right),\qquad c^3=1.
\]
The three branch points are cyclically permuted.  The other fixed
points have $x^2-x+1=0$.  For $\rho+\rho^{-1}=1$, the curve and
fixed-point equations become $y^{37}=y^9=-1$, whose unique common
solution is $y=-1$.  Thus there are exactly two fixed points,
$(\rho,-1)$ and $(\rho^{-1},-1)$, and
\[
 34=3(2g(B)-2)+4.
\]
This proves $g(B)=6$.  Norm and pullback identify the image of
$(1+c+c^2)/3$ with $A$ up to isogeny and give the stated complement.

The Jacobi-sum CM type on $K$ is
\[
 S=\{1,2,3,4,5,6,8,9,10,12,13,15,16,19,20,23,26,30\};
\]
equivalently, $b\in S$ if $b+(10b\bmod37)<37$.
Its multiplicative stabilizer is exactly $H$.  With generator
$2\bmod37$, the descended type on the cyclic degree-twelve field
$M$ is
\[
 T=\{0,1,2,3,4,11\}\subset\mathbb Z/12\mathbb Z,
\]
whose translation stabilizer is trivial: a preserving translation
must belong to $T$, and each of $1,2,3,4,11$ sends an element of
$T$ outside $T$. The CM decomposition of
\citet[Lemma~1.6 and the following remark, pp.~206--207]{GrossRohrlich1978}
therefore gives $J_{\overline{\mathbb Q}}\sim D^3$, with $D$
absolutely simple of dimension six and
$\operatorname{End}_{\overline{\mathbb Q}}^0(D)=M$.
If $\tau(x,y)=(x,\zeta_{37}y)$, then $c\tau c^{-1}=\tau^{10}$.
The crossed product generated by $K$ and $c$ is
\[
 K\rtimes H=\operatorname{End}_M(K)=M_3(M)
          =\operatorname{End}_{\overline{\mathbb Q}}^0(J).
\]
Galois acts on the unique expansion $\sum_{j=0}^2a_jc^j$ through
its usual action on the $a_j\in K$ and fixes $c$.  The rational
fixed algebra is $\mathbb Q[c]=\mathbb Q\times\mathbb Q(\zeta_3)$.
Its central factors prove the endomorphism assertions and rational
simplicity.  The same argument over $F_0$ gives $F_0[c]$, so
$A_{F_0}$ has real multiplication by $F_0$.

On $X^{37}+Y^{37}=1$, the quotient map is
$(X,Y)\mapsto(X^{37},XY^{10})$.  It takes
$(\rho,\rho^{-1})$ to $(\rho,-1)$.
\citet[Theorem~2.1, pp.~208--209]{GrossRohrlich1978} applies because
$10\notin\{1,18,35\}$, and proves that
\[
 D_0=[(\rho,-1)+(\rho^{-1},-1)-2R_\infty]\in J(\mathbb Q)
\]
has infinite order.  Changing the normalizing branch cusp changes
this class only by $37$-torsion.  The two displayed points are fixed
by $c$, whereas differences of branch points are killed by $37$.
Consequently, for $\pi:C\to B$,
\[
 \pi^*\pi_*D_0=(1+c+c^2)D_0=3D_0+\text{torsion}
\]
is non-torsion.  Thus $A(\mathbb Q)$ has positive rank.
The $\mathbb Q(\zeta_3)$-action makes the rank of $P(\mathbb Q)$ even.

By \citet[Theorem~1.1, p.~202]{GrossRohrlich1978},
$J(\mathbb Q)_{\mathrm{tors}}\simeq C_{37}$.
Put $D=[R_0-R_\infty]$ and $E=[R_1-R_\infty]$.
The divisor of $x$ gives $37D=0$, and $D\ne0$, since a principal
divisor $R_0-R_\infty$ would give a degree-one map to the projective
line.  The divisor of $y$ gives $D+10E=0$, so
$cD=E-D=10D$ and $c^2D=26D$.
Hence $(2-c-c^2)D=3D$, which generates the order-thirty-seven
torsion and belongs to $P$.  The inclusion $P\subset J$ proves
the torsion assertion.
\end{proof}

For an auxiliary prime $\ell\ne37$ and an abelian variety $B'$,
write $V(B')=T_\ell B'\otimes_{\mathbb Z_\ell}\mathbb Q_\ell$.
Write $N$ for the arithmetic conductor and $W$ for the global root number.

\begin{proposition}[Conductor, signs, and the order-nine character]
\label{prop:thirtyseven-character}
The conductors and root numbers are
\[
 N(J)=37^{36},\quad N(A)=37^{12},\quad N(P)=37^{24},
 \qquad W(J)=W(A)=-1,\quad W(P)=+1.
\]
Let $\chi$ be the character modulo $37$ with $\chi(2)=\zeta_9$.
Over a coefficient field containing its values,
\[
 V(J)\simeq V(A)\oplus V(A)\otimes\chi
                   \oplus V(A)\otimes\chi^2.
\]
The two complementary weight-two Hilbert constituents over the real
sextic field $F_0$ have nontrivial cubic central characters.
\end{proposition}
\begin{proof}
In the parameters $(r,s,t,\delta)=(1,10,26,1)$,
\[
 10^{10}(-11)^{26}\equiv581\pmod{37^2},\qquad
 581^{36}\equiv1\pmod{37^2}.
\]
Thus \citet[Proposition~3.12, p.~16; Theorem~3.13,
pp.~17--18]{ShuExpanded2021} applies in the tame case.
The Jacobi-sum character on $K$ has conductor the prime above $37$
and nontrivial quadratic tame inertia.  Its induction to $\mathbb Q$
is tame and has no inertia invariants, so its conductor exponent is
$36$.  Each rational direct summand is also tame and has no inertia
invariants.  Its exponent is consequently its representation
dimension, giving $12$ and $24$.  The same root-number formula gives
$W(J)=(2/37)=-1$.

The equalities $26=2^{12}$ and $10=2^{24}\pmod{37}$ show that every
cubic character modulo $37$ is trivial on $H$.  The chosen
order-nine $\chi$ instead restricts nontrivially to $H$.
Let $\varphi$ be the CM Hecke character of $M$ associated with $A$
in one fixed embedding, and put $\nu=\chi\circ N_{M/\mathbb Q}$.
The latter cuts out $K/M$.  The $c$-invariant line extends the
Jacobi-sum character from $G_K$ to $G_M$; the three extensions are
$\varphi,\varphi\nu,\varphi\nu^2$.  Transitivity of induction gives
the asserted representation identity and
\[
 L(P/\mathbb Q,s)=L_M(\varphi\nu,s)L_M(\varphi\nu^2,s),
 \qquad L(A/\mathbb Q,s)=L_M(\varphi,s).
\]
Since $\chi^3|_{G_M}=1$, it is a self-twist of $V(A)$.
Thus $V(A)\otimes\chi^{-1}\simeq V(A)\otimes\chi^2$.
Polarization makes the two normalized complementary constituents a
dual pair.  Their completed functional equations have reciprocal
root numbers, so $W(P)=+1$ and $W(A)=W(J)=-1$.

Quadratic induction from $M/F_0$ gives a weight-two Hilbert form
$\pi$ with trivial normalized central character.  Indeed the real
multiplication on $A_{F_0}$ is fixed by Rosati, so every
two-dimensional constituent has cyclotomic determinant.
The character $\theta=\chi\circ N_{F_0/\mathbb Q}$ has order three,
since $G_{F_0}$ has image $\langle2^6\rangle$ in
$\operatorname{Gal}(K/\mathbb Q)$.  The complementary forms
$\pi\otimes\theta$ and $\pi\otimes\theta^2$ therefore have central
characters $\theta^2$ and $\theta$, respectively.
\end{proof}

The arithmetic conclusion needed for a rank-zero finite-descent
application remains unproved.  In particular, no assertion is made
that $L(P/\mathbb Q,1)\ne0$, that $\operatorname{rank}P(\mathbb Q)=0$,
or that $\Sha(P/\mathbb Q)$ is finite.  The following applicability
restrictions explain why the previous arithmetic inputs do not
immediately settle it.

The direct central-value statements of
\citet[Theorem~A, p.~28]{Zhang2001} and
\citet[Theorems~B and~B', pp.~593--594]{Nekovar2012} require trivial
central character.  The two complementary forms violate this
condition.  Nekov\'a\v r's more general Theorem~A permits a
compensating character over an auxiliary CM extension, subject to
equation~(0.3.1), but this changes the Selmer representation.
Taking that character trivial again requires trivial central
character.  Killing the cubic twists by passing to $K^+$ makes
$P$ isogenous to $A^2$ and imports its positive-rank point.

The quadratic-Hilbert construction of \citet{LoefflerZerbes2025},
including the rational refinement used above, does not give a theorem
for this real sextic base.  The primitive type cannot be reduced to
the quartic CM subfield: its counts above the four embeddings are
$(2,1,1,2)$, whereas an induced type would have counts zero or three.
There is no imaginary quadratic subfield of the cyclic degree-twelve
field $M$; its unique quadratic subfield is real.

In the September 2025 author version of \citet{Hsieh2014},
Theorem~2 and Corollaries~2--3, pp.~3--4, the general inclusion
$(F_{\psi,\Sigma})\subset(L_{\psi,\Sigma})$ has the lower-bound
direction for the Selmer module.  The anticyclotomic equality case
fails here at $37$: complex conjugation fixes $\nu$, and
\[
 (\varphi\nu)^c(\varphi\nu)=\operatorname{Norm}\,\nu^2.
\]
The nontrivial cubic tame inertia of $\nu^2$ cannot be removed by
cyclotomic or Tate twists at any auxiliary coefficient prime other
than $3$ and $37$.  The equality case for a type induced from an
imaginary quadratic field is also unavailable.  The contemporary
completion of the Eisenstein-congruence argument in
\citet[Theorem~1.6 and Corollary~1.8, pp.~4--5]{BurungaleHeTianYe2026}
preserves this divisibility direction.

Finally, $37$ is irregular, with irregular index $32$ and
$B_{32}=-7709321041217/510$.  Thus the regular-prime Selmer
equalities in \citet[\S4.4.2]{ShuExpanded2021} do not apply.
These are failures of the specified hypotheses, not a claim that no
other arithmetic theorem could apply.

The image of the Fermat curve in $P$ is nonconstant and generates
$P$: the cyclic quotient and the projection above induce a surjection
of Jacobians.  It is therefore not a torsion translate of an elliptic
subvariety.  Proving rank zero and whole-$\Sha$ finiteness for $P$
would supply a proper finite-descent input.  The earlier split-prime
criterion needs only rank zero together with vanishing of
$T_\ell\Sha(P/\mathbb Q)$ for all but finitely many primes completely
split in a fixed finite Galois extension.  Neither conclusion follows
from the present structural calculation.

\section{Frey representations in solvable quotients}
\subsection{A prosolvable Frey-three criterion}
\label{subsec:frey-three}
Let $n\ge3$ be odd. Identify $X_n$ with
$F_n:X^n+Y^n=Z^n$ and put $U=F_n-\{XYZ=0\}$.
For a point-or-cusp Selmer section $s$ of the geometrically
prosolvable quotient, use the Legendre family
$E_\lambda:y^2=x(x-1)(x-\lambda)$, where $\lambda=(X/Z)^n$.
All local statements below concern representatives of this section;
no lift to the full arithmetic fundamental group is assumed.
Write $\Pi_U^{\mathrm{sol}}$ for its arithmetic quotient.
\par\medskip\noindent\textit{The representation on the stated quotient.}\quad 

The geometric determinant of the Tate-module representation is trivial.
Every finite geometric image on \(E[3^m]\) is therefore a subgroup of
\(\mathrm{SL}_2(\mathbb Z/3^m)\). This group is solvable: its congruence
kernel over \(\mathbb F_3\) is a finite 3-group, while
\(\mathrm{SL}_2(\mathbb F_3)\) has a normal quaternion subgroup with
cyclic quotient and is solvable. Thus the representation factors through
\(\Pi_U^{\mathrm{sol}}\) at every level, and hence on \(T_3E\).

A quotient section \(s\) consequently supplies

\[
\rho_s:G_{\mathbb Q}\longrightarrow
\mathrm{GL}_2(\mathbb Z_3),\qquad \det\rho_s=\chi_3.
\]

No section of the full arithmetic fundamental group is assumed or
constructed. The action on a finite cyclic-subgroup scheme below is
likewise already defined on the prosolvable quotient.

\par\medskip\noindent\textit{The topology of solvable covers.}\quad 

\begin{lemma}[The full prosolvable topology on a solvable cover]\label{lem:solvable-open-topology} Let \(\Delta\) be a profinite group and
\(H\subseteq\Delta\) an open subgroup such that
\(\Delta/\operatorname{core}_\Delta(H)\) is solvable. If \(R_\Delta\)
and \(R_H\) are the respective kernels of the maximal prosolvable
quotients, then

\[
R_H=R_\Delta\subseteq H.
\]

In particular \(H^{\mathrm{sol}}\) is the open subgroup \(H/R_\Delta\)
of \(\Delta^{\mathrm{sol}}\).

\end{lemma}\begin{proof} The solvable permutation quotient first gives
\(R_\Delta\subseteq H\). Let \(K\triangleleft H\) be open with \(H/K\)
solvable. The core \(C=\bigcap_{\delta\in\Delta}\delta K\delta^{-1}\) is
a finite intersection because \(K\) is open and normal in \(H\). Set
\(L=\operatorname{core}_\Delta(H)\). Then

\[
L/C\hookrightarrow
\prod_{\delta\in\Delta/H}
 \delta H\delta^{-1}/\delta K\delta^{-1}.
\]

The right side is solvable, as is \(\Delta/L\). Hence \(\Delta/C\) is
solvable and \(R_\Delta\subseteq C\subseteq K\). Intersecting over \(K\)
proves \(R_\Delta\subseteq R_H\).

Conversely, for every open normal \(M\triangleleft\Delta\) with
\(\Delta/M\) solvable, the quotient \(H/(H\cap M)\) is solvable.
Therefore \(R_H\subseteq\bigcap_M(H\cap M)=R_\Delta\). This proves the
lemma.

\end{proof}

For a finite cover with solvable geometric Galois closure, the lemma
applies to geometric fundamental groups. Quotienting its full arithmetic
fundamental group by this common kernel gives exactly the corresponding
open subgroup of the base's geometrically prosolvable arithmetic group.
This is an identification of entire prosolvable quotients, not just
their particular finite quotients.

\par\medskip\noindent\textit{A finite sufficient condition.}\quad 

\begin{proposition}[A stable cyclic subgroup modulo nine]\label{prop:frey-three-mod-nine} If the representation \(\rho_s\bmod9\) has a
\(G_{\mathbb Q}\)-stable cyclic subgroup of order 9, then \(s\) is
globally cuspidal.
\end{proposition}

The hypothesis is equivalent to a stable rank-one direct summand of
\((\mathbb Z/9)^2\). It is stronger than a stable line modulo 3. For
example, the inverse image of a Borel modulo 3 in
\(\mathrm{GL}_2(\mathbb Z/9)\) preserves the residual line, but the
congruence element \(1+3E_{21}\) moves each of its lifts. This is a
finite arithmetic sufficient condition, not a reduction of the
difficulty of proving mod-3 reducibility.

\begin{proof}\par\medskip\noindent\textit{The correct neighborhood.}\quad 

Let \(\mathcal S\) be the finite étale scheme over \(U\) parametrizing
cyclic order-9 subgroups of \(E[9]\). Choose the geometric point \(C\)
fixed by the pulled-back action of \(s(G_{\mathbb Q})\). Its stabilizer
is an open subgroup of \(\Pi_U^{\mathrm{sol}}\) containing
\(s(G_{\mathbb Q})\).

Take the corresponding component \(V\to U\). It is defined over
\(\mathbb Q\) and geometrically connected: the arithmetic orbit of \(C\)
is its geometric orbit, since the arithmetic quotient is generated by
its geometric subgroup together with the section and the latter fixes
\(C\). Equivalently, the inverse-image stabilizer in \(\pi_1(U)\)
surjects onto \(G_{\mathbb Q}\).

The geometric permutation group of \(\mathcal S\) is a quotient of a
subgroup of \(\mathrm{SL}_2(\mathbb Z/9)\), so its Galois closure is
solvable. Lemma~\ref{lem:solvable-open-topology} identifies the open arithmetic stabilizer with
\(\Pi_V^{\mathrm{sol}}\). The given section therefore lifts to a section
of that precise group. A stable subgroup in the pulled-back
representation has thus been converted into a compatible lift to the
level curve; this conversion is justified by a stabilizer argument, not
by assuming a full-section lift.

The lifted section is locally point-or-cusp geometric. For an interior
local point, fixedness of the corresponding finite fibre gives a
rational point of the cover. At a boundary point, the same statement
uses the finite set of prolongations and their stabilizers. One can also
lift the local cusp section before passing to the finite cover: cusp
inertia injects into the prosolvable quotient here, since each puncture
has a nonzero primitive class in the geometric abelianization and the
curve has more than one puncture. Its intersection with the stabilizer
projects onto the local Galois group and therefore gives a rational
boundary lift. This argument does not require that the entire cusp
decomposition group be contained in the neighborhood.

\par\medskip\noindent\textit{The map to the modular curve.}\quad 

Over \(V\), the elliptic curve has its original full rational 2-torsion
and the tautological cyclic subgroup \(C_9\). Choose independent nonzero
2-torsion sections \(P,Q\). Let

\[
\phi:E\longrightarrow E'=E/\langle P\rangle
\]

be the 2-isogeny, with dual \(\widehat\phi\). The subgroup

\[
C_4=\widehat\phi^{-1}(\langle Q\rangle)
\]

is cyclic of order 4. Indeed, for \(\widehat\phi(R)=Q\) one has
\(2R=\phi(Q)\ne0\), so \(R\) has order 4. The image \(C'_9=\phi(C_9)\)
is cyclic of order 9. Their sum is a cyclic subgroup \(C_{36}\) of order
36 because their orders are coprime.

This gives an algebraic map \(V\to Y_0(36)\), and thus a map from the
smooth proper compactification \(\overline V\) to \(X_0(36)\). It is
nonconstant. If \(j(E')\) were constant, the degree-2 modular relation
would force \(j(E)\) algebraic over the constant field, contrary to the
nonconstant Legendre parameter.

The modular curve has genus one: the standard index is 72, the number of
cusps is 12, and both elliptic-point contributions vanish, giving
\(1+72/12-12/2=1\). \citet[Corollary~8.8, p.~385]{Stoll2007} proves
that every positive-genus \(X_0(N)\) over \(\mathbb Q\) is excellent for
finite abelian descent. Proposition~8.5, p.~384, of the same source consequently
applies to every proper curve mapping nonconstantly to \(X_0(36)\). Its
real-place proof correction is in \citet[\S3]{StollErrata2017} and does not
alter the assertion.

\par\medskip\noindent\textit{Cusp containment.}\quad 

Let \(W\to V\) run through finite neighborhoods of the lifted
prosolvable section. One can use a cofinal family of open characteristic
subgroups of the finitely generated geometric prosolvable group and
adjoin the section image. Their geometric Galois closures are solvable.
Lemma~\ref{lem:solvable-open-topology} applies again, so each carries its induced entire prosolvable
section. The compactification \(\overline W\) maps nonconstantly to
\(X_0(36)\); hence

\[
\overline W(\mathbb A_{\mathbb Q})_\bullet^{f\text{-ab}}
=\overline W(\mathbb Q).
\]

The adelic point attached to the lifted local sections survives finite
abelian descent. This assertion only uses the given prosolvable section:
any finite abelian torsor on \(\overline W\), pulled back to \(W\),
factors through that quotient, and the section supplies its common
global twist with local lifts. No non-solvable torsor or full
fundamental-group section is required.

Consequently every \(\overline W\) has a rational point compatible with
the chosen local section. To make compatibility explicit, fix one finite
place and a point or cusp representative of the local section of \(V\).
Its lifts through all section neighborhoods give compatible local
points. Excellence gives rational points with those localizations;
injectivity of rational-point localization then makes the rational
points compatible under every transition map.

Fix their image \(z\in\overline V(\mathbb Q)\). Above \(z\), the
compactified finite geometric fibres form a projective system of finite
sets. A rational lift is a section-fixed orbit at the corresponding
level. The compatible rational points thus determine a fixed
prolongation in the inverse-limit fibre. Its stabilizer is a
decomposition group. Hence the lifted section image is contained in that
decomposition group.

The direction is

\[
s(G_{\mathbb Q})\subseteq D_{\widetilde z},
\]

not the reverse; cusp inertia need not lie in a section image. This
finite-orbit argument remains valid at ramified boundary fibres.

Finally the image of \(z\) on \(F_n\) is a rational coordinate cusp,
by the classical Fermat theorem \citep[Theorem~0.5, p.~448]{Wiles1995}. Thus \(z\) belongs to the boundary of \(V\), and
its image section on \(U\) is cuspidal. This proves the proposition.\end{proof}

\begin{corollary}\label{cor:frey-three-irreducible} If \(\rho_s\) is reducible over \(\mathbb Q_3\),
then \(s\) is cuspidal.

\end{corollary}\begin{proof}
Indeed, an invariant \(\mathbb Q_3\)-line intersects the stable Tate
lattice in a saturated rank-one \(\mathbb Z_3\)-submodule. Its reduction
modulo 9 is a stable direct summand, so Proposition~\ref{prop:frey-three-mod-nine} applies. Therefore
every hypothetical noncuspidal section has \(\mathbb Q_3\)-irreducible
\(\rho_s\).
\end{proof}

\par\medskip\noindent\textit{Density-one integrality.}\quad 

The next argument explicitly handles disconnected algebraic monodromy
and representations that are irreducible over \(\mathbb Q_3\) but
reducible over its algebraic closure.

\begin{lemma}[The squared-trace zero set]\label{lem:frey-three-trace-density} Let \(H\) be a compact subgroup of
\(\mathrm{GL}_2(\mathbb Q_3)\), acting irreducibly over \(\mathbb Q_3\),
and suppose \(\det H\) is infinite. Then

\[
T=\{h\in H:\operatorname{tr}(h)^2=(1+\det h)^2\}
\]

has Haar measure zero.

\end{lemma}\begin{proof} Let \(G\) be the algebraic Zariski closure of \(H\), and
\(G^\circ\) its connected identity component. Its determinant is
\(\mathbb G_m\): the determinant of \(G\) is infinite, and a finite
component group cannot carry that dimension.

The unipotent radical of \(G^\circ\) is trivial. Otherwise a nontrivial
unipotent subgroup in dimension two has a unique fixed line, which is
preserved by all of \(G\), since the radical is characteristic in
\(G^\circ\). That line descends to \(\mathbb Q_3\) and contradicts
irreducibility.

If the connected reductive group \(G^\circ\) is not a torus, its derived
group contains \(\mathrm{SL}_2\). On every coset \(gG^\circ\), the trace
varies along \(g\mathrm{SL}_2\) while the determinant is constant.
Therefore the polynomial

\[
F(g)=\operatorname{tr}(g)^2-(1+\det g)^2
\]

cannot vanish identically on that component.

It remains to treat a torus, after extension to
\(\overline{\mathbb Q}_3\). If it acts by scalars, the scalar parameter
\(t\) gives

\[
F(tg)=t^2\operatorname{tr}(g)^2-(1+t^2\det g)^2,
\]

a nonzero polynomial. If it has two distinct weights, each \(g\in G\)
either preserves or exchanges their weight lines. In an exchanging
component the trace is zero, whereas the determinant varies, so \(F\)
again is not identically zero. In a preserving component let the
eigenvalues be \(a\alpha(t)\) and \(b\beta(t)\). Then

\[
F(gt)=-
\bigl(a^2\alpha(t)^2-1\bigr)
\bigl(b^2\beta(t)^2-1\bigr).
\]

Identical vanishing on the connected torus would force one weight to be
trivial. Since the determinant varies, precisely one weight is trivial.
Its weight space is the intrinsic invariant subspace of \(G^\circ\); it
is preserved by all of \(G\) and defined over \(\mathbb Q_3\). This
contradicts irreducibility. This last observation also covers a nonsplit
torus and the absolutely reducible case; no absolute-irreducibility
assumption has been made.

Thus \(F\) is nonzero on every algebraic component. Every open subgroup
of \(H\) has Zariski closure containing \(G^\circ\), since it has finite
index. Hence \(F\) cannot vanish identically on an analytic open subset
of any component of \(H\). A compact linear 3-adic group is a 3-adic
analytic group, and a nonzero analytic function on a chart has a zero
locus of Haar measure zero. The latter fact follows by induction on the
chart dimension, applying the one-variable zero theorem to fibres and
then Fubini. A finite chart cover proves the result. \citet[\S2.4, author pp.~4--5]{KhareRajan2001} use the same squared-trace hypersurface and
finite-level density mechanism in the open-image case; the component
analysis above supplies the extra cases needed here.

\end{proof}

\begin{proposition}[Density-one integrality]\label{prop:frey-three-density} If \(\rho_s\) is irreducible over
\(\mathbb Q_3\), the set of primes at which its associated proper local
point is cuspidal or reduces to the boundary of \(U\) has upper
Dirichlet density zero.

\end{proposition}\begin{proof} Discard the finite set of primes where the chosen
curve/family model is bad, together with 3. At a boundary or degenerate
place \(q\), the local Tate representation, possibly with a quadratic
twist, is upper triangular with diagonal characters
\(\delta_q\chi_3,\delta_q\). Thus every arithmetic Frobenius lift has
squared trace \((1+q)^2\). The cusp calculation has the same diagonal
form. These are the local calculations in \citet[\S4.3, author pp.~7--8]{Stix2015};
their validity is local and does not require a full global section. At the Legendre cusp at infinity, the quadratic twist may be ramified. No unramifiedness of that character is required: its square is one, so the squared-trace identity is unchanged, and the finite-level argument below allows its ramification.

For each \(m\), the representation modulo \(3^m\) has finite image and
is ramified at only finitely many primes. At a degenerate place outside
this finite exceptional set its Frobenius belongs to

\[
T_m=\{h\in \rho_s(G_{\mathbb Q}) :
 F(h)\equiv0\pmod {3^m}\}.
\]

Finite-level Chebotarev bounds the upper density of all degenerate
places by the Haar measure of \(T_m\). The closed sets \(T_m\) decrease
to the set \(T\) of Lemma~\ref{lem:frey-three-trace-density}, whose Haar measure is zero. Letting \(m\)
increase proves the proposition.
\end{proof}

The finite exceptional set is allowed to depend on \(m\). No finite
exceptional set for the entire 3-adic representation is used. Combining
Propositions~\ref{prop:frey-three-mod-nine} and~\ref{prop:frey-three-density} gives the proved reduction:

A noncuspidal prosolvable Selmer section has irreducible $\rho_s$
and is integral at a density-one set of primes. These are the modular
and density assertions used in Theorem~\ref{thm:prosolvable-frey-open-intro}.

Density-one integrality is not membership in
\(U(\mathbb A_{\mathbb Q})\), which would require integrality outside a
finite set.

\begin{remark}[The local parameter at seventeen]\label{rem:frey-three-seventeen-parameter}

Near the coordinate cusp \(\lambda=0\), put \(t=X/Z\). The Legendre
\(j\)-invariant and its Tate parameter give

\[
j=256\frac{(1-\lambda+\lambda^2)^3}
 {\lambda^2(1-\lambda)^2},
\qquad
q_E=2^{-8}t^{34}(1+O(t^{17})).
\]

Consequently, up to the fixed quadratic twist, the 3-adic extension
parameter for a tangential section is

\[
\kappa_3(2^{-8})+34\alpha_3.
\]

The coefficient 34 is a unit in \(\mathbb Z_3\). Thus the Fermat
exponent does not annihilate residual 3-adic tame monodromy. In the
usual exponent-prime Frey argument the corresponding divisibility kills
the residual inertia; that step has no analogue here.

At \(\mathbb Q_{17}\), the pro-3 completion of the unit group is zero,
since its finite part has order 16 and its principal-unit group is
pro-17. Thus \(H^1(\mathbb Q_{17},\mathbb Z_3(1))\cong\mathbb Z_3\) by
the valuation, and the constant term vanishes. Multiplication by 34 is
an automorphism. Local cuspidality at 17 therefore permits every 3-adic
Tate extension parameter; it gives a local invariant line, with no
global invariant line and no vanishing of local monodromy forced by this
calculation.

There is a concrete illustration of this distinction at the
representation level. The rational Legendre curve with \(\lambda=17\)
has full rational 2-torsion and split multiplicative reduction at 17.
Its 3-adic local representation there occurs in a cusp packet by the
preceding surjectivity. Yet its mod-3 representation has no global
stable line: its 3-division polynomial is

\[
3x^4-72x^3+102x^2-289.
\]

Modulo 5 its values at \(0,1,2,3,4\) are \(1,4,1,3,3\), so it has no
rational root, hence no rational cyclic 3-isogeny. This is not claimed
to arise from a Selmer section on \(U_{17}\); it only refutes the
implication from the local Tate shape at 17, determinant, and full
2-torsion alone to global residual reducibility.

Continuity gives finite ramification separately for each \(3^m\)-torsion
representation. It does not give a common finite ramification set for
the inverse limit. \citet[Theorem~1]{KhareRajan2001} prove density zero, rather than
finiteness, for ramification of semisimple \(p\)-adic representations.
\citet[Main Theorem and Remark~1, author p.~3]{KhareLarsenRamakrishna2005} construct, for coefficient primes at least 5,
infinitely ramified representations that are potentially semistable at
the coefficient prime and pure and rational at almost all unramified
primes. Their result is a counterexample to that general shortcut; its
stated prime restriction means it is not itself a 3-adic or
section-theoretic counterexample.

Nor does the finite-level modular curve associated only with a stable
line modulo 3 finish the argument. With full 2-torsion it gives level
12, whose modular curve has genus zero. The order-9 subgroup used above
gives level 36 and is a strictly stronger condition.

\end{remark}

\par\medskip\noindent\textit{The remaining sufficient condition.}\quad 

The sufficient conditions are global mod-9 Borel reduction, or
the stronger assumption of reducibility of the entire 3-adic
representation. Neither has been proved for an arbitrary prosolvable
Selmer section. The local theorem at seventeen by itself gives only a local line at
17. The next subsection proves that the remaining case has jointly open
two-adic and three-adic images. It has density-one integrality, but no
finite ramification set and no modularity theorem applicable solely from
the verified hypotheses.

This Frey-curve argument does not establish either sufficient condition
for an arbitrary prosolvable Selmer section. Exponent seventeen is
classified separately by the regular-prime theorem,
Theorem~\ref{thm:regular-solvable-intro}. Outside that theorem's range,
the simultaneous open-image case remains unresolved; no all-prime
finite-descent equality or global lifting theorem follows from the
Frey-curve argument alone.

\subsection{Simultaneous open Frey images}
\label{subsec:frey-open-images}
\paragraph{A rank-one weight obstruction.}

For $\ell=2,3$, the geometric image of the Legendre Tate module is
prosolvable: $\operatorname{SL}_2(\F_\ell)$ is solvable and the
successive congruence kernels are abelian. Thus both representations
are defined on the quotient under consideration.

\begin{lemma}[Common pure Frobenius polynomials]\label{lem:frey-common-frobenius}
For a noncuspidal section in Theorem~\ref{thm:prosolvable-frey-open-intro},
there is a Dirichlet-density-one set of rational primes at which both
Frey representations are unramified and have the same polynomial
$T^2-a_qT+q$, $a_q\in\Z$, with complex roots of modulus $q^{1/2}$.
\end{lemma}
\begin{proof}
Proposition~\ref{prop:frey-three-density} makes the local representative
interior with integral Legendre parameter away from the discriminant
on a density-one set, after removing the fixed finite bad-model set.
Both representations are then the Tate modules of the same elliptic
curve with good reduction at $q$. Its Frobenius polynomial and Weil
bound give the assertion.
\end{proof}

\begin{lemma}\label{lem:frey-rank-one-weights}
Let \(\psi:G_L\to E^\times\) be a continuous character of a number field,
where \(E/\Q_\ell\) is finite. If it is Hodge--Tate at all places over
\(\ell\), then it has finite ramification and is the \(\ell\)-adic
realization of an algebraic character of a Serre group
\(S_{\mathfrak m}\), defined over a number field.
It has realizations at every coefficient prime.
For some integer \(b\), at every place \(v\) outside its modulus and
different from the coefficient prime its algebraic Frobenius value has
absolute value \((Nv)^{b/2}\) under every complex embedding.
\end{lemma}
\begin{proof}
Tate's local theorem gives local algebraicity:
\cite[III \S1.2, III-7]{Serre1989}.
Finite ramification for abelian continuous representations, and the
existence of a modulus, are proved in
\cite[III \S2.2, III-10--11]{Serre1989}.
These statements are only being applied to a character.

There is a rationality qualification in the stated global theorem
\cite[III \S2.3, Theorem~2, III-13]{Serre1989}.
Its proof, III-14--15, first constructs a morphism from
\(S_{\mathfrak m}\) over the coefficient field, and uses rationality only
in the final descent to \(\Q\).
In the rank-one case that first morphism is a character after extending
coefficients to \(\overline\Q_\ell\).
Character groups of groups of multiplicative type do not acquire new
elements upon extending algebraically closed fields: the group-like
elements of a diagonalizable group's coordinate group algebra are its
character monomials. Hence this character descends through
\(\overline\Q\hookrightarrow\overline\Q_\ell\) to a finite number field.

Composition with the canonical representation
\(\varepsilon_{\ell'}:G_L^{\mathrm{ab}}\to
S_{\mathfrak m}(\Q_{\ell'})\) gives each companion. The common algebraic
Frobenius values are supplied by
\cite[II \S2.5, proposition, II-18--19]{Serre1989}.
The absolute-value formula is exactly
\cite[II \S3.4, Proposition~2, II-37--38]{Serre1989}.
\end{proof}

\begin{proposition}\label{prop:frey-weight-open-image}
Let \(\rho:G_\Q\to\operatorname{GL}_2(\Q_\ell)\) be continuous, with
\(\det\rho=\chi_\ell\), and Hodge--Tate at \(\ell\).
Assume that at a density-one set of unramified rational primes \(q\)
its polynomial is \(T^2-a_qT+q\), with \(a_q\in\Z\) and both complex
roots of absolute value \(q^{1/2}\).
If, for one \(r\ne\ell\),
\[
 (\rho|_{G_{\Q_r}})^{\mathrm{ss}}\simeq\delta\chi_\ell\oplus\delta
 \quad\text{with }\delta^2=1,
\]
then \(\rho\) has open image.
\end{proposition}
\begin{proof}
Let \(G\) be the algebraic monodromy group of \(\rho^{\mathrm{ss}}\).
Its connected component is reductive: the fixed space of its unipotent
radical is nonzero and stable under \(G\), and semisimplicity makes
the radical act trivially on every simple summand.
If \(G^\circ\) were a torus, after a finite Galois extension \(L/\Q\)
and coefficient extension its representation would be
\(\psi_1\oplus\psi_2\).
Enlarge \(L\) by a quadratic extension of \(\Q\) killing \(\delta\)
at \(r\). Such an extension is obtained by choosing a rational
representative of the local squareclass.
At each \(w\mid r\), the two local characters are now
\(\chi_\ell\) and \(1\), so both are unramified there.
Choose their moduli without these places.

Lemma~\ref{lem:frey-rank-one-weights} gives constant weights \(b_1,b_2\).
Choose a good rational prime \(q\) splitting completely in \(L\) and
avoiding the finite conductors. Such a prime exists by Chebotarev and
the density-one hypothesis.
Both character values at \(v\mid q\) are roots of
\(T^2-a_qT+q\). Under every complex embedding they remain roots of
this rational polynomial. Thus \(b_1=b_2=1\).
At \(w\mid r\), of residue degree \(f\), the local character values
are \(1\) and \(r^f\). Their absolute values contradict the requirement
that both be \(r^{f/2}\).

Therefore \(G^\circ\) is not a torus. A connected reductive nontorus
subgroup of \(\operatorname{GL}_2\) contains \(\operatorname{SL}_2\); the cyclotomic determinant
then gives \(G^\circ=\operatorname{GL}_2\).
The semisimplification is absolutely irreducible, so \(\rho\) itself
is absolutely irreducible.

For \(H=\rho(G_\Q)\), its Lie algebra \(\mathfrak h\subseteq
\mathfrak{gl}_2\) is \(\operatorname{Ad}(H)\)-stable and therefore
\(\operatorname{Ad}(\operatorname{GL}_2)\)-stable. It is an ideal with nonzero trace because the
determinant image is open. Such an ideal is the center or all of
\(\mathfrak{gl}_2\).
The center would make an open subgroup of \(H\) scalar, contradicting
its full algebraic monodromy. Thus \(\mathfrak h=\mathfrak{gl}_2\),
and \(H\) is open. This Lie argument also works at \(2\), by taking
logarithms on a sufficiently deep congruence subgroup.
\end{proof}

\paragraph{Local geometry and transfer between coefficient primes.}

\begin{lemma}\label{lem:frey-local-two}
For odd \(n\ge5\), every local point or cusp witness at \(2\) gives
\[
 (\rho_{s,3}|_{G_{\Q_2}})^{\mathrm{ss}}
 \simeq\delta\chi_3\oplus\delta,\qquad \delta^2=1.
\]
At either coefficient prime \(\ell=2,3\), the local Frey representation
is potentially semistable.
\end{lemma}
\begin{proof}
For an interior point at \(2\), normalize its projective coordinates
to be integral with at least one unit. Exactly two are units.
After an anharmonic substitution preserving \(j\), the parameter
has valuation \(nm\), with \(m\ge1\). Consequently
\[
 v_2(j)=v_2\!\left(
 2^8\frac{(\lambda^2-\lambda+1)^3}{\lambda^2(\lambda-1)^2}
 \right)=8-2nm<0.
\]
This calculation for prime exponent is already
\cite[Lemmas~5.11--5.12, pp.~21--22]{Steklov2026}; the same calculation
applies to every odd \(n\ge5\).
Negative \(j\)-valuation gives a Tate curve after at most a quadratic
extension \cite[IV Appendix A.1.1--A.1.2, IV-29--32]{Serre1989}.
Its local semisimplification has the asserted form. At a cusp the
Tate description gives the same conclusion.

At a coefficient prime an interior point gives the Tate module of a
local elliptic curve, which is potentially semistable.
A cusp representation is a finite quadratic twist of an extension of
\(\Q_\ell\) by \(\Q_\ell(1)\).
Every such extension is semistable, including at \(\ell=2\), by
\cite[Lemma~8.3.9, p.~125]{BrinonConrad2009}.
One can also see this by spanning \(H^1(K,\Q_\ell(1))\) with Kummer
classes of a uniformizer and of its products with a basis of principal
units. They are Tate-curve classes; sums and scalar multiples are
obtained through direct sums, subobjects and quotients.
\end{proof}

\begin{lemma}\label{lem:frey-companion-comparison}
Suppose continuous representations \(\rho_2\) and \(\rho_3\) of
\(G_\Q\), over \(\Q_2\) and \(\Q_3\), have equal rational Frobenius
polynomials at a density-one set \(\mathcal P\) of primes where both
are unramified. Assume \(\rho_2\) is Hodge--Tate at \(2\), and
\(\rho_3\) has open image. Then the connected algebraic monodromy
of \(\rho_2^{\mathrm{ss}}\) is not a torus.
\end{lemma}
\begin{proof}
If it were a torus, choose a finite Galois \(L/\Q\) and a finite
coefficient extension splitting
\(\rho_2^{\mathrm{ss}}|_{G_L}=\psi_{1,2}\oplus\psi_{2,2}\).
Lemma~\ref{lem:frey-rank-one-weights} supplies algebraic characters over a number
field \(F\) and their realizations \(\psi_{1,3},\psi_{2,3}\), for any
fixed embedding \(F\hookrightarrow\overline\Q_3\).

The prime ideals of \(L\) lying above \(q\in\mathcal P\) which split
completely in \(L\), after finitely many exclusions, have density one
among prime ideals of \(L\).
Indeed, prime ideals of residue degree at least two have a convergent
prime sum at \(s=1\); all unramified degree-one primes lie above
completely split rational primes; and the prime sum above
\(q\notin\mathcal P\) is at most \([L:\Q]\sum_{q\notin\mathcal P}q^{-s}\).
At this density-one set, the Frobenius polynomials of
\(\rho_3|_{G_L}\) and \(\psi_{1,3}\oplus\psi_{2,3}\) agree, because the
algebraic polynomial over \(F\) maps to the given rational polynomial
under the injective two-adic embedding. It therefore equals that
polynomial already in \(F[T]\), and is unchanged by the three-adic embedding.

Their traces agree on all of \(G_L\).
Otherwise, a finite congruence quotient detects a trace difference on
a conjugacy-invariant open subset, to which finite-level Chebotarev
assigns positive density. It must meet the preceding density-one set,
a contradiction.
This argument only discards the finite ramification set of the chosen
finite quotient; it assumes no finite ramification set for \(\rho_3\).
Brauer--Nesbitt gives
\[
 (\rho_3|_{G_L})^{\mathrm{ss}}\simeq\psi_{1,3}\oplus\psi_{2,3}.
\]
Restriction to \(G_L\) preserves openness of the compact image.
An open subgroup has full algebraic monodromy and acts absolutely
irreducibly, contradicting this decomposition.
\end{proof}

\begin{proof}[Proof of Theorem~\ref{thm:prosolvable-frey-open-intro}]
The first assertion is Proposition~\ref{prop:frey-three-mod-nine}.
For a noncuspidal section Lemma~\ref{lem:frey-common-frobenius} gives the common
density-one pure rational polynomials.
Lemmas~\ref{lem:frey-local-two} and Proposition~\ref{prop:frey-weight-open-image},
with \(\ell=3,r=2\), make \(\rho_{s,3}\) open.
Lemma~\ref{lem:frey-companion-comparison} then excludes toroidal connected
monodromy for \(\rho_{s,2}^{\mathrm{ss}}\).
The determinant and Lie-algebra argument in
Proposition~\ref{prop:frey-weight-open-image} makes \(\rho_{s,2}\) open as well.

Let \(H_\ell\) be the individual images and \(H\subseteq H_2\times H_3\)
the joint image. Compact Goursat gives a common profinite quotient.
The images in that quotient of open pro-2 and open pro-3 subgroups
are open. Their intersection is both pro-2 and pro-3, hence trivial.
The quotient is finite, so \(H\) contains a product of open subgroups.
After any finite field extension the individual images remain open,
hence absolutely irreducible with algebraic monodromy \(\operatorname{GL}_2\).
\end{proof}

\paragraph{The remaining arithmetic question.}

The theorem removes potentially abelian semisimplifications, including
reducible, finite-projective, toroidal and dihedral cases, from the
possible noncuspidal quotient sections. It does not prove finite
ramification or modularity in the remaining simultaneous open-image
case.

Khare--Larsen--Ramakrishna construct density-one compatible pure rational
pairs, potentially semistable at their coefficient primes, with infinite
ramification permitted
\cite[\S2.4, Corollary~14, manuscript pp.~23--24]{KhareLarsenRamakrishna2005}.
Their hypotheses require \(p>q\ge5\). This shows why two coefficients
alone do not give a general finite-ramification theorem; it is neither
a counterexample at coefficients \(2,3\) nor a section counterexample.
The usual finite exceptional set in a compatible system cannot be
introduced merely by renaming density-one compatibility.


\subsection{The small-prime hypothesis in the exponent-thirteen argument}
\label{subsec:thirteen-small-prime-gate}

Let $A_B$ be the genus-two Jacobian in the exponent-thirteen
decomposition, let $\chi$ be a cubic Dirichlet character modulo
thirteen, and put $\mathcal O=\mathbb Z_3[\zeta_3]$ with maximal
ideal $\mathfrak m=(1-\zeta_3)$.  The large split-prime argument
does not apply at three merely because three splits in the quartic
CM field.

\begin{proposition}[A small-prime hypothesis obstruction]
\label{prop:thirteen-small-prime-mr}
For $a=1,2$, the natural lattice
\[
 T=(T_3A_B\otimes_{\mathbb Z_3}\mathcal O)\otimes\chi^a
\]
does not satisfy hypothesis (H.4) of
\citet[\S3.5, p.~27]{MazurRubin2004}.  The same holds for its
dual and for every integral Tate twist of either lattice.
This is a limitation on that theorem's hypotheses, not a
nonvanishing statement for any Selmer group.
\end{proposition}
\begin{proof}
The residue field of $\mathcal O$ is $\mathbb F_3$, and the
cubic character becomes trivial modulo $\mathfrak m$.  Consequently
$T/\mathfrak mT\simeq A_B[3]$ as $\mathbb F_3[G_{\mathbb Q}]$-modules.
The principal polarization of the Jacobian gives a perfect
Galois-equivariant alternating Weil pairing
\[
 A_B[3]\times A_B[3]\longrightarrow\mu_3.
\]
It identifies this residual module with its Cartier dual and
therefore gives a nonzero element of
\[
 \operatorname{Hom}_{\mathbb F_3[G_{\mathbb Q}]}
       (T/\mathfrak mT,T^*[\mathfrak m]).
\]
Thus (H.4a) fails.  The alternative (H.4b) requires $p>4$, so
it fails at $p=3$ as well.

The mod-three cyclotomic character has order two.  Multiplying
a residual representation by any of its powers preserves
Cartier self-duality, since the square of the twisting character
is one.  Moreover the ordinary dual of $A_B[3]$ is its
cyclotomic twist.  This proves the assertion for duals and
integral Tate twists.
\end{proof}

The proposition concerns the displayed natural lattices.  It does
not assert that every stable lattice in the rational representation
has the same extension structure.  It does show why the verified
large-prime proof cannot simply be specialized to three.  A
different small-prime argument, or a theorem with different
hypotheses, would be needed to obtain the missing three-primary
Selmer input by this route.



\section{The remaining irregular-prime conditions}
\label{sec:irreg-gates}

The finite sets above do not by themselves identify the exact
centres in the cusp residue discs. We describe a sufficient
additional input for that local step, and then the distinct
restriction-map issue.

\subsection{A prescribed supported pairing}

In this subsection assume \(p>5\), fix a wild quotient \(C=C_{a,b,c}\),
and put \(A=\Jac(C)\), \(K=\Q(\zeta_p)\), and
\(\pi=1-\zeta_p\), acting on \(A/K\) by complex multiplication.
Let \(\mathfrak p\) be the unique prime of \(K\) above \(p\), and set
\begin{equation}
 \Lambda=
 \bigl(A(K_{\mathfrak p})/\pi A(K_{\mathfrak p})\bigr)^\Delta.
 \label{eq:irreg-Lambda}
\end{equation}
The local calculation in
\cite[Corollary 2 and Proposition 5, pp.~574--575]{McCallum1994}
shows that \(\Lambda\) is one-dimensional over \(\F_p\), and that
the natural map
\begin{equation}
 \chi:A(\Q_p)/pA(\Q_p)\longrightarrow\Lambda
 \label{eq:irreg-chi}
\end{equation}
is nonzero. Under the Faddeev descent map, the invariant line is
generated by the class of \(1+p\).

Consider the following additional condition:
\begin{equation}
 \begin{gathered}
 d_n\in H^1(K,A)[\pi^n],\qquad
 \pi d_{n+1}=d_n \quad(n\geq1),\\
 d_{n,v}=0\quad(v\neq\mathfrak p),\qquad
 d_{1,\mathfrak p}\text{ pairs nontrivially with }\Lambda.
 \end{gathered}
 \label{eq:irreg-prescribed-tower}
\end{equation}
The dual of \(\pi\) is \(\overline\pi=1-\zeta_p^{-1}\), a unit
multiple of \(\pi\), so the last pairing is the usual finite local
Tate pairing on \(A(K_{\mathfrak p})/\pi A(K_{\mathfrak p})\).
Condition \eqref{eq:irreg-prescribed-tower} is stronger than the
mere existence of a nonzero supported \(p\)-primary tower.

\begin{proposition}\label{prop:irreg-prescribed-isolation}
If a wild quotient satisfies
\eqref{eq:irreg-prescribed-tower}, then any compact \(p\)-Selmer
datum on \(F_p\) whose \(p\)-local centre lies in a coordinate cusp
residue disc has that cusp as its centre.
\end{proposition}

\begin{proof}
The proof of Lemma~\ref{lem:irreg-reciprocity} applies over the
number field \(K\), with support at \(\mathfrak p\). Restrict a
compact datum to \(K\), push it to \(A\), and apply that argument
to every \(d_n\), using a finite \(p\)-power level killing \(d_n\).
The point \(j(f(R_p))\in A(\Q_p)\) lies in the common annihilator
\[
 B_K=\{P\in A(K_{\mathfrak p}):
             \langle P,d_{n,\mathfrak p}\rangle=0
             \text{ for every }n\}.
\]

Compatibility of the tower implies
\[
 \langle\overline\pi P,d_{n+1,\mathfrak p}\rangle
       =\langle P,d_{n,\mathfrak p}\rangle.
\]
Hence \(B_K\) is stable under \(\overline\pi\) and saturated for
multiplication by \(\overline\pi\). It is stable under
\(\Z[\zeta_p]\), and \(\overline\pi^{p-1}\) is \(p\) times a unit
of this ring. It is therefore saturated for multiplication by
\(p\). Multiplication by an integer prime to \(p\) is invertible
on every pairing target, giving saturation for all nonzero
integers. The same compatibility shows that \(B_K\) contains
all torsion: use a sufficiently large power of \(\overline\pi\)
for the \(p\)-primary part, and the order for the prime-to-\(p\) part.

Put \(B=B_K\cap A(\Q_p)\). Since the last condition in
\eqref{eq:irreg-prescribed-tower} is a nonzero functional on the
one-dimensional space \(\Lambda\), one has \(B\subseteq\ker\chi\).
In particular \(\chi\) vanishes on torsion.
The quotient \(A(\Q_p)/A(\Q_p)_{\tors}\) is a free
\(\Z_p\)-module, and the image of \(B\) is a saturated submodule.
Consequently there is a saturated subgroup \(M\) containing \(B\),
of corank one modulo torsion, whose image modulo \(p\) is
\(\ker\chi\). To see this directly, choose a basis adapted to the
direct summand defined by \(B\), and lift the nonzero linear
functional \(\chi\) on the complementary quotient to a primitive
\(\Z_p\)-linear functional; take its kernel.

The integral local calculation of
\cite[Section 4, Propositions 6--7 and Theorem 4,
 pp.~577--580]{McCallum1994}
now applies to \(M\). It gives a nonzero holomorphic differential
on \(C\), annihilating \(M\) under the abelian logarithm, whose
reduction is nonzero at at least one of \(O_0,O_1,O_\infty\).
This part of the argument uses the wild model and the specified
reduction \(M\bmod p=\ker\chi\), not regularity of \(p\).
Its Coleman integral vanishes at every allowed centre and at the
three torsion cusp images.

At a cusp where the reduced differential is nonzero, the local
zero estimate in
\cite[Lemma 7, pp.~579--580]{McCallum1994}
shows that the integral has at most one zero in that residue disc.
Since the cusp is already a zero, it is the only one there.
The quotient map respects these coordinate residue classes
\cite[Lemma 8 and the proof of Theorem 1, p.~580]{McCallum1994},
and is injective on \(\Q_p\)-points by
Lemma~\ref{lem:irreg-quotient-free}. Thus one coordinate cusp disc
on \(F_p\) contains no other permitted centre.

Finally the rational automorphism
\[
 \tau(X:Y:Z)=(-Y:Z:X)
\]
cycles the three coordinate cusps. It carries compact data to
compact data: on cohomology, replace \(t\) by
\[
 \tau_*t+
 \kappa_p([\tau(Q_0)-Q_0]),
\]
and replace every \(R_v\) by \(\tau(R_v)\). Thus the same isolated
disc may be applied after either cyclic permutation. All three
coordinate cusp discs are isolated for every compact datum.
\end{proof}

The hypotheses \((\mathrm H1)\), \((\mathrm H2)\) in
\cite[p.~575]{McCallum1994} ask for supported classes at arbitrarily
large \(\pi\)-levels with the specified nontrivial pairing after
multiplication down to the first level. Condition
\eqref{eq:irreg-prescribed-tower} is a compatible version sufficient
for the compact Selmer argument above. A compatible lift of the
initial Kummer class corresponding to \(\zeta_p\) is one sufficient
way to construct it; equivalence with that particular initial
class is not asserted.

\subsection{The obstruction space does not decide the specified lift}

Let \(\mathcal U=\Spec\mathcal O_K[1/p]\). Over \(K\), choose the
usual isomorphism \(A[\pi]\simeq\mu_p\). The Kummer sequence gives
\begin{equation}
 H^2(\mathcal U,A[\pi])
 \simeq H^2(\mathcal U,\mu_p)
 \simeq\Cl(K)/p\Cl(K).
 \label{eq:irreg-obstruction-space}
\end{equation}
Here \(\mathfrak p\) is principal, so
\(\operatorname{Pic}(\mathcal U)=\Cl(K)\). Also
\(\operatorname{Br}(\mathcal U)=0\): a class in this Brauer group
has zero invariant away from \(\mathfrak p\), all infinite places
of \(K\) are complex, and global reciprocity forces its remaining
invariant to vanish. These are the calculations in
\cite[Section 5, Theorem 5, pp.~581--582]{McCallum1994} and
\cite[Lemma 6, pp.~401--402]{McCallum2001Brauer}.

The regular-prime lifting construction uses the vanishing of
\eqref{eq:irreg-obstruction-space} to obtain the successive
surjections in cohomology associated with
\begin{equation}
 0\longrightarrow A[\pi]\longrightarrow A[\pi^{n+1}]
   \xrightarrow{\pi}A[\pi^n]\longrightarrow0.
 \label{eq:irreg-lift-sequence}
\end{equation}
At \(p=37,59,67\), the group in
\eqref{eq:irreg-obstruction-space} has dimension one.
This identifies a possible obstruction space. It does not compute
the obstruction to lifting the specified initial class, and does
not prove that such a lift fails.

In fact the first local pairing direction itself is present:
\cite[Lemma 10, pp.~581--582]{McCallum1994}
identifies the initial class with \(\zeta_p\) and its pairing with
the invariant line with the nonzero Hilbert symbol of
\(\zeta_p\) and \(1+p\).
The unresolved issue in this approach is a compatible lift through
arbitrarily high levels preserving the required support and
pairing direction.

Proposition~\ref{prop:irreg-McCallum}(1) constructs some nonzero
supported tower under \eqref{eq:irreg-Mc-bound}. Its statement
does not prescribe the first-layer pairing in
\eqref{eq:irreg-prescribed-tower}, so it cannot be substituted for
that stronger hypothesis. Likewise, the final observation in
\cite[pp.~405--406]{McCallum2001Brauer} that an individual complex
multiplication eigendifferential has only the three cusp zeros
does not show that such an eigendifferential annihilates all the
Selmer centres. The differential used in
Theorem~\ref{thm:irreg-fixed} is not asserted to have that property.

\subsection{Failure of the regular local-source surjectivity}

\begin{proposition}\label{prop:irreg-source-failure}
Suppose the \(p\)-primary part of \(\Cl(K)\) is nonzero. Let
\(K_{\{p\}}(p)\) be the maximal pro-\(p\) extension of \(K\)
unramified outside \(\mathfrak p\). Then the decomposition map
\begin{equation}
 G_{K_{\mathfrak p}}\longrightarrow
 \Gal(K_{\{p\}}(p)/K)
 \label{eq:irreg-source-map}
\end{equation}
is not surjective.
\end{proposition}

\begin{proof}
The maximal unramified abelian \(p\)-extension \(H/K\) is nontrivial
by class field theory and is contained in \(K_{\{p\}}(p)\).
The prime \(\mathfrak p=(1-\zeta_p)\) is principal, so its Artin
class in \(\Gal(H/K)\) is trivial. Since \(H/K\) is unramified,
\(\mathfrak p\) splits completely in \(H\). Its decomposition
group therefore maps trivially to the nontrivial quotient
\(\Gal(H/K)\), proving the assertion.
\end{proof}

For \(p=37,59,67\), the maximal unramified abelian \(p\)-extension
has degree \(p\). It is stable under \(\Delta\) and hence Galois
over \(\Q\); a decomposition group at \(p\) in \(\Gal(H/\Q)\) has
order \(p-1\), while the full group has order \(p(p-1)\).
Retaining the full arithmetic quotient therefore does not restore
this particular surjectivity input.

Proposition~\ref{prop:irreg-source-failure} is a statement about the
Galois source. It neither constructs a kernel for localization on
a particular geometric Tate module nor disproves an integral
restriction theorem for that module. A replacement restriction
argument, if needed for a global classification, must be supplied
separately. Even exact identification of all proper local centres
would not by itself classify sections in the remaining affine
geometric kernel.



\section*{Reproducibility and attribution}
The mathematics, text, and programs were produced by GPT-6 Astra Pro
in ChatGPT Pro mode; the reasoning effort setting was not exposed.
Separate instances of the same model attempted selected principal
statements from their statements before comparing proofs. Other
passes were targeted audits of proposed arguments, and numerical
certificates were repeated with different parameters or methods.
These are same-model checks, not independent human verification.
Earlier global arguments were examined by Claude Opus~5.5 at maximal
effort on 25 November 2026, and selected earlier computations by that
model at high effort on 24 November 2026. Those historical checks
have their recorded scope and do not certify the extensions to
smaller geometric quotients proved here. No human has checked this
paper.

The programs \emph{Quintic finite-depth certificate}, \emph{Quintic class-number certificate},
\emph{Quintic residual-group certificate},
\emph{Quintic interval certificate},
\emph{Quintic rational lower bound},
\emph{Septic principal-prime certificate}, \emph{Septic class-number analytic certificate}, \emph{Septic auxiliary class-number certificate}, \emph{Septic residual-field polynomial check},
\emph{Septic central-value certificate},
\emph{Exponent-eleven central-value certificate},
\emph{Exponent-eleven rank-one derivative certificate},
\emph{Exponent-eleven rank-one coefficient enumerator},
\emph{Jacobi-sum coefficient enumerator},
\emph{Exponent-thirteen partial quotient checks},
\emph{Exponent-thirteen central-value certificate},
\emph{Exponent-thirteen coefficient enumerator},
\emph{Exponent-seventeen arithmetic checks},
\emph{Regular-prime seventeen residue certificate},
\emph{Nineteenth CM arithmetic checks},
\emph{Nineteenth central-derivative certificate},
\emph{Nineteenth canonical-model check},
\emph{Exponent-thirty-one arithmetic checks},
\emph{Cyclotomic relative-class-number checks},
\emph{First-case character and Bernoulli checks}, and
\emph{Auxiliary level-twenty-seven Kummer certificate}, together with their
outputs and exact build instructions, are included in the
\href{\workbenchurl}{workbench}. The exponent-thirteen programs verify
the point counts and analytic certificate for the explicit Prym.
Its split-prime Selmer consequence is proved separately in
Section~\ref{subsec:lm-lz-rational} from the cited Euler-system
constructions and the explicit representation-theoretic checks.
The companion \href{\recordurl}{record repository} records the project.
The mathematical dependence on the named published theorems is
explicit; those theorems are not claimed as new results. No priority
claim is made for the specializations or transfer arguments.
The original text, arguments, and programs are dedicated to the public
domain under
\href{https://creativecommons.org/publicdomain/zero/1.0/}{CC0~1.0}.
Published sources retain their own rights and are not redistributed.

\IfFileExists{identity.tex}{% Generated source identity; exclude from SOURCE_MANIFEST.sha256.
\par\smallskip\begingroup\footnotesize\raggedright
\noindent Source identity. Commit \texttt{580bd7ef71676ad031eeda31e70f55265ab4bfc3}.\par
\noindent Source-manifest SHA-256:\\
\texttt{7b0e6aa49b797f322a9fd549bfe8bd04}\allowbreak\texttt{d5562f7dfad14b0df74356868d68585c}.\par\endgroup
}{}

\bibliographystyle{plainnat}
\bibliography{references}
\end{document}
